% === DEBUT DU CHUNK preamble+ch1-2 ===
\documentclass[12pt,a4paper,openany]{book}

% ============================================================
%  ENCODAGE ET LANGUE
% ============================================================
\usepackage[utf8]{inputenc}
\usepackage[T1]{fontenc}
\usepackage[french]{babel}
\usepackage[margin=2.5cm]{geometry}

% ============================================================
%  MATHÉMATIQUES
% ============================================================
\usepackage{amsmath,amssymb,amsthm,mathrsfs,mathtools}

% ============================================================
%  GRAPHIQUES ET DIAGRAMMES
% ============================================================
\usepackage{tikz}
\usetikzlibrary{arrows.meta,calc,positioning,patterns,decorations.markings,decorations.pathmorphing,cd}
\usepackage{pgfplots}
\pgfplotsset{compat=1.18}
\usepackage{tikz-cd}

% ============================================================
%  BOÎTES COLORÉES
% ============================================================
\usepackage[most]{tcolorbox}
\tcbuselibrary{theorems,skins,breakable}

% ============================================================
%  HYPERLIENS ET RÉFÉRENCES
% ============================================================
\usepackage[colorlinks=true,linkcolor=blue!60!black,citecolor=green!50!black]{hyperref}

% ============================================================
%  INDEX
% ============================================================
\usepackage{makeidx}
\makeindex

% ============================================================
%  MISE EN PAGE ET TYPOGRAPHIE
% ============================================================
\usepackage{enumitem}
\usepackage{fancyhdr}
\usepackage{caption}
\usepackage{booktabs}
\usepackage{array}
\usepackage{longtable}
\usepackage{microtype}
\usepackage{xcolor}
\usepackage{minitoc}

% ============================================================
%  STYLE DE PAGE
% ============================================================
\pagestyle{fancy}
\fancyhf{}
\fancyhead[LE]{\textsl{\leftmark}}
\fancyhead[RO]{\textsl{\rightmark}}
\fancyfoot[C]{\thepage}
\renewcommand{\headrulewidth}{0.4pt}
\renewcommand{\footrulewidth}{0pt}

% ============================================================
%  ENVIRONNEMENTS DE THÉORÈMES (amsthm)
% ============================================================
\theoremstyle{definition}
\newtheorem{definition}{Définition}[chapter]
\newtheorem{exemple}[definition]{Exemple}
\newtheorem{exercice}{Exercice}[chapter]

\theoremstyle{plain}
\newtheorem{theoreme}[definition]{Théorème}
\newtheorem{proposition}[definition]{Proposition}
\newtheorem{lemme}[definition]{Lemme}
\newtheorem{corollaire}[definition]{Corollaire}

\theoremstyle{remark}
\newtheorem{remarque}[definition]{Remarque}
\newtheorem{notation}[definition]{Notation}

% ============================================================
%  HABILLAGE tcolorbox DES ENVIRONNEMENTS
% ============================================================
\tcolorboxenvironment{definition}{enhanced,breakable,colback=blue!5,colframe=blue!70!black,fonttitle=\bfseries,before skip=10pt,after skip=10pt,left=8pt,right=8pt}
\tcolorboxenvironment{theoreme}{enhanced,breakable,colback=red!5,colframe=red!70!black,fonttitle=\bfseries,before skip=10pt,after skip=10pt,left=8pt,right=8pt}
\tcolorboxenvironment{proposition}{enhanced,breakable,colback=orange!5,colframe=orange!70!black,fonttitle=\bfseries,before skip=10pt,after skip=10pt,left=8pt,right=8pt}
\tcolorboxenvironment{lemme}{enhanced,breakable,colback=yellow!5,colframe=yellow!50!black,fonttitle=\bfseries,before skip=10pt,after skip=10pt,left=8pt,right=8pt}
\tcolorboxenvironment{corollaire}{enhanced,breakable,colback=green!5,colframe=green!50!black,fonttitle=\bfseries,before skip=10pt,after skip=10pt,left=8pt,right=8pt}
\tcolorboxenvironment{remarque}{enhanced,breakable,colback=gray!5,colframe=gray!50!black,before skip=8pt,after skip=8pt,left=6pt,right=6pt}
\tcolorboxenvironment{exemple}{enhanced,breakable,colback=purple!5,colframe=purple!50!black,before skip=8pt,after skip=8pt,left=6pt,right=6pt}

% ============================================================
%  COMMANDES PERSONNALISÉES
% ============================================================
\newcommand{\N}{\mathbb{N}}
\newcommand{\Z}{\mathbb{Z}}
\newcommand{\Q}{\mathbb{Q}}
\newcommand{\R}{\mathbb{R}}
\newcommand{\C}{\mathbb{C}}
\newcommand{\RP}{\mathbb{R}\mathrm{P}}
\newcommand{\CP}{\mathbb{C}\mathrm{P}}
\DeclareMathOperator{\Hom}{Hom}
\DeclareMathOperator{\im}{im}
\DeclareMathOperator{\coker}{coker}
\DeclareMathOperator{\Aut}{Aut}
\DeclareMathOperator{\Id}{Id}
\newcommand{\abs}[1]{\left|#1\right|}
\newcommand{\set}[1]{\left\{#1\right\}}
\newcommand{\rel}{\operatorname{rel}}

% ============================================================
%  PAGE DE TITRE
% ============================================================
% ============================================================
\begin{document}
% ============================================================
\dominitoc

\begin{titlepage}
\begin{tikzpicture}[remember picture, overlay]
  \fill[left color=blue!15!white, right color=blue!5!white]
    (current page.south west) rectangle (current page.north east);
  \fill[color=orange!70!yellow]
    ([yshift=-2cm]current page.north west) rectangle ([yshift=-2.3cm]current page.north east);
  \fill[color=blue!80!black, rounded corners=8pt]
    ([xshift=3cm, yshift=-4cm]current page.north west) rectangle
    ([xshift=-3cm, yshift=-10cm]current page.north east);
  \node[anchor=center, text=white, font=\Huge\bfseries]
    at ([yshift=-6cm]current page.north) {Topologie Alg\'ebrique};
  \node[anchor=center, text=orange!80!yellow, font=\Large]
    at ([yshift=-7.5cm]current page.north) {Notes de Cours};
  \node[anchor=center, text=white!80, font=\large]
    at ([yshift=-8.8cm]current page.north) {Master M1 --- 2025--2026};
  \node[anchor=center, text=blue!80!black, font=\Large\itshape]
    at ([yshift=-12cm]current page.north) {Ya\'e Ulrich Gaba};
  \draw[orange!70!yellow, line width=2pt]
    ([xshift=5cm, yshift=-13.5cm]current page.north west) --
    ([xshift=-5cm, yshift=-13.5cm]current page.north east);
  \node[anchor=center, text width=12cm, align=center, text=blue!60!black, font=\itshape]
    at ([yshift=-16cm]current page.north) {<<\,La topologie alg\'ebrique est l'art de transformer des probl\`emes de g\'eom\'etrie en probl\`emes d'alg\`ebre.\,>>\\[4pt]--- Henri Poincar\'e};
  \node[anchor=south, text=gray, font=\small]
    at ([yshift=1.5cm]current page.south) {\today};
\end{tikzpicture}
\end{titlepage}
\tableofcontents

% ============================================================
%  TABLE DE NOTATIONS
% ============================================================
\chapter*{Table des notations}
\addcontentsline{toc}{chapter}{Table des notations}
\markboth{TABLE DES NOTATIONS}{TABLE DES NOTATIONS}

\begin{longtable}{>{\raggedright}p{4.5cm} p{9cm}}
\toprule
\textbf{Notation} & \textbf{Signification} \\
\midrule
\endfirsthead
\toprule
\textbf{Notation} & \textbf{Signification} \\
\midrule
\endhead
\midrule
\multicolumn{2}{r}{\emph{Suite page suivante}} \\
\bottomrule
\endfoot
\bottomrule
\endlastfoot
$X, Y, Z$ & Espaces topologiques \\
$\mathbf{Top}$ & Catégorie des espaces topologiques et applications continues \\
$\mathbf{Top}_*$ & Catégorie des espaces topologiques pointés \\
$\mathbf{hTop}$ & Catégorie homotopique \\
$\mathbf{Grp}$ & Catégorie des groupes \\
$\mathbf{Ab}$ & Catégorie des groupes abéliens \\
$[X,Y]$ & Ensemble des classes d'homotopie d'applications $X \to Y$ \\
$f \simeq g$ & $f$ est homotope à $g$ \\
$X \simeq Y$ & $X$ et $Y$ ont le même type d'homotopie \\
$\pi_1(X, x_0)$ & Groupe fondamental\index{groupe fondamental} de $X$ en $x_0$ \\
$\pi_n(X, x_0)$ & $n$-ième groupe d'homotopie de $X$ en $x_0$ \\
$H_n(X)$ & $n$-ième groupe d'homologie singulière\index{homologie singulière} de $X$ \\
$H_n(X, A)$ & Homologie singulière relative du couple $(X, A)$ \\
$H^n(X)$ & $n$-ième groupe de cohomologie singulière\index{cohomologie singulière} de $X$ \\
$H^n(X; G)$ & Cohomologie à coefficients dans $G$ \\
$\partial_n$ & Opérateur bord\index{opérateur bord} $C_n(X) \to C_{n-1}(X)$ \\
$\delta^n$ & Opérateur cobord $C^n(X) \to C^{n+1}(X)$ \\
$\Delta^n$ & Simplexe standard\index{simplexe standard} de dimension $n$ \\
$\sigma : \Delta^n \to X$ & Simplexe singulier\index{simplexe singulier} \\
$C_n(X)$ & Groupe des chaînes singulières de dimension $n$ \\
$Z_n(X)$ & Groupe des cycles de dimension $n$ \\
$B_n(X)$ & Groupe des bords de dimension $n$ \\
$\smile$ & Produit cup\index{produit cup} en cohomologie \\
$\frown$ & Produit cap\index{produit cap} \\
$\chi(X)$ & Caractéristique d'Euler\index{caractéristique d'Euler} de $X$ \\
$S^n$ & Sphère\index{sphère} de dimension $n$ \\
$D^n$ & Disque (boule fermée) de dimension $n$ \\
$T^n$ & Tore\index{tore} de dimension $n$ \\
$\RP^n$ & Espace projectif réel\index{espace projectif réel} de dimension $n$ \\
$\CP^n$ & Espace projectif complexe\index{espace projectif complexe} de dimension $n$ \\
$X \vee Y$ & Bouquet\index{bouquet} (wedge sum) de $X$ et $Y$ \\
$X \times Y$ & Produit cartésien \\
$X / A$ & Espace quotient\index{espace quotient} de $X$ par $A$ \\
$\widetilde{X}$ & Revêtement universel\index{revêtement universel} de $X$ \\
$f_*$ & Morphisme induit en homologie ou sur $\pi_1$ \\
$f^*$ & Morphisme induit en cohomologie \\
$\Hom(G, H)$ & Groupe des homomorphismes de $G$ dans $H$ \\
$\im f$ & Image de $f$ \\
$\coker f$ & Conoyau de $f$ \\
$\Aut(G)$ & Groupe des automorphismes de $G$ \\
$\Id$ & Application identité \\
\end{longtable}

% ============================================================
%  PRÉFACE
% ============================================================
\chapter*{Préface}
\addcontentsline{toc}{chapter}{Préface}
\markboth{PRÉFACE}{PRÉFACE}

La topologie algébrique constitue l'une des branches les plus fécondes des
mathématiques du vingtième siècle. Son ambition fondamentale est d'associer à
chaque espace topologique des invariants algébriques --- groupes, anneaux,
modules --- qui permettent de distinguer des espaces que la topologie ensembliste
seule ne parvient pas à séparer. Ce passage du continu au discret, du
géométrique à l'algébrique, est au c\oe ur de la pensée mathématique moderne.

Les origines de cette discipline remontent aux travaux pionniers d'\textbf{Henri
Poincaré}\index{Poincaré, Henri}, qui introduisit en 1895 dans son
\emph{Analysis Situs} la notion de groupe fondamental et les premiers invariants
homologiques. Poincaré comprit que pour classifier les variétés, il fallait leur
associer des quantités algébriques préservées par déformation continue. Ses
intuitions, parfois formulées de manière informelle, ont guidé un siècle entier
de développements.

La formalisation rigoureuse de l'homologie singulière est due en grande partie à
\textbf{Emmy Noether}\index{Noether, Emmy}, qui, dans les années 1920--1930,
reconnut que les \og nombres de Betti \fg{} de Poincaré devaient être compris
comme les rangs de \emph{groupes} d'homologie, et non comme de simples entiers.
Cette vision algébrique a transformé la topologie en une théorie structurelle
profonde.

L'axiomatisation de la théorie fut achevée par \textbf{Eilenberg} et
\textbf{Steenrod}\index{Eilenberg, Samuel}\index{Steenrod, Norman} dans leur
ouvrage fondamental de 1952, \emph{Foundations of Algebraic Topology}. Ils
identifièrent un ensemble d'axiomes --- exactitude, excision, homotopie,
dimension --- qui caractérisent essentiellement de façon unique la théorie de
l'homologie sur la catégorie des couples de CW-complexes. Cette approche
axiomatique a permis d'unifier les différentes constructions (simpliciale,
singulière, cellulaire, de \v{C}ech) en un cadre conceptuel cohérent.

Le présent cours s'adresse aux étudiants de Master et de Doctorat. Il suppose
acquises les bases de la topologie générale (espaces métriques, compacité,
connexité) ainsi qu'une familiarité avec l'algèbre (groupes, anneaux, modules,
suites exactes). Nous développerons la théorie selon le plan suivant :

\begin{enumerate}[label=\textbullet]
\item Le \textbf{groupe fondamental} $\pi_1(X, x_0)$\index{groupe fondamental},
premier invariant algébrique, avec ses applications classiques (théorème de
Brouwer, théorème de Borsuk--Ulam en dimension 2).
\item La théorie des \textbf{revêtements}\index{revêtement}, intimement liée au
groupe fondamental par la correspondance de Galois topologique.
\item L'\textbf{homologie singulière}\index{homologie singulière} $H_n(X)$,
invariant plus puissant, avec les axiomes d'Eilenberg--Steenrod, la suite exacte
longue du couple et le théorème d'excision.
\item La \textbf{cohomologie}\index{cohomologie} $H^n(X; G)$ et la structure
d'anneau donnée par le produit cup $\smile$.
\item Les outils avancés : dualité de Poincaré, formule de
Künneth, coefficients universels.
\end{enumerate}

Notre approche sera résolument rigoureuse, dans la tradition de Godement, Serre
et Spanier, tout en accordant une place importante aux exemples et aux exercices.
Les diagrammes commutatifs et les illustrations géométriques jalonnent le texte
afin de cultiver l'intuition autant que la rigueur.

\bigskip
\hfill \emph{Bonne lecture.}

% ============================================================
% ============================================================
%  CHAPITRE 1 : RAPPELS DE TOPOLOGIE ET MOTIVATIONS
% ============================================================
% ============================================================
\chapter{Rappels de topologie et motivations}\label{chap:rappels}

\section{Espaces topologiques, continuité, homéomorphismes}\label{sec:espaces-top}

Nous commençons par rappeler les notions fondamentales de la topologie générale
qui seront utilisées de manière constante dans la suite.

\begin{definition}[Espace topologique]\label{def:espace-top}
\index{espace topologique}
Un \textbf{espace topologique} est un couple $(X, \mathcal{T})$ où $X$ est un
ensemble et $\mathcal{T} \subseteq \mathcal{P}(X)$ est une famille de
sous-ensembles de $X$, appelés \textbf{ouverts}\index{ouvert}, vérifiant :
\begin{enumerate}[label=(\roman*)]
\item $\varnothing \in \mathcal{T}$ et $X \in \mathcal{T}$.
\item $\mathcal{T}$ est stable par réunion quelconque : si $(U_i)_{i \in I}$ est
une famille d'éléments de $\mathcal{T}$, alors $\bigcup_{i \in I} U_i \in
\mathcal{T}$.
\item $\mathcal{T}$ est stable par intersection finie : si $U_1, \ldots, U_n \in
\mathcal{T}$, alors $U_1 \cap \cdots \cap U_n \in \mathcal{T}$.
\end{enumerate}
\end{definition}

\begin{definition}[Application continue]\label{def:continuite}
\index{application continue}
Soient $(X, \mathcal{T}_X)$ et $(Y, \mathcal{T}_Y)$ deux espaces topologiques.
Une application $f : X \to Y$ est dite \textbf{continue} si pour tout ouvert
$V \in \mathcal{T}_Y$, l'image réciproque $f^{-1}(V)$ est un ouvert de $X$ :
\[
\forall\, V \in \mathcal{T}_Y, \quad f^{-1}(V) \in \mathcal{T}_X.
\]
\end{definition}

\begin{definition}[Homéomorphisme]\label{def:homeomorphisme}
\index{homéomorphisme}
Une application $f : X \to Y$ est un \textbf{homéomorphisme} si $f$ est
bijective, continue, et si son inverse $f^{-1} : Y \to X$ est également
continue. On écrit alors $X \cong Y$.
\end{definition}

\begin{remarque}\label{rem:homeo-equivalence}
La relation \og être homéomorphe \fg{} est une relation d'équivalence sur la
classe des espaces topologiques. L'un des objectifs centraux de la topologie est
de classifier les espaces à homéomorphisme près. Cependant, cette classification
est en général inaccessible directement, ce qui motive l'introduction
d'invariants algébriques.
\end{remarque}

\begin{exemple}\label{ex:homeo-intervalles}
L'application $f : {]}{-1},1{[} \to \R$ définie par $f(t) = \tan(\pi t / 2)$
est un homéomorphisme. Ainsi l'intervalle ouvert borné ${]}{-1},1{[}$ est
homéomorphe à $\R$ tout entier. De même, toute boule ouverte de $\R^n$ est
homéomorphe à $\R^n$.
\end{exemple}

\begin{exemple}[Topologie quotient]\label{ex:topologie-quotient}
\index{topologie quotient}
Soit $X$ un espace topologique et $\sim$ une relation d'équivalence sur $X$.
L'\textbf{espace quotient} $X/{\sim}$ est l'ensemble des classes d'équivalence
muni de la topologie la plus fine rendant continue la projection canonique
$\pi : X \to X/{\sim}$. Concrètement, $U \subseteq X/{\sim}$ est ouvert si et
seulement si $\pi^{-1}(U)$ est ouvert dans $X$.

Voici quelques exemples fondamentaux :
\begin{enumerate}[label=(\alph*)]
\item $[0,1]/(0 \sim 1) \cong S^1$ : on identifie les extrémités de l'intervalle.
\item $D^n / S^{n-1} \cong S^n$ : on écrase le bord du disque en un point.
\item $S^n / (x \sim -x) \cong \RP^n$ : on identifie les points antipodaux.
\end{enumerate}
\end{exemple}

\begin{definition}[Plongement topologique]\label{def:plongement}
\index{plongement}
Une application continue $f : X \to Y$ est un \textbf{plongement} si $f$ induit
un homéomorphisme de $X$ sur $f(X)$ muni de la topologie induite.
\end{definition}

\begin{proposition}\label{prop:plongement-compact}
Soit $f : X \to Y$ une application continue injective, avec $X$ compact et $Y$
de Hausdorff. Alors $f$ est un plongement.
\end{proposition}

\begin{proof}
L'application $f : X \to f(X)$ est continue, bijective, $X$ est compact et
$f(X)$ est de Hausdorff (comme sous-espace d'un espace de Hausdorff). Par la

ef{prop:compact-hausdorff}, $f$ est un homéomorphisme sur son image.
\end{proof}

% ============================================================
\section{Connexité et connexité par arcs}\label{sec:connexite}

\begin{definition}[Espace connexe]\label{def:connexe}
\index{connexité}
Un espace topologique $X$ est dit \textbf{connexe} s'il ne peut pas s'écrire
comme réunion disjointe de deux ouverts non vides. Autrement dit, les seules
parties de $X$ qui sont à la fois ouvertes et fermées sont $\varnothing$ et $X$.
\end{definition}

\begin{definition}[Connexité par arcs]\label{def:connexe-arcs}
\index{connexité par arcs}
Un espace $X$ est \textbf{connexe par arcs} si pour tous $x, y \in X$, il existe
un chemin continu $\gamma : [0,1] \to X$ tel que $\gamma(0) = x$ et
$\gamma(1) = y$.
\end{definition}

\begin{proposition}\label{prop:arc-connexe-implique-connexe}
Tout espace connexe par arcs est connexe.
\end{proposition}

\begin{proof}
Supposons $X$ connexe par arcs et supposons par l'absurde que $X = U \cup V$
avec $U, V$ ouverts non vides et $U \cap V = \varnothing$. Choisissons
$x \in U$ et $y \in V$. Il existe un chemin $\gamma : [0,1] \to X$ tel que
$\gamma(0) = x$ et $\gamma(1) = y$. Les ensembles $\gamma^{-1}(U)$ et
$\gamma^{-1}(V)$ forment une partition de $[0,1]$ en deux ouverts non vides,
ce qui contredit la connexité de $[0,1]$.
\end{proof}

\begin{remarque}\label{rem:connexe-pas-arcs}
La réciproque est fausse en général. Le \og peigne du topologiste\fg{}
\index{peigne du topologiste} fournit un exemple classique d'espace connexe qui
n'est pas connexe par arcs. Cependant, pour les espaces localement connexes par
arcs (en particulier les variétés topologiques et les CW-complexes), les deux
notions coïncident.
\end{remarque}

\begin{definition}[Composantes connexes par arcs]\label{def:composantes-arcs}
\index{composantes connexes par arcs}
Soit $X$ un espace topologique. La relation \og $x$ et $y$ sont reliés par un
chemin \fg{} est une relation d'équivalence. Les classes d'équivalence sont
appelées \textbf{composantes connexes par arcs} de $X$. On note $\pi_0(X)$
\index{$\pi_0(X)$} l'ensemble des composantes connexes par arcs de $X$.
\end{definition}

\begin{remarque}\label{rem:pi0-foncteur}
L'ensemble $\pi_0(X)$ est un invariant homotopique : si $f : X \to Y$ est une
équivalence d'homotopie, alors $f$ induit une bijection $\pi_0(X) \cong \pi_0(Y)$.
On peut voir $\pi_0$ comme un foncteur $\pi_0 : \mathbf{Top} \to \mathbf{Ens}$.
\end{remarque}

\begin{exemple}\label{ex:composantes-On}
Le groupe orthogonal $O(n) = \set{A \in M_n(\R) : A^T A = I_n}$ a exactement
deux composantes connexes par arcs : $SO(n) = \set{A \in O(n) : \det A = 1}$
et son complémentaire. Ainsi $\pi_0(O(n)) \cong \Z/2\Z$.
\end{exemple}

% ============================================================
\section{Compacité et espaces de Hausdorff}\label{sec:compacite}

\begin{definition}[Espace compact]\label{def:compact}
\index{compacité}
Un espace topologique $X$ est \textbf{compact} si de tout recouvrement ouvert de
$X$ on peut extraire un sous-recouvrement fini : pour toute famille
$(U_i)_{i \in I}$ d'ouverts telle que $X = \bigcup_{i \in I} U_i$, il existe
$i_1, \ldots, i_n \in I$ tels que $X = U_{i_1} \cup \cdots \cup U_{i_n}$.
\end{definition}

\begin{definition}[Espace de Hausdorff]\label{def:hausdorff}
\index{espace de Hausdorff}
Un espace topologique $X$ est dit \textbf{de Hausdorff} (ou \textbf{séparé}) si
pour tous points distincts $x, y \in X$, il existe des ouverts disjoints $U$ et
$V$ tels que $x \in U$ et $y \in V$.
\end{definition}

\begin{proposition}\label{prop:compact-hausdorff}
Soit $f : X \to Y$ une application continue bijective, avec $X$ compact et $Y$
de Hausdorff. Alors $f$ est un homéomorphisme.
\end{proposition}

\begin{proof}
Il suffit de montrer que $f$ est fermée. Soit $F \subseteq X$ un fermé. Comme
$X$ est compact, $F$ est compact. L'image continue d'un compact étant compacte,
$f(F)$ est compact dans $Y$. Or un compact dans un espace de Hausdorff est
fermé, donc $f(F)$ est fermé dans $Y$. Ainsi $f$ est une application fermée
bijective, donc $f^{-1}$ est continue.
\end{proof}

\begin{remarque}\label{rem:compact-hausdorff-pratique}
Cette proposition est un outil extrêmement utile en pratique : pour montrer
qu'une bijection continue est un homéomorphisme, il suffit souvent de vérifier
les conditions de compacité et de séparation.
\end{remarque}

\begin{theoreme}[Tychonov]\label{thm:tychonov}
\index{théorème de Tychonov}
Le produit $\prod_{i \in I} X_i$ d'espaces compacts (muni de la topologie
produit) est compact.
\end{theoreme}

Nous ne redémontrons pas ce théorème (qui requiert l'axiome du choix) mais
l'utiliserons librement.

\begin{exemple}[Espaces compacts fondamentaux]\label{ex:compacts-fondamentaux}
Voici les espaces compacts que nous rencontrerons constamment :
\begin{enumerate}[label=(\alph*)]
\item Le disque fermé $D^n = \set{x \in \R^n : \|x\| \leqslant 1}$.
\item La sphère $S^n = \set{x \in \R^{n+1} : \|x\| = 1} = \partial D^{n+1}$.
\item Le tore $T^n = (S^1)^n$ (produit fini de compacts).
\item L'espace projectif $\RP^n = S^n/(x \sim -x)$ (quotient d'un compact).
\item L'espace projectif complexe $\CP^n$ (quotient de $S^{2n+1}$).
\end{enumerate}
Tous ces espaces sont compacts et de Hausdorff.
\end{exemple}

% ============================================================
\section{Homotopie entre applications}\label{sec:homotopie}

La notion d'homotopie est la pierre angulaire de la topologie algébrique. Elle
formalise l'idée intuitive de \og déformation continue \fg{} d'une application en
une autre.

\begin{definition}[Homotopie]\label{def:homotopie}
\index{homotopie}
Soient $X, Y$ des espaces topologiques et $f, g : X \to Y$ deux applications
continues. Une \textbf{homotopie} de $f$ à $g$ est une application continue
$H : X \times [0,1] \to Y$ telle que
\[
H(x, 0) = f(x) \quad \text{et} \quad H(x, 1) = g(x) \quad \text{pour tout } x \in X.
\]
On note $f \simeq g$\index{$\simeq$} et on dit que $f$ et $g$ sont
\textbf{homotopes}. Si de plus il existe un sous-espace $A \subseteq X$ tel que
$H(a, t) = f(a) = g(a)$ pour tout $a \in A$ et tout $t \in [0,1]$, on parle
d'\textbf{homotopie relative à $A$} et on note $f \simeq g \pmod{A}$ ou
$f \simeq g \rel A$.
\end{definition}

\begin{proposition}\label{prop:homotopie-equivalence}
La relation d'homotopie $\simeq$ est une relation d'équivalence sur l'ensemble
$\mathcal{C}(X, Y)$ des applications continues de $X$ dans $Y$.
\end{proposition}

\begin{proof}
\emph{Réflexivité.} L'homotopie constante $H(x,t) = f(x)$ montre $f \simeq f$.

\emph{Symétrie.} Si $H : f \simeq g$, alors $\overline{H}(x,t) = H(x, 1-t)$
réalise $g \simeq f$.

\emph{Transitivité.} Si $H : f \simeq g$ et $K : g \simeq h$, alors
\[
L(x,t) = \begin{cases} H(x, 2t) & \text{si } 0 \leqslant t \leqslant \tfrac{1}{2}, \\
K(x, 2t-1) & \text{si } \tfrac{1}{2} \leqslant t \leqslant 1,
\end{cases}
\]
est continue par le lemme de recollement et réalise $f \simeq h$.
\end{proof}

\begin{notation}\label{not:classes-homotopie}
On note $[X, Y]$\index{$[X,Y]$} l'ensemble des classes d'homotopie
d'applications continues de $X$ dans $Y$. C'est le quotient
$\mathcal{C}(X,Y)/{\simeq}$.
\end{notation}

\begin{definition}[Équivalence d'homotopie]\label{def:equiv-homotopie}
\index{équivalence d'homotopie}
Une application continue $f : X \to Y$ est une \textbf{équivalence d'homotopie}
s'il existe une application continue $g : Y \to X$ telle que
$g \circ f \simeq \Id_X$ et $f \circ g \simeq \Id_Y$. On dit alors que $X$ et
$Y$ ont le même \textbf{type d'homotopie}\index{type d'homotopie} et on note
$X \simeq Y$.
\end{definition}

\begin{definition}[Espace contractile]\label{def:contractile}
\index{espace contractile}
Un espace topologique $X$ est \textbf{contractile} s'il a le type d'homotopie
d'un point, c'est-à-dire si $\Id_X$ est homotope à une application constante.
\end{definition}

\begin{exemple}\label{ex:Rn-contractile}
L'espace $\R^n$ est contractile : l'homotopie $H(x,t) = (1-t)x$ réalise une
contraction vers l'origine. De même, tout espace convexe de $\R^n$ est
contractile.
\end{exemple}

\begin{exemple}[Espaces non contractiles]\label{ex:non-contractile}
La sphère $S^n$ (pour $n \geqslant 1$) n'est pas contractile. En effet, nous
montrerons que $\pi_1(S^1) \cong \Z \neq 0$ et, pour $n \geqslant 2$, que
$H_n(S^n) \cong \Z \neq 0$. Si $S^n$ était contractile, ces invariants seraient
triviaux, d'où contradiction.
\end{exemple}

\begin{proposition}[Homotopie et composition]\label{prop:homotopie-composition}
\index{homotopie!et composition}
Soient $f_0, f_1 : X \to Y$ et $g_0, g_1 : Y \to Z$ des applications continues
avec $f_0 \simeq f_1$ et $g_0 \simeq g_1$. Alors $g_0 \circ f_0 \simeq g_1 \circ f_1$.
\end{proposition}

\begin{proof}
Par transitivité, il suffit de montrer : (a) $g_0 \circ f_0 \simeq g_0 \circ f_1$,
et (b) $g_0 \circ f_1 \simeq g_1 \circ f_1$.

Pour (a), si $H : X \times [0,1] \to Y$ est une homotopie de $f_0$ à $f_1$,
alors $g_0 \circ H : X \times [0,1] \to Z$ est une homotopie de $g_0 \circ f_0$
à $g_0 \circ f_1$.

Pour (b), si $K : Y \times [0,1] \to Z$ est une homotopie de $g_0$ à $g_1$,
alors $L(x,t) = K(f_1(x), t)$ est une homotopie de $g_0 \circ f_1$ à
$g_1 \circ f_1$.
\end{proof}

% ============================================================
\section{Rétractes et rétractes par déformation}\label{sec:retractes}

\begin{definition}[Rétraction et rétracte]\label{def:retracte}
\index{rétraction}\index{rétracte}
Soit $A \subseteq X$ un sous-espace. Une \textbf{rétraction} de $X$ sur $A$ est
une application continue $r : X \to A$ telle que $r|_A = \Id_A$, c'est-à-dire
$r(a) = a$ pour tout $a \in A$. On dit alors que $A$ est un \textbf{rétracte}
de $X$.
\end{definition}

\begin{definition}[Rétracte par déformation]\label{def:retracte-deformation}
\index{rétracte par déformation}
Un sous-espace $A \subseteq X$ est un \textbf{rétracte par déformation} de $X$
s'il existe une application continue $H : X \times [0,1] \to X$ telle que :
\begin{enumerate}[label=(\roman*)]
\item $H(x, 0) = x$ pour tout $x \in X$,
\item $H(x, 1) \in A$ pour tout $x \in X$,
\item $H(a, t) = a$ pour tout $a \in A$ et tout $t \in [0,1]$.
\end{enumerate}
L'application $r : x \mapsto H(x,1)$ est alors une rétraction, et l'inclusion
$\iota : A \hookrightarrow X$ est une équivalence d'homotopie d'inverse
homotopique $r$.
\end{definition}

\begin{exemple}\label{ex:retracte-Rn-0}
Le cercle $S^1$ est un rétracte par déformation de $\R^2 \setminus \{0\}$ via
$H(x,t) = (1-t)x + t \cdot x/\|x\|$. Plus généralement, $S^{n-1}$ est un
rétracte par déformation de $\R^n \setminus \{0\}$.
\end{exemple}

\begin{proposition}\label{prop:retracte-def-equiv-homotopie}
Si $A$ est un rétracte par déformation de $X$, alors l'inclusion
$\iota : A \hookrightarrow X$ est une équivalence d'homotopie.
\end{proposition}

\begin{proof}
Soit $r : X \to A$ la rétraction associée. On a $r \circ \iota = \Id_A$ et
$\iota \circ r \simeq \Id_X$ via l'homotopie de déformation $H$.
\end{proof}

% ============================================================
\section{CW-complexes}\label{sec:CW}

Les CW-complexes\index{CW-complexe}, introduits par J.~H.~C.
Whitehead\index{Whitehead, J.H.C.}, forment la classe d'espaces topologiques la
mieux adaptée aux méthodes de la topologie algébrique.

\begin{definition}[CW-complexe]\label{def:CW-complexe}
Un \textbf{CW-complexe} est un espace topologique $X$ muni d'une filtration
\[
\varnothing = X^{-1} \subseteq X^0 \subseteq X^1 \subseteq X^2 \subseteq \cdots \subseteq X = \bigcup_{n \geqslant 0} X^n,
\]
où chaque \textbf{squelette}\index{squelette} $X^n$ est obtenu à partir de
$X^{n-1}$ par attachement de cellules\index{cellule} de dimension $n$ : il
existe un ensemble d'indices $J_n$ et pour chaque $\alpha \in J_n$ une
\textbf{application caractéristique}\index{application caractéristique}
$\Phi_\alpha : D^n \to X^n$ telle que :
\begin{enumerate}[label=(\roman*)]
\item $\Phi_\alpha|_{\mathring{D}^n}$ est un homéomorphisme sur son image
$e_\alpha^n$ (la cellule ouverte),
\item $\Phi_\alpha(S^{n-1}) \subseteq X^{n-1}$ (l'application d'attachement),
\item $X^n$ est l'espace quotient de $X^{n-1} \sqcup \bigsqcup_{\alpha \in
J_n} D^n$ obtenu en identifiant chaque $x \in S^{n-1} \subseteq D^n_\alpha$ avec
$\Phi_\alpha(x) \in X^{n-1}$.
\end{enumerate}
L'espace $X$ est muni de la topologie faible : un ensemble $A \subseteq X$ est
fermé si et seulement si $A \cap X^n$ est fermé dans $X^n$ pour tout $n$.
\end{definition}

\begin{exemple}[CW-structure de $S^2$]\label{ex:CW-S2}
\index{sphère!CW-structure}
La sphère $S^2$ admet une CW-structure minimale avec une $0$-cellule $e^0$ et
une $2$-cellule $e^2$ attachée via l'application constante
$\varphi : S^1 \to e^0$. On a donc $S^2 = e^0 \cup e^2$.

\begin{center}
\begin{tikzpicture}[scale=1.2]
  % 0-squelette
  \fill (0,0) circle (2pt) node[below] {$e^0$};
  \draw[-{Stealth},thick] (1,0) -- (2.5,0) node[midway,above] {$\Phi$};
  % sphère
  \begin{scope}[xshift=4.5cm]
    \shade[ball color=red!15!white,opacity=0.8] (0,0) circle (1);
    \draw (0,0) circle (1);
    \fill (0,-1) circle (2pt) node[below] {$e^0$};
    \node at (0,0.3) {$e^2$};
  \end{scope}
\end{tikzpicture}
\captionof{figure}{CW-structure minimale de $S^2$ : un point et un disque.}\label{fig:CW-S2}
\end{center}
\end{exemple}

\begin{exemple}[CW-structure du tore $T^2$]\label{ex:CW-T2}
\index{tore!CW-structure}
Le tore $T^2 = S^1 \times S^1$ admet une CW-structure avec une $0$-cellule
$e^0$, deux $1$-cellules $a$ et $b$, et une $2$-cellule $e^2$ attachée via le
mot $aba^{-1}b^{-1}$.

\begin{center}
\begin{tikzpicture}[scale=1.8]
  % carré fondamental
  \draw[thick,blue!70!black,
        decoration={markings, mark=at position 0.55 with {\arrow{Stealth}}},
        postaction={decorate}]
    (0,0) -- (2,0) node[midway,below] {$a$};
  \draw[thick,red!70!black,
        decoration={markings, mark=at position 0.55 with {\arrow{Stealth}}},
        postaction={decorate}]
    (2,0) -- (2,2) node[midway,right] {$b$};
  \draw[thick,blue!70!black,
        decoration={markings, mark=at position 0.55 with {\arrow{Stealth}}},
        postaction={decorate}]
    (0,2) -- (2,2) node[midway,above] {$a$};
  \draw[thick,red!70!black,
        decoration={markings, mark=at position 0.55 with {\arrow{Stealth}}},
        postaction={decorate}]
    (0,0) -- (0,2) node[midway,left] {$b$};
  \fill (0,0) circle (1.5pt);
  \fill (2,0) circle (1.5pt);
  \fill (0,2) circle (1.5pt);
  \fill (2,2) circle (1.5pt);
  \node at (1,1) {$e^2$};
  \node[below left] at (0,0) {$e^0$};
\end{tikzpicture}
\captionof{figure}{CW-structure du tore : identification des bords du carré.}\label{fig:CW-T2}
\end{center}
\end{exemple}

\begin{exemple}[CW-structure de $\RP^2$]\label{ex:CW-RP2}
\index{espace projectif réel!CW-structure}
L'espace projectif réel $\RP^2$ s'obtient en identifiant les points antipodaux de
$S^1$ sur le bord de $D^2$. Sa CW-structure est : une $0$-cellule $e^0$, une
$1$-cellule $e^1$ (cercle obtenu par identification), et une $2$-cellule $e^2$
dont l'application d'attachement $\varphi : S^1 \to X^1$ fait le tour de la
$1$-cellule deux fois.

\begin{center}
\begin{tikzpicture}[scale=1.5]
  \draw[thick,fill=orange!10] (0,0) circle (1);
  \draw[thick,blue!70!black,
        decoration={markings,
          mark=at position 0.25 with {\arrow{Stealth}},
          mark=at position 0.75 with {\arrow{Stealth}}},
        postaction={decorate}]
    (0,0) circle (1);
  \node at (0,0) {$e^2$};
  \node[blue!70!black] at (1.3,0.4) {$e^1$};
  \fill (1,0) circle (1.5pt);
  \fill (-1,0) circle (1.5pt);
  \draw[dashed,gray] (1,0) -- (-1,0);
  \node[below] at (0,-0.05) {\small antipodaux identifiés};
  \node[below right] at (1,0) {$e^0$};
\end{tikzpicture}
\captionof{figure}{CW-structure de $\RP^2$ : le bord du disque est identifié par l'antipode.}\label{fig:CW-RP2}
\end{center}
\end{exemple}

\begin{proposition}\label{prop:CW-proprietes}
Tout CW-complexe est un espace de Hausdorff, paracompact et localement
contractile. En particulier, tout CW-complexe est localement connexe par arcs.
\end{proposition}

\begin{exemple}[CW-structure de $\CP^n$]\label{ex:CW-CPn}
\index{espace projectif complexe!CW-structure}
L'espace projectif complexe $\CP^n$ admet une CW-structure avec exactement une
cellule en chaque dimension paire $0, 2, 4, \ldots, 2n$. En effet, la
décomposition
\[
\CP^n = \C^0 \cup \C^1 \cup \cdots \cup \C^n
\]
(où chaque $\C^k$ est plongé via $[z_0 : \cdots : z_k : 0 : \cdots : 0]$ avec
$z_k = 1$) fournit une cellule $e^{2k} \cong \C^k \cong \R^{2k}$ pour chaque
$0 \leqslant k \leqslant n$. On a donc $\CP^n = e^0 \cup e^2 \cup e^4 \cup
\cdots \cup e^{2n}$.
\end{exemple}

\begin{remarque}[Avantages des CW-complexes]\label{rem:avantages-CW}
Les CW-complexes présentent plusieurs avantages considérables :
\begin{enumerate}[label=(\roman*)]
\item L'\textbf{homologie cellulaire}\index{homologie cellulaire} offre un moyen
de calcul très efficace : le complexe de chaînes $C_n^{\mathrm{CW}}(X)$ est le
groupe abélien libre engendré par les $n$-cellules.
\item Le théorème de Whitehead affirme qu'une application continue entre
CW-complexes qui induit un isomorphisme sur tous les groupes d'homotopie est une
équivalence d'homotopie.
\item Le théorème d'approximation CW affirme que tout espace topologique a le
type d'homotopie d'un CW-complexe (sous des hypothèses légères).
\end{enumerate}
\end{remarque}

\begin{definition}[Sous-complexe]\label{def:sous-complexe}
\index{sous-complexe}
Un \textbf{sous-complexe} d'un CW-complexe $X$ est un sous-espace fermé $A
\subseteq X$ qui est réunion de cellules de $X$ et tel que la fermeture de
chaque cellule de $A$ soit contenue dans $A$.
\end{definition}

% ============================================================
\section{La catégorie homotopique}\label{sec:hTop}

\begin{definition}[Catégorie homotopique $\mathbf{hTop}$]\label{def:hTop}
\index{catégorie homotopique}\index{$\mathbf{hTop}$}
La \textbf{catégorie homotopique} $\mathbf{hTop}$ a pour objets les espaces
topologiques et pour morphismes les classes d'homotopie d'applications continues :
\[
\Hom_{\mathbf{hTop}}(X, Y) = [X, Y].
\]
La composition est induite par la composition des applications : si
$[f] \in [X,Y]$ et $[g] \in [Y,Z]$, alors $[g] \circ [f] := [g \circ f]$.
Cette opération est bien définie car la composition est compatible avec la
relation d'homotopie.
\end{definition}

\begin{remarque}\label{rem:hTop-isomorphismes}
Les isomorphismes dans $\mathbf{hTop}$ sont précisément les classes
d'équivalences d'homotopie. Autrement dit, $X$ et $Y$ sont isomorphes dans
$\mathbf{hTop}$ si et seulement si $X \simeq Y$. Tout invariant de la
topologie algébrique (groupe fondamental, groupes d'homologie, etc.) est un
foncteur défini sur $\mathbf{hTop}$ ou sur une catégorie apparentée.
\end{remarque}

De même, on dispose de la catégorie pointée.

\begin{definition}[Catégorie $\mathbf{hTop}_*$]\label{def:hTopstar}
\index{catégorie homotopique pointée}
La catégorie $\mathbf{hTop}_*$ a pour objets les espaces topologiques pointés
$(X, x_0)$ et pour morphismes les classes d'homotopie pointée
$[(X,x_0),(Y,y_0)]$, où les homotopies préservent le point base.
\end{definition}

% ============================================================
\section{Motivations : pourquoi des invariants algébriques ?}\label{sec:motivations}

La question fondamentale de la topologie est la suivante : étant donnés deux
espaces topologiques $X$ et $Y$, sont-ils homéomorphes ? Ou, de manière plus
souple, ont-ils le même type d'homotopie ? Cette question est en général très
difficile, et l'on a besoin d'outils permettant de \emph{distinguer} les espaces.

\begin{exemple}[Une question simple, une réponse non triviale]\label{ex:motivation-R2-R3}
Les espaces $\R^2$ et $\R^3$ ne sont pas homéomorphes. Bien que l'énoncé
paraisse évident, sa démonstration rigoureuse est délicate et repose sur des
outils de topologie algébrique. En effet, si $f : \R^2 \to \R^3$ était un
homéomorphisme, en ôtant un point on obtiendrait un homéomorphisme
$\R^2 \setminus \{p\} \cong \R^3 \setminus \{f(p)\}$. Or
$\R^2 \setminus \{p\} \simeq S^1$ et $\R^3 \setminus \{q\} \simeq S^2$, et ces
espaces n'ont pas le même type d'homotopie (leurs groupes fondamentaux diffèrent :
$\pi_1(S^1) \cong \Z$ tandis que $\pi_1(S^2) = 0$).
\end{exemple}

\begin{exemple}[Sphère et tore]\label{ex:motivation-S2-T2}
La sphère $S^2$ et le tore $T^2$ sont des surfaces compactes connexes orientées
de genres respectifs $0$ et $1$. Elles ne sont pas homéomorphes, mais comment le
prouver rigoureusement ? L'homologie fournit une réponse nette :
\[
H_1(S^2) = 0 \qquad \text{tandis que} \qquad H_1(T^2) \cong \Z^2.
\]
\end{exemple}

\begin{exemple}[Bande de Möbius et cylindre]\label{ex:motivation-mobius}
\index{bande de Möbius}\index{cylindre}
La bande de Möbius $M$ et le cylindre $C = S^1 \times [0,1]$ ont le même type
d'homotopie (tous deux se rétractent par déformation sur $S^1$), donc
$\pi_1(M) \cong \pi_1(C) \cong \Z$ et $H_*(M) \cong H_*(C)$. Pour les
distinguer, il faut utiliser des invariants plus fins, par exemple l'orientabilité
ou la cohomologie à coefficients dans $\Z/2\Z$ des compactifications.
\end{exemple}

\begin{remarque}\label{rem:strategie-invariants}
La stratégie générale de la topologie algébrique est la suivante :
\begin{enumerate}[label=(\arabic*)]
\item Associer à chaque espace topologique $X$ un objet algébrique $F(X)$
(groupe, anneau, module\ldots).
\item S'assurer que cette construction est \textbf{fonctorielle} : une application
continue $f : X \to Y$ induit un morphisme $F(f) : F(X) \to F(Y)$ (ou dans le
sens opposé pour la cohomologie), de manière compatible avec la composition.
\item Montrer que cet invariant est \textbf{homotopique} : si $f \simeq g$, alors
$F(f) = F(g)$.
\item Utiliser ces invariants pour distinguer des espaces non homéomorphes ou pour
démontrer des résultats de non-existence (pas de rétraction, pas de point fixe,
etc.).
\end{enumerate}
\end{remarque}

Le diagramme fonctoriel fondamental est le suivant :
\[
\begin{tikzcd}[column sep=large, row sep=large]
\mathbf{Top}_* \arrow[r, "\pi_1"] \arrow[d, "\text{quotient}"'] & \mathbf{Grp} \\
\mathbf{hTop}_* \arrow[ru, "\pi_1"', dashed] &
\end{tikzcd}
\]
Le foncteur $\pi_1$ passe au quotient par la relation d'homotopie, ce qui sera
démontré au 
ef{chap:pi1}.

% ============================================================
\section{Exercices}\label{sec:exos-ch1}

\begin{exercice}[$\star\star$]\label{exo:homeo-sphere-bord}
\index{sphère}
Montrer que $S^n$ est homéomorphe au bord $\partial D^{n+1}$ du disque fermé
$D^{n+1} = \set{x \in \R^{n+1} : \|x\| \leqslant 1}$.
\end{exercice}

\begin{exercice}[$\star\star$]\label{exo:produit-connexes}
Montrer que le produit $X \times Y$ de deux espaces connexes par arcs est connexe
par arcs. En déduire que le tore $T^n = (S^1)^n$ est connexe par arcs.
\end{exercice}

\begin{exercice}[$\star\star$]\label{exo:retracte-S1}
Montrer que $S^1$ est un rétracte par déformation de
$\R^2 \setminus \{0\}$. En déduire que $S^1$ n'est pas un rétracte de $D^2$.
\emph{(On pourra admettre que $\pi_1(S^1) \cong \Z$ et $\pi_1(D^2) = 0$.)}
\end{exercice}

\begin{exercice}[$\star\star$]\label{exo:contractile-pi1}
Montrer que tout espace contractile est connexe par arcs. En déduire que si $X$
est contractile, alors $[Y, X]$ est réduit à un seul élément pour tout espace $Y$.
\end{exercice}

\begin{exercice}[$\star\star\star$]\label{exo:CW-Sn}
Décrire une CW-structure sur $S^n$ avec exactement deux cellules (une de dimension
$0$ et une de dimension $n$). Décrire une CW-structure sur $S^n$ avec exactement
$2(n+1)$ cellules (deux en chaque dimension $0, 1, \ldots, n$), obtenue en
voyant $S^n$ comme le bord du $(n+1)$-cube $[-1,1]^{n+1}$.
\end{exercice}

\begin{exercice}[$\star\star\star$]\label{exo:CW-RPn}
Montrer que $\RP^n$ admet une CW-structure avec exactement une cellule en chaque
dimension $0, 1, \ldots, n$. \emph{Indication : considérer la décomposition
$\RP^n = \RP^{n-1} \cup e^n$ et procéder par récurrence.}
\end{exercice}

\begin{exercice}[$\star\star$]\label{exo:homotopie-composition}
Soient $f_0 \simeq f_1 : X \to Y$ et $g_0 \simeq g_1 : Y \to Z$. Montrer que
$g_0 \circ f_0 \simeq g_1 \circ f_1 : X \to Z$. En déduire que la composition
dans $\mathbf{hTop}$ est bien définie.
\end{exercice}

\begin{exercice}[$\star\star\star$]\label{exo:equiv-homotopie-transitive}
Montrer que la relation \og avoir le même type d'homotopie \fg{} est une relation
d'équivalence. \emph{Attention : la transitivité n'est pas triviale.}
\end{exercice}

\begin{exercice}[$\star\star$]\label{exo:cone-contractile}
\index{cône}
Le \textbf{cône} sur un espace $X$ est l'espace $CX = (X \times [0,1]) / (X \times \{1\})$.
Montrer que $CX$ est contractile pour tout espace $X$.
\end{exercice}

\begin{exercice}[$\star\star\star$]\label{exo:suspension}
\index{suspension}
La \textbf{suspension} de $X$ est $\Sigma X = CX / (X \times \{0\})$, ou de
façon équivalente $\Sigma X = (X \times [0,1]) / {\sim}$ où $(x,0) \sim (x',0)$
et $(x,1) \sim (x',1)$ pour tous $x, x'$. Montrer que $\Sigma S^n \cong S^{n+1}$.
\end{exercice}


% ============================================================
% ============================================================
%  CHAPITRE 2 : LE GROUPE FONDAMENTAL
% ============================================================
% ============================================================
\chapter{Le groupe fondamental --- définition et premières propriétés}\label{chap:pi1}

Le groupe fondamental est le premier et le plus accessible des invariants
algébriques associés à un espace topologique. Introduit par Poincaré en 1895, il
encode l'information sur les lacets que l'on peut tracer dans un espace, à
déformation continue près.

\section{Chemins et lacets}\label{sec:chemins}

\begin{definition}[Chemin]\label{def:chemin}
\index{chemin}
Un \textbf{chemin} dans un espace topologique $X$ est une application continue
$\gamma : [0,1] \to X$. Le point $\gamma(0)$ est l'\textbf{origine} et
$\gamma(1)$ est l'\textbf{extrémité} du chemin. On dit que $\gamma$ est un
chemin \textbf{de $\gamma(0)$ à $\gamma(1)$}.
\end{definition}

\begin{definition}[Lacet]\label{def:lacet}
\index{lacet}
Un \textbf{lacet} basé en $x_0 \in X$ est un chemin $\gamma : [0,1] \to X$ tel
que $\gamma(0) = \gamma(1) = x_0$. L'ensemble des lacets basés en $x_0$ est
noté $\Omega(X, x_0)$\index{$\Omega(X,x_0)$}.
\end{definition}

\begin{definition}[Chemin inverse et chemin constant]\label{def:chemin-inverse}
\index{chemin!inverse}\index{chemin!constant}
Soit $\gamma : [0,1] \to X$ un chemin.
\begin{enumerate}[label=(\roman*)]
\item Le \textbf{chemin inverse} de $\gamma$ est $\overline{\gamma} : [0,1] \to X$
défini par $\overline{\gamma}(t) = \gamma(1-t)$.
\item Le \textbf{chemin constant} en $x_0$ est $c_{x_0} : [0,1] \to X$ défini
par $c_{x_0}(t) = x_0$ pour tout $t$.
\end{enumerate}
\end{definition}

\begin{definition}[Concaténation de chemins]\label{def:concatenation}
\index{concaténation de chemins}
Soient $\gamma, \delta : [0,1] \to X$ deux chemins tels que
$\gamma(1) = \delta(0)$. La \textbf{concaténation} (ou \textbf{produit}) de
$\gamma$ et $\delta$ est le chemin $\gamma \cdot \delta : [0,1] \to X$ défini par
\[
(\gamma \cdot \delta)(t) =
\begin{cases}
\gamma(2t) & \text{si } 0 \leqslant t \leqslant \tfrac{1}{2}, \\
\delta(2t - 1) & \text{si } \tfrac{1}{2} \leqslant t \leqslant 1.
\end{cases}
\]
Ce chemin est continu par le lemme de recollement.
\end{definition}

\begin{remarque}\label{rem:concatenation-convention}
Attention à la convention : $\gamma \cdot \delta$ signifie \og parcourir d'abord
$\gamma$, puis $\delta$ \fg{}. Certains auteurs utilisent la convention opposée.
Notre convention est celle de Hatcher et de la plupart des ouvrages modernes.
\end{remarque}

\begin{lemme}[Lemme de recollement]\label{lem:recollement}
\index{lemme de recollement}
Soient $X = A \cup B$ où $A$ et $B$ sont des sous-ensembles fermés de $X$. Si
$f : A \to Y$ et $g : B \to Y$ sont des applications continues telles que
$f|_{A \cap B} = g|_{A \cap B}$, alors l'application $h : X \to Y$ définie par
$h(x) = f(x)$ si $x \in A$ et $h(x) = g(x)$ si $x \in B$ est continue.
\end{lemme}

\begin{proof}
Soit $F \subseteq Y$ un fermé. Alors
$h^{-1}(F) = f^{-1}(F) \cup g^{-1}(F)$. Or $f^{-1}(F)$ est fermé dans $A$, et
comme $A$ est fermé dans $X$, $f^{-1}(F)$ est fermé dans $X$. De même pour
$g^{-1}(F)$. La réunion de deux fermés étant fermée, $h^{-1}(F)$ est fermé.
\end{proof}

Ce lemme justifie rigoureusement la continuité de la concaténation de chemins
dans la 
ef{def:concatenation}, puisque $[0, \tfrac{1}{2}]$ et
$[\tfrac{1}{2}, 1]$ sont fermés dans $[0,1]$ et les deux morceaux coïncident en
$t = \tfrac{1}{2}$.

% ============================================================
\section{Homotopie de chemins}\label{sec:homotopie-chemins}

\begin{definition}[Homotopie de chemins à extrémités fixées]\label{def:homotopie-chemins}
\index{homotopie de chemins}
Soient $\gamma, \delta : [0,1] \to X$ deux chemins ayant mêmes extrémités :
$\gamma(0) = \delta(0) = x_0$ et $\gamma(1) = \delta(1) = x_1$. On dit que
$\gamma$ et $\delta$ sont \textbf{homotopes à extrémités fixées}, et on note
$\gamma \simeq \delta \rel \partial$, s'il existe une application continue
$H : [0,1] \times [0,1] \to X$ vérifiant :
\begin{enumerate}[label=(\roman*)]
\item $H(s, 0) = \gamma(s)$ et $H(s, 1) = \delta(s)$ pour tout $s \in [0,1]$,
\item $H(0, t) = x_0$ et $H(1, t) = x_1$ pour tout $t \in [0,1]$.
\end{enumerate}
On note $[\gamma]$ la classe d'homotopie de $\gamma$ à extrémités fixées.
\end{definition}

\begin{proposition}\label{prop:homotopie-chemins-equiv}
La relation $\simeq_{\partial}$ (homotopie de chemins à extrémités fixées) est
une relation d'équivalence.
\end{proposition}

\begin{proof}
La vérification est analogue à celle de la 
ef{prop:homotopie-equivalence} :
la réflexivité utilise l'homotopie constante $H(s,t) = \gamma(s)$, la symétrie
utilise $\overline{H}(s,t) = H(s,1-t)$, et la transitivité utilise la
concaténation des homotopies en reparamétrant le paramètre $t$. Dans chaque cas,
les conditions aux bords $H(0,t) = x_0$ et $H(1,t) = x_1$ sont préservées.
\end{proof}

% ============================================================
\section{Structure de groupe sur les classes de lacets}\label{sec:structure-groupe}

\begin{lemme}[Associativité à homotopie près]\label{lem:associativite}
\index{associativité des chemins}
Soient $\gamma$, $\delta$, $\epsilon$ trois chemins composables dans $X$. Alors
\[
(\gamma \cdot \delta) \cdot \epsilon \simeq \gamma \cdot (\delta \cdot \epsilon) \rel \partial.
\]
\end{lemme}

\begin{proof}
Posons $\alpha = (\gamma \cdot \delta) \cdot \epsilon$ et
$\beta = \gamma \cdot (\delta \cdot \epsilon)$. On calcule explicitement :
\[
\alpha(s) = \begin{cases}
\gamma(4s) & \text{si } 0 \leqslant s \leqslant \tfrac{1}{4}, \\
\delta(4s - 1) & \text{si } \tfrac{1}{4} \leqslant s \leqslant \tfrac{1}{2}, \\
\epsilon(2s - 1) & \text{si } \tfrac{1}{2} \leqslant s \leqslant 1,
\end{cases}
\qquad
\beta(s) = \begin{cases}
\gamma(2s) & \text{si } 0 \leqslant s \leqslant \tfrac{1}{2}, \\
\delta(4s - 2) & \text{si } \tfrac{1}{2} \leqslant s \leqslant \tfrac{3}{4}, \\
\epsilon(4s - 3) & \text{si } \tfrac{3}{4} \leqslant s \leqslant 1.
\end{cases}
\]
L'homotopie est donnée par une reparamétrisation linéaire. Définissons
$H : [0,1]^2 \to X$ par
\[
H(s,t) = \begin{cases}
\gamma\!\left(\dfrac{4s}{t+1}\right) & \text{si } 0 \leqslant s \leqslant \dfrac{t+1}{4}, \\[6pt]
\delta(4s - t - 1) & \text{si } \dfrac{t+1}{4} \leqslant s \leqslant \dfrac{t+2}{4}, \\[6pt]
\epsilon\!\left(\dfrac{4s - t - 2}{2 - t}\right) & \text{si } \dfrac{t+2}{4} \leqslant s \leqslant 1.
\end{cases}
\]
On vérifie que $H$ est continue (par le lemme de recollement, les trois
morceaux coïncidant sur les segments de jonction), que $H(s,0) = \alpha(s)$,
$H(s,1) = \beta(s)$, et que les extrémités sont fixées : $H(0,t) = \gamma(0)$
et $H(1,t) = \epsilon(1)$.
\end{proof}

\begin{lemme}[Élément neutre à homotopie près]\label{lem:neutre}
\index{chemin constant}
Pour tout chemin $\gamma : [0,1] \to X$ d'origine $x_0$ et d'extrémité $x_1$,
\[
c_{x_0} \cdot \gamma \simeq \gamma \simeq \gamma \cdot c_{x_1} \rel \partial.
\]
\end{lemme}

\begin{proof}
L'homotopie pour $c_{x_0} \cdot \gamma \simeq \gamma$ est
\[
H(s,t) = \begin{cases}
x_0 & \text{si } 0 \leqslant s \leqslant \dfrac{1-t}{2}, \\[4pt]
\gamma\!\left(\dfrac{2s - 1 + t}{1 + t}\right) & \text{si } \dfrac{1-t}{2} \leqslant s \leqslant 1.
\end{cases}
\]
La construction pour $\gamma \cdot c_{x_1} \simeq \gamma$ est analogue.
\end{proof}

\begin{lemme}[Inverse à homotopie près]\label{lem:inverse}
\index{chemin inverse}
Pour tout chemin $\gamma : [0,1] \to X$ d'origine $x_0$ et d'extrémité $x_1$,
\[
\gamma \cdot \overline{\gamma} \simeq c_{x_0} \rel \partial \qquad \text{et} \qquad \overline{\gamma} \cdot \gamma \simeq c_{x_1} \rel \partial.
\]
\end{lemme}

\begin{proof}
L'homotopie pour $\gamma \cdot \overline{\gamma} \simeq c_{x_0}$ est définie par
\[
H(s,t) = \begin{cases}
\gamma(2s) & \text{si } 0 \leqslant s \leqslant \dfrac{t}{2}, \\[4pt]
\gamma(t) & \text{si } \dfrac{t}{2} \leqslant s \leqslant 1 - \dfrac{t}{2}, \\[4pt]
\gamma(2 - 2s) & \text{si } 1 - \dfrac{t}{2} \leqslant s \leqslant 1.
\end{cases}
\]
Au temps $t = 0$, $H(s,0) = \gamma(0) = x_0$ pour tout $s$ (chemin constant).
Au temps $t = 1$, $H(s,1)$ parcourt $\gamma$ sur $[0, \tfrac{1}{2}]$ puis
$\overline{\gamma}$ sur $[\tfrac{1}{2}, 1]$. La seconde assertion se traite de
manière symétrique.
\end{proof}

\begin{theoreme}[Structure de groupe du groupe fondamental]\label{thm:pi1-groupe}
\index{groupe fondamental!structure de groupe}
Soit $(X, x_0)$ un espace topologique pointé. L'ensemble des classes d'homotopie
(à extrémités fixées) de lacets basés en $x_0$, muni de l'opération
\[
[\gamma] \cdot [\delta] := [\gamma \cdot \delta],
\]
est un groupe, appelé \textbf{groupe fondamental} de $X$ en $x_0$ et noté
$\pi_1(X, x_0)$\index{$\pi_1(X,x_0)$}.
\end{theoreme}

\begin{proof}
Nous devons vérifier les quatre axiomes de groupe.

\emph{Bonne définition.} Si $\gamma \simeq \gamma'$ et $\delta \simeq \delta'$
(à extrémités fixées), alors $\gamma \cdot \delta \simeq \gamma' \cdot \delta'$
à extrémités fixées. En effet, si $H$ est une homotopie de $\gamma$ à $\gamma'$
et $K$ une homotopie de $\delta$ à $\delta'$, alors
\[
L(s,t) = \begin{cases}
H(2s, t) & \text{si } 0 \leqslant s \leqslant \tfrac{1}{2}, \\
K(2s-1, t) & \text{si } \tfrac{1}{2} \leqslant s \leqslant 1,
\end{cases}
\]
est une homotopie de $\gamma \cdot \delta$ à $\gamma' \cdot \delta'$ (les deux
morceaux se recollent en $s = \tfrac{1}{2}$ car $H(1,t) = x_0 = K(0,t)$ pour
les lacets basés en $x_0$).

\emph{Associativité.} C'est le 
ef{lem:associativite}.

\emph{Élément neutre.} La classe $[c_{x_0}]$ du lacet constant est l'élément
neutre d'après le 
ef{lem:neutre}.

\emph{Inverse.} Pour tout lacet $\gamma$ basé en $x_0$, le lacet inverse
$\overline{\gamma}$ vérifie $[\gamma] \cdot [\overline{\gamma}] = [c_{x_0}]$ et
$[\overline{\gamma}] \cdot [\gamma] = [c_{x_0}]$ d'après le 
ef{lem:inverse}.
\end{proof}

\begin{remarque}\label{rem:pi1-non-abelien}
Le groupe $\pi_1(X, x_0)$ n'est en général \emph{pas} abélien. Nous verrons que
$\pi_1(S^1 \vee S^1)$ est le groupe libre à deux générateurs $F_2$, qui est non
abélien. En revanche, $\pi_1(T^2) \cong \Z^2$ est abélien.
\end{remarque}

\begin{remarque}[Interprétation géométrique des axiomes]\label{rem:interpretation-geometrique}
La preuve de la structure de groupe sur $\pi_1(X, x_0)$ repose sur des
reparamétrisations du paramètre $s \in [0,1]$. Il est instructif de visualiser
ces homotopies comme des déformations dans le carré $[0,1]^2$ :

\begin{center}
\begin{tikzpicture}[scale=2.5]
  % Associativité
  \draw[thick] (0,0) rectangle (1,1);
  \node[below] at (0,0) {$0$};
  \node[below] at (0.25,0) {$\frac{1}{4}$};
  \node[below] at (0.5,0) {$\frac{1}{2}$};
  \node[below] at (1,0) {$1$};
  \node[left] at (0,0) {$0$};
  \node[left] at (0,1) {$1$};
  % Lignes de séparation en bas
  \draw[dashed,blue] (0.25,0) -- (0.5,1);
  \draw[dashed,red] (0.5,0) -- (0.75,1);
  % Labels
  \node at (0.1,0.5) {\small $\gamma$};
  \node at (0.4,0.5) {\small $\delta$};
  \node at (0.8,0.5) {\small $\epsilon$};
  % Titre
  \node[above] at (0.5,1.05) {\small Associativité : reparamétrisation};
  % Flèches de bord
  \node[below] at (0.5,-0.15) {\small $(\gamma\cdot\delta)\cdot\epsilon \xrightarrow{\simeq} \gamma\cdot(\delta\cdot\epsilon)$};
\end{tikzpicture}
\end{center}

Le bord inférieur du carré correspond au découpage $(\gamma \cdot \delta) \cdot
\epsilon$ et le bord supérieur au découpage $\gamma \cdot (\delta \cdot \epsilon)$.
L'homotopie consiste à \og glisser \fg{} continûment les points de coupure.
\end{remarque}

% ============================================================
\section{Fonctorialité du groupe fondamental}\label{sec:fonctorialite}

\begin{definition}[Morphisme induit]\label{def:morphisme-induit-pi1}
\index{morphisme induit!groupe fondamental}
Soit $f : (X, x_0) \to (Y, y_0)$ une application continue pointée (i.e.\
$f(x_0) = y_0$). Le \textbf{morphisme induit} est l'application
\[
f_* : \pi_1(X, x_0) \longrightarrow \pi_1(Y, y_0), \qquad [\gamma] \longmapsto [f \circ \gamma].
\]
\end{definition}

\begin{proposition}\label{prop:f-star-bien-defini}
L'application $f_*$ est bien définie et est un homomorphisme de groupes.
\end{proposition}

\begin{proof}
\emph{Bonne définition.} Si $\gamma \simeq \delta \rel \partial$ via une
homotopie $H$, alors $f \circ H$ est une homotopie de $f \circ \gamma$ à
$f \circ \delta$ à extrémités fixées.

\emph{Homomorphisme.} Soient $[\gamma], [\delta] \in \pi_1(X, x_0)$. On a
\[
f_*([\gamma] \cdot [\delta]) = f_*([\gamma \cdot \delta]) = [f \circ (\gamma \cdot \delta)].
\]
Or $(f \circ (\gamma \cdot \delta))(s) = f((\gamma \cdot \delta)(s))$, et l'on
vérifie immédiatement que $f \circ (\gamma \cdot \delta) = (f \circ \gamma) \cdot
(f \circ \delta)$, ce qui donne
$f_*([\gamma] \cdot [\delta]) = [f \circ \gamma] \cdot [f \circ \delta] =
f_*([\gamma]) \cdot f_*([\delta])$.
\end{proof}

\begin{theoreme}[Fonctorialité de $\pi_1$]\label{thm:pi1-foncteur}
\index{foncteur!$\pi_1$}
Le groupe fondamental $\pi_1$ définit un foncteur
\[
\pi_1 : \mathbf{Top}_* \longrightarrow \mathbf{Grp}
\]
de la catégorie des espaces topologiques pointés vers la catégorie des groupes,
au sens suivant :
\begin{enumerate}[label=(\roman*)]
\item $(\Id_X)_* = \Id_{\pi_1(X, x_0)}$ pour tout $(X, x_0)$.
\item $(g \circ f)_* = g_* \circ f_*$ pour toutes applications continues
pointées $f : (X, x_0) \to (Y, y_0)$ et $g : (Y, y_0) \to (Z, z_0)$.
\end{enumerate}
\end{theoreme}

\begin{proof}
(i) Pour tout lacet $\gamma$ basé en $x_0$, $(\Id_X)_*([\gamma]) =
[\Id_X \circ \gamma] = [\gamma]$.

(ii) Pour tout $[\gamma] \in \pi_1(X, x_0)$,
\[
(g \circ f)_*([\gamma]) = [(g \circ f) \circ \gamma] = [g \circ (f \circ \gamma)]
= g_*([f \circ \gamma]) = g_*(f_*([\gamma])) = (g_* \circ f_*)([\gamma]).
\qedhere
\]
\end{proof}

Le diagramme commutatif suivant résume la fonctorialité :
\[
\begin{tikzcd}[column sep=large, row sep=large]
(X, x_0) \arrow[r, "f"] \arrow[rr, bend left=20, "g \circ f"] & (Y, y_0) \arrow[r, "g"] & (Z, z_0)
\end{tikzcd}
\quad \longmapsto \quad
\begin{tikzcd}[column sep=large, row sep=large]
\pi_1(X, x_0) \arrow[r, "f_*"] \arrow[rr, bend left=20, "(g \circ f)_*"] & \pi_1(Y, y_0) \arrow[r, "g_*"] & \pi_1(Z, z_0)
\end{tikzcd}
\]

\begin{corollaire}\label{cor:homeo-iso-pi1}
Si $f : (X, x_0) \to (Y, y_0)$ est un homéomorphisme, alors
$f_* : \pi_1(X, x_0) \to \pi_1(Y, y_0)$ est un isomorphisme de groupes.
\end{corollaire}

\begin{proof}
Si $f$ est un homéomorphisme, $f^{-1}$ est continue et
$(f^{-1})_* \circ f_* = (f^{-1} \circ f)_* = (\Id_X)_* = \Id$. De même
$f_* \circ (f^{-1})_* = \Id$.
\end{proof}

% ============================================================
\section{Changement de point base}\label{sec:changement-base}

Que se passe-t-il lorsqu'on change de point base ? Si $X$ est connexe par arcs,
les groupes fondamentaux en différents points base sont isomorphes, mais
l'isomorphisme n'est pas canonique --- il dépend du choix d'un chemin.

\begin{theoreme}[Isomorphisme de changement de point base]\label{thm:changement-base}
\index{changement de point base}
Soit $X$ un espace connexe par arcs et soient $x_0, x_1 \in X$. Pour tout
chemin $\alpha : [0,1] \to X$ de $x_0$ à $x_1$, l'application
\[
\beta_\alpha : \pi_1(X, x_1) \longrightarrow \pi_1(X, x_0), \qquad [\gamma] \longmapsto [\alpha \cdot \gamma \cdot \overline{\alpha}],
\]
est un isomorphisme de groupes.
\end{theoreme}

\begin{proof}
\emph{Homomorphisme.} Soient $[\gamma], [\delta] \in \pi_1(X, x_1)$. On a
\begin{align*}
\beta_\alpha([\gamma] \cdot [\delta])
&= [\alpha \cdot (\gamma \cdot \delta) \cdot \overline{\alpha}] \\
&= [\alpha \cdot \gamma \cdot c_{x_1} \cdot \delta \cdot \overline{\alpha}] \\
&= [\alpha \cdot \gamma \cdot (\overline{\alpha} \cdot \alpha) \cdot \delta \cdot \overline{\alpha}] \\
&= [(\alpha \cdot \gamma \cdot \overline{\alpha}) \cdot (\alpha \cdot \delta \cdot \overline{\alpha})] \\
&= \beta_\alpha([\gamma]) \cdot \beta_\alpha([\delta]).
\end{align*}

\emph{Isomorphisme.} L'application $\beta_{\overline{\alpha}} : \pi_1(X, x_0) \to \pi_1(X, x_1)$ est un inverse : pour tout lacet $\gamma$ basé en $x_1$,
\[
\beta_{\overline{\alpha}}(\beta_\alpha([\gamma])) = [\overline{\alpha} \cdot (\alpha \cdot \gamma \cdot \overline{\alpha}) \cdot \alpha] = [(\overline{\alpha} \cdot \alpha) \cdot \gamma \cdot (\overline{\alpha} \cdot \alpha)] = [c_{x_1} \cdot \gamma \cdot c_{x_1}] = [\gamma].
\]
Le calcul symétrique montre que $\beta_\alpha \circ \beta_{\overline{\alpha}} = \Id$.
\end{proof}

\begin{remarque}\label{rem:changement-base-non-canonique}
L'isomorphisme $\beta_\alpha$ dépend du choix du chemin $\alpha$. Si
$\alpha' : x_0 \to x_1$ est un autre chemin, alors
$\beta_{\alpha'} = \iota_{[\alpha \cdot \overline{\alpha'}]} \circ \beta_\alpha$
où $\iota_g$ désigne la conjugaison par $g$ dans $\pi_1(X, x_0)$. En
particulier, $\beta_\alpha = \beta_{\alpha'}$ si et seulement si
$[\alpha \cdot \overline{\alpha'}]$ est dans le centre de $\pi_1(X, x_0)$. Si
$\pi_1(X, x_0)$ est abélien, l'isomorphisme est donc indépendant du chemin, et
l'on peut parler de \og $\pi_1(X)$ \fg{} sans ambiguïté.
\end{remarque}

Le diagramme suivant illustre le changement de point base :
\[
\begin{tikzcd}[row sep=large, column sep=large]
\pi_1(X, x_0) \arrow[rr, "\beta_{\overline{\alpha}}"] & & \pi_1(X, x_1) \arrow[ll, bend left=15, "\beta_\alpha"]
\end{tikzcd}
\]

% ============================================================
\section{Invariance homotopique de $\pi_1$}\label{sec:invariance-homotopique}

\begin{theoreme}[Invariance homotopique]\label{thm:pi1-invariance-homotopique}
\index{invariance homotopique!$\pi_1$}
Soit $f : (X, x_0) \to (Y, y_0)$ une application continue telle que $f$ est une
équivalence d'homotopie. Alors
$f_* : \pi_1(X, x_0) \to \pi_1(Y, y_0)$ est un isomorphisme de groupes (où
$y_0 = f(x_0)$ et le point base au but est ajusté par un isomorphisme de
changement de point base si nécessaire).
\end{theoreme}

\begin{proof}
Soit $g : Y \to X$ une inverse homotopique de $f$, de sorte que
$g \circ f \simeq \Id_X$ et $f \circ g \simeq \Id_Y$.

\emph{Étape 1.} Posons $x_0' = g(f(x_0)) = g(y_0)$. L'homotopie
$H : X \times [0,1] \to X$ avec $H(\cdot, 0) = g \circ f$ et
$H(\cdot, 1) = \Id_X$ fournit un chemin $\alpha : t \mapsto H(x_0, t)$ de
$x_0' = g(y_0)$ à $x_0$. Pour tout lacet $\gamma$ basé en $x_0$, le lacet
$g \circ f \circ \gamma$ est basé en $x_0'$, et la relation
$g \circ f \simeq \Id_X$ implique
\[
[\alpha \cdot (g \circ f \circ \gamma) \cdot \overline{\alpha}] = [\gamma]
\quad \text{dans } \pi_1(X, x_0).
\]
Autrement dit, $\beta_\alpha \circ (g \circ f)_* = \Id_{\pi_1(X, x_0)}$, soit
encore $\beta_\alpha \circ g_* \circ f_* = \Id$.

\emph{Étape 2.} Un raisonnement identique appliqué à l'homotopie
$f \circ g \simeq \Id_Y$ fournit un chemin $\alpha'$ et l'égalité
$\beta_{\alpha'} \circ f_* \circ g_* = \Id_{\pi_1(Y, y_0)}$.

\emph{Étape 3.} De l'étape 1, $f_*$ est injective (car elle admet un inverse à
gauche $\beta_\alpha \circ g_*$). De l'étape 2, $f_*$ est surjective (car
$g_*$ admet un inverse à gauche $\beta_{\alpha'} \circ f_*$, donc $g_*$ est
injective, et l'étape 1 donne $f_*$ surjective). Ainsi $f_*$ est un
isomorphisme.
\end{proof}

\begin{remarque}\label{rem:invariance-homotopique-pratique}
Le 
ef{thm:pi1-invariance-homotopique} est un résultat extrêmement puissant.
Il signifie que le groupe fondamental est un invariant du \emph{type d'homotopie},
et pas seulement du type topologique. Pour calculer $\pi_1(X)$, on peut donc
remplacer $X$ par n'importe quel espace ayant le même type d'homotopie. Par
exemple :
\begin{itemize}
\item $\pi_1(\R^n \setminus \{0\}) \cong \pi_1(S^{n-1})$, car
$S^{n-1}$ est un rétracte par déformation de $\R^n \setminus \{0\}$.
\item $\pi_1(\text{cylindre plein}) = 0$, car le cylindre plein est contractile.
\item $\pi_1(\text{bande de Möbius}) \cong \Z$, car la bande de Möbius se
rétracte par déformation sur son cercle central.
\end{itemize}
\end{remarque}

\begin{corollaire}\label{cor:contractile-pi1-trivial}
Si $X$ est contractile, alors $\pi_1(X, x_0) = 0$ pour tout $x_0 \in X$.
\end{corollaire}

\begin{proof}
Un espace contractile a le type d'homotopie d'un point $\{*\}$. Or
$\pi_1(\{*\}) = 0$. Par le 
ef{thm:pi1-invariance-homotopique},
$\pi_1(X, x_0) \cong 0$.
\end{proof}

% ============================================================
\section{Premiers calculs}\label{sec:premiers-calculs}

\subsection{Espaces simplement connexes}\label{subsec:simplement-connexe}

\begin{definition}[Espace simplement connexe]\label{def:simplement-connexe}
\index{simplement connexe}
Un espace topologique $X$ est dit \textbf{simplement connexe} s'il est connexe
par arcs et si $\pi_1(X, x_0) = 0$ pour un (donc pour tout) point $x_0 \in X$.
\end{definition}

\begin{proposition}\label{prop:pi1-Rn}
Pour tout $n \geqslant 1$, $\pi_1(\R^n) = 0$. Plus généralement, tout espace
convexe de $\R^n$ est simplement connexe.
\end{proposition}

\begin{proof}
$\R^n$ est contractile (
ef{ex:Rn-contractile}), donc
$\pi_1(\R^n) = 0$ par le 
ef{cor:contractile-pi1-trivial}.
\end{proof}

\begin{proposition}\label{prop:pi1-Sn}
Pour $n \geqslant 2$, la sphère $S^n$ est simplement connexe :
$\pi_1(S^n) = 0$.
\end{proposition}

\begin{proof}
Soit $\gamma : [0,1] \to S^n$ un lacet basé en le pôle nord $N$.
Pour $n \geqslant 2$, l'image de $\gamma$ ne peut remplir $S^n$ tout entier
(un argument de dimension ou de mesure le montre). Plus rigoureusement, on
utilise le résultat suivant : si $U_1, U_2$ sont deux ouverts simplement connexes
de $X$ tels que $X = U_1 \cup U_2$ et $U_1 \cap U_2$ est connexe par arcs,
alors $X$ est simplement connexe (conséquence du théorème de van
Kampen\index{van Kampen}). On prend $U_1 = S^n \setminus \{S\}$ et
$U_2 = S^n \setminus \{N\}$ (où $S$ est le pôle sud). Chacun est homéomorphe à
$\R^n$ (par projection stéréographique), donc simplement connexe. Leur
intersection $S^n \setminus \{N, S\}$ est homéomorphe à
$\R^n \setminus \{0\} \simeq S^{n-1}$, qui est connexe par arcs pour
$n \geqslant 2$. Par van Kampen, $\pi_1(S^n) = 0$.
\end{proof}

\subsection{Le groupe fondamental du cercle}\label{subsec:pi1-S1}

Le résultat suivant est le premier calcul non trivial de groupe fondamental et
l'un des théorèmes fondateurs de la topologie algébrique.

\begin{theoreme}[Groupe fondamental du cercle]\label{thm:pi1-S1}
\index{cercle!groupe fondamental}\index{$\pi_1(S^1)$}
Le groupe fondamental du cercle est infini cyclique :
\[
\pi_1(S^1, 1) \cong \Z.
\]
L'isomorphisme envoie la classe $[\omega_n]$ du lacet
$\omega_n(t) = e^{2\pi i n t}$ (qui parcourt $n$ fois le cercle) sur l'entier $n$.
\end{theoreme}

La démonstration repose sur la théorie des revêtements, dont nous utilisons ici
les éléments essentiels.

\begin{definition}[Revêtement]\label{def:revetement-apercu}
\index{revêtement}
Un \textbf{revêtement} d'un espace $X$ est une application continue
$p : \widetilde{X} \to X$ telle que tout point $x \in X$ possède un voisinage
ouvert $U$ dont l'image réciproque $p^{-1}(U)$ est une réunion disjointe
d'ouverts, chacun envoyé homéomorphiquement sur $U$ par $p$.
\end{definition}

\begin{exemple}\label{ex:revetement-R-S1}
L'application $p : \R \to S^1$ définie par $p(t) = e^{2\pi i t}$ est un
revêtement. La fibre $p^{-1}(z)$ au-dessus de chaque point $z \in S^1$ est
isomorphe à $\Z$.
\end{exemple}

\begin{lemme}[Relèvement des chemins]\label{lem:relevement-chemins}
\index{relèvement des chemins}
Soient $p : \widetilde{X} \to X$ un revêtement et $\gamma : [0,1] \to X$ un
chemin d'origine $x_0$. Pour tout $\widetilde{x}_0 \in p^{-1}(x_0)$, il existe
un unique chemin $\widetilde{\gamma} : [0,1] \to \widetilde{X}$ tel que
$\widetilde{\gamma}(0) = \widetilde{x}_0$ et
$p \circ \widetilde{\gamma} = \gamma$.
\end{lemme}

\begin{proof}[Idée de la preuve]
Par compacité de $[0,1]$, on peut trouver une subdivision
$0 = t_0 < t_1 < \cdots < t_k = 1$ telle que chaque $\gamma([t_i, t_{i+1}])$
est contenu dans un ouvert trivialement recouvert. On relève alors le chemin
morceau par morceau, en utilisant l'unicité sur chaque intervalle pour garantir
la cohérence du recollement.
\end{proof}

\begin{lemme}[Relèvement des homotopies]\label{lem:relevement-homotopies}
\index{relèvement des homotopies}
Soient $p : \widetilde{X} \to X$ un revêtement et $H : [0,1]^2 \to X$ une
homotopie à extrémités fixées entre deux chemins. Si
$\widetilde{x}_0 \in p^{-1}(H(0,0))$, alors il existe un unique relèvement
$\widetilde{H} : [0,1]^2 \to \widetilde{X}$ avec
$\widetilde{H}(0,0) = \widetilde{x}_0$ et
$p \circ \widetilde{H} = H$. De plus, $\widetilde{H}$ est une homotopie à
extrémités fixées.
\end{lemme}

\begin{proof}[Preuve du 
ef{thm:pi1-S1}]
Considérons le revêtement $p : \R \to S^1$, $p(t) = e^{2\pi i t}$, et posons
$\widetilde{x}_0 = 0 \in p^{-1}(1)$.

\emph{Construction de l'homomorphisme.} Soit $[\gamma] \in \pi_1(S^1, 1)$. Par le

ef{lem:relevement-chemins}, il existe un unique relèvement
$\widetilde{\gamma} : [0,1] \to \R$ avec $\widetilde{\gamma}(0) = 0$ et
$p \circ \widetilde{\gamma} = \gamma$. Comme $\gamma$ est un lacet,
$\widetilde{\gamma}(1) \in p^{-1}(1) = \Z$. On définit
\[
\deg : \pi_1(S^1, 1) \longrightarrow \Z, \qquad [\gamma] \longmapsto \widetilde{\gamma}(1).
\]

\emph{Bonne définition.} Si $\gamma \simeq \delta \rel \partial$ via $H$, le

ef{lem:relevement-homotopies} fournit un relèvement $\widetilde{H}$ qui est
une homotopie à extrémités fixées entre $\widetilde{\gamma}$ et
$\widetilde{\delta}$. En particulier,
$\widetilde{\gamma}(1) = \widetilde{H}(1,0) = \widetilde{H}(1,1) = \widetilde{\delta}(1)$.

\emph{Homomorphisme.} Si $\widetilde{\gamma}$ relève $\gamma$ avec
$\widetilde{\gamma}(0) = 0$ et $\widetilde{\gamma}(1) = m$, et
$\widetilde{\delta}$ relève $\delta$ avec $\widetilde{\delta}(0) = 0$ et
$\widetilde{\delta}(1) = n$, alors le chemin
\[
\widetilde{\gamma \cdot \delta}(t) = \begin{cases}
\widetilde{\gamma}(2t) & \text{si } 0 \leqslant t \leqslant \tfrac{1}{2}, \\
m + \widetilde{\delta}(2t-1) & \text{si } \tfrac{1}{2} \leqslant t \leqslant 1,
\end{cases}
\]
relève $\gamma \cdot \delta$ avec $\widetilde{\gamma \cdot \delta}(0) = 0$ et
$\widetilde{\gamma \cdot \delta}(1) = m + n$. Donc $\deg([\gamma] \cdot [\delta]) = m + n = \deg([\gamma]) + \deg([\delta])$.

\emph{Surjectivité.} Pour $n \in \Z$, le lacet $\omega_n(t) = e^{2\pi i n t}$
se relève en $\widetilde{\omega}_n(t) = nt$, et $\deg([\omega_n]) = n$.

\emph{Injectivité.} Supposons $\deg([\gamma]) = 0$, i.e.\
$\widetilde{\gamma}(1) = 0 = \widetilde{\gamma}(0)$. Le relèvement
$\widetilde{\gamma}$ est donc un lacet dans $\R$ basé en $0$. Comme $\R$ est
simplement connexe, $\widetilde{\gamma} \simeq c_0 \rel \partial$ dans $\R$.
En composant avec $p$, on obtient $\gamma = p \circ \widetilde{\gamma} \simeq
p \circ c_0 = c_1 \rel \partial$ dans $S^1$. Donc $[\gamma] = [c_1]$ est
l'élément neutre.
\end{proof}

\begin{remarque}[Signification géométrique]\label{rem:pi1-S1-geometrique}
L'isomorphisme $\pi_1(S^1) \cong \Z$ a une signification géométrique limpide :
un lacet sur le cercle est entièrement déterminé (à homotopie près) par son
\textbf{nombre d'enroulement}\index{nombre d'enroulement}, c'est-à-dire le nombre
de fois qu'il fait le tour du cercle (compté avec signe). Le générateur
$[\omega_1]$ correspond à un tour dans le sens trigonométrique, et
$[\omega_{-1}]$ à un tour dans le sens horaire.

\begin{center}
\begin{tikzpicture}[scale=0.9]
  % n = 0
  \begin{scope}[xshift=-4cm]
    \draw[gray!50] (0,0) circle (1);
    \fill (1,0) circle (2pt) node[right] {$1$};
    \node[below] at (0,-1.5) {$n = 0$};
    \node[below] at (0,-1.9) {(lacet constant)};
  \end{scope}
  % n = 1
  \begin{scope}
    \draw[gray!50] (0,0) circle (1);
    \fill (1,0) circle (2pt);
    \draw[thick,blue!70!black,
          decoration={markings, mark=at position 0.3 with {\arrow{Stealth}},
                                mark=at position 0.7 with {\arrow{Stealth}}},
          postaction={decorate}]
      (1,0) arc[start angle=0,end angle=360,radius=1];
    \node[below] at (0,-1.5) {$n = 1$};
  \end{scope}
  % n = 2
  \begin{scope}[xshift=4cm]
    \draw[gray!50] (0,0) circle (1);
    \fill (1,0) circle (2pt);
    \draw[thick,red!70!black,
          decoration={markings, mark=at position 0.15 with {\arrow{Stealth}},
                                mark=at position 0.65 with {\arrow{Stealth}}},
          postaction={decorate}]
      (0.95,0) arc[start angle=0,end angle=720,radius=0.95];
    \node[below] at (0,-1.5) {$n = 2$};
    \node[below] at (0,-1.9) {(deux tours)};
  \end{scope}
\end{tikzpicture}
\captionof{figure}{Lacets sur le cercle classifiés par le nombre d'enroulement.}\label{fig:enroulement}
\end{center}
\end{remarque}

\begin{corollaire}\label{cor:S1-pas-simplement-connexe}
Le cercle $S^1$ n'est pas simplement connexe. En particulier, $S^1$ n'est pas
contractile.
\end{corollaire}

\begin{proof}
On a $\pi_1(S^1) \cong \Z \neq 0$, donc $S^1$ n'est pas simplement connexe.
Un espace contractile étant simplement connexe (
ef{cor:contractile-pi1-trivial}),
$S^1$ n'est pas contractile.
\end{proof}

% ============================================================
\section{Applications du groupe fondamental}\label{sec:applications-pi1}

\subsection{Théorème de Brouwer en dimension 2}\label{subsec:brouwer}

\begin{theoreme}[Pas de rétraction de $D^2$ sur $S^1$]\label{thm:pas-retraction}
\index{rétraction!$D^2$ sur $S^1$}
Il n'existe pas de rétraction continue $r : D^2 \to S^1$.
\end{theoreme}

\begin{proof}
Supposons par l'absurde qu'il existe une rétraction $r : D^2 \to S^1$, i.e.\ une
application continue telle que $r \circ \iota = \Id_{S^1}$, où
$\iota : S^1 \hookrightarrow D^2$ est l'inclusion. En appliquant le foncteur
$\pi_1$ (avec point base $1 \in S^1$), on obtient le diagramme commutatif
\[
\begin{tikzcd}
\pi_1(S^1, 1) \arrow[r, "\iota_*"] \arrow[rr, bend left=20, "\Id"] & \pi_1(D^2, 1) \arrow[r, "r_*"] & \pi_1(S^1, 1)
\end{tikzcd}
\]
avec $r_* \circ \iota_* = (\Id_{S^1})_* = \Id$. Or $\pi_1(D^2, 1) = 0$ (car
$D^2$ est contractile) et $\pi_1(S^1, 1) \cong \Z$. L'identité de $\Z$ se
factorise donc à travers le groupe trivial, ce qui est impossible.
\end{proof}

\begin{theoreme}[Théorème du point fixe de Brouwer en dimension 2]\label{thm:brouwer-dim2}
\index{théorème du point fixe de Brouwer}
Toute application continue $f : D^2 \to D^2$ possède un point fixe.
\end{theoreme}

\begin{proof}
Supposons par l'absurde que $f(x) \neq x$ pour tout $x \in D^2$. On construit
une rétraction $r : D^2 \to S^1$ de la manière suivante : pour chaque
$x \in D^2$, on considère la demi-droite partant de $f(x)$ et passant par $x$;
cette demi-droite coupe le cercle $S^1$ en un unique point $r(x)$.

Plus explicitement, $r(x)$ est l'unique point de $S^1$ de la forme
$x + t(x - f(x))$ avec $t \geqslant 0$ et $\|x + t(x - f(x))\| = 1$. Cette
application est bien définie et continue (car $f(x) \neq x$), et si $x \in S^1$,
la demi-droite sort immédiatement du disque, donc $r(x) = x$. Ainsi $r$ est une
rétraction de $D^2$ sur $S^1$, ce qui contredit le 
ef{thm:pas-retraction}.
\end{proof}

\begin{center}
\begin{tikzpicture}[scale=1.8]
  \draw[thick] (0,0) circle (1);
  \fill[blue!10] (0,0) circle (1);
  \draw[thick] (0,0) circle (1);
  % point x
  \coordinate (x) at (0.3, 0.2);
  \fill (x) circle (1pt) node[below right] {$x$};
  % point f(x)
  \coordinate (fx) at (-0.3, 0.5);
  \fill[red] (fx) circle (1pt) node[above left] {$f(x)$};
  % demi-droite
  \draw[dashed, -{Stealth}] (fx) -- ($(fx)!2.5!(x)$);
  % point r(x) sur le cercle
  \coordinate (rx) at (0.82, -0.57);
  \fill[green!60!black] (rx) circle (1.5pt) node[below right] {$r(x)$};
  % label
  \node at (0, -1.4) {Construction de la rétraction};
\end{tikzpicture}
\captionof{figure}{Démonstration par l'absurde du théorème de Brouwer.}\label{fig:brouwer}
\end{center}

\subsection{Théorème de Borsuk--Ulam en dimension 1}\label{subsec:borsuk-ulam}

\begin{theoreme}[Borsuk--Ulam en dimension 1]\label{thm:borsuk-ulam-1}
\index{théorème de Borsuk--Ulam}
Pour toute application continue $f : S^1 \to \R$, il existe un point $x \in S^1$
tel que $f(x) = f(-x)$.
\end{theoreme}

\begin{proof}
Posons $g(x) = f(x) - f(-x)$. Alors $g$ est continue et
$g(-x) = f(-x) - f(x) = -g(x)$. Si $g(x_0) = 0$ pour un $x_0$, on a terminé.
Sinon, $g$ ne s'annule pas, et le signe de $g$ est opposé en deux points
antipodaux. Par le théorème des valeurs intermédiaires appliqué à un arc reliant
$x_0$ à $-x_0$, $g$ doit s'annuler, contradiction.
\end{proof}

\begin{corollaire}[Ham sandwich en dimension 1]\label{cor:ham-sandwich-1}
Pour toute application continue $f : S^1 \to \R$, il existe un diamètre du
cercle dont les deux extrémités ont la même image par $f$. Autrement dit, il
est toujours possible de couper un \og sandwich unidimensionnel \fg{} en deux
parts égales.
\end{corollaire}

\begin{remarque}\label{rem:borsuk-ulam-generalisation}
Le théorème de Borsuk--Ulam se généralise en toute dimension : pour toute
application continue $f : S^n \to \R^n$, il existe $x \in S^n$ tel que
$f(x) = f(-x)$. La preuve en dimension supérieure fait appel à l'homologie ou à
la cohomologie.
\end{remarque}

\subsection{Application à l'algèbre : théorème fondamental de l'algèbre}\label{subsec:TFA}

\begin{theoreme}[Théorème fondamental de l'algèbre]\label{thm:TFA}
\index{théorème fondamental de l'algèbre}
Tout polynôme $P \in \C[z]$ de degré $n \geqslant 1$ possède au moins une
racine dans $\C$.
\end{theoreme}

\begin{proof}[Esquisse de preuve via $\pi_1$]
On peut supposer $P(z) = z^n + a_{n-1}z^{n-1} + \cdots + a_0$ unitaire.
Supposons par l'absurde que $P$ ne s'annule pas. Pour $r > 0$, considérons le
lacet $\gamma_r : [0,1] \to \C \setminus \{0\}$ défini par
$\gamma_r(t) = P(re^{2\pi it})$.

Pour $r$ suffisamment grand, le terme dominant $z^n$ l'emporte, et $\gamma_r$
est homotope au lacet $t \mapsto r^n e^{2\pi i n t}$, qui a pour classe
$n \in \pi_1(\C \setminus \{0\}, P(r)) \cong \Z$ (non triviale car $n \geqslant 1$).

Pour $r = 0$, $\gamma_0$ est le lacet constant en $a_0 \neq 0$, de classe
$0 \in \Z$.

Or la famille $r \mapsto \gamma_r$ fournit une homotopie continue (dans
$\C \setminus \{0\}$ par hypothèse), ce qui impliquerait $n = 0$ dans $\Z$,
contradiction.
\end{proof}

% ============================================================
\section{Le groupe fondamental du produit}\label{sec:pi1-produit}

\begin{theoreme}\label{thm:pi1-produit}
\index{groupe fondamental!produit}
Soient $(X, x_0)$ et $(Y, y_0)$ deux espaces topologiques pointés. Alors
\[
\pi_1(X \times Y, (x_0, y_0)) \cong \pi_1(X, x_0) \times \pi_1(Y, y_0).
\]
\end{theoreme}

\begin{proof}
Les projections $p_X : X \times Y \to X$ et $p_Y : X \times Y \to Y$ induisent
un homomorphisme
\[
\Phi = ((p_X)_*, (p_Y)_*) : \pi_1(X \times Y, (x_0, y_0)) \longrightarrow \pi_1(X, x_0) \times \pi_1(Y, y_0).
\]
Réciproquement, si $\gamma : [0,1] \to X$ est un lacet basé en $x_0$ et
$\delta : [0,1] \to Y$ est un lacet basé en $y_0$, le lacet
$(\gamma, \delta) : t \mapsto (\gamma(t), \delta(t))$ est un lacet basé en
$(x_0, y_0)$ dans $X \times Y$. L'application $\Psi : ([\gamma], [\delta]) \mapsto [(\gamma, \delta)]$ est un homomorphisme et on vérifie aisément que $\Phi$ et $\Psi$ sont inverses l'un de l'autre.
\end{proof}

\begin{corollaire}\label{cor:pi1-tore}
\index{tore!groupe fondamental}
Le groupe fondamental du tore $T^2 = S^1 \times S^1$ est
\[
\pi_1(T^2) \cong \Z \times \Z \cong \Z^2.
\]
Plus généralement, $\pi_1(T^n) \cong \Z^n$.
\end{corollaire}

Le diagramme commutatif suivant résume la situation pour le tore :
\[
\begin{tikzcd}[column sep=huge, row sep=large]
\pi_1(T^2, (1,1)) \arrow[r, "\sim"] \arrow[d, "(p_1)_*"'] & \Z \times \Z \arrow[d, "\mathrm{pr}_1"] \\
\pi_1(S^1, 1) \arrow[r, "\sim"] & \Z
\end{tikzcd}
\]

\begin{remarque}\label{rem:pi1-produit-infini}
Pour un produit infini, on a encore
$\pi_1\!\left(\prod_{i \in I} X_i\right) \cong \prod_{i \in I} \pi_1(X_i)$,
pourvu que les espaces soient pointés et le produit muni de la topologie produit.
Cependant, pour un bouquet infini $\bigvee_{i \in I} X_i$, le résultat est plus
délicat et fait intervenir le produit libre (éventuellement non dénombrable) des
groupes fondamentaux.
\end{remarque}

\begin{exemple}[Groupe fondamental de la bouteille de Klein]\label{ex:pi1-klein}
\index{bouteille de Klein}
La bouteille de Klein $K$ est obtenue en identifiant les bords d'un carré via
le mot $aba^{-1}b$ (au lieu de $aba^{-1}b^{-1}$ pour le tore). Par le théorème
de van Kampen, on obtient la présentation
\[
\pi_1(K) \cong \langle a, b \mid abab^{-1} = 1 \rangle.
\]
Ce groupe est non abélien et contient $\Z$ et $\Z/2\Z$ comme quotients. La
bouteille de Klein n'est donc homéomorphe ni au tore (dont le groupe fondamental
est $\Z^2$) ni à $\RP^2$ (dont le groupe fondamental est $\Z/2\Z$).
\end{exemple}

\begin{exemple}[Résumé des groupes fondamentaux classiques]\label{ex:resume-pi1}
Voici un tableau récapitulatif des groupes fondamentaux des espaces les plus
courants :

\begin{center}
\begin{tabular}{lll}
\toprule
\textbf{Espace $X$} & $\pi_1(X)$ & \textbf{Commentaire} \\
\midrule
$\R^n$ & $0$ & contractile \\
$D^n$ & $0$ & contractile \\
$S^1$ & $\Z$ & engendré par $[\omega_1]$ \\
$S^n$ ($n \geqslant 2$) & $0$ & simplement connexe \\
$T^n = (S^1)^n$ & $\Z^n$ & produit \\
$\RP^n$ ($n \geqslant 2$) & $\Z/2\Z$ & revêtement par $S^n$ \\
$S^1 \vee S^1$ & $\Z * \Z = F_2$ & bouquet, groupe libre \\
$\Sigma_g$ (genre $g$) & $\langle a_i, b_i \mid \prod [a_i, b_i]\rangle$ & surface orientable \\
\bottomrule
\end{tabular}
\end{center}
\end{exemple}

% ============================================================
\section{Espaces simplement connexes et relèvement}\label{sec:simplement-connexe-relevement}

\begin{proposition}[Critère de simple connexité]\label{prop:critere-simpl-connexe}
\index{simplement connexe!critère}
Soit $X$ un espace connexe par arcs. Les conditions suivantes sont équivalentes :
\begin{enumerate}[label=(\roman*)]
\item $X$ est simplement connexe.
\item Deux chemins quelconques de mêmes extrémités dans $X$ sont homotopes à
extrémités fixées.
\item Tout lacet dans $X$ est homotope à un point.
\item Pour tout revêtement $p : \widetilde{X} \to X$, toute application continue
$f : [0,1] \to X$ admet un relèvement (et celui-ci est unique une fois le point
de départ fixé).
\end{enumerate}
\end{proposition}

\begin{proof}
(i) $\Leftrightarrow$ (iii) est la définition (un lacet homotope à un point
est précisément un élément trivial de $\pi_1$).

(ii) $\Rightarrow$ (iii) : Un lacet $\gamma$ et le chemin constant $c_{x_0}$
sont deux chemins de mêmes extrémités $x_0$, donc $\gamma \simeq c_{x_0}$.

(iii) $\Rightarrow$ (ii) : Si $\gamma$ et $\delta$ vont de $x_0$ à $x_1$, le
lacet $\gamma \cdot \overline{\delta}$ est homotope au point, i.e.\
$\gamma \cdot \overline{\delta} \simeq c_{x_0}$. En composant à droite par
$\delta$, on obtient $\gamma \simeq \delta$.

(i) $\Rightarrow$ (iv) : C'est une conséquence du théorème de relèvement des
chemins (
ef{lem:relevement-chemins}).

(iv) $\Rightarrow$ (i) : On considère le revêtement universel et on montre que
la trivialité du relèvement implique $\pi_1 = 0$.
\end{proof}

% ============================================================
\section{Exercices}\label{sec:exos-ch2}

\begin{exercice}[$\star\star$]\label{exo:pi1-trivial-convexe}
Soit $C \subseteq \R^n$ un sous-ensemble convexe. Montrer directement (sans
utiliser la contractilité) que $\pi_1(C, x_0) = 0$ pour tout $x_0 \in C$.
\emph{Indication : utiliser l'homotopie rectiligne $H(s,t) = (1-t)\gamma(s) + t\delta(s)$.}
\end{exercice}

\begin{exercice}[$\star\star$]\label{exo:pi1-homeo-invariant}
Soient $X$ et $Y$ deux espaces homéomorphes et $f : X \to Y$ un homéomorphisme.
Montrer que $f_* : \pi_1(X, x_0) \to \pi_1(Y, f(x_0))$ est un isomorphisme.
En déduire que $\R$ et $S^1$ ne sont pas homéomorphes.
\end{exercice}

\begin{exercice}[$\star\star$]\label{exo:R2-R3-non-homeo}
Montrer que $\R^2$ et $\R^3$ ne sont pas homéomorphes en utilisant le groupe
fondamental. \emph{Indication : utiliser la technique de l'espace épointé
(
ef{ex:motivation-R2-R3}).}
\end{exercice}

\begin{exercice}[$\star\star$]\label{exo:pi1-retracte}
Soit $A$ un rétracte de $X$. Montrer que
$\iota_* : \pi_1(A, a_0) \to \pi_1(X, a_0)$ est injectif. En déduire que si
$\pi_1(X, x_0) = 0$, alors $\pi_1(A, a_0) = 0$ pour tout rétracte $A$.
\end{exercice}

\begin{exercice}[$\star\star\star$]\label{exo:pi1-cercle-winding}
Soit $\gamma : [0,1] \to \C \setminus \{0\}$ un lacet basé en $1$. Le
\textbf{nombre d'enroulement}\index{nombre d'enroulement} de $\gamma$ autour de
$0$ est l'entier
\[
n(\gamma, 0) = \frac{1}{2\pi i} \int_0^1 \frac{\gamma'(t)}{\gamma(t)} \, dt
\]
(lorsque $\gamma$ est de classe $\mathcal{C}^1$). Montrer que
$n(\gamma, 0) = \deg([\gamma])$ où $\deg$ est l'isomorphisme du

ef{thm:pi1-S1} (via la rétraction de $\C \setminus \{0\}$ sur $S^1$).
\end{exercice}

\begin{exercice}[$\star\star\star$]\label{exo:pi1-produit-preuve}
Donner une preuve détaillée du 
ef{thm:pi1-produit}. Vérifier soigneusement
que $\Phi$ et $\Psi$ sont des homomorphismes et qu'ils sont inverses l'un de
l'autre.
\end{exercice}

\begin{exercice}[$\star\star$]\label{exo:pi1-bouquet}
\index{bouquet}
Soit $X \vee Y = (X \sqcup Y) / (x_0 \sim y_0)$ le bouquet de deux espaces
pointés. Montrer que si $X$ et $Y$ sont des CW-complexes, le morphisme naturel
\[
\pi_1(X, x_0) * \pi_1(Y, y_0) \longrightarrow \pi_1(X \vee Y, *)
\]
est un isomorphisme (théorème de van Kampen appliqué au bouquet). En déduire que
$\pi_1(S^1 \vee S^1) \cong \Z * \Z$, le groupe libre\index{groupe libre} à
deux générateurs.
\end{exercice}

\begin{exercice}[$\star\star\star$]\label{exo:lacet-surface-genre}
\index{surface de genre $g$}
La surface orientable compacte $\Sigma_g$ de genre $g \geqslant 1$ admet une
CW-structure avec une $0$-cellule, $2g$ arêtes $a_1, b_1, \ldots, a_g, b_g$, et
une $2$-cellule attachée via le mot
$a_1 b_1 a_1^{-1} b_1^{-1} \cdots a_g b_g a_g^{-1} b_g^{-1}$.
En utilisant le théorème de van Kampen, montrer que
\[
\pi_1(\Sigma_g) \cong \left\langle a_1, b_1, \ldots, a_g, b_g \;\middle|\; \prod_{i=1}^g [a_i, b_i] = 1 \right\rangle.
\]
En déduire que $\pi_1(\Sigma_g)^{\mathrm{ab}} \cong \Z^{2g}$, et que les
surfaces $\Sigma_g$ et $\Sigma_h$ ne sont pas homéomorphes si $g \neq h$.
\end{exercice}

\begin{exercice}[$\star\star\star$]\label{exo:pi1-RP2}
\index{espace projectif réel!groupe fondamental}
En utilisant la CW-structure de $\RP^2$ décrite à l'
ef{ex:CW-RP2} et le
théorème de van Kampen, montrer que $\pi_1(\RP^2) \cong \Z/2\Z$. En déduire
que $\RP^2$ n'est pas simplement connexe et que la sphère $S^2$ est son
revêtement universel.
\end{exercice}

\begin{exercice}[$\star\star$]\label{exo:pi1-application-degree}
Soit $f : S^1 \to S^1$ une application continue. Le \textbf{degré}\index{degré}
de $f$ est l'entier $d \in \Z$ tel que
$f_* : \pi_1(S^1) \to \pi_1(S^1)$ soit la multiplication par $d$ (via
l'identification $\pi_1(S^1) \cong \Z$). Montrer que :
\begin{enumerate}[label=(\alph*)]
\item Si $f$ est un homéomorphisme, alors $\deg(f) = \pm 1$.
\item Si $f$ n'est pas surjective, alors $\deg(f) = 0$.
\item $\deg(f \circ g) = \deg(f) \cdot \deg(g)$.
\item L'application antipodale $a : S^1 \to S^1$, $z \mapsto -z$, est de
degré $1$.
\end{enumerate}
\end{exercice}

% === FIN DU CHUNK preamble+ch1-2 ===
% === DEBUT DU CHUNK ch3-4 ===

%%%%%%%%%%%%%%%%%%%%%%%%%%%%%%%%%%%%%%%%%%%%%%%%%%%%%%%%%%%%%%%%%%%%%
%% CHAPITRE 3 : THÉORÈME DE VAN KAMPEN
%%%%%%%%%%%%%%%%%%%%%%%%%%%%%%%%%%%%%%%%%%%%%%%%%%%%%%%%%%%%%%%%%%%%%
\chapter{Théorème de Van Kampen}\label{ch:van-kampen}

\minitoc

\vspace{1em}

Le théorème de Seifert--Van Kampen\index{théorème!de Van Kampen} constitue l'outil fondamental
pour le calcul effectif du groupe fondamental d'un espace topologique. Il exprime le groupe
fondamental d'une réunion $X = U_1 \cup U_2$ d'ouverts en termes des groupes fondamentaux
de $U_1$, $U_2$ et $U_1 \cap U_2$. Pour énoncer ce théorème dans toute sa généralité, nous
aurons besoin des notions de produit libre et de produit libre amalgamé de groupes.

%---------------------------------------------------------------
\section{Produit libre de groupes}\label{sec:produit-libre}
%---------------------------------------------------------------

\begin{definition}[Produit libre]\label{def:produit-libre}
\index{produit libre!de groupes}
Soient $(G_\alpha)_{\alpha \in A}$ une famille de groupes. Le \emph{produit libre}
$\displaystyle\mathop{\scalebox{1.3}{$\ast$}}_{\alpha \in A} G_\alpha$ est un groupe $G$ muni
de morphismes $\iota_\alpha \colon G_\alpha \to G$ satisfaisant la \emph{propriété universelle}
suivante : pour tout groupe $H$ et toute famille de morphismes $\varphi_\alpha \colon G_\alpha \to H$,
il existe un unique morphisme $\Phi \colon G \to H$ tel que $\Phi \circ \iota_\alpha = \varphi_\alpha$
pour tout $\alpha \in A$.
\end{definition}

La propriété universelle se résume par le diagramme commutatif suivant :
\[
\begin{tikzcd}
G_\alpha \arrow[r, "\iota_\alpha"] \arrow[dr, "\varphi_\alpha"'] & \displaystyle\mathop{\scalebox{1.3}{$\ast$}}_{\alpha \in A} G_\alpha \arrow[d, dashed, "\exists!\,\Phi"] \\
& H
\end{tikzcd}
\]

\begin{proposition}[Existence et unicité du produit libre]\label{prop:existence-produit-libre}
\index{produit libre!existence}
Le produit libre existe et est unique à isomorphisme unique près.
\end{proposition}

\begin{proof}
\textbf{Unicité.} Supposons que $(G, (\iota_\alpha))$ et $(G', (\iota'_\alpha))$
satisfassent tous deux la propriété universelle. En appliquant la propriété universelle de $G$
à la famille $(\iota'_\alpha)$, on obtient $\Phi \colon G \to G'$. Symétriquement, on obtient
$\Psi \colon G' \to G$. Alors $\Psi \circ \Phi \colon G \to G$ vérifie
$(\Psi \circ \Phi) \circ \iota_\alpha = \iota_\alpha$ pour tout $\alpha$. Par unicité dans
la propriété universelle (appliquée à $H = G$ et $\varphi_\alpha = \iota_\alpha$), on a
$\Psi \circ \Phi = \Id_G$. De même $\Phi \circ \Psi = \Id_{G'}$.

\textbf{Existence.} On considère l'ensemble des \emph{mots réduits}\index{mot réduit}
en les éléments des $G_\alpha$. Un mot est une suite finie
$(g_1, g_2, \ldots, g_n)$ où chaque $g_i \in G_{\alpha_i} \setminus \{e_{\alpha_i}\}$
et $\alpha_i \neq \alpha_{i+1}$ pour tout $i$. Le mot vide représente l'élément neutre.
La multiplication est donnée par concaténation suivie de réduction (si les dernières lettres
du premier mot et les premières du second proviennent du même facteur, on les multiplie dans
ce facteur, et on élimine le résultat s'il est trivial). On vérifie que cela définit une
structure de groupe satisfaisant la propriété universelle.
\end{proof}

\begin{notation}
Pour deux groupes $G$ et $H$, on note $G * H$ leur produit libre.\index{produit libre!notation}
Plus généralement, pour une famille finie, on écrit $G_1 * G_2 * \cdots * G_n$.
\end{notation}

\begin{exemple}[Produit libre $\Z * \Z$]\label{ex:Z-star-Z}
\index{produit libre!exemples}
Le produit libre $\Z * \Z$ est le \emph{groupe libre à deux générateurs}\index{groupe libre},
noté $F_2$. Ses éléments sont les mots réduits en $a^{\pm 1}$ et $b^{\pm 1}$, par exemple
$a^2 b^{-1} a b^3$. Ce groupe est non abélien : $ab \neq ba$.
\end{exemple}

\begin{exemple}[Produit libre $\Z/2\Z * \Z/3\Z$]\label{ex:Z2-star-Z3}
Le produit libre $\Z/2\Z * \Z/3\Z$ est isomorphe au groupe modulaire
$\mathrm{PSL}(2, \Z)$.\index{groupe modulaire} Si $a$ et $b$ désignent les générateurs
respectifs de $\Z/2\Z$ et $\Z/3\Z$, les éléments de $\Z/2\Z * \Z/3\Z$ sont les
mots réduits alternant $a$ et des puissances de $b$ (c'est-à-dire $b$ ou $b^2 = b^{-1}$).
\end{exemple}

\begin{remarque}[Produit libre et abélianisation]\label{rem:abelianise-produit-libre}
L'abélianisé du produit libre est le produit direct :
$\left(\displaystyle\mathop{\scalebox{1.3}{$\ast$}}_\alpha G_\alpha\right)^{\mathrm{ab}}
\cong \bigoplus_\alpha G_\alpha^{\mathrm{ab}}$.
\end{remarque}

%---------------------------------------------------------------
\section{Produit libre amalgamé}\label{sec:produit-amalgame}
%---------------------------------------------------------------

\begin{definition}[Produit libre amalgamé]\label{def:produit-amalgame}
\index{produit libre amalgamé}
Soient $G_1$, $G_2$ des groupes et $K$ un groupe muni de morphismes
$j_1 \colon K \to G_1$ et $j_2 \colon K \to G_2$. Le \emph{produit libre amalgamé}
de $G_1$ et $G_2$ au-dessus de $K$, noté $G_1 *_K G_2$, est le groupe défini par
la propriété universelle suivante : c'est un groupe $G$ muni de morphismes
$\iota_1 \colon G_1 \to G$ et $\iota_2 \colon G_2 \to G$ tels que
$\iota_1 \circ j_1 = \iota_2 \circ j_2$, et universel pour cette propriété.

Autrement dit, pour tout groupe $H$ et tous morphismes $\varphi_1 \colon G_1 \to H$,
$\varphi_2 \colon G_2 \to H$ vérifiant $\varphi_1 \circ j_1 = \varphi_2 \circ j_2$,
il existe un unique morphisme $\Phi \colon G \to H$ tel que $\Phi \circ \iota_i = \varphi_i$.
\end{definition}

En langage catégorique, $G_1 *_K G_2$ est le \emph{pushout}\index{pushout!de groupes}
(colimite cocartésienne) dans la catégorie des groupes :
\[
\begin{tikzcd}
K \arrow[r, "j_1"] \arrow[d, "j_2"'] & G_1 \arrow[d, "\iota_1"] \\
G_2 \arrow[r, "\iota_2"'] & G_1 *_K G_2 \arrow[ul, phantom, "\ulcorner" very near start]
\end{tikzcd}
\]

\begin{proposition}[Construction du produit libre amalgamé]\label{prop:construction-amalgame}
Le produit libre amalgamé $G_1 *_K G_2$ est isomorphe au quotient
\[
G_1 *_K G_2 \;\cong\; (G_1 * G_2) \,\big/\, N
\]
où $N$ est le sous-groupe normal engendré par les éléments
$\set{j_1(k)\, j_2(k)^{-1} \mid k \in K}$.
\end{proposition}

\begin{proof}
Notons $G = (G_1 * G_2)/N$ et $\pi \colon G_1 * G_2 \to G$ la projection canonique.
Posons $\iota_i = \pi \circ \iota_i'$ où $\iota_i' \colon G_i \hookrightarrow G_1 * G_2$
est l'inclusion canonique. Par construction, pour tout $k \in K$, on a
$\iota_1(j_1(k)) = \iota_2(j_2(k))$ dans $G$.

Soit $H$ un groupe et $\varphi_i \colon G_i \to H$ des morphismes avec
$\varphi_1 \circ j_1 = \varphi_2 \circ j_2$. Par la propriété universelle du produit libre,
il existe un unique $\Phi' \colon G_1 * G_2 \to H$ avec $\Phi' \circ \iota_i' = \varphi_i$.
Alors $\Phi'(j_1(k)\, j_2(k)^{-1}) = \varphi_1(j_1(k))\, \varphi_2(j_2(k))^{-1} = e$,
donc $N \subset \ker \Phi'$. Par passage au quotient, on obtient
$\Phi \colon G \to H$ avec les propriétés requises. L'unicité découle de celle
pour le produit libre.
\end{proof}

\begin{exemple}[Présentation par générateurs et relations]\label{ex:amalgame-presentation}
\index{présentation!de groupe}
Si $G_1 = \langle S_1 \mid R_1 \rangle$ et $G_2 = \langle S_2 \mid R_2 \rangle$
(avec $S_1 \cap S_2 = \varnothing$), et $K$ est engendré par des éléments dont les images
dans $G_1$ et $G_2$ sont $w_1^{(i)}$ et $w_2^{(i)}$ respectivement, alors
\[
G_1 *_K G_2 = \langle S_1 \cup S_2 \mid R_1 \cup R_2 \cup \set{w_1^{(i)} = w_2^{(i)}} \rangle.
\]
\end{exemple}

%---------------------------------------------------------------
\section{Énoncé du théorème de Van Kampen}\label{sec:enonce-vk}
%---------------------------------------------------------------

\begin{theoreme}[{Seifert--Van Kampen}]\label{thm:van-kampen}
\index{théorème!de Van Kampen}
Soit $X$ un espace topologique, et soient $U_1, U_2 \subset X$ des ouverts connexes par arcs
tels que $X = U_1 \cup U_2$ et $U_0 = U_1 \cap U_2$ est connexe par arcs et non vide.
Fixons un point base $x_0 \in U_0$, et notons $j_k \colon \pi_1(U_0, x_0) \to \pi_1(U_k, x_0)$
et $\iota_k \colon \pi_1(U_k, x_0) \to \pi_1(X, x_0)$ les morphismes induits par les inclusions.
Alors le diagramme
\[
\begin{tikzcd}
\pi_1(U_0, x_0) \arrow[r, "j_1"] \arrow[d, "j_2"'] & \pi_1(U_1, x_0) \arrow[d, "\iota_1"] \\
\pi_1(U_2, x_0) \arrow[r, "\iota_2"'] & \pi_1(X, x_0)
\arrow[ul, phantom, "\ulcorner" very near start]
\end{tikzcd}
\]
est un carré cocartésien (pushout) dans la catégorie des groupes. Autrement dit,
\[
\pi_1(X, x_0) \;\cong\; \pi_1(U_1, x_0) *_{\pi_1(U_0, x_0)} \pi_1(U_2, x_0).
\]
\end{theoreme}

\begin{remarque}[Version groupoïde]\label{rem:vk-groupoide}
\index{théorème!de Van Kampen!version groupoïde}
R.~Brown a donné une version du théorème en termes de groupoïdes fondamentaux\index{groupoïde fondamental}.
Si $A \subset X$ est un sous-ensemble qui rencontre chaque composante connexe par arcs de
$U_0$, $U_1$, $U_2$ et $X$, alors le diagramme
\[
\begin{tikzcd}
\Pi_1(U_0, A \cap U_0) \arrow[r] \arrow[d] & \Pi_1(U_1, A \cap U_1) \arrow[d] \\
\Pi_1(U_2, A \cap U_2) \arrow[r] & \Pi_1(X, A)
\arrow[ul, phantom, "\ulcorner" very near start]
\end{tikzcd}
\]
est un pushout dans la catégorie des groupoïdes. Cette version permet de traiter le cas
où $U_0$ n'est pas connexe par arcs.
\end{remarque}

\begin{corollaire}\label{cor:vk-trivial}
Si $U_0$ est simplement connexe, alors
$\pi_1(X, x_0) \cong \pi_1(U_1, x_0) * \pi_1(U_2, x_0)$.
\end{corollaire}

\begin{proof}
Si $\pi_1(U_0, x_0) = \{e\}$, les morphismes $j_1, j_2$ sont triviaux,
et le produit libre amalgamé au-dessus du groupe trivial est simplement le produit libre.
\end{proof}

%---------------------------------------------------------------
\section{Preuve du théorème de Van Kampen}\label{sec:preuve-vk}
%---------------------------------------------------------------

\begin{proof}[Preuve du théorème~\ref{thm:van-kampen}]
Notons $G = \pi_1(U_1, x_0) *_{\pi_1(U_0, x_0)} \pi_1(U_2, x_0)$ le produit libre amalgamé.
Par la propriété universelle du pushout, les morphismes $\iota_k$ induisent un unique morphisme
\[
\Phi \colon G \longrightarrow \pi_1(X, x_0).
\]
Nous devons montrer que $\Phi$ est un isomorphisme. La preuve comporte deux étapes :
la surjectivité et l'injectivité.

\medskip
\textbf{Étape 1 : Surjectivité de $\Phi$.}
Soit $[\gamma] \in \pi_1(X, x_0)$, où $\gamma \colon [0,1] \to X$ est un lacet basé en $x_0$.
Comme $\{U_1, U_2\}$ est un recouvrement ouvert de $X$ et $[0,1]$ est compact, par le
lemme du nombre de Lebesgue\index{lemme!du nombre de Lebesgue}, il existe une subdivision
\[
0 = t_0 < t_1 < \cdots < t_n = 1
\]
telle que, pour chaque $i$, l'image $\gamma([t_{i-1}, t_i])$ est contenue dans $U_1$ ou dans $U_2$.

Pour chaque $i = 1, \ldots, n-1$, le point $\gamma(t_i)$ appartient à $U_1$ ou $U_2$.
S'il n'appartient pas à $U_0$, les segments adjacents $\gamma|_{[t_{i-1},t_i]}$ et
$\gamma|_{[t_i, t_{i+1}]}$ sont dans le même ouvert $U_k$, et on peut fusionner les
intervalles. Quitte à raffiner, on peut donc supposer que $\gamma(t_i) \in U_0$ pour tout $i$.

Comme $U_0$ est connexe par arcs, on choisit un chemin $\alpha_i$ dans $U_0$
de $x_0$ à $\gamma(t_i)$ (avec $\alpha_0 = \alpha_n$ le chemin constant en $x_0$).
Posons
\[
\gamma_i = \alpha_{i-1}^{-1} \cdot \gamma|_{[t_{i-1}, t_i]} \cdot \alpha_i^{-1}.
\]
Alors $\gamma_i$ est un lacet basé en $x_0$ contenu dans $U_1$ ou dans $U_2$, et
$[\gamma] = [\gamma_1] \cdot [\gamma_2] \cdots [\gamma_n]$ dans $\pi_1(X, x_0)$.
Chaque $[\gamma_i]$ est dans l'image de $\iota_1$ ou $\iota_2$, donc
$[\gamma] \in \im(\Phi)$.

\medskip
\textbf{Étape 2 : Injectivité de $\Phi$.}
Supposons que $g = g_1 g_2 \cdots g_m \in G$ (écrit comme mot en des éléments de
$\pi_1(U_1, x_0)$ et $\pi_1(U_2, x_0)$) satisfasse $\Phi(g) = e$ dans $\pi_1(X, x_0)$.
Chaque $g_i = [\gamma_i]$ est un lacet dans $U_1$ ou $U_2$. Alors $\gamma_1 \cdot \gamma_2
\cdots \gamma_m$ est homotope au lacet constant dans $X$, disons via une homotopie
$H \colon [0,1]^2 \to X$.

Par le lemme du nombre de Lebesgue appliqué au recouvrement ouvert
$\{H^{-1}(U_1), H^{-1}(U_2)\}$ du carré compact $[0,1]^2$, il existe un quadrillage
du carré en petits rectangles dont chacun a son image dans $U_1$ ou dans $U_2$.

On peut alors, par récurrence sur le nombre de petits carrés, décomposer l'homotopie
en une suite de modifications élémentaires. Chaque modification élémentaire est de l'une
des formes suivantes :
\begin{enumerate}[label=(\roman*)]
\item remplacement d'un facteur $g_i \in \pi_1(U_k, x_0)$ par un produit
$g_i' g_i''$ dans $\pi_1(U_k, x_0)$ (avec $g_i' g_i'' = g_i$ dans $\pi_1(U_k, x_0)$) ;
\item insertion de $j_1(h) \cdot j_2(h)^{-1}$ (ou son inverse) pour $h \in \pi_1(U_0, x_0)$,
ce qui est trivial dans le produit amalgamé par construction ;
\item suppression d'un facteur trivial.
\end{enumerate}

Chacune de ces opérations est triviale dans le produit libre amalgamé $G$.
Donc $g = e$ dans $G$, ce qui prouve l'injectivité.
\end{proof}

%---------------------------------------------------------------
\section{Applications}\label{sec:applications-vk}
%---------------------------------------------------------------

\subsection{Bouquet de cercles}\label{subsec:bouquet-cercles}

\begin{definition}[Bouquet]\label{def:bouquet}
\index{bouquet!de cercles}
Le \emph{bouquet} (ou \emph{wedge sum}\index{wedge sum}) de deux espaces pointés $(X, x_0)$ et $(Y, y_0)$ est
\[
X \vee Y = (X \sqcup Y) / (x_0 \sim y_0).
\]
\end{definition}

\begin{center}
\begin{tikzpicture}[scale=1.2]
  % Bouquet de deux cercles
  \draw[thick, blue] (0,0) circle (0.8);
  \draw[thick, red] (1.6,0) circle (0.8);
  \fill (0.8,0) circle (2pt) node[below=3pt] {$x_0$};
  \node at (0, -1.3) {$S^1 \vee S^1$};
  % Bouquet de trois cercles
  \begin{scope}[xshift=5cm]
    \draw[thick, blue] (0,0) circle (0.7);
    \draw[thick, red] (1.4,0) circle (0.7);
    \draw[thick, green!60!black] (0.7, 1.2) circle (0.7);
    \fill (0.7,0) circle (2pt) node[below=3pt] {$x_0$};
    \node at (0.7, -1.3) {$S^1 \vee S^1 \vee S^1$};
  \end{scope}
\end{tikzpicture}
\end{center}

\begin{theoreme}[{Groupe fondamental de $S^1 \vee S^1$}]\label{thm:pi1-bouquet}
\index{groupe fondamental!du bouquet de cercles}
$\pi_1(S^1 \vee S^1) \cong \Z * \Z$, le groupe libre à deux générateurs.
Plus généralement, $\pi_1\!\left(\bigvee_{i=1}^n S^1\right) \cong F_n$, le groupe libre
de rang $n$.
\end{theoreme}

\begin{proof}
Identifions $S^1 \vee S^1$ avec deux cercles $C_1, C_2$ se touchant en un point $x_0$.
Posons $U_1 = X \setminus \{p_2\}$ et $U_2 = X \setminus \{p_1\}$, où $p_i$ est le point
antipodal à $x_0$ sur $C_i$. Alors :
\begin{itemize}
\item $U_1$ se rétracte par déformation sur $C_1$, donc $\pi_1(U_1, x_0) \cong \Z$ ;
\item $U_2$ se rétracte par déformation sur $C_2$, donc $\pi_1(U_2, x_0) \cong \Z$ ;
\item $U_0 = U_1 \cap U_2$ est une croix (réunion de deux intervalles ouverts), qui est contractile :
$\pi_1(U_0, x_0) = \{e\}$.
\end{itemize}
Par le corollaire~\ref{cor:vk-trivial},
$\pi_1(S^1 \vee S^1, x_0) \cong \Z * \Z = F_2$.
Le cas de $n$ cercles se traite par récurrence.
\end{proof}

\subsection{Plan projectif réel}\label{subsec:rp2}

\begin{theoreme}[{Groupe fondamental de $\RP^2$}]\label{thm:pi1-rp2}
\index{groupe fondamental!du plan projectif réel}
\index{plan projectif réel}
$\pi_1(\RP^2) \cong \Z/2\Z$.
\end{theoreme}

\begin{proof}
On réalise $\RP^2$ comme quotient du disque $D^2$ par l'identification antipodale sur le bord :
$x \sim -x$ pour $x \in S^1 = \partial D^2$.

\begin{center}
\begin{tikzpicture}[scale=1.5]
  \draw[thick] (0,0) circle (1);
  \draw[->, thick, blue] (90:1) arc (90:-90:1);
  \draw[->, thick, blue] (270:1) arc (270:90:1);
  \node[blue, right] at (1.1, 0) {$a$};
  \node[blue, left] at (-1.1, 0) {$a$};
  \draw[->] (70:1.15) -- (70:1.15) node[above right] {};
  \fill (90:1) circle (1.5pt);
  \fill (270:1) circle (1.5pt);
  \node at (0, -1.5) {$\RP^2$ : identification $a = a$};
\end{tikzpicture}
\end{center}

On peut aussi voir $\RP^2$ comme le CW-complexe obtenu en attachant une $2$-cellule à $S^1$
par l'application $\varphi \colon S^1 \to S^1$ de degré $2$ (l'application $z \mapsto z^2$).
L'espace $X^{(1)} = S^1$ a pour groupe fondamental $\Z$, engendré par $a$.
L'attachement de la $2$-cellule ajoute la relation $a^2 = e$ (le bord de la cellule parcourt
le générateur deux fois). Par Van Kampen (voir le théorème~\ref{thm:attachement-cellule} ci-dessous),
\[
\pi_1(\RP^2) = \langle a \mid a^2 = 1 \rangle \cong \Z/2\Z. \qedhere
\]
\end{proof}

\subsection{Surfaces orientables de genre \texorpdfstring{$g$}{g}}\label{subsec:surfaces-genre}

\begin{definition}[Surface orientable de genre $g$]\label{def:surface-genre}
\index{surface!orientable de genre $g$}
La surface orientable fermée de genre $g$, notée $\Sigma_g$, est obtenue en identifiant les
côtés d'un $4g$-gone selon le mot
\[
a_1 b_1 a_1^{-1} b_1^{-1}\, a_2 b_2 a_2^{-1} b_2^{-1} \,\cdots\, a_g b_g a_g^{-1} b_g^{-1}.
\]
\end{definition}

\begin{center}
\begin{tikzpicture}[scale=2]
  % Octogone pour g=2
  \foreach \i in {0,...,7} {
    \coordinate (P\i) at ({90+\i*45}:1.2);
  }
  \draw[thick] (P0) -- (P1) -- (P2) -- (P3) -- (P4) -- (P5) -- (P6) -- (P7) -- cycle;
  % Étiquettes
  \node[above] at ($(P0)!0.5!(P1)$) {\color{blue}$a_1$};
  \node[above left] at ($(P1)!0.5!(P2)$) {\color{red}$b_1$};
  \node[left] at ($(P2)!0.5!(P3)$) {\color{blue}$a_1^{-1}$};
  \node[below left] at ($(P3)!0.5!(P4)$) {\color{red}$b_1^{-1}$};
  \node[below] at ($(P4)!0.5!(P5)$) {\color{blue}$a_2$};
  \node[below right] at ($(P5)!0.5!(P6)$) {\color{red}$b_2$};
  \node[right] at ($(P6)!0.5!(P7)$) {\color{blue}$a_2^{-1}$};
  \node[above right] at ($(P7)!0.5!(P0)$) {\color{red}$b_2^{-1}$};
  % Flèches sur les côtés
  \foreach \i/\nexti in {0/1, 1/2, 2/3, 3/4, 4/5, 5/6, 6/7, 7/0} {
    \draw[->, thick] ($(P\i)!0.4!(P\nexti)$) -- ($(P\i)!0.6!(P\nexti)$);
  }
  \node at (0, -1.7) {Présentation polygonale de $\Sigma_2$};
\end{tikzpicture}
\end{center}

\begin{theoreme}[{Groupe fondamental de $\Sigma_g$}]\label{thm:pi1-sigma-g}
\index{groupe fondamental!des surfaces orientables}
Pour $g \geq 1$,
\[
\pi_1(\Sigma_g) \;\cong\; \left\langle
a_1, b_1, \ldots, a_g, b_g \;\middle|\;
\prod_{i=1}^g [a_i, b_i] = 1
\right\rangle
\]
où $[a_i, b_i] = a_i b_i a_i^{-1} b_i^{-1}$ est le commutateur.
\end{theoreme}

\begin{proof}
On voit $\Sigma_g$ comme un CW-complexe avec un $0$-squelette réduit à un point $x_0$,
un $1$-squelette $X^{(1)} = \bigvee_{i=1}^{2g} S^1$ (bouquet de $2g$ cercles étiquetés
$a_1, b_1, \ldots, a_g, b_g$) et une unique $2$-cellule $e^2$ attachée par l'application
$\varphi \colon S^1 \to X^{(1)}$ qui parcourt le mot $\prod_{i=1}^g [a_i, b_i]$.

Par le théorème~\ref{thm:pi1-bouquet}, $\pi_1(X^{(1)}) \cong F_{2g}$. L'attachement de la
$2$-cellule (théorème~\ref{thm:attachement-cellule}) ajoute la relation
$\prod_{i=1}^g [a_i, b_i] = 1$.
\end{proof}

\begin{corollaire}\label{cor:abelianise-sigma}
$H_1(\Sigma_g; \Z) \cong \pi_1(\Sigma_g)^{\mathrm{ab}} \cong \Z^{2g}$.
\end{corollaire}

\begin{proof}
L'abélianisé du groupe $\pi_1(\Sigma_g)$ est obtenu en rendant triviale la relation
$\prod [a_i, b_i] = 1$, qui est déjà triviale dans un groupe abélien. Donc l'abélianisé
est le groupe abélien libre de rang $2g$.
\end{proof}

\subsection{Bouteille de Klein}\label{subsec:klein}

\begin{definition}[Bouteille de Klein]\label{def:bouteille-klein}
\index{bouteille de Klein}
La \emph{bouteille de Klein} $K$ est la surface non orientable obtenue en identifiant les
côtés du carré selon le mot $aba^{-1}b$.
\end{definition}

\begin{center}
\begin{tikzpicture}[scale=2]
  \draw[thick] (0,0) rectangle (1.5, 1.5);
  % Flèches
  \draw[->, thick, blue] (0, 0.4) -- (0, 0.9) node[midway, left] {$b$};
  \draw[->, thick, blue] (1.5, 0.9) -- (1.5, 0.4) node[midway, right] {$b$};
  \draw[->, thick, red] (0.4, 0) -- (0.9, 0) node[midway, below] {$a$};
  \draw[->, thick, red] (0.4, 1.5) -- (0.9, 1.5) node[midway, above] {$a$};
  \node at (0.75, -0.5) {Bouteille de Klein : $aba^{-1}b$};
\end{tikzpicture}
\end{center}

\begin{theoreme}[{Groupe fondamental de la bouteille de Klein}]\label{thm:pi1-klein}
\index{groupe fondamental!de la bouteille de Klein}
$\pi_1(K) \cong \langle a, b \mid abab^{-1} = 1 \rangle$.
\end{theoreme}

\begin{proof}
La bouteille de Klein admet une structure de CW-complexe avec un $0$-sommet, deux $1$-cellules
$a, b$ et une $2$-cellule attachée selon le mot $aba^{-1}b$. Le $1$-squelette est un bouquet
de deux cercles, de groupe fondamental $F_2 = \langle a, b \rangle$. L'attachement de la
$2$-cellule ajoute la relation $aba^{-1}b = 1$, soit $abab^{-1} = 1$ (en réécrivant
$aba^{-1}b = 1$ comme $ab = b^{-1}a$).
\end{proof}

\begin{remarque}\label{rem:klein-abelianise}
L'abélianisé est $\pi_1(K)^{\mathrm{ab}} \cong \Z \oplus \Z/2\Z$.
En effet, dans un groupe abélien la relation $aba^{-1}b = 1$ devient $b^2 = 1$.
\end{remarque}

\subsection{Attachement de cellules et CW-complexes}\label{subsec:attachement-cellules}

\begin{theoreme}[Effet de l'attachement d'une $2$-cellule]\label{thm:attachement-cellule}
\index{CW-complexe!attachement de cellule}
\index{attachement de cellule}
Soit $X$ un espace connexe par arcs, $\varphi \colon S^1 \to X$ une application continue, et
$Y = X \cup_\varphi D^2$ l'espace obtenu en attachant une $2$-cellule via $\varphi$.
Alors, pour tout point base $x_0 \in X$,
\[
\pi_1(Y, x_0) \;\cong\; \pi_1(X, x_0) \,\big/\, N([\varphi])
\]
où $N([\varphi])$ est le sous-groupe normal engendré par la classe de $\varphi$
dans $\pi_1(X, x_0)$ (à conjugaison près, dépendant du choix d'un chemin de $x_0$ à
$\varphi(1)$).
\end{theoreme}

\begin{proof}
On applique Van Kampen. Soit $p$ le centre de la cellule $D^2$ dans $Y$.
Posons $U_1 = X \cup_\varphi (D^2 \setminus \{p\})$ et $U_2 = \mathring{D}^2$ (l'intérieur
de la cellule, vu dans $Y$). Alors $U_1 \cup U_2 = Y$, $U_2$ est contractile, et
$U_1 \cap U_2 \cong S^1 \times (0,1)$ se rétracte sur un cercle.
De plus, $U_1$ se rétracte par déformation sur $X$.

Par Van Kampen :
$\pi_1(Y) \cong \pi_1(U_1) *_{\pi_1(U_0)} \pi_1(U_2) \cong \pi_1(X) *_\Z \{e\}$.
Le morphisme $\Z = \pi_1(U_0) \to \pi_1(U_1) \cong \pi_1(X)$ envoie le générateur
sur $[\varphi]$. Le produit libre amalgamé $\pi_1(X) *_\Z \{e\}$ est le quotient de
$\pi_1(X)$ par le sous-groupe normal engendré par l'image de $\Z$, soit $N([\varphi])$.
\end{proof}

\begin{corollaire}\label{cor:cellules-dim-sup}
L'attachement de cellules de dimension $n \geq 3$ ne modifie pas le groupe fondamental.
\end{corollaire}

\begin{proof}
Si $Y = X \cup_\varphi D^n$ avec $n \geq 3$, on applique Van Kampen de la même manière.
L'intersection $U_0 \simeq S^{n-1}$ avec $n-1 \geq 2$ est simplement connexe
(car $\pi_1(S^m) = 0$ pour $m \geq 2$). Donc $j_1$ et $j_2$ sont triviaux,
et $\pi_1(Y) \cong \pi_1(X) * \{e\} = \pi_1(X)$.
\end{proof}

\begin{proposition}[Groupe fondamental d'un CW-complexe]\label{prop:pi1-cw}
\index{CW-complexe!groupe fondamental}
Soit $X$ un CW-complexe connexe. Alors $\pi_1(X)$ ne dépend que du $2$-squelette $X^{(2)}$,
et admet la présentation
\[
\pi_1(X) \cong \langle\, \text{$1$-cellules} \mid \text{relations données par les $2$-cellules} \,\rangle.
\]
\end{proposition}

\begin{proof}
Notons $X^{(n)}$ le $n$-squelette de $X$. Par le théorème~\ref{thm:pi1-bouquet},
après avoir choisi un arbre maximal dans le $1$-squelette $X^{(1)}$ et contracté celui-ci,
on obtient un bouquet de cercles, un par $1$-cellule hors de l'arbre maximal. D'où
$\pi_1(X^{(1)})$ est un groupe libre.

Le passage de $X^{(1)}$ à $X^{(2)}$ se fait par attachement de $2$-cellules.
Par le théorème~\ref{thm:attachement-cellule}, chaque $2$-cellule ajoute une relation
(donnée par le mot de recollement lu sur les $1$-cellules).

Le passage de $X^{(2)}$ à $X^{(3)}$ se fait par attachement de $3$-cellules. Par le
corollaire~\ref{cor:cellules-dim-sup}, cela ne modifie pas $\pi_1$. Par récurrence,
l'attachement des $n$-cellules pour $n \geq 3$ ne modifie pas le groupe fondamental.
\end{proof}

\begin{exemple}[Tout groupe de présentation finie est un $\pi_1$]\label{ex:tout-groupe-pi1}
\index{groupe!de présentation finie}
Soit $G = \langle g_1, \ldots, g_m \mid r_1, \ldots, r_k \rangle$ un groupe de présentation
finie. On construit un CW-complexe $X_G$ avec :
\begin{itemize}
\item un unique $0$-sommet $x_0$;
\item $m$ boucles $e_1^1, \ldots, e_m^1$ (une par générateur), formant un bouquet $X^{(1)} = \bigvee_m S^1$;
\item $k$ cellules de dimension $2$, la $j$-ème étant attachée selon le mot $r_j$ lu dans les générateurs.
\end{itemize}
Alors $\pi_1(X_G) \cong G$. Ceci montre que tout groupe de présentation finie est le groupe
fondamental d'un CW-complexe de dimension $2$.
\end{exemple}

%---------------------------------------------------------------
\section{Exercices}\label{sec:exo-vk}
%---------------------------------------------------------------

\begin{exercice}\label{exo:pi1-tore-vk}
Retrouver $\pi_1(T^2) \cong \Z^2$ en utilisant le théorème de Van Kampen appliqué à la
décomposition standard du tore en carré identifié. \textbf{Difficulté : \(\bigstar\bigstar\)}
\end{exercice}

\begin{exercice}\label{exo:pi1-sphere}
Montrer que $\pi_1(S^n) = 0$ pour $n \geq 2$ en utilisant Van Kampen.
\textbf{Difficulté : \(\bigstar\bigstar\)}
\end{exercice}

\begin{exercice}\label{exo:pi1-bouquet-n}
Montrer par récurrence que $\pi_1\!\left(\bigvee_{i=1}^n S^1\right) \cong F_n$.
\textbf{Difficulté : \(\bigstar\bigstar\)}
\end{exercice}

\begin{exercice}\label{exo:amalgame-calcul}
Calculer $\Z *_\Z \Z$ où les deux morphismes $\Z \to \Z$ sont donnés par
$n \mapsto 2n$ et $n \mapsto 3n$ respectivement. Montrer que le groupe obtenu admet la
présentation $\langle a, b \mid a^2 = b^3 \rangle$.
\textbf{Difficulté : \(\bigstar\bigstar\)}
\end{exercice}

\begin{exercice}\label{exo:pi1-surface-non-orientable}
\index{surface!non orientable}
La surface non orientable fermée $N_k$ de genre $k$ est obtenue en identifiant les côtés d'un
$2k$-gone selon le mot $a_1^2 a_2^2 \cdots a_k^2$. Montrer que
\[
\pi_1(N_k) \cong \langle a_1, \ldots, a_k \mid a_1^2 a_2^2 \cdots a_k^2 = 1 \rangle.
\]
Vérifier que $N_1 = \RP^2$ et $N_2 = K$ (bouteille de Klein).
\textbf{Difficulté : \(\bigstar\bigstar\bigstar\)}
\end{exercice}

\begin{exercice}\label{exo:vk-trois-ouverts}
Soit $X = U_1 \cup U_2 \cup U_3$ avec $U_i$ ouverts connexes par arcs, $U_i \cap U_j$
connexes par arcs, et $U_1 \cap U_2 \cap U_3$ connexe par arcs. Énoncer et démontrer
une généralisation du théorème de Van Kampen à ce cas.
\textbf{Difficulté : \(\bigstar\bigstar\bigstar\)}
\end{exercice}

\begin{exercice}\label{exo:dunce-hat}
\index{dunce hat}
Le \emph{dunce hat} est le CW-complexe obtenu à partir d'un triangle en identifiant les
trois arêtes selon le schéma $a \cdot a \cdot a^{-1}$. Calculer son groupe fondamental.
\textbf{Difficulté : \(\bigstar\bigstar\bigstar\)}
\end{exercice}

\begin{exercice}\label{exo:espace-lenticulaire}
\index{espace lenticulaire}
L'espace lenticulaire $L(p, q)$ est le quotient de $S^3 \subset \C^2$ par l'action de
$\Z/p\Z$ donnée par $[k] \cdot (z_1, z_2) = (e^{2\pi i k/p} z_1, e^{2\pi i kq/p} z_2)$.
Montrer que $\pi_1(L(p,q)) \cong \Z/p\Z$.
\textbf{Difficulté : \(\bigstar\bigstar\bigstar\)}
\end{exercice}

\begin{exercice}\label{exo:complement-noeud}
Soit $K \subset S^3$ un n\oe ud (plongement de $S^1$ dans $S^3$). Montrer, en utilisant
Van Kampen, que l'abélianisé du groupe du n\oe ud $\pi_1(S^3 \setminus K)$ est isomorphe à $\Z$.
\textbf{Difficulté : \(\bigstar\bigstar\bigstar\)}
\end{exercice}

%---------------------------------------------------------------
\section*{Résumé du chapitre}\label{sec:resume-ch3}
\addcontentsline{toc}{section}{Résumé du chapitre}
%---------------------------------------------------------------

\begin{center}
\renewcommand{\arraystretch}{1.4}
\begin{tabular}{p{5cm} p{9cm}}
\toprule
\textbf{Concept} & \textbf{Résultat principal} \\
\midrule
Produit libre $G_1 * G_2$ &
Groupe universel pour les morphismes depuis $G_1$ et $G_2$. \\
Produit amalgamé $G_1 *_K G_2$ &
Pushout dans la catégorie des groupes; quotient de $G_1 * G_2$ par les relations d'amalgamation. \\
Théorème de Van Kampen &
$\pi_1(U_1 \cup U_2) \cong \pi_1(U_1) *_{\pi_1(U_0)} \pi_1(U_2)$
quand $U_0 = U_1 \cap U_2$ est connexe par arcs. \\
Bouquet de cercles &
$\pi_1\!\left(\bigvee_n S^1\right) \cong F_n$. \\
Surfaces &
$\pi_1(\Sigma_g) = \langle a_i, b_i \mid \prod [a_i, b_i] = 1 \rangle$. \\
Attachement de $2$-cellule &
$\pi_1(X \cup_\varphi D^2) \cong \pi_1(X) / N([\varphi])$. \\
\bottomrule
\end{tabular}
\end{center}


%%%%%%%%%%%%%%%%%%%%%%%%%%%%%%%%%%%%%%%%%%%%%%%%%%%%%%%%%%%%%%%%%%%%%
%% CHAPITRE 4 : REVÊTEMENTS ET GROUPE FONDAMENTAL
%%%%%%%%%%%%%%%%%%%%%%%%%%%%%%%%%%%%%%%%%%%%%%%%%%%%%%%%%%%%%%%%%%%%%
\chapter{Revêtements et groupe fondamental}\label{ch:revetements}

\minitoc

\vspace{1em}

La théorie des revêtements\index{revêtement} établit un lien profond entre la topologie d'un espace
et l'algèbre de son groupe fondamental. Dans ce chapitre, nous développons cette théorie en
détail : propriétés de relèvement, action de monodromie, classification des revêtements
connexes, revêtement universel, et transformations de revêtement.

%---------------------------------------------------------------
\section{Définitions et premiers exemples}\label{sec:rev-def}
%---------------------------------------------------------------

\begin{definition}[Revêtement]\label{def:revetement}
\index{revêtement!définition}
Soit $X$ un espace topologique connexe et localement connexe par arcs.
Un \emph{revêtement} de $X$ est la donnée d'un espace topologique $\widetilde{X}$
et d'une application continue surjective $p \colon \widetilde{X} \to X$ telle que,
pour tout $x \in X$, il existe un voisinage ouvert $U$ de $x$ (appelé \emph{voisinage
trivialisé}\index{voisinage trivialisé}) vérifiant
\[
p^{-1}(U) = \bigsqcup_{\alpha \in A} V_\alpha
\]
où chaque $V_\alpha$ est un ouvert de $\widetilde{X}$ et $p|_{V_\alpha} \colon V_\alpha \xrightarrow{\;\sim\;} U$
est un homéomorphisme.
\end{definition}

\begin{definition}[Fibre, degré]\label{def:fibre-degre}
\index{fibre d'un revêtement}\index{revêtement!degré}
Pour $x \in X$, l'ensemble $p^{-1}(x)$ est appelé la \emph{fibre} au-dessus de $x$.
Si toutes les fibres ont le même cardinal $n \in \N \cup \{\infty\}$, on dit que $p$ est un
revêtement de \emph{degré} (ou \emph{nombre de feuillets}) $n$, noté $\deg(p) = n$.
Lorsque $X$ est connexe, le cardinal des fibres est constant.
\end{definition}

\begin{exemple}[{Revêtement exponentiel $\R \to S^1$}]\label{ex:rev-exp}
\index{revêtement!exponentiel}
L'application $p \colon \R \to S^1$ définie par $p(t) = e^{2\pi i t}$ est un revêtement.
Pour tout $z_0 = e^{2\pi i t_0} \in S^1$, un voisinage trivialisé est
$U = S^1 \setminus \{-z_0\}$, et $p^{-1}(U) = \bigsqcup_{n \in \Z} (t_0 - \tfrac{1}{2} + n,\, t_0 + \tfrac{1}{2} + n)$.
C'est un revêtement de degré infini.

\begin{center}
\begin{tikzpicture}[scale=0.9]
  % Hélice
  \draw[thick, ->] (0, -0.5) -- (0, 5) node[right] {$\R$};
  \foreach \k in {0, 1, 2, 3} {
    \draw[blue, thick] ({-1.2*cos(0)}, \k) arc[start angle=180, end angle=0,
      x radius=1.2, y radius=0.3];
    \draw[blue, thick, dashed] ({1.2*cos(0)}, \k) arc[start angle=0, end angle=-180,
      x radius=1.2, y radius=0.3];
  }
  % Cercle en bas
  \draw[red, thick] (5, 1.5) circle (1.2);
  \fill[red] (5, 2.7) circle (2pt) node[above] {$1$};
  \node[red] at (5, -0.2) {$S^1$};
  % Flèche p
  \draw[->, thick] (1.8, 2) -- (3.5, 2) node[midway, above] {$p$};
\end{tikzpicture}
\end{center}
\end{exemple}

\begin{exemple}[{Revêtement $S^n \to \RP^n$}]\label{ex:rev-sn-rpn}
\index{revêtement!de l'espace projectif}
L'application quotient $p \colon S^n \to \RP^n = S^n / \{x \sim -x\}$ est un revêtement
à deux feuillets. Chaque fibre est $\{x, -x\}$.
\end{exemple}

\begin{exemple}[{Revêtement $\R^2 \to T^2$}]\label{ex:rev-plan-tore}
\index{revêtement!du tore}
L'application quotient $p \colon \R^2 \to \R^2/\Z^2 \cong T^2$ est un revêtement de degré
infini. La fibre au-dessus de $[x, y]$ est le réseau $\{(x+m, y+n) \mid (m,n) \in \Z^2\}$.

\begin{center}
\begin{tikzpicture}[scale=0.7]
  % Grille
  \draw[step=1, gray!50, thin] (-0.5, -0.5) grid (4.5, 4.5);
  \foreach \i in {0,...,4} {
    \foreach \j in {0,...,4} {
      \fill[blue] (\i, \j) circle (2pt);
    }
  }
  \draw[thick, ->] (-0.5, 0) -- (4.8, 0);
  \draw[thick, ->] (0, -0.5) -- (0, 4.8);
  \draw[very thick, red, fill=red!10] (0,0) rectangle (1,1);
  \node at (0.5, 0.5) {\small$F$};
  \node at (2, -1) {$\R^2$, domaine fondamental $F$};
  % Flèche
  \draw[->, thick] (5.5, 2) -- (7.5, 2) node[midway, above] {$p$};
  % Tore schématique
  \begin{scope}[xshift=10cm, yshift=1.5cm]
    \draw[thick, red, fill=red!10] (0,0) rectangle (2, 2);
    \draw[->, thick, blue] (0, 0.5) -- (0, 1.3);
    \draw[->, thick, blue] (2, 0.5) -- (2, 1.3);
    \draw[->, thick, green!60!black] (0.5, 0) -- (1.3, 0);
    \draw[->, thick, green!60!black] (0.5, 2) -- (1.3, 2);
    \node at (1, -0.7) {$T^2$};
  \end{scope}
\end{tikzpicture}
\end{center}
\end{exemple}

\begin{exemple}[{Revêtement cyclique $S^1 \to S^1$}]\label{ex:rev-cyclique}
\index{revêtement!cyclique}
L'application $p_n \colon S^1 \to S^1$ définie par $p_n(z) = z^n$ est un revêtement
de degré $n$. La fibre au-dessus de $w \in S^1$ est l'ensemble des racines $n$-ièmes de $w$.
\end{exemple}

%---------------------------------------------------------------
\section{Relèvement des chemins}\label{sec:relevement-chemins}
%---------------------------------------------------------------

\begin{theoreme}[Relèvement des chemins]\label{thm:relevement-chemins}
\index{relèvement!des chemins}
Soit $p \colon \widetilde{X} \to X$ un revêtement et $\gamma \colon [0,1] \to X$ un chemin.
Pour tout $\tilde{x}_0 \in p^{-1}(\gamma(0))$, il existe un unique chemin
$\tilde{\gamma} \colon [0,1] \to \widetilde{X}$ tel que $p \circ \tilde{\gamma} = \gamma$
et $\tilde{\gamma}(0) = \tilde{x}_0$.
\[
\begin{tikzcd}
& \widetilde{X} \arrow[d, "p"] \\
{[0,1]} \arrow[r, "\gamma"'] \arrow[ur, dashed, "\exists!\,\tilde{\gamma}"] & X
\end{tikzcd}
\]
\end{theoreme}

\begin{proof}
\textbf{Existence.}
Par le lemme du nombre de Lebesgue\index{lemme!du nombre de Lebesgue}, il existe une
subdivision $0 = t_0 < t_1 < \cdots < t_n = 1$ telle que chaque $\gamma([t_{i-1}, t_i])$
est contenu dans un voisinage trivialisé $U_i$ de $\gamma(t_{i-1})$.

On construit $\tilde{\gamma}$ par récurrence sur $i$. Pour $i = 0$, on part de
$\tilde{x}_0$. Supposons $\tilde{\gamma}$ construit sur $[0, t_i]$, avec
$\tilde{\gamma}(t_i) = \tilde{x}_i$. L'ouvert $U_{i+1}$ est trivialisé, donc
$p^{-1}(U_{i+1}) = \bigsqcup_\alpha V_\alpha$, et $\tilde{x}_i$ appartient à un unique
$V_{\alpha_0}$. On définit $\tilde{\gamma}|_{[t_i, t_{i+1}]}$ comme
$(p|_{V_{\alpha_0}})^{-1} \circ \gamma|_{[t_i, t_{i+1}]}$. Ceci est bien défini, continu,
et prolonge $\tilde{\gamma}$ continûment car $\tilde{\gamma}(t_i) = \tilde{x}_i \in V_{\alpha_0}$.

\textbf{Unicité.}
Soient $\tilde{\gamma}$ et $\tilde{\gamma}'$ deux relèvements avec $\tilde{\gamma}(0)
= \tilde{\gamma}'(0)$. L'ensemble $S = \{t \in [0,1] \mid \tilde{\gamma}(t) = \tilde{\gamma}'(t)\}$
est fermé par continuité. Montrons qu'il est ouvert. Si $t_0 \in S$, alors
$\tilde{\gamma}(t_0) = \tilde{\gamma}'(t_0) \in V_\alpha$ pour un certain feuillet $V_\alpha$
au-dessus d'un voisinage trivialisé. Par continuité, $\tilde{\gamma}(t)$ et
$\tilde{\gamma}'(t)$ sont dans $V_\alpha$ pour $t$ proche de $t_0$. Puisque
$p|_{V_\alpha}$ est un homéomorphisme et que $p \circ \tilde{\gamma} = p \circ \tilde{\gamma}'$,
on a $\tilde{\gamma}(t) = \tilde{\gamma}'(t)$ au voisinage de $t_0$.
Comme $[0,1]$ est connexe et $0 \in S$, on a $S = [0,1]$.
\end{proof}

%---------------------------------------------------------------
\section{Relèvement des homotopies}\label{sec:relevement-homotopies}
%---------------------------------------------------------------

\begin{theoreme}[Relèvement des homotopies]\label{thm:relevement-homotopies}
\index{relèvement!des homotopies}
Soit $p \colon \widetilde{X} \to X$ un revêtement et
$H \colon [0,1]^2 \to X$ une homotopie (c'est-à-dire une application continue).
Pour tout relèvement $\tilde{H}_0 \colon [0,1] \to \widetilde{X}$ de $t \mapsto H(t, 0)$
(c'est-à-dire $p \circ \tilde{H}_0 = H(-,0)$), il existe une unique application continue
$\widetilde{H} \colon [0,1]^2 \to \widetilde{X}$ telle que $p \circ \widetilde{H} = H$
et $\widetilde{H}(t, 0) = \tilde{H}_0(t)$ pour tout $t$.
\end{theoreme}

\begin{proof}
Par le lemme du nombre de Lebesgue, il existe des subdivisions
$0 = s_0 < s_1 < \cdots < s_m = 1$ et $0 = t_0 < t_1 < \cdots < t_n = 1$ telles que
chaque rectangle $[s_{i-1}, s_i] \times [t_{j-1}, t_j]$ est envoyé par $H$ dans un voisinage
trivialisé $U_{ij}$.

On construit $\widetilde{H}$ par récurrence, en balayant les rectangles de gauche à droite,
puis de bas en haut. Pour le rectangle $R_{11} = [s_0, s_1] \times [t_0, t_1]$, on connaît
$\widetilde{H}$ sur le bord inférieur $[s_0, s_1] \times \{t_0\}$ (donné par $\tilde{H}_0$)
et sur le bord gauche $\{s_0\} \times [t_0, t_1]$ (relèvement du chemin $s \mapsto H(s_0, s)$
partant de $\tilde{H}_0(s_0)$). Puisque $H(R_{11}) \subset U_{11}$ et $p^{-1}(U_{11})$
est une union disjointe de feuillets, la continuité force $\widetilde{H}(R_{11})$ à rester
dans un unique feuillet $V_\alpha$, et on définit
$\widetilde{H}|_{R_{11}} = (p|_{V_\alpha})^{-1} \circ H|_{R_{11}}$.

On procède de même pour les rectangles suivants, en utilisant les valeurs déjà définies
sur les bords communs pour déterminer le feuillet. La compatibilité aux bords est garantie
par la continuité et la discrétion des fibres.

\textbf{Unicité.} L'argument est identique à celui du théorème~\ref{thm:relevement-chemins} :
deux relèvements coïncidant sur $[0,1] \times \{0\}$ coïncident partout, par un argument
d'ouvert-fermé dans le connexe $[0,1]^2$.
\end{proof}

\begin{corollaire}[Invariance homotopique du relèvement]\label{cor:invariance-relevement}
\index{relèvement!invariance homotopique}
Soit $p \colon \widetilde{X} \to X$ un revêtement, et soient $\gamma_0, \gamma_1 \colon [0,1]
\to X$ des chemins homotopes à extrémités fixées : $\gamma_0 \simeq \gamma_1 \rel \{0,1\}$.
Si $\tilde{\gamma}_0, \tilde{\gamma}_1$ sont les relèvements partant du même point
$\tilde{x}_0 \in p^{-1}(\gamma_0(0))$, alors
$\tilde{\gamma}_0(1) = \tilde{\gamma}_1(1)$ et
$\tilde{\gamma}_0 \simeq \tilde{\gamma}_1 \rel \{0,1\}$.
\end{corollaire}

\begin{proof}
Soit $H$ une homotopie de $\gamma_0$ à $\gamma_1$ relative à $\{0,1\}$.
Par le théorème~\ref{thm:relevement-homotopies}, $H$ se relève en
$\widetilde{H} \colon [0,1]^2 \to \widetilde{X}$ avec $\widetilde{H}(t, 0) = \tilde{\gamma}_0(t)$.
Puisque $H(0, s) = x_0$ pour tout $s$ et $\widetilde{H}(0, -) \colon [0,1] \to p^{-1}(x_0)$
est un chemin continu dans la fibre discrète $p^{-1}(x_0)$, c'est le chemin constant $\tilde{x}_0$.
De même, $s \mapsto \widetilde{H}(1, s)$ est constant, donc $\tilde{\gamma}_0(1) = \widetilde{H}(1, 0)
= \widetilde{H}(1, 1) = \tilde{\gamma}_1(1)$.

De plus, $\widetilde{H}(-, 1) = \tilde{\gamma}_1$ par unicité du relèvement, et $\widetilde{H}$
est une homotopie relative à $\{0,1\}$ de $\tilde{\gamma}_0$ à $\tilde{\gamma}_1$.
\end{proof}

\begin{corollaire}\label{cor:injectivite-p-star}
\index{revêtement!morphisme induit}
Le morphisme induit $p_* \colon \pi_1(\widetilde{X}, \tilde{x}_0) \to \pi_1(X, x_0)$
est injectif.
\end{corollaire}

\begin{proof}
Soit $[\tilde{\gamma}] \in \ker(p_*)$. Alors $p \circ \tilde{\gamma}$ est un lacet
homotope au lacet constant en $x_0$. Par le corollaire~\ref{cor:invariance-relevement},
le relèvement de cette homotopie montre que $\tilde{\gamma}$ est homotope au relèvement
du lacet constant, qui est le lacet constant en $\tilde{x}_0$.
\end{proof}

%---------------------------------------------------------------
\section{Action de monodromie}\label{sec:monodromie}
%---------------------------------------------------------------

\begin{definition}[Action de monodromie]\label{def:monodromie}
\index{monodromie!action de}
Soit $p \colon \widetilde{X} \to X$ un revêtement et $x_0 \in X$. L'\emph{action
de monodromie} est l'action à droite de $\pi_1(X, x_0)$ sur la fibre $F = p^{-1}(x_0)$
définie comme suit : pour $\tilde{x} \in F$ et $[\gamma] \in \pi_1(X, x_0)$, on pose
\[
\tilde{x} \cdot [\gamma] = \tilde{\gamma}(1)
\]
où $\tilde{\gamma}$ est l'unique relèvement de $\gamma$ partant de $\tilde{x}$.
\end{definition}

\begin{proposition}\label{prop:monodromie-bien-definie}
L'action de monodromie est bien définie (ne dépend que de la classe d'homotopie de $\gamma$)
et définit une action à droite de $\pi_1(X, x_0)$ sur $F$.
\end{proposition}

\begin{proof}
Le fait que l'action ne dépende que de $[\gamma]$ résulte du corollaire~\ref{cor:invariance-relevement}.
Vérifions que c'est une action à droite. Pour le lacet constant $c_{x_0}$, son relèvement
partant de $\tilde{x}$ est le lacet constant en $\tilde{x}$, donc $\tilde{x} \cdot [c_{x_0}] = \tilde{x}$.
Pour $[\gamma], [\delta] \in \pi_1(X, x_0)$, le relèvement de $\gamma \cdot \delta$
partant de $\tilde{x}$ est la concaténation du relèvement de $\gamma$ partant de $\tilde{x}$
(dont l'extrémité est $\tilde{x} \cdot [\gamma]$) et du relèvement de $\delta$
partant de $\tilde{x} \cdot [\gamma]$. D'où
$\tilde{x} \cdot ([\gamma][\delta]) = (\tilde{x} \cdot [\gamma]) \cdot [\delta]$.
\end{proof}

\begin{theoreme}[Transitivité de la monodromie]\label{thm:transitivite-monodromie}
\index{monodromie!transitivité}
Soit $p \colon \widetilde{X} \to X$ un revêtement avec $\widetilde{X}$ connexe par arcs.
Alors l'action de monodromie de $\pi_1(X, x_0)$ sur $F = p^{-1}(x_0)$ est \emph{transitive}.
\end{theoreme}

\begin{proof}
Soient $\tilde{x}_0, \tilde{x}_1 \in F$. Puisque $\widetilde{X}$ est connexe par arcs,
il existe un chemin $\tilde{\alpha}$ de $\tilde{x}_0$ à $\tilde{x}_1$ dans $\widetilde{X}$.
Alors $\gamma = p \circ \tilde{\alpha}$ est un lacet en $x_0$ (car $p(\tilde{x}_0) = p(\tilde{x}_1) = x_0$),
et $\tilde{x}_0 \cdot [\gamma] = \tilde{\alpha}(1) = \tilde{x}_1$.
\end{proof}

\begin{corollaire}[Stabilisateur et image de $p_*$]\label{cor:stabilisateur-image}
\index{monodromie!stabilisateur}
Soit $p \colon \widetilde{X} \to X$ un revêtement connexe avec $p(\tilde{x}_0) = x_0$.
Le stabilisateur de $\tilde{x}_0$ sous l'action de monodromie est exactement
$p_*(\pi_1(\widetilde{X}, \tilde{x}_0))$. Autrement dit :
\[
\mathrm{Stab}(\tilde{x}_0) = \set{[\gamma] \in \pi_1(X, x_0) \mid \tilde{x}_0 \cdot [\gamma] = \tilde{x}_0}
= p_*\!\big(\pi_1(\widetilde{X}, \tilde{x}_0)\big).
\]
\end{corollaire}

\begin{proof}
Le lacet $\gamma$ se relève en un lacet basé en $\tilde{x}_0$ si et seulement si
$\tilde{\gamma}(1) = \tilde{x}_0$, c'est-à-dire si et seulement si $[\gamma]$ appartient
au stabilisateur de $\tilde{x}_0$. Or $[\gamma] \in p_*(\pi_1(\widetilde{X}, \tilde{x}_0))$
signifie exactement que $\gamma = p \circ \tilde{\gamma}$ pour un lacet $\tilde{\gamma}$
basé en $\tilde{x}_0$.
\end{proof}

\begin{exemple}[Monodromie du revêtement $z \mapsto z^3$]\label{ex:monodromie-z3}
\index{monodromie!exemple}
Considérons le revêtement $p_3 \colon S^1 \to S^1$ défini par $p_3(z) = z^3$. La fibre
au-dessus de $1 \in S^1$ est $F = \{1, \omega, \omega^2\}$ où $\omega = e^{2\pi i/3}$.
Le générateur $[\gamma]$ de $\pi_1(S^1, 1) \cong \Z$ (où $\gamma(t) = e^{2\pi it}$)
agit sur $F$ par permutation cyclique :
\[
1 \cdot [\gamma] = \omega, \qquad
\omega \cdot [\gamma] = \omega^2, \qquad
\omega^2 \cdot [\gamma] = 1.
\]
L'action de monodromie est donc la permutation cyclique $(1\;\omega\;\omega^2)$.
Le stabilisateur de chaque point est $3\Z \leq \Z$, ce qui confirme que
$p_{3*}(\pi_1(S^1)) = 3\Z$.
\end{exemple}

\begin{corollaire}[Formule de l'indice]\label{cor:formule-indice}
Si $p \colon \widetilde{X} \to X$ est un revêtement connexe, alors
\[
\deg(p) = \abs{p^{-1}(x_0)} = [\pi_1(X, x_0) : p_*(\pi_1(\widetilde{X}, \tilde{x}_0))].
\]
\end{corollaire}

\begin{proof}
Par transitivité de l'action de monodromie, $F \cong \pi_1(X, x_0) / \mathrm{Stab}(\tilde{x}_0)$
en tant qu'ensembles. Le cardinal de $F$ est donc l'indice du stabilisateur.
\end{proof}

%---------------------------------------------------------------
\section{Classification des revêtements connexes}\label{sec:classification-rev}
%---------------------------------------------------------------

\begin{definition}[Morphisme de revêtements]\label{def:morphisme-rev}
\index{revêtement!morphisme de}
Soient $p_1 \colon \widetilde{X}_1 \to X$ et $p_2 \colon \widetilde{X}_2 \to X$ deux
revêtements. Un \emph{morphisme de revêtements} est une application continue
$f \colon \widetilde{X}_1 \to \widetilde{X}_2$ telle que $p_2 \circ f = p_1$ :
\[
\begin{tikzcd}
\widetilde{X}_1 \arrow[rr, "f"] \arrow[dr, "p_1"'] & & \widetilde{X}_2 \arrow[dl, "p_2"] \\
& X &
\end{tikzcd}
\]
Un tel morphisme est automatiquement un revêtement de $\widetilde{X}_2$.
\end{definition}

\begin{lemme}[Critère de relèvement]\label{lem:critere-relevement}
\index{relèvement!critère de}
Soient $p \colon \widetilde{X} \to X$ un revêtement, $Y$ un espace connexe et localement
connexe par arcs, et $f \colon (Y, y_0) \to (X, x_0)$ une application continue. Il existe
un relèvement $\tilde{f} \colon (Y, y_0) \to (\widetilde{X}, \tilde{x}_0)$ (c'est-à-dire
$p \circ \tilde{f} = f$ et $\tilde{f}(y_0) = \tilde{x}_0$) si et seulement si
\[
f_*\!\big(\pi_1(Y, y_0)\big) \subset p_*\!\big(\pi_1(\widetilde{X}, \tilde{x}_0)\big).
\]
De plus, si ce relèvement existe, il est unique.
\[
\begin{tikzcd}
& (\widetilde{X}, \tilde{x}_0) \arrow[d, "p"] \\
(Y, y_0) \arrow[r, "f"'] \arrow[ur, dashed, "\exists!\,\tilde{f}"] & (X, x_0)
\end{tikzcd}
\]
\end{lemme}

\begin{proof}
\textbf{Nécessité.} Si $\tilde{f}$ existe, alors
$f_* = (p \circ \tilde{f})_* = p_* \circ \tilde{f}_*$, d'où
$f_*(\pi_1(Y, y_0)) = p_*(\tilde{f}_*(\pi_1(Y, y_0))) \subset p_*(\pi_1(\widetilde{X}, \tilde{x}_0))$.

\textbf{Suffisance et construction.}
Pour $y \in Y$, choisissons un chemin $\alpha$ de $y_0$ à $y$ dans $Y$. Le chemin
$f \circ \alpha$ dans $X$ se relève de façon unique en $\widetilde{f \circ \alpha}$
partant de $\tilde{x}_0$. On pose $\tilde{f}(y) = \widetilde{f \circ \alpha}(1)$.

Il faut vérifier que cette définition ne dépend pas du choix de $\alpha$. Si $\beta$ est
un autre chemin de $y_0$ à $y$, alors $\alpha \cdot \beta^{-1}$ est un lacet en $y_0$
et $f_*([\alpha \cdot \beta^{-1}]) \in f_*(\pi_1(Y, y_0)) \subset p_*(\pi_1(\widetilde{X}, \tilde{x}_0))$.
Donc $f \circ (\alpha \cdot \beta^{-1})$ se relève en un lacet basé en $\tilde{x}_0$,
ce qui implique que les relèvements de $f \circ \alpha$ et $f \circ \beta$ ont la même extrémité.

La continuité de $\tilde{f}$ s'obtient par la connexité locale par arcs de $Y$ et la trivialité
locale du revêtement. L'unicité découle du même argument que pour l'unicité du relèvement des chemins.
\end{proof}

\begin{theoreme}[{Correspondance de Galois des revêtements}]\label{thm:galois-revetements}
\index{correspondance de Galois!des revêtements}
\index{classification!des revêtements}
Soit $X$ un espace connexe, localement connexe par arcs, et semi-localement simplement connexe,
avec un point base $x_0$. Il existe une bijection
\[
\left\{
\begin{array}{c}
\text{Revêtements connexes de } X \\
\text{à isomorphisme de revêtements près}
\end{array}
\right\}
\;\xleftrightarrow{\;1:1\;}\;
\left\{
\begin{array}{c}
\text{Sous-groupes de } \pi_1(X, x_0) \\
\text{à conjugaison près}
\end{array}
\right\}
\]
donnée par $\big(p \colon \widetilde{X} \to X,\; \tilde{x}_0\big) \;\longmapsto\; p_*(\pi_1(\widetilde{X}, \tilde{x}_0))$.

Si l'on fixe les points base, la correspondance devient : les revêtements connexes pointés
$(p, \tilde{x}_0)$ de $(X, x_0)$ sont en bijection avec les sous-groupes de $\pi_1(X, x_0)$.
\end{theoreme}

\begin{proof}
\textbf{Injectivité.} Supposons que deux revêtements pointés
$(p_1, \tilde{x}_0^{(1)})$ et $(p_2, \tilde{x}_0^{(2)})$ satisfassent
$H_1 = p_{1*}(\pi_1(\widetilde{X}_1, \tilde{x}_0^{(1)})) = p_{2*}(\pi_1(\widetilde{X}_2, \tilde{x}_0^{(2)})) = H_2$.
Par le lemme~\ref{lem:critere-relevement}, l'identité de $X$ se relève en
$f \colon \widetilde{X}_1 \to \widetilde{X}_2$ (car $H_1 \subset H_2$) et
$g \colon \widetilde{X}_2 \to \widetilde{X}_1$ (car $H_2 \subset H_1$).
Les composées $g \circ f$ et $f \circ g$ sont des relèvements de l'identité, donc
égales à l'identité par unicité. Donc $f$ est un isomorphisme de revêtements.

\textbf{Surjectivité.} Soit $H \leq \pi_1(X, x_0)$. On construit
$\widetilde{X}_H$ comme suit. Considérons l'ensemble des chemins de $x_0$ dans $X$,
muni de la relation d'équivalence : $\gamma \sim \gamma'$ si $\gamma(1) = \gamma'(1)$ et
$[\gamma \cdot (\gamma')^{-1}] \in H$. On pose $\widetilde{X}_H = \{\text{classes d'équivalence}\}$
et $p([\gamma]) = \gamma(1)$. Le point base est $\tilde{x}_0 = [c_{x_0}]$ (classe du chemin constant).

On munit $\widetilde{X}_H$ d'une topologie : pour tout ouvert $U \subset X$ connexe par arcs
et tout chemin $\gamma$ de $x_0$ à un point de $U$, on définit
$\widetilde{U}_{[\gamma]} = \set{[\gamma \cdot \alpha] \mid \alpha \text{ chemin dans } U}$.
Ces ensembles forment une base de topologie, et on vérifie que $p$ est un revêtement
avec $p_*(\pi_1(\widetilde{X}_H, \tilde{x}_0)) = H$.

Le fait que $X$ soit semi-localement simplement connexe garantit que $\widetilde{X}_H$
est bien un revêtement (la trivialisation locale est assurée car les lacets locaux se
relèvent en lacets).
\end{proof}

Sous la correspondance de Galois, les propriétés des revêtements se traduisent en propriétés
des sous-groupes :

\begin{center}
\renewcommand{\arraystretch}{1.4}
\begin{tabular}{ll}
\toprule
\textbf{Revêtement} & \textbf{Sous-groupe} \\
\midrule
Degré du revêtement & Indice du sous-groupe \\
Revêtement connexe & --- \\
Revêtement universel & Sous-groupe trivial $\{e\}$ \\
Revêtement régulier (normal) & Sous-groupe normal \\
$\widetilde{X}_1 \to \widetilde{X}_2$ existe & $H_1 \leq H_2$ \\
\bottomrule
\end{tabular}
\end{center}

\begin{remarque}[Analogie galoisienne]\label{rem:analogie-galoisienne}
\index{correspondance de Galois!analogie avec la théorie de Galois}
La correspondance ci-dessus est l'analogue topologique de la correspondance de Galois en théorie
des corps : les extensions intermédiaires d'une extension galoisienne $L/K$ correspondent aux
sous-groupes du groupe de Galois $\mathrm{Gal}(L/K)$, les extensions normales correspondant aux
sous-groupes normaux. Le revêtement universel joue le rôle de la clôture algébrique.
\end{remarque}

\begin{remarque}[Point de vue fonctoriel]\label{rem:foncteur-fibre}
\index{foncteur fibre}
La correspondance de Galois s'interprète en termes de la catégorie
$\mathsf{Rev}(X)$ des revêtements de $X$. Le \emph{foncteur fibre}
$\mathrm{Fib}_{x_0} \colon \mathsf{Rev}(X) \to \mathsf{Ens}$, qui à un revêtement $p$
associe la fibre $p^{-1}(x_0)$, est une équivalence entre $\mathsf{Rev}(X)$ et la catégorie
des $\pi_1(X, x_0)$-ensembles (ensembles munis d'une action de $\pi_1(X, x_0)$).
Les revêtements connexes correspondent aux actions transitives, et les classes
d'isomorphisme de revêtements connexes aux classes de conjugaison de sous-groupes de $\pi_1(X, x_0)$.
\end{remarque}

%---------------------------------------------------------------
\section{Revêtement universel}\label{sec:rev-universel}
%---------------------------------------------------------------

\begin{definition}[Espace semi-localement simplement connexe]\label{def:semi-loc-simp-connexe}
\index{semi-localement simplement connexe}
Un espace $X$ est \emph{semi-localement simplement connexe} si tout point $x \in X$ admet
un voisinage ouvert $U$ tel que le morphisme $\pi_1(U, x) \to \pi_1(X, x)$ induit par
l'inclusion est trivial. Cela signifie que tout lacet dans $U$ est contractile \emph{dans $X$}
(pas nécessairement dans $U$).
\end{definition}

\begin{remarque}\label{rem:localement-vs-semi}
Localement simplement connexe implique semi-localement simplement connexe, mais la réciproque est fausse.
Le peigne hawaïen (réunion de cercles de rayon $1/n$ tangents à l'origine) n'est pas semi-localement
simplement connexe à l'origine.\index{peigne hawaïen}
\end{remarque}

\begin{definition}[Revêtement universel]\label{def:rev-universel}
\index{revêtement!universel}
Un revêtement $p \colon \widetilde{X} \to X$ est dit \emph{universel} si $\widetilde{X}$
est simplement connexe ($\pi_1(\widetilde{X}) = \{e\}$).
\end{definition}

\begin{theoreme}[Existence et unicité du revêtement universel]\label{thm:rev-universel}
\index{revêtement!universel!existence}
Soit $X$ un espace connexe, localement connexe par arcs, et semi-localement simplement connexe.
Alors $X$ admet un revêtement universel $p \colon \widetilde{X} \to X$, unique à isomorphisme
de revêtements près.
\end{theoreme}

\begin{proof}
Par le théorème~\ref{thm:galois-revetements}, le sous-groupe trivial $H = \{e\} \leq \pi_1(X, x_0)$
correspond à un revêtement $\widetilde{X}$ dont le groupe fondamental vérifie
$p_*(\pi_1(\widetilde{X}, \tilde{x}_0)) = \{e\}$. Puisque $p_*$ est injective
(corollaire~\ref{cor:injectivite-p-star}), on a $\pi_1(\widetilde{X}, \tilde{x}_0) = \{e\}$,
donc $\widetilde{X}$ est simplement connexe. L'unicité découle de l'injectivité de la
correspondance de Galois.
\end{proof}

\begin{proposition}[Propriété universelle du revêtement universel]\label{prop:propriete-universelle-rev}
\index{revêtement!universel!propriété universelle}
Le revêtement universel $\widetilde{X}$ est initial dans la catégorie des revêtements connexes
de $X$ : pour tout revêtement connexe $q \colon Y \to X$ et tout choix de points base
$\tilde{x}_0 \in \widetilde{X}$, $\tilde{y}_0 \in Y$ avec $p(\tilde{x}_0) = q(\tilde{y}_0)$,
il existe un unique morphisme de revêtements $f \colon \widetilde{X} \to Y$ avec
$f(\tilde{x}_0) = \tilde{y}_0$.
\[
\begin{tikzcd}
\widetilde{X} \arrow[rr, dashed, "f"] \arrow[dr, "p"'] & & Y \arrow[dl, "q"] \\
& X &
\end{tikzcd}
\]
\end{proposition}

\begin{proof}
On a $p_*(\pi_1(\widetilde{X}, \tilde{x}_0)) = \{e\} \subset q_*(\pi_1(Y, \tilde{y}_0))$.
Par le lemme~\ref{lem:critere-relevement}, $p$ se relève en $f \colon \widetilde{X} \to Y$,
et ce relèvement est unique.
\end{proof}

\begin{exemple}[Revêtements universels classiques]\label{ex:rev-universels}
\index{revêtement!universel!exemples}
\begin{enumerate}[label=(\roman*)]
\item $\R \to S^1$ est le revêtement universel de $S^1$ (car $\R$ est contractile).
\item $S^n \to \RP^n$ est le revêtement universel de $\RP^n$ pour $n \geq 2$
(car $\pi_1(S^n) = 0$ pour $n \geq 2$).
\item $\R^2 \to T^2$ est le revêtement universel du tore.
\item Le revêtement universel du bouquet $S^1 \vee S^1$ est l'arbre de Cayley\index{arbre de Cayley}
du groupe libre $F_2$ : un arbre infini $4$-régulier.
\item $\mathrm{SU}(2) \cong S^3 \to \mathrm{SO}(3) \cong \RP^3$ est un revêtement
universel à deux feuillets.
\end{enumerate}
\end{exemple}

%---------------------------------------------------------------
\section{Transformations de revêtement}\label{sec:deck-transformations}
%---------------------------------------------------------------

\begin{definition}[Transformation de revêtement]\label{def:deck}
\index{transformation de revêtement}\index{deck transformation|see{transformation de revêtement}}
Soit $p \colon \widetilde{X} \to X$ un revêtement. Une \emph{transformation de revêtement}
(ou \emph{deck transformation}, ou \emph{automorphisme de revêtement}) est un homéomorphisme
$\varphi \colon \widetilde{X} \to \widetilde{X}$ tel que $p \circ \varphi = p$.
\[
\begin{tikzcd}
\widetilde{X} \arrow[rr, "\varphi"] \arrow[dr, "p"'] & & \widetilde{X} \arrow[dl, "p"] \\
& X &
\end{tikzcd}
\]
L'ensemble des transformations de revêtement forme un groupe pour la composition, noté
$\Aut(\widetilde{X}/X)$ ou $\mathrm{Deck}(\widetilde{X}/X)$.
\end{definition}

\begin{proposition}\label{prop:deck-libre}
\index{transformation de revêtement!action libre}
Le groupe $\Aut(\widetilde{X}/X)$ agit librement sur chaque fibre $p^{-1}(x)$.
En particulier, toute transformation de revêtement non triviale est sans point fixe.
\end{proposition}

\begin{proof}
Si $\varphi \in \Aut(\widetilde{X}/X)$ fixe un point $\tilde{x}$, alors $\varphi$ et
$\Id_{\widetilde{X}}$ sont deux relèvements de $p$ coïncidant en $\tilde{x}$.
Par unicité des relèvements (pour $Y = \widetilde{X}$ connexe dans le
lemme~\ref{lem:critere-relevement}), $\varphi = \Id$.
\end{proof}

\begin{theoreme}\label{thm:deck-normaliser}
\index{transformation de revêtement!et sous-groupes}
Soit $p \colon \widetilde{X} \to X$ un revêtement connexe avec
$H = p_*(\pi_1(\widetilde{X}, \tilde{x}_0)) \leq \pi_1(X, x_0)$. Alors
\[
\Aut(\widetilde{X}/X) \;\cong\; N(H)/H
\]
où $N(H) = \set{g \in \pi_1(X, x_0) \mid gHg^{-1} = H}$ est le normalisateur de $H$.
\end{theoreme}

\begin{proof}
Soit $\varphi \in \Aut(\widetilde{X}/X)$. Posons $\tilde{x}_1 = \varphi(\tilde{x}_0)$.
Comme $p(\tilde{x}_1) = x_0$, il existe un chemin $\tilde{\alpha}$ de $\tilde{x}_0$ à $\tilde{x}_1$
dans $\widetilde{X}$. Posons $\gamma = p \circ \tilde{\alpha}$, un lacet en $x_0$.

On vérifie que la classe $[\gamma] \in \pi_1(X, x_0)$ appartient à $N(H)$. En effet,
$p_*(\pi_1(\widetilde{X}, \tilde{x}_1)) = [\gamma] \cdot H \cdot [\gamma]^{-1}$
(changement de point base dans le revêtement). Mais $\varphi$ induit un isomorphisme
$\pi_1(\widetilde{X}, \tilde{x}_0) \xrightarrow{\sim} \pi_1(\widetilde{X}, \tilde{x}_1)$,
donc $p_*(\pi_1(\widetilde{X}, \tilde{x}_1)) = H$ aussi, d'où $[\gamma] H [\gamma]^{-1} = H$.

L'application $\varphi \mapsto [\gamma] H \in N(H)/H$ est bien définie
(indépendante du choix de $\tilde{\alpha}$), et c'est un isomorphisme de groupes. L'injectivité
provient de la liberté de l'action sur la fibre. La surjectivité utilise le critère de relèvement :
si $[\gamma] \in N(H)$, le point $\tilde{x}_1 = \tilde{x}_0 \cdot [\gamma]$ vérifie
$p_*(\pi_1(\widetilde{X}, \tilde{x}_1)) = H$, et le lemme~\ref{lem:critere-relevement}
fournit un automorphisme $\varphi$ envoyant $\tilde{x}_0$ sur $\tilde{x}_1$.
\end{proof}

%---------------------------------------------------------------
\section{Revêtements réguliers}\label{sec:rev-reguliers}
%---------------------------------------------------------------

\begin{definition}[Revêtement régulier]\label{def:rev-regulier}
\index{revêtement!régulier (normal)}
Un revêtement connexe $p \colon \widetilde{X} \to X$ est dit \emph{régulier}
(ou \emph{normal}, ou \emph{galoisien}) si $p_*(\pi_1(\widetilde{X}, \tilde{x}_0))$
est un sous-groupe normal de $\pi_1(X, x_0)$.
\end{definition}

\begin{proposition}[Caractérisations des revêtements réguliers]\label{prop:caract-regulier}
\index{revêtement!régulier!caractérisations}
Soit $p \colon \widetilde{X} \to X$ un revêtement connexe. Les conditions suivantes
sont équivalentes :
\begin{enumerate}[label=(\roman*)]
\item $p$ est régulier ;
\item $\Aut(\widetilde{X}/X)$ agit transitivement sur chaque fibre ;
\item $\Aut(\widetilde{X}/X)$ agit transitivement sur $p^{-1}(x_0)$ ;
\item pour tout lacet $\gamma$ en $x_0$, soit tous les relèvements de $\gamma$ sont des lacets,
soit aucun ne l'est.
\end{enumerate}
\end{proposition}

\begin{proof}
$\mathrm{(i) \Leftrightarrow (ii) \Leftrightarrow (iii)}$ :
Par le théorème~\ref{thm:deck-normaliser}, $\Aut(\widetilde{X}/X) \cong N(H)/H$ où $H = p_*(\pi_1(\widetilde{X}, \tilde{x}_0))$.
Le revêtement est régulier si et seulement si $N(H) = \pi_1(X, x_0)$, auquel cas
$\Aut(\widetilde{X}/X) \cong \pi_1(X, x_0)/H$. L'action de $\Aut(\widetilde{X}/X)$ sur
la fibre est libre et de cardinal $[G:H] = \abs{F}$ si et seulement si elle est transitive.

$\mathrm{(i) \Leftrightarrow (iv)}$ :
Le relèvement de $\gamma$ partant de $\tilde{x}$ est un lacet si et seulement si
$[\gamma] \in p_*(\pi_1(\widetilde{X}, \tilde{x}))$. Quand on change de point de départ
$\tilde{x}$, le sous-groupe $p_*(\pi_1(\widetilde{X}, \tilde{x}))$ est conjugué de $H$.
La condition (iv) signifie que tous les conjugués de $H$ coïncident, c'est-à-dire que $H$ est normal.
\end{proof}

\begin{theoreme}[Revêtement universel et groupe fondamental]\label{thm:rev-universel-pi1}
\index{revêtement!universel!et groupe fondamental}
Soit $X$ un espace connexe, localement connexe par arcs, et semi-localement simplement connexe,
et soit $p \colon \widetilde{X} \to X$ son revêtement universel. Alors :
\begin{enumerate}[label=(\roman*)]
\item Le revêtement universel est régulier.
\item $\Aut(\widetilde{X}/X) \cong \pi_1(X, x_0)$.
\item $X \cong \widetilde{X} / \Aut(\widetilde{X}/X)$, où l'action est proprement discontinue.
\end{enumerate}
\end{theoreme}

\begin{proof}
\begin{enumerate}[label=(\roman*)]
\item Le sous-groupe $H = p_*(\pi_1(\widetilde{X}, \tilde{x}_0)) = \{e\}$ est normal
dans tout groupe.
\item Par le théorème~\ref{thm:deck-normaliser},
$\Aut(\widetilde{X}/X) \cong N(\{e\})/\{e\} = \pi_1(X, x_0)$.
\item L'action de $\Aut(\widetilde{X}/X)$ sur $\widetilde{X}$ est libre et proprement
discontinue (car le revêtement est un homéomorphisme local). Le quotient est $X$.
\qedhere
\end{enumerate}
\end{proof}

\begin{corollaire}[Quotient par action de groupe]\label{cor:quotient-action}
\index{action de groupe!et revêtement}
Si un groupe $G$ agit librement et proprement discontinûment sur un espace $Y$
connexe et simplement connexe (localement connexe par arcs), alors la projection
$Y \to Y/G$ est le revêtement universel et $\pi_1(Y/G) \cong G$.
\end{corollaire}

\begin{exemple}\label{ex:applications-quotient}
\begin{enumerate}[label=(\roman*)]
\item $\Z$ agit sur $\R$ par translations : $\R / \Z \cong S^1$, d'où $\pi_1(S^1) \cong \Z$.
\item $\Z^2$ agit sur $\R^2$ par translations : $\R^2 / \Z^2 \cong T^2$, d'où $\pi_1(T^2) \cong \Z^2$.
\item $\Z/2\Z$ agit sur $S^n$ par l'antipodie : $S^n / (\Z/2\Z) \cong \RP^n$,
d'où $\pi_1(\RP^n) \cong \Z/2\Z$ pour $n \geq 2$.
\end{enumerate}
\end{exemple}

\begin{theoreme}[Correspondance pour les revêtements intermédiaires]\label{thm:rev-intermediaire}
\index{revêtement!intermédiaire}
Soit $p \colon \widetilde{X} \to X$ le revêtement universel. Pour tout sous-groupe
$H \leq \pi_1(X, x_0) \cong \Aut(\widetilde{X}/X)$, l'espace $\widetilde{X}/H$ est
un revêtement connexe de $X$, et
\[
\pi_1(\widetilde{X}/H) \cong H.
\]
Le revêtement $\widetilde{X}/H \to X$ est régulier si et seulement si $H \trianglelefteq \pi_1(X, x_0)$,
auquel cas
\[
\Aut((\widetilde{X}/H) / X) \cong \pi_1(X, x_0) / H.
\]
\end{theoreme}

\begin{proof}
Puisque $H$ agit librement et proprement discontinûment sur $\widetilde{X}$ (en tant
que sous-groupe du groupe des deck transformations), la projection
$\widetilde{X} \to \widetilde{X}/H$ est un revêtement. La composée
$\widetilde{X} \to \widetilde{X}/H \to X$ est $p$, donc $\widetilde{X}/H \to X$
est un revêtement. Le revêtement $\widetilde{X} \to \widetilde{X}/H$ est universel
(car $\widetilde{X}$ est simplement connexe), donc
$\pi_1(\widetilde{X}/H) \cong \Aut(\widetilde{X}/(\widetilde{X}/H)) = H$.

Le revêtement est régulier si et seulement si $H$ est normal dans $\pi_1(X, x_0)$,
par le théorème~\ref{thm:deck-normaliser}, et dans ce cas
$\Aut((\widetilde{X}/H)/X) \cong N(H)/H = \pi_1(X,x_0)/H$.
\end{proof}

%---------------------------------------------------------------
\section{Applications complémentaires}\label{sec:applications-rev}
%---------------------------------------------------------------

\begin{proposition}[Revêtements du cercle]\label{prop:rev-cercle}
\index{revêtement!du cercle}
Les revêtements connexes de $S^1$ sont (à isomorphisme près) :
\begin{enumerate}[label=(\roman*)]
\item $p_n \colon S^1 \to S^1$, $z \mapsto z^n$, pour $n \geq 1$ (degré $n$,
correspondant au sous-groupe $n\Z \leq \Z$) ;
\item $p \colon \R \to S^1$, $t \mapsto e^{2\pi it}$ (revêtement universel,
correspondant au sous-groupe $\{0\} \leq \Z$).
\end{enumerate}
\end{proposition}

\begin{proof}
Les sous-groupes de $\Z$ sont $n\Z$ pour $n \geq 0$. Par la correspondance de Galois,
$n\Z$ avec $n \geq 1$ donne un revêtement de degré $[\Z : n\Z] = n$, qui est $z \mapsto z^n$.
Le sous-groupe trivial $\{0\}$ donne le revêtement universel $\R \to S^1$.
\end{proof}

\begin{proposition}[Revêtements du tore]\label{prop:rev-tore}
\index{revêtement!du tore}
Les revêtements connexes de $T^2$ correspondent aux sous-groupes de $\Z^2$, qui sont
de la forme $H \cong \Z^k$ avec $k = 0, 1, 2$ (car les sous-groupes de $\Z^2$ sont libres
de rang $\leq 2$). Les cas sont :
\begin{enumerate}[label=(\roman*)]
\item $k=0$ : le revêtement universel $\R^2 \to T^2$ ;
\item $k=1$ : $H \cong \Z$, donnant un revêtement par un cylindre $\R \times S^1 \to T^2$ ;
\item $k=2$ : $H \cong \Z^2$ d'indice fini $n$ dans $\Z^2$, donnant un revêtement $T^2 \to T^2$
de degré $n$.
\end{enumerate}
\end{proposition}

\begin{theoreme}[Borsuk--Ulam en dimension $2$]\label{thm:borsuk-ulam-2}
\index{théorème!de Borsuk--Ulam}
Pour toute application continue $f \colon S^2 \to \R^2$, il existe $x \in S^2$ tel que
$f(x) = f(-x)$.
\end{theoreme}

\begin{proof}
Supposons par l'absurde que $f(x) \neq f(-x)$ pour tout $x$. Définissons
$g \colon S^2 \to S^1$ par $g(x) = \frac{f(x) - f(-x)}{\|f(x) - f(-x)\|}$.
Alors $g(-x) = -g(x)$ pour tout $x$.

La restriction de $g$ à l'équateur $S^1 \subset S^2$ donne une application
$h \colon S^1 \to S^1$ vérifiant $h(-z) = -h(z)$. Cette application a un degré impair
(car elle est équivariante pour l'action antipodale), en particulier de degré non nul.

Mais $h$ se prolonge au disque supérieur $D^2_+$ (hémisphère nord de $S^2$), donc $h$ est
homotopiquement triviale. Ceci contredit le fait que son degré est non nul.
\end{proof}

%---------------------------------------------------------------
\section{Exercices}\label{sec:exo-rev}
%---------------------------------------------------------------

\begin{exercice}\label{exo:rev-degre-2}
Montrer que tout revêtement de degré $2$ est régulier.
\textbf{Difficulté : \(\bigstar\bigstar\)}
\end{exercice}

\begin{exercice}\label{exo:rev-sous-groupe}
Soit $p \colon \widetilde{X} \to X$ un revêtement avec $\widetilde{X}$ connexe par arcs.
Montrer que si $\pi_1(X, x_0)$ est abélien, alors tout revêtement connexe de $X$ est régulier.
\textbf{Difficulté : \(\bigstar\bigstar\)}
\end{exercice}

\begin{exercice}\label{exo:relevement-applications}
Soit $p \colon \R \to S^1$ le revêtement exponentiel. Montrer qu'une application continue
$f \colon S^1 \to S^1$ se relève en $\tilde{f} \colon S^1 \to \R$ si et seulement si
$f$ est homotope à une application constante.
\textbf{Difficulté : \(\bigstar\bigstar\)}
\end{exercice}

\begin{exercice}\label{exo:rev-bouquet}
Décrire (par un dessin) le revêtement de $S^1 \vee S^1$ correspondant au sous-groupe
$\langle aba^{-1}b^{-1} \rangle$ de $F_2 = \langle a, b \rangle$. Quel est le degré de ce
revêtement ? Est-il régulier ?
\textbf{Difficulté : \(\bigstar\bigstar\bigstar\)}
\end{exercice}

\begin{exercice}\label{exo:pi1-so3}
En utilisant le revêtement $\mathrm{SU}(2) \to \mathrm{SO}(3)$, montrer que
$\pi_1(\mathrm{SO}(3)) \cong \Z/2\Z$.
\textbf{Difficulté : \(\bigstar\bigstar\)}
\end{exercice}

\begin{exercice}\label{exo:deck-action-propre}
\index{action!proprement discontinue}
Soit $G$ un groupe agissant par homéomorphismes sur un espace topologique $Y$.
\begin{enumerate}[label=(\alph*)]
\item Montrer que si l'action est libre et proprement discontinue, alors la projection
$Y \to Y/G$ est un revêtement.
\item Montrer que la réciproque est vraie : si $Y \to Y/G$ est un revêtement, alors l'action
est libre et proprement discontinue.
\end{enumerate}
\textbf{Difficulté : \(\bigstar\bigstar\bigstar\)}
\end{exercice}

\begin{exercice}\label{exo:rev-surface}
Soit $\Sigma_g$ la surface orientable de genre $g \geq 1$.
\begin{enumerate}[label=(\alph*)]
\item Montrer que $\Sigma_g$ admet un revêtement connexe de degré $n$ pour tout $n \geq 1$.
\item Pour $g = 1$ (le tore), décrire explicitement les revêtements connexes de degré $2$
et montrer qu'il y en a exactement $3$ à isomorphisme près.
\item Montrer que le revêtement universel de $\Sigma_g$ (pour $g \geq 2$) est homéomorphe
au plan $\R^2$. \emph{Indication : utiliser le théorème d'uniformisation.}
\end{enumerate}
\textbf{Difficulté : \(\bigstar\bigstar\bigstar\)}
\end{exercice}

\begin{exercice}\label{exo:functorialite-rev}
\index{fonctorialité!des revêtements}
Soient $X, Y$ des espaces admettant des revêtements universels $\widetilde{X}, \widetilde{Y}$.
Montrer que toute application continue $f \colon X \to Y$ se relève en
$\tilde{f} \colon \widetilde{X} \to \widetilde{Y}$ (non nécessairement de façon unique).
Montrer que le diagramme suivant commute :
\[
\begin{tikzcd}
\widetilde{X} \arrow[r, "\tilde{f}"] \arrow[d, "p_X"'] & \widetilde{Y} \arrow[d, "p_Y"] \\
X \arrow[r, "f"'] & Y
\end{tikzcd}
\]
En déduire que $f$ est $\pi_1$-compatible :
$\tilde{f} \circ \varphi = f_*(\varphi) \circ \tilde{f}$ pour tout $\varphi \in \Aut(\widetilde{X}/X) \cong \pi_1(X)$,
où l'on identifie $f_*(\varphi)$ à l'élément correspondant de $\Aut(\widetilde{Y}/Y)$.
\textbf{Difficulté : \(\bigstar\bigstar\bigstar\)}
\end{exercice}

\begin{exercice}\label{exo:rev-produit}
Soient $p_1 \colon \widetilde{X}_1 \to X_1$ et $p_2 \colon \widetilde{X}_2 \to X_2$ des revêtements.
Montrer que $p_1 \times p_2 \colon \widetilde{X}_1 \times \widetilde{X}_2 \to X_1 \times X_2$ est un revêtement.
En déduire que si $\widetilde{X}_1$ et $\widetilde{X}_2$ sont les revêtements universels,
alors $\widetilde{X}_1 \times \widetilde{X}_2$ est le revêtement universel de $X_1 \times X_2$.
\textbf{Difficulté : \(\bigstar\bigstar\)}
\end{exercice}

%---------------------------------------------------------------
\section*{Résumé du chapitre}\label{sec:resume-ch4}
\addcontentsline{toc}{section}{Résumé du chapitre}
%---------------------------------------------------------------

\begin{center}
\renewcommand{\arraystretch}{1.4}
\begin{tabular}{p{5cm} p{9cm}}
\toprule
\textbf{Concept} & \textbf{Résultat principal} \\
\midrule
Revêtement $p \colon \widetilde{X} \to X$ &
Homéomorphisme local, fibres discrètes de cardinal constant. \\
Relèvement des chemins &
Existence et unicité du relèvement, invariance par homotopie. \\
Injectivité de $p_*$ &
$p_* \colon \pi_1(\widetilde{X}, \tilde{x}_0) \hookrightarrow \pi_1(X, x_0)$. \\
Action de monodromie &
$\pi_1(X, x_0) \curvearrowright p^{-1}(x_0)$, transitive si $\widetilde{X}$ connexe. \\
Correspondance de Galois &
Revêtements connexes $\leftrightarrow$ sous-groupes de $\pi_1$ (à conjugaison près). \\
Revêtement universel &
$\widetilde{X}$ simplement connexe; $\Aut(\widetilde{X}/X) \cong \pi_1(X, x_0)$. \\
Revêtement régulier &
Sous-groupe normal; $\Aut(\widetilde{X}/X) \cong \pi_1(X)/H$. \\
Quotient par action &
$G \curvearrowright Y$ libre proprement discontinue $\Rightarrow$ $Y \to Y/G$ revêtement. \\
\bottomrule
\end{tabular}
\end{center}

% === FIN DU CHUNK ch3-4 ===
% === DEBUT DU CHUNK ch5-6 ===

%%%%%%%%%%%%%%%%%%%%%%%%%%%%%%%%%%%%%%%%%%%%%%%%%%%%%%%%%%%%%%%%%%%%%
%% CHAPITRE 5 : HOMOLOGIE SINGULIÈRE — DÉFINITION ET PROPRIÉTÉS   %%
%%%%%%%%%%%%%%%%%%%%%%%%%%%%%%%%%%%%%%%%%%%%%%%%%%%%%%%%%%%%%%%%%%%%%

\chapter{Homologie singulière --- Définition et propriétés}\label{chap:homologie-singuliere}
\minitoc\vspace{1cm}

\noindent
L'homologie singulière constitue l'un des invariants algébriques les plus puissants de la topologie.
Contrairement au groupe fondamental, qui ne capture que des phénomènes essentiellement unidimensionnels,
l'homologie singulière associe à chaque espace topologique $X$ une suite de groupes abéliens
$H_n(X)$, $n \geq 0$, qui détectent les «~trous~» de toute dimension.
L'idée maîtresse, remontant à Poincaré et systématisée par Eilenberg et Steenrod, consiste à
sonder l'espace $X$ au moyen de simplexes continus, puis à extraire l'information
homologique par un procédé algébrique canonique.\index{homologie singulière}

% ===================================================================
\section{Le simplexe standard et les simplexes singuliers}\label{sec:simplexe-standard}
% ===================================================================

\begin{definition}[Simplexe standard]\label{def:simplexe-standard}
\index{simplexe standard}
Pour tout entier $n \geq 0$, le \emph{$n$-simplexe standard} est le sous-espace de $\R^{n+1}$
défini par
\[
  \Delta^n = \set{(t_0, t_1, \ldots, t_n) \in \R^{n+1} \;\Big|\; t_i \geq 0 \text{ pour tout } i, \;\sum_{i=0}^{n} t_i = 1}.
\]
On note $e_0, e_1, \ldots, e_n$ les sommets de $\Delta^n$, où $e_i$ est le $(i+1)$-ième vecteur de
la base canonique de $\R^{n+1}$.
\end{definition}

\begin{remarque}\label{rem:simplexes-bas}
En petite dimension : $\Delta^0$ est un point, $\Delta^1$ est un segment, $\Delta^2$ est un
triangle plein, $\Delta^3$ est un tétraèdre plein.
\end{remarque}

\begin{exemple}[Coordonnées barycentriques]\label{ex:coordonnees-barycentriques}
\index{coordonnées barycentriques}
Les coordonnées $(t_0, \ldots, t_n)$ d'un point de $\Delta^n$ sont ses \emph{coordonnées
barycentriques} par rapport aux sommets $e_0, \ldots, e_n$. Le point de coordonnées
$(t_0, \ldots, t_n)$ est le barycentre $\sum_{i=0}^n t_i e_i$, et la condition $\sum t_i = 1$
avec $t_i \geq 0$ exprime que ce point est une combinaison convexe des sommets.
Par exemple, le barycentre (centre de gravité) de $\Delta^n$ a pour coordonnées
$(\frac{1}{n+1}, \ldots, \frac{1}{n+1})$.
\end{exemple}

\begin{proposition}[Relations entre applications face]\label{prop:relations-face}
\index{application face!relations}
Pour $i < j$, on a la relation de commutation :
\[
  \delta_j \circ \delta_i = \delta_i \circ \delta_{j-1} \colon \Delta^{n-2} \longrightarrow \Delta^n.
\]
Ces relations, dites \emph{relations cosimpliciales}\index{relations cosimpliciales}, jouent un
rôle essentiel dans la preuve que $\partial^2 = 0$.
\end{proposition}

\begin{proof}
Les deux membres sont des applications affines $\Delta^{n-2} \to \Delta^n$, et il suffit de
vérifier qu'elles coïncident sur les sommets. Pour $e_k \in \Delta^{n-2}$ :
\begin{itemize}
  \item Si $k < i$ : $\delta_j(\delta_i(e_k)) = \delta_j(e_k) = e_k$ (car $k < i < j$)
  et $\delta_i(\delta_{j-1}(e_k)) = \delta_i(e_k) = e_k$ (car $k < i \leq j-1$).
  \item Si $i \leq k < j-1$ : $\delta_j(\delta_i(e_k)) = \delta_j(e_{k+1}) = e_{k+1}$ (car $k+1 < j$)
  et $\delta_i(\delta_{j-1}(e_k)) = \delta_i(e_{k+1}) = e_{k+1}$ (car $k \geq i$, puis $k+1 \leq j-1$, donc pas de décalage par $\delta_i$ ni $\delta_{j-1}$ au-delà).

  Un examen cas par cas complet sur les sommets confirme l'égalité.
\end{itemize}
\end{proof}

\begin{definition}[Applications face]\label{def:applications-face}
\index{application face}
Pour $0 \leq i \leq n$, la \emph{$i$-ième application face} est l'inclusion affine
\[
  \delta_i \colon \Delta^{n-1} \longrightarrow \Delta^n
\]
définie par $\delta_i(e_j) = e_j$ si $j < i$ et $\delta_i(e_j) = e_{j+1}$ si $j \geq i$.
Autrement dit, $\delta_i$ est l'unique application affine qui envoie $\Delta^{n-1}$ sur la face
de $\Delta^n$ opposée au sommet $e_i$.
\end{definition}

\begin{definition}[Simplexe singulier]\label{def:simplexe-singulier}
\index{simplexe singulier}
Soit $X$ un espace topologique. Un \emph{$n$-simplexe singulier} dans $X$ est une application
continue
\[
  \sigma \colon \Delta^n \longrightarrow X.
\]
On note $\operatorname{Sing}_n(X)$ l'ensemble de tous les $n$-simplexes singuliers dans $X$.
\end{definition}

\begin{remarque}\label{rem:singulier-pas-injectif}
Le terme «~singulier~» souligne qu'aucune condition d'injectivité n'est imposée : l'application
$\sigma$ peut écraser $\Delta^n$ sur un point, ou identifier des faces entre elles. Cette liberté
totale est la clef de la fonctorialité et de la puissance de la théorie.
\end{remarque}

\begin{exemple}[Simplexes singuliers en petite dimension]\label{ex:simplexes-petite-dim}
\begin{enumerate}[label=(\roman*)]
  \item Un $0$-simplexe singulier dans $X$ est simplement un point de $X$ (une application
  $\Delta^0 = \{e_0\} \to X$).
  \item Un $1$-simplexe singulier $\sigma \colon \Delta^1 \to X$ est un chemin dans $X$,
  d'origine $\sigma(e_0)$ et d'extrémité $\sigma(e_1)$. La face $\sigma \circ \delta_0 = \sigma(e_1)$
  est le point terminal et $\sigma \circ \delta_1 = \sigma(e_0)$ est le point initial.
  \item Un $2$-simplexe singulier $\sigma \colon \Delta^2 \to X$ est une application continue
  d'un triangle plein dans $X$, dont les trois côtés sont les $1$-simplexes singuliers
  $\sigma \circ \delta_0$, $\sigma \circ \delta_1$, $\sigma \circ \delta_2$.
\end{enumerate}
\end{exemple}

\begin{definition}[Applications de dégénérescence]\label{def:degenerescence}
\index{application de dégénérescence}
Pour $0 \leq i \leq n$, l'\emph{application de dégénérescence} $s_i \colon \Delta^{n+1} \to \Delta^n$
est la surjection affine définie par $s_i(e_j) = e_j$ si $j \leq i$ et $s_i(e_j) = e_{j-1}$ si $j > i$.
Intuitivement, $s_i$ «~écrase~» l'arête $[e_i, e_{i+1}]$ en le sommet $e_i$.
Un simplexe singulier de la forme $\sigma \circ s_i$ est dit \emph{dégénéré}\index{simplexe dégénéré}.
\end{definition}

% ===================================================================
\section{Chaînes singulières et opérateur bord}\label{sec:chaines-singulieres}
% ===================================================================

\begin{definition}[Groupe des chaînes singulières]\label{def:chaines-singulieres}
\index{chaîne singulière}\index{groupe des chaînes}
Pour $n \geq 0$, le \emph{groupe des $n$-chaînes singulières} de $X$ est le groupe abélien libre
engendré par l'ensemble $\operatorname{Sing}_n(X)$ :
\[
  C_n(X) = \Z[\operatorname{Sing}_n(X)] = \bigoplus_{\sigma \in \operatorname{Sing}_n(X)} \Z \cdot \sigma.
\]
Un élément de $C_n(X)$ est donc une combinaison linéaire formelle finie
$\sum_{i} n_i \sigma_i$ avec $n_i \in \Z$ et $\sigma_i \colon \Delta^n \to X$ continues.
Pour $n < 0$, on pose $C_n(X) = 0$.
\end{definition}

\begin{remarque}\label{rem:chaines-enormes}
Le groupe $C_n(X)$ est en général énormément grand, même pour des espaces simples : dès que
$X$ n'est pas discret, l'ensemble $\operatorname{Sing}_n(X)$ est de cardinal au moins celui du
continu pour $n \geq 1$. Ce n'est pas un obstacle, car l'information essentielle est contenue
dans le quotient $H_n(X) = \ker \partial_n / \im \partial_{n+1}$, qui s'avère de taille
raisonnable pour les espaces «~ordinaires~».
\end{remarque}

\begin{exemple}[Le bord d'un $1$-simplexe]\label{ex:bord-1-simplexe}
Soit $\sigma \colon \Delta^1 \to X$ un chemin. Son bord est
\[
  \partial_1(\sigma) = \sigma \circ \delta_0 - \sigma \circ \delta_1 = \sigma(e_1) - \sigma(e_0),
\]
c'est-à-dire le point terminal moins le point initial (vus comme $0$-chaînes).
\end{exemple}

\begin{exemple}[Le bord d'un $2$-simplexe]\label{ex:bord-2-simplexe}
Soit $\sigma \colon \Delta^2 \to X$ un $2$-simplexe singulier. Son bord est
\[
  \partial_2(\sigma) = \sigma \circ \delta_0 - \sigma \circ \delta_1 + \sigma \circ \delta_2.
\]
Géométriquement, c'est la somme alternée des trois côtés du triangle : le côté opposé à $e_0$
(de $e_1$ à $e_2$), moins le côté opposé à $e_1$ (de $e_0$ à $e_2$), plus le côté opposé à
$e_2$ (de $e_0$ à $e_1$). Les signes sont choisis pour que ces chemins «~s'enchaînent~»
de manière cohérente, formant un cycle parcouru dans le sens positif.
\end{exemple}

\begin{definition}[Opérateur bord]\label{def:operateur-bord}
\index{opérateur bord}\index{bord!opérateur}
L'\emph{opérateur bord} $\partial_n \colon C_n(X) \to C_{n-1}(X)$ est l'unique homomorphisme de
groupes abéliens défini sur les générateurs par
\[
  \partial_n(\sigma) = \sum_{i=0}^{n} (-1)^i \, \sigma \circ \delta_i,
\]
où $\delta_i \colon \Delta^{n-1} \to \Delta^n$ est la $i$-ième application face. On étend $\partial_n$
par linéarité à tout $C_n(X)$.
\end{definition}

\begin{theoreme}[$\partial^2 = 0$]\label{thm:bord-carre-nul}
\index{bord!carré nul}
Pour tout $n \geq 1$, on a $\partial_{n-1} \circ \partial_n = 0$, c'est-à-dire
$\partial^2 = 0$.
\end{theoreme}

\begin{proof}
Soit $\sigma \colon \Delta^n \to X$ un $n$-simplexe singulier. On calcule :
\begin{align*}
  \partial_{n-1}(\partial_n(\sigma))
  &= \partial_{n-1}\!\left(\sum_{i=0}^{n} (-1)^i\, \sigma \circ \delta_i\right) \\
  &= \sum_{i=0}^{n} (-1)^i \sum_{j=0}^{n-1} (-1)^j\, \sigma \circ \delta_i \circ \delta_j.
\end{align*}
La relation fondamentale entre les applications face est $\delta_i \circ \delta_j = \delta_{j+1} \circ \delta_i$
pour $i \leq j$. En effet, ces deux applications affines coïncident sur les sommets de $\Delta^{n-2}$.

Décomposons la somme double en deux parties selon que $j \geq i$ ou $j < i$.

Pour les termes avec $j \geq i$, on applique $\delta_i \circ \delta_j = \delta_{j+1} \circ \delta_i$ :
\[
  \sum_{0 \leq i \leq j \leq n-1} (-1)^{i+j}\, \sigma \circ \delta_{j+1} \circ \delta_i.
\]
En posant $k = j+1$ (de sorte que $k$ va de $i+1$ à $n$) :
\[
  \sum_{0 \leq i < k \leq n} (-1)^{i+k-1}\, \sigma \circ \delta_k \circ \delta_i.
\]
Pour les termes avec $j < i$, on a directement :
\[
  \sum_{0 \leq j < i \leq n} (-1)^{i+j}\, \sigma \circ \delta_i \circ \delta_j.
\]
En renommant $(i,j)$ en $(k,i)$ dans cette dernière somme (c'est-à-dire $k = i$, $i = j$), on obtient :
\[
  \sum_{0 \leq i < k \leq n} (-1)^{i+k}\, \sigma \circ \delta_k \circ \delta_i.
\]
Les deux expressions sont opposées (elles diffèrent par un facteur $(-1)^{-1} = -1$), donc leur somme est nulle.
Par linéarité, on conclut $\partial_{n-1} \circ \partial_n = 0$.
\end{proof}

% ===================================================================
\section{Le complexe de chaînes singulières}\label{sec:complexe-chaines}
% ===================================================================

\begin{definition}[Complexe de chaînes]\label{def:complexe-chaines}
\index{complexe de chaînes}
Le \emph{complexe de chaînes singulières} de $X$ est la suite de groupes abéliens et d'homomorphismes
\[
  \cdots \xrightarrow{\partial_{n+1}} C_n(X) \xrightarrow{\partial_n} C_{n-1}(X) \xrightarrow{\partial_{n-1}} \cdots
  \xrightarrow{\partial_1} C_0(X) \xrightarrow{\partial_0} 0,
\]
que l'on note $(C_*(X), \partial)$. La condition $\partial^2 = 0$ assure que $\im \partial_{n+1} \subset \ker \partial_n$
pour tout $n$.
\end{definition}

On représente schématiquement un complexe de chaînes par le diagramme :
\[
\begin{tikzcd}
  \cdots \ar[r, "\partial_{n+2}"] & C_{n+1}(X) \ar[r, "\partial_{n+1}"] & C_n(X) \ar[r, "\partial_n"]
  & C_{n-1}(X) \ar[r, "\partial_{n-1}"] & \cdots \ar[r, "\partial_1"] & C_0(X) \ar[r] & 0
\end{tikzcd}
\]

% ===================================================================
\section{Cycles, bords et groupes d'homologie}\label{sec:groupes-homologie}
% ===================================================================

\begin{definition}[Cycles et bords]\label{def:cycles-bords}
\index{cycle}\index{bord}
Soit $(C_*(X), \partial)$ le complexe de chaînes singulières de $X$. Pour tout $n \geq 0$ :
\begin{enumerate}[label=(\roman*)]
  \item Le groupe des \emph{$n$-cycles} est $Z_n(X) = \ker(\partial_n \colon C_n(X) \to C_{n-1}(X))$.
  \item Le groupe des \emph{$n$-bords} est $B_n(X) = \im(\partial_{n+1} \colon C_{n+1}(X) \to C_n(X))$.
  \item Puisque $\partial^2 = 0$, on a $B_n(X) \subset Z_n(X)$.
\end{enumerate}
\end{definition}

\begin{definition}[Groupes d'homologie singulière]\label{def:groupes-homologie}
\index{homologie singulière!groupes}
Le \emph{$n$-ième groupe d'homologie singulière} de $X$ est le groupe quotient
\[
  H_n(X) = \frac{Z_n(X)}{B_n(X)} = \frac{\ker \partial_n}{\im \partial_{n+1}}.
\]
La classe d'un cycle $z \in Z_n(X)$ dans $H_n(X)$ est notée $[z]$ et s'appelle la
\emph{classe d'homologie}\index{classe d'homologie} de $z$. Deux cycles $z, z'$ sont
\emph{homologues}\index{homologue} si $[z] = [z']$, c'est-à-dire si $z - z' \in B_n(X)$.
\end{definition}

\begin{remarque}\label{rem:homologie-mesure-trous}
Intuitivement, $H_n(X)$ mesure les «~trous $n$-dimensionnels~» de $X$ : un $n$-cycle est une
chaîne fermée (sans bord), et il est trivial en homologie s'il borde une $(n+1)$-chaîne.
Le rang de $H_n(X)$ (en tant que groupe abélien) s'appelle le \emph{$n$-ième nombre de Betti}
\index{nombre de Betti} $\beta_n(X)$.
\end{remarque}

% ===================================================================
\section{Fonctorialité de l'homologie}\label{sec:fonctorialite}
% ===================================================================

\begin{definition}[Morphisme de chaînes induit]\label{def:morphisme-chaines-induit}
\index{morphisme induit}
Soit $f \colon X \to Y$ une application continue. Pour tout $n$-simplexe singulier
$\sigma \colon \Delta^n \to X$, on pose $f_\sharp(\sigma) = f \circ \sigma \colon \Delta^n \to Y$.
On étend par linéarité pour obtenir un homomorphisme
\[
  f_\sharp \colon C_n(X) \longrightarrow C_n(Y).
\]
\end{definition}

\begin{proposition}[Morphisme de complexes de chaînes]\label{prop:morphisme-complexes}
Le morphisme $f_\sharp$ commute avec les opérateurs bord, c'est-à-dire
$\partial_n \circ f_\sharp = f_\sharp \circ \partial_n$ pour tout $n$. Autrement dit, $f_\sharp$
est un \emph{morphisme de complexes de chaînes}\index{morphisme de complexes de chaînes}.
\end{proposition}

\begin{proof}
Pour un générateur $\sigma \colon \Delta^n \to X$ :
\[
  \partial_n(f_\sharp(\sigma)) = \partial_n(f \circ \sigma) = \sum_{i=0}^{n} (-1)^i (f \circ \sigma) \circ \delta_i
  = \sum_{i=0}^{n} (-1)^i f \circ (\sigma \circ \delta_i) = f_\sharp\!\left(\sum_{i=0}^{n} (-1)^i \sigma \circ \delta_i\right)
  = f_\sharp(\partial_n(\sigma)).
\]
Le résultat s'étend par linéarité.
\end{proof}

Le diagramme suivant commute :
\[
\begin{tikzcd}
  \cdots \ar[r] & C_n(X) \ar[r, "\partial_n"] \ar[d, "f_\sharp"] & C_{n-1}(X) \ar[r] \ar[d, "f_\sharp"] & \cdots \\
  \cdots \ar[r] & C_n(Y) \ar[r, "\partial_n"'] & C_{n-1}(Y) \ar[r] & \cdots
\end{tikzcd}
\]

\begin{corollaire}[Morphisme induit en homologie]\label{cor:morphisme-induit-homologie}
Toute application continue $f \colon X \to Y$ induit un homomorphisme
\[
  f_* \colon H_n(X) \longrightarrow H_n(Y), \qquad [z] \longmapsto [f_\sharp(z)],
\]
bien défini pour tout $n \geq 0$.
\end{corollaire}

\begin{proof}
Le morphisme $f_\sharp$ envoie cycles sur cycles (car il commute au bord) et bords sur bords
(car $f_\sharp(\partial c) = \partial(f_\sharp(c))$), donc il passe au quotient.
\end{proof}

\begin{theoreme}[Fonctorialité de l'homologie]\label{thm:fonctorialite-homologie}
\index{fonctorialité!homologie}
L'homologie singulière définit un foncteur covariant
\[
  H_n \colon \mathbf{Top} \longrightarrow \mathbf{Ab}
\]
de la catégorie des espaces topologiques vers la catégorie des groupes abéliens, pour chaque $n \geq 0$.
Précisément :
\begin{enumerate}[label=(\roman*)]
  \item $(\Id_X)_* = \Id_{H_n(X)}$.
  \item Si $f \colon X \to Y$ et $g \colon Y \to Z$ sont continues, alors $(g \circ f)_* = g_* \circ f_*$.
\end{enumerate}
\end{theoreme}

\begin{proof}
(i) est immédiat car $(\Id_X)_\sharp = \Id_{C_n(X)}$.
(ii) découle de $(g \circ f)_\sharp(\sigma) = g \circ f \circ \sigma = g_\sharp(f_\sharp(\sigma))$,
donc $(g \circ f)_\sharp = g_\sharp \circ f_\sharp$, et le passage au quotient est compatible.
\end{proof}

\begin{corollaire}\label{cor:homeo-iso-homologie}
Si $f \colon X \to Y$ est un homéomorphisme, alors $f_* \colon H_n(X) \to H_n(Y)$ est un
isomorphisme pour tout $n$.
\end{corollaire}

\begin{proof}
Si $g = f^{-1}$, alors $g_* \circ f_* = (g \circ f)_* = (\Id_X)_* = \Id$ et de même
$f_* \circ g_* = \Id$.
\end{proof}

% ===================================================================
\section{Invariance homotopique}\label{sec:invariance-homotopique}
% ===================================================================

Le résultat fondamental suivant renforce considérablement le 
ef{cor:homeo-iso-homologie} : non seulement
les homéomorphismes, mais même les équivalences d'homotopie, induisent des isomorphismes en homologie.

\begin{theoreme}[Invariance homotopique de l'homologie]\label{thm:invariance-homotopique}
\index{invariance homotopique}\index{homologie singulière!invariance homotopique}
Si $f, g \colon X \to Y$ sont deux applications continues homotopes ($f \simeq g$), alors
$f_* = g_* \colon H_n(X) \to H_n(Y)$ pour tout $n \geq 0$.

En conséquence, si $X$ et $Y$ ont le même type d'homotopie, alors $H_n(X) \cong H_n(Y)$ pour tout $n$.
\end{theoreme}

La preuve repose sur la construction d'un \emph{opérateur de prisme}, dont voici le mécanisme.

\begin{definition}[Opérateur de prisme]\label{def:operateur-prisme}
\index{opérateur de prisme}\index{homotopie de chaînes}
Soit $F \colon X \times [0,1] \to Y$ une homotopie entre $f = F(-,0)$ et $g = F(-,1)$.
On définit un homomorphisme $P_n \colon C_n(X) \to C_{n+1}(Y)$, appelé \emph{opérateur de prisme},
de la manière suivante. Pour un $n$-simplexe singulier $\sigma \colon \Delta^n \to X$,
le «~prisme~» $\Delta^n \times [0,1]$ se décompose en $(n+1)$-simplexes. Plus précisément,
on pose
\[
  P_n(\sigma) = \sum_{i=0}^{n} (-1)^i F \circ (\sigma \times \Id) \circ p_i,
\]
où $p_i \colon \Delta^{n+1} \to \Delta^n \times [0,1]$ est l'application affine définie par
\[
  p_i(e_j) = \begin{cases}
    (e_j, 0) & \text{si } j \leq i, \\
    (e_{j-1}, 1) & \text{si } j > i.
  \end{cases}
\]
\end{definition}

\begin{lemme}[Relation fondamentale du prisme]\label{lem:relation-prisme}
\index{relation du prisme}
L'opérateur de prisme satisfait la \emph{relation d'homotopie de chaînes} :
\[
  \partial_{n+1} \circ P_n + P_{n-1} \circ \partial_n = g_\sharp - f_\sharp \colon C_n(X) \to C_n(Y).
\]
\end{lemme}

\begin{proof}
On vérifie la relation sur un générateur $\sigma \colon \Delta^n \to X$. Le calcul de
$\partial_{n+1}(P_n(\sigma))$ fait apparaître, par développement, des termes correspondant aux
faces du prisme $\Delta^n \times [0,1]$. Les faces «~verticales~» (celles qui n'impliquent ni
$t=0$ ni $t=1$) se simplifient en paires de signes opposés avec les termes de $P_{n-1}(\partial_n(\sigma))$.
Il ne subsiste que les faces «~supérieure~» ($t=1$) et «~inférieure~» ($t=0$), qui donnent
respectivement $g_\sharp(\sigma)$ et $-f_\sharp(\sigma)$.

En détail, écrivons $\partial_{n+1}(P_n(\sigma)) = \sum_{i=0}^{n}\sum_{k=0}^{n+1} (-1)^{i+k} F \circ (\sigma \times \Id) \circ p_i \circ \delta_k$.
L'analyse des cas $k \leq i$ et $k > i$ combinée aux identités entre les $p_i$ et les $\delta_k$
fournit, après télescopage, exactement $g_\sharp(\sigma) - f_\sharp(\sigma) - P_{n-1}(\partial_n(\sigma))$.
\end{proof}

\begin{proof}[Preuve du 
ef{thm:invariance-homotopique}]
Soit $[z] \in H_n(X)$ avec $z \in Z_n(X)$, c'est-à-dire $\partial_n(z) = 0$. Alors :
\[
  g_\sharp(z) - f_\sharp(z) = \partial_{n+1}(P_n(z)) + P_{n-1}(\partial_n(z)) = \partial_{n+1}(P_n(z)).
\]
Ainsi $g_\sharp(z) - f_\sharp(z) \in B_n(Y)$, d'où $[g_\sharp(z)] = [f_\sharp(z)]$, c'est-à-dire $g_*[z] = f_*[z]$.
\end{proof}

\begin{corollaire}\label{cor:contractile-homologie}
\index{espace contractile}
Si $X$ est contractile, alors $H_n(X) = 0$ pour tout $n \geq 1$ et $H_0(X) \cong \Z$.
\end{corollaire}

\begin{proof}
Un espace contractile $X$ est homotopiquement équivalent à un point $\{*\}$.
Par invariance homotopique, $H_n(X) \cong H_n(\{*\})$ pour tout $n$. D'après le

ef{thm:homologie-point}, $H_0(\{*\}) = \Z$ et $H_n(\{*\}) = 0$ pour $n \geq 1$.
\end{proof}

\begin{corollaire}\label{cor:equivalence-homotopie-iso}
Si $f \colon X \to Y$ est une équivalence d'homotopie, alors $f_* \colon H_n(X) \to H_n(Y)$ est
un isomorphisme pour tout $n \geq 0$.
\end{corollaire}

\begin{proof}
Soit $g \colon Y \to X$ une inverse homotopique de $f$, c'est-à-dire $g \circ f \simeq \Id_X$
et $f \circ g \simeq \Id_Y$. Par le 
ef{thm:invariance-homotopique},
$g_* \circ f_* = (g \circ f)_* = (\Id_X)_* = \Id_{H_n(X)}$ et de même
$f_* \circ g_* = \Id_{H_n(Y)}$. Donc $f_*$ est un isomorphisme d'inverse $g_*$.
\end{proof}

\begin{exemple}[Espaces contractiles classiques]\label{ex:espaces-contractiles}
\index{espace contractile!exemples}
Les espaces suivants sont contractiles et ont donc la même homologie qu'un point :
\begin{enumerate}[label=(\roman*)]
  \item L'espace euclidien $\R^n$ (homotopie de contraction $F(x,t) = (1-t)x$).
  \item Le disque fermé $D^n$ (même argument).
  \item Tout sous-espace convexe de $\R^n$.
  \item Le cône $CX = (X \times [0,1])/(X \times \{1\})$ sur un espace quelconque $X$.
\end{enumerate}
\end{exemple}

\begin{remarque}[Homotopie de chaînes et algèbre homologique]\label{rem:homotopie-algebrique}
\index{homotopie de chaînes}
La notion d'homotopie de chaînes (un opérateur $P$ tel que $\partial P + P \partial = g_\sharp - f_\sharp$)
est purement algébrique et ne fait intervenir aucune topologie. Elle joue un rôle central en
algèbre homologique, bien au-delà du cadre topologique. On dit que $P$ est une
\emph{homotopie de chaînes} entre $f_\sharp$ et $g_\sharp$, et que ces deux morphismes de
complexes sont \emph{homotopes}\index{morphismes homotopes} au sens de l'algèbre homologique.
Le résultat clef est qu'une homotopie de chaînes implique l'égalité des morphismes induits en
homologie.
\end{remarque}

% ===================================================================
\section{Homologie réduite}\label{sec:homologie-reduite}
% ===================================================================

\begin{definition}[Morphisme d'augmentation]\label{def:augmentation}
\index{augmentation}\index{morphisme d'augmentation}
Le \emph{morphisme d'augmentation} est l'homomorphisme $\varepsilon \colon C_0(X) \to \Z$ défini par
$\varepsilon(\sigma) = 1$ pour tout $0$-simplexe singulier $\sigma$ et étendu par linéarité. On a
$\varepsilon \circ \partial_1 = 0$, de sorte que le \emph{complexe augmenté} est
\[
\begin{tikzcd}
  \cdots \ar[r, "\partial_2"] & C_1(X) \ar[r, "\partial_1"] & C_0(X) \ar[r, "\varepsilon"] & \Z \ar[r] & 0.
\end{tikzcd}
\]
\end{definition}

\begin{definition}[Homologie réduite]\label{def:homologie-reduite}
\index{homologie réduite}
Les \emph{groupes d'homologie réduite} de $X$ (pour $X \neq \varnothing$) sont les groupes
d'homologie du complexe augmenté :
\[
  \widetilde{H}_n(X) = \begin{cases}
    H_n(X) & \text{si } n \geq 1, \\
    \ker \varepsilon / \im \partial_1 & \text{si } n = 0.
  \end{cases}
\]
On a une suite exacte courte
$0 \to \widetilde{H}_0(X) \to H_0(X) \xrightarrow{\varepsilon_*} \Z \to 0$,
de sorte que $H_0(X) \cong \widetilde{H}_0(X) \oplus \Z$.
\end{definition}

\begin{remarque}\label{rem:homologie-reduite-motivation}
La motivation de l'homologie réduite est d'éliminer la «~copie parasite~» de $\Z$ en degré $0$
due au fait que tout espace non vide a $H_0 \neq 0$. L'homologie réduite est plus naturelle
pour formuler certains résultats, notamment la suspension et le théorème de Mayer--Vietoris.
En particulier, $\widetilde{H}_n(\{*\}) = 0$ pour tout $n$.
\end{remarque}

\begin{remarque}[Homologie réduite via le complexe augmenté]\label{rem:reduite-augmente}
\index{complexe augmenté}
On peut voir le complexe augmenté comme un complexe de chaînes indexé en degré $\geq -1$,
avec $C_{-1}(X) = \Z$ et $\partial_0 = \varepsilon$. Les groupes d'homologie de ce complexe
étendu sont exactement les $\widetilde{H}_n(X)$ pour $n \geq 0$, et $\widetilde{H}_{-1}(X) = 0$
pour tout espace non vide.
\end{remarque}

\begin{proposition}\label{prop:H0-composantes}
\index{composantes connexes par arcs}
Si $X$ possède $k$ composantes connexes par arcs $X_1, \ldots, X_k$, alors
$H_0(X) \cong \Z^k$ et $\widetilde{H}_0(X) \cong \Z^{k-1}$.
\end{proposition}

\begin{proof}
Deux $0$-simplexes $\sigma, \tau \colon \Delta^0 \to X$ (c'est-à-dire deux points de $X$) sont
homologues si et seulement s'ils sont dans la même composante connexe par arcs, car un chemin
$\gamma$ de $\sigma$ à $\tau$ fournit un $1$-simplexe $\gamma$ avec $\partial_1(\gamma) = \tau - \sigma$.
Ainsi $H_0(X)$ est le groupe abélien libre engendré par les composantes connexes par arcs.
\end{proof}

% ===================================================================
\section{Calculs fondamentaux}\label{sec:calculs-fondamentaux}
% ===================================================================

\begin{theoreme}[Homologie d'un point]\label{thm:homologie-point}
\index{homologie!d'un point}
Soit $\{*\}$ l'espace réduit à un point. Alors
\[
  H_n(\{*\}) = \begin{cases} \Z & \text{si } n = 0, \\ 0 & \text{si } n \geq 1. \end{cases}
\]
\end{theoreme}

\begin{proof}
Pour chaque $n \geq 0$, il existe un unique $n$-simplexe singulier $\sigma_n \colon \Delta^n \to \{*\}$
(l'application constante). Ainsi $C_n(\{*\}) \cong \Z$ pour tout $n \geq 0$, engendré par $\sigma_n$.
L'opérateur bord donne
\[
  \partial_n(\sigma_n) = \sum_{i=0}^{n} (-1)^i \sigma_{n-1} = \begin{cases} 0 & \text{si $n$ est pair,} \\ \sigma_{n-1} & \text{si $n$ est impair.} \end{cases}
\]
Le complexe de chaînes est donc
\[
\begin{tikzcd}
  \cdots \ar[r, "0"] & \Z \ar[r, "\Id"] & \Z \ar[r, "0"] & \Z \ar[r, "\Id"] & \Z \ar[r, "0"] & \Z \ar[r] & 0,
\end{tikzcd}
\]
d'où $H_0 = \Z/0 = \Z$ et $H_n = 0$ pour $n \geq 1$ (car les morphismes alternent entre $0$ et $\Id$).
\end{proof}

\begin{theoreme}[Homologie des sphères]\label{thm:homologie-spheres}
\index{homologie!des sphères}\index{sphère!homologie}
Pour tout $n \geq 0$, on a
\[
  H_k(S^n) \cong \begin{cases} \Z & \text{si } k = 0 \text{ ou } k = n, \\ 0 & \text{sinon.} \end{cases}
\]
En homologie réduite : $\widetilde{H}_k(S^n) \cong \Z$ si $k = n$, et $0$ sinon.
\end{theoreme}

Nous donnons la preuve par récurrence, en utilisant la suite exacte de Mayer--Vietoris
\index{Mayer--Vietoris}
(démontrée au 
ef{chap:suites-exactes}).

\begin{proof}
\textbf{Initialisation.} $S^0 = \{-1, +1\}$ a deux composantes connexes par arcs, donc
$H_0(S^0) \cong \Z^2$ et $H_k(S^0) = 0$ pour $k \geq 1$. En homologie réduite, $\widetilde{H}_0(S^0) \cong \Z$.

\textbf{Hérédité.} Supposons le résultat pour $S^{n-1}$. Décomposons $S^n$ en deux hémisphères
ouverts $U = S^n \setminus \{N\}$ et $V = S^n \setminus \{S\}$ (pôle nord et pôle sud).
Les deux ouverts sont contractiles (homéomorphes à $\R^n$ par projection stéréographique),
et $U \cap V$ est homotopiquement équivalent à $S^{n-1}$ (c'est l'équateur épaissi).

Par la suite de Mayer--Vietoris (
ef{thm:mayer-vietoris-mention}) :
\[
\begin{tikzcd}
  \cdots \ar[r] & \widetilde{H}_k(U) \oplus \widetilde{H}_k(V) \ar[r] & \widetilde{H}_k(S^n) \ar[r, "\partial_*"] & \widetilde{H}_{k-1}(U \cap V) \ar[r] & \cdots
\end{tikzcd}
\]
Puisque $\widetilde{H}_k(U) = \widetilde{H}_k(V) = 0$ pour tout $k$, la suite exacte fournit
des isomorphismes $\widetilde{H}_k(S^n) \cong \widetilde{H}_{k-1}(U \cap V) \cong \widetilde{H}_{k-1}(S^{n-1})$
pour tout $k$. L'hypothèse de récurrence donne le résultat.
\end{proof}

% ===================================================================
\section{Homologie relative}\label{sec:homologie-relative}
% ===================================================================

\begin{definition}[Chaînes relatives]\label{def:chaines-relatives}
\index{homologie relative}\index{chaînes relatives}
Soit $A \subset X$ un sous-espace. L'inclusion $\iota \colon A \hookrightarrow X$ induit une
injection $\iota_\sharp \colon C_n(A) \hookrightarrow C_n(X)$. Le \emph{groupe des $n$-chaînes relatives}
est le quotient
\[
  C_n(X, A) = C_n(X) / C_n(A).
\]
\end{definition}

\begin{proposition}\label{prop:complexe-relatif}
L'opérateur bord $\partial_n$ de $C_*(X)$ passe au quotient en un opérateur bord
$\bar{\partial}_n \colon C_n(X, A) \to C_{n-1}(X, A)$ satisfaisant $\bar{\partial}^2 = 0$.
\end{proposition}

\begin{proof}
Comme $\partial_n(C_n(A)) \subset C_{n-1}(A)$ (car l'inclusion $A \hookrightarrow X$ est un
morphisme de complexes), l'opérateur bord passe au quotient. La relation $\bar{\partial}^2 = 0$
découle de $\partial^2 = 0$ par passage au quotient.
\end{proof}

\begin{definition}[Groupes d'homologie relative]\label{def:homologie-relative}
\index{homologie relative!groupes}
Les \emph{groupes d'homologie relative} de la paire $(X, A)$ sont
\[
  H_n(X, A) = \ker \bar{\partial}_n / \im \bar{\partial}_{n+1}, \qquad n \geq 0.
\]
Un élément de $H_n(X, A)$ est représenté par une chaîne $c \in C_n(X)$ dont le bord
$\partial_n(c)$ est dans $C_{n-1}(A)$ (un «~cycle relatif~»), modulo les chaînes de la forme
$\partial_{n+1}(d) + a$ avec $d \in C_{n+1}(X)$ et $a \in C_n(A)$.
\end{definition}

\begin{remarque}\label{rem:paire-ses}
La suite exacte courte de complexes de chaînes
\[
\begin{tikzcd}
  0 \ar[r] & C_*(A) \ar[r, "\iota_\sharp"] & C_*(X) \ar[r, "\pi"] & C_*(X,A) \ar[r] & 0
\end{tikzcd}
\]
engendre, par le théorème algébrique de la suite exacte longue (
ef{thm:suite-exacte-longue}),
la suite exacte longue de la paire $(X, A)$ (
ef{sec:suite-exacte-paire}).
\end{remarque}

\begin{proposition}[Fonctorialité de l'homologie relative]\label{prop:fonctorialite-relative}
\index{fonctorialité!homologie relative}
Si $f \colon (X, A) \to (Y, B)$ est une application de paires (c'est-à-dire $f \colon X \to Y$
continue avec $f(A) \subset B$), alors $f$ induit un homomorphisme
$f_* \colon H_n(X, A) \to H_n(Y, B)$ pour tout $n$, et l'homologie relative est fonctorielle :
$(\Id)_* = \Id$ et $(g \circ f)_* = g_* \circ f_*$.
\end{proposition}

\begin{proof}
L'application $f_\sharp \colon C_n(X) \to C_n(Y)$ envoie $C_n(A)$ dans $C_n(B)$ (car $f(A) \subset B$),
donc passe au quotient en $\bar{f}_\sharp \colon C_n(X,A) \to C_n(Y,B)$. Ce morphisme commute
aux bords et induit donc un morphisme en homologie. Les propriétés de fonctorialité se vérifient
par passage au quotient.
\end{proof}

\begin{exemple}[Homologie relative de $(\R^n, \R^n \setminus \{0\})$]\label{ex:homologie-relative-Rn}
En utilisant l'excision et le fait que $(\R^n, \R^n \setminus \{0\})$ est une bonne paire,
on obtient (cf.\ 
ef{thm:relative-quotient}) :
\[
  H_k(\R^n, \R^n \setminus \{0\}) \cong \widetilde{H}_k(\R^n / (\R^n \setminus \{0\})) \cong \widetilde{H}_k(S^n) \cong \begin{cases} \Z & \text{si } k = n, \\ 0 & \text{sinon.} \end{cases}
\]
Ce groupe est appelé le \emph{groupe d'orientation locale}\index{orientation locale} de $\R^n$ en l'origine.
\end{exemple}

% ===================================================================
\section{Commentaires historiques et perspectives}\label{sec:perspectives-ch5}
% ===================================================================

\begin{remarque}[Perspective historique]\label{rem:histoire-homologie}
\index{homologie!histoire}
La notion d'homologie a été introduite par Henri Poincaré dans son \emph{Analysis Situs} (1895),
sous une forme géométrique assez différente de la présentation algébrique moderne. Poincaré
travaillait avec des variétés triangulées et définissait les «~cycles~» comme des combinaisons
de cellules sans bord. La formulation en termes de simplexes singuliers est due à Samuel Eilenberg
(1944), qui a montré que cette approche produit les mêmes groupes d'homologie que les définitions
simpliciale et cellulaire pour les espaces raisonnables (CW-complexes), tout en étant
définie pour \emph{tout} espace topologique.

La caractérisation axiomatique de l'homologie par Eilenberg et Steenrod (1952) a ensuite montré
que toute théorie d'homologie satisfaisant un petit nombre d'axiomes (fonctorialité, invariance
homotopique, excision, axiome de la dimension) coïncide avec l'homologie singulière sur les
CW-complexes finis.\index{Eilenberg--Steenrod}
\end{remarque}

\begin{remarque}[Homologie à coefficients]\label{rem:homologie-coefficients}
\index{homologie!à coefficients}
La construction s'étend immédiatement à des coefficients dans un groupe abélien $G$ quelconque :
on pose $C_n(X; G) = C_n(X) \otimes_\Z G$ et on étend $\partial$ par $\partial \otimes \Id_G$.
Les groupes d'homologie $H_n(X; G)$ sont particulièrement utiles pour $G = \Z/p\Z$
(homologie modulo $p$), $G = \Q$ (homologie rationnelle) et $G = \R$ (lié à la cohomologie
de de Rham pour les variétés).
\end{remarque}

% ===================================================================
\section{Exercices}\label{sec:exo-ch5}
% ===================================================================

\begin{exercice}[$\star\star$]\label{exo:bord-explicite}
Soit $X = \R^2$ et $\sigma \colon \Delta^2 \to \R^2$ le $2$-simplexe singulier défini par
$\sigma(t_0, t_1, t_2) = (t_1, t_2)$. Calculer explicitement $\partial_2(\sigma)$ comme
combinaison de $1$-simplexes singuliers. Vérifier que $\partial_1(\partial_2(\sigma)) = 0$.
\end{exercice}

\begin{exercice}[$\star\star$]\label{exo:H0-connexe}
Montrer directement (sans utiliser la 
ef{prop:H0-composantes}) que si $X$ est connexe par arcs,
alors $H_0(X) \cong \Z$. \emph{Indication : montrer que l'augmentation $\varepsilon \colon C_0(X) \to \Z$
induit un isomorphisme $H_0(X) \xrightarrow{\sim} \Z$.}
\end{exercice}

\begin{exercice}[$\star\star$]\label{exo:homologie-espace-discret}
Soit $X$ un ensemble discret de cardinal $m$. Calculer $H_n(X)$ pour tout $n \geq 0$.
\end{exercice}

\begin{exercice}[$\star\star$]\label{exo:fonctorialite-retraction}
Soit $A \subset X$ un rétracte, c'est-à-dire qu'il existe une application continue $r \colon X \to A$
telle que $r \circ \iota = \Id_A$ où $\iota \colon A \hookrightarrow X$.
Montrer que $\iota_* \colon H_n(A) \to H_n(X)$ est injective pour tout $n$.
\end{exercice}

\begin{exercice}[$\star\star\star$]\label{exo:homotopie-chaines}
Soient $f_\sharp, g_\sharp \colon C_*(X) \to C_*(Y)$ deux morphismes de complexes de chaînes.
Montrer que s'il existe une \emph{homotopie de chaînes} $P_n \colon C_n(X) \to C_{n+1}(Y)$
satisfaisant $\partial_{n+1} P_n + P_{n-1} \partial_n = g_\sharp - f_\sharp$ pour tout $n$,
alors $f_*= g_* \colon H_n(X) \to H_n(Y)$.
\end{exercice}

\begin{exercice}[$\star\star\star$]\label{exo:homologie-produit}
Soient $X$ et $Y$ deux espaces topologiques. Montrer que
$H_0(X \times Y) \cong H_0(X) \otimes_\Z H_0(Y)$.
\emph{Indication : relier les composantes connexes par arcs de $X \times Y$ à celles de $X$ et $Y$.}
\end{exercice}

\begin{exercice}[$\star\star$]\label{exo:homologie-reduite-point}
Montrer que $\widetilde{H}_n(\{*\}) = 0$ pour tout $n \geq 0$.
\end{exercice}

\begin{exercice}[$\star\star\star$]\label{exo:suite-exacte-augmentation}
Soit $x_0 \in X$. Montrer qu'il existe une suite exacte courte
\[
  0 \longrightarrow \widetilde{H}_n(X) \longrightarrow H_n(X) \longrightarrow H_n(\{x_0\}) \longrightarrow 0
\]
pour tout $n \geq 0$, et que cette suite est scindée. En déduire que $H_n(X) \cong \widetilde{H}_n(X) \oplus H_n(\{x_0\})$.
\end{exercice}

\begin{exercice}[$\star\star\star$]\label{exo:prisme-detail}
Compléter les détails de la preuve du 
ef{lem:relation-prisme} en vérifiant l'annulation
des termes correspondant aux faces verticales du prisme $\Delta^n \times [0,1]$.
\end{exercice}

% ===================================================================
\section*{Résumé du chapitre}\label{sec:resume-ch5}
\addcontentsline{toc}{section}{Résumé du chapitre}
% ===================================================================

\begin{enumerate}[label=\textbf{\arabic*.}]
  \item Le $n$-simplexe standard $\Delta^n$ est l'enveloppe convexe de $n+1$ sommets dans $\R^{n+1}$.
  Un $n$-simplexe singulier dans $X$ est une application continue $\sigma \colon \Delta^n \to X$.
  \item Le groupe des $n$-chaînes singulières $C_n(X)$ est le groupe abélien libre engendré par les
  $n$-simplexes singuliers. L'opérateur bord $\partial_n$ est défini par somme alternée des faces.
  \item La propriété $\partial^2 = 0$ fait de $(C_*(X), \partial)$ un complexe de chaînes.
  Les groupes d'homologie $H_n(X) = \ker \partial_n / \im \partial_{n+1}$ mesurent les trous
  $n$-dimensionnels.
  \item L'homologie est un foncteur $H_n \colon \mathbf{Top} \to \mathbf{Ab}$, invariant par homotopie.
  \item $H_*(\{*\})$ est $\Z$ en degré $0$ et trivial ailleurs. Pour les sphères,
  $\widetilde{H}_k(S^n) \cong \Z$ si $k=n$ et $0$ sinon.
  \item L'homologie relative $H_n(X,A)$ provient du complexe quotient $C_*(X)/C_*(A)$
  et intervient dans la suite exacte longue de la paire.
\end{enumerate}


%%%%%%%%%%%%%%%%%%%%%%%%%%%%%%%%%%%%%%%%%%%%%%%%%%%%%%%%%%%%%%%%%%%%%
%% CHAPITRE 6 : SUITES EXACTES LONGUES ET EXCISION                %%
%%%%%%%%%%%%%%%%%%%%%%%%%%%%%%%%%%%%%%%%%%%%%%%%%%%%%%%%%%%%%%%%%%%%%

\chapter{Suites exactes longues et excision}\label{chap:suites-exactes}
\minitoc\vspace{1cm}

\noindent
Le présent chapitre développe la machinerie algébrique et les résultats géométriques qui
confèrent à l'homologie singulière sa puissance de calcul. La suite exacte longue associée
à une paire $(X, A)$ relie l'homologie de $X$, de $A$ et de la paire, tandis que le théorème
d'excision permet de localiser les calculs homologiques. Combinés, ces outils conduisent à des
applications remarquables : calcul de l'homologie des sphères, notion de degré, théorèmes de
point fixe de Brouwer et invariance du domaine.\index{suite exacte longue}\index{excision}

% ===================================================================
\section{Suites exactes courtes de complexes de chaînes}\label{sec:ses-complexes}
% ===================================================================

\begin{definition}[Suite exacte courte de complexes]\label{def:ses-complexes}
\index{suite exacte courte!de complexes de chaînes}
Soient $(A_*, \partial^A)$, $(B_*, \partial^B)$, $(C_*, \partial^C)$ trois complexes de chaînes
de groupes abéliens. Une \emph{suite exacte courte de complexes de chaînes} est une suite
\[
\begin{tikzcd}
  0 \ar[r] & A_* \ar[r, "i"] & B_* \ar[r, "p"] & C_* \ar[r] & 0,
\end{tikzcd}
\]
où $i$ et $p$ sont des morphismes de complexes de chaînes tels que, pour chaque $n$,
\[
\begin{tikzcd}
  0 \ar[r] & A_n \ar[r, "i_n"] & B_n \ar[r, "p_n"] & C_n \ar[r] & 0
\end{tikzcd}
\]
est une suite exacte courte de groupes abéliens.
\end{definition}

Le diagramme complet de la situation est le suivant :
\[
\begin{tikzcd}
  & \vdots \ar[d] & \vdots \ar[d] & \vdots \ar[d] & \\
  0 \ar[r] & A_{n+1} \ar[r, "i"] \ar[d, "\partial^A"] & B_{n+1} \ar[r, "p"] \ar[d, "\partial^B"] & C_{n+1} \ar[r] \ar[d, "\partial^C"] & 0 \\
  0 \ar[r] & A_n \ar[r, "i"] \ar[d, "\partial^A"] & B_n \ar[r, "p"] \ar[d, "\partial^B"] & C_n \ar[r] \ar[d, "\partial^C"] & 0 \\
  0 \ar[r] & A_{n-1} \ar[r, "i"] \ar[d] & B_{n-1} \ar[r, "p"] \ar[d] & C_{n-1} \ar[r] \ar[d] & 0 \\
  & \vdots & \vdots & \vdots &
\end{tikzcd}
\]
où tous les carrés commutent.

% ===================================================================
\section{Le morphisme de connexion}\label{sec:morphisme-connexion}
% ===================================================================

\begin{definition}[Morphisme de connexion]\label{def:morphisme-connexion}
\index{morphisme de connexion}\index{opérateur de connexion}
Soit $0 \to A_* \xrightarrow{i} B_* \xrightarrow{p} C_* \to 0$ une suite exacte courte de
complexes de chaînes. Le \emph{morphisme de connexion} (ou \emph{opérateur bord})
\[
  \partial_* \colon H_n(C_*) \longrightarrow H_{n-1}(A_*)
\]
est construit par la procédure suivante de \emph{chasse au diagramme}\index{chasse au diagramme}.
\end{definition}

\begin{proposition}[Construction de $\partial_*$]\label{prop:construction-connexion}
Soit $[c] \in H_n(C_*)$ représenté par un cycle $c \in C_n$ (c'est-à-dire $\partial^C(c) = 0$).
\begin{enumerate}[label=(\roman*)]
  \item Comme $p_n$ est surjectif, il existe $b \in B_n$ tel que $p_n(b) = c$.
  \item On a $p_{n-1}(\partial^B(b)) = \partial^C(p_n(b)) = \partial^C(c) = 0$, donc
  $\partial^B(b) \in \ker p_{n-1} = \im i_{n-1}$.
  \item Il existe un unique $a \in A_{n-1}$ tel que $i_{n-1}(a) = \partial^B(b)$.
  \item Cet élément $a$ est un cycle : $i_{n-2}(\partial^A(a)) = \partial^B(i_{n-1}(a)) = \partial^B(\partial^B(b)) = 0$,
  et comme $i_{n-2}$ est injectif, $\partial^A(a) = 0$.
  \item On pose $\partial_*[c] = [a] \in H_{n-1}(A_*)$.
\end{enumerate}
\end{proposition}

Le diagramme ci-dessous illustre la chasse :
\[
\begin{tikzcd}[column sep=large]
  & b \ar[r, mapsto, "p_n"] \ar[d, mapsto, "\partial^B"] & c \ar[d, mapsto, "\partial^C"] \\
  a \ar[r, mapsto, "i_{n-1}"] & \partial^B(b) \ar[r, mapsto, "p_{n-1}"] & 0
\end{tikzcd}
\]

\begin{lemme}[Bonne définition de $\partial_*$]\label{lem:connexion-bien-defini}
Le morphisme $\partial_*$ est bien défini, c'est-à-dire que la classe $[a]$ ne dépend ni du
choix du relèvement $b$ de $c$, ni du choix du représentant $c$ de la classe $[c]$.
\end{lemme}

\begin{proof}
\textbf{Indépendance du relèvement.} Si $b'$ est un autre relèvement, $p_n(b - b') = 0$,
donc $b - b' = i_n(a')$ pour un certain $a' \in A_n$. Alors
$i_{n-1}(\partial^A(a')) = \partial^B(i_n(a')) = \partial^B(b) - \partial^B(b')$, d'où
les deux éléments $a$ obtenus diffèrent par $\partial^A(a')$, un bord.

\textbf{Indépendance du représentant.} Si $c' = c + \partial^C(c'')$ avec $c'' \in C_{n+1}$,
choisissons $b'' \in B_{n+1}$ avec $p_{n+1}(b'') = c''$. Alors $b + \partial^B(b'')$ est un
relèvement de $c'$, et $\partial^B(b + \partial^B(b'')) = \partial^B(b)$, d'où le même $a$.
\end{proof}

% ===================================================================
\section{Théorème de la suite exacte longue}\label{sec:sel-homologie}
% ===================================================================

\begin{theoreme}[Suite exacte longue en homologie]\label{thm:suite-exacte-longue}
\index{suite exacte longue!en homologie}\index{lemme du serpent}
Soit $0 \to A_* \xrightarrow{i} B_* \xrightarrow{p} C_* \to 0$ une suite exacte courte de
complexes de chaînes. Il existe une suite exacte longue de groupes abéliens :
\[
\begin{tikzcd}[column sep=1.8em]
  \cdots \ar[r] & H_n(A_*) \ar[r, "i_*"] & H_n(B_*) \ar[r, "p_*"] & H_n(C_*) \ar[dll, "\partial_*"', out=0, in=180, looseness=1.5, overlay] \\
  & H_{n-1}(A_*) \ar[r, "i_*"] & H_{n-1}(B_*) \ar[r, "p_*"] & H_{n-1}(C_*) \ar[r] & \cdots
\end{tikzcd}
\]
c'est-à-dire
\[
  \cdots \xrightarrow{\partial_*} H_n(A_*) \xrightarrow{i_*} H_n(B_*) \xrightarrow{p_*} H_n(C_*) \xrightarrow{\partial_*} H_{n-1}(A_*) \xrightarrow{i_*} \cdots
\]
De plus, la suite exacte longue est \emph{naturelle} : un morphisme de suites exactes courtes
de complexes induit un morphisme de suites exactes longues.
\end{theoreme}

\begin{proof}
Il y a six exactitudes à vérifier. Nous les démontrons toutes.

\textbf{1. Exactitude en $H_n(B_*)$ : $\im i_* \subset \ker p_*$.}
Pour $[a] \in H_n(A_*)$, on a $p_*(i_*[a]) = [p(i(a))] = [0] = 0$ car $p \circ i = 0$.

\textbf{2. Exactitude en $H_n(B_*)$ : $\ker p_* \subset \im i_*$.}
Soit $[b] \in H_n(B_*)$ avec $p_*[b] = 0$. Alors $p(b) = \partial^C(c')$ pour un certain
$c' \in C_{n+1}$. On choisit $b' \in B_{n+1}$ avec $p(b') = c'$. Posons $b'' = b - \partial^B(b')$.
Alors $p(b'') = p(b) - \partial^C(p(b')) = p(b) - \partial^C(c') = 0$, donc $b'' \in \ker p = \im i$.
Il existe $a \in A_n$ avec $i(a) = b''$, et $a$ est un cycle car $i(\partial^A(a)) = \partial^B(i(a)) = \partial^B(b'') = \partial^B(b) - \partial^B\partial^B(b') = 0$
et $i$ est injectif. Alors $i_*[a] = [b''] = [b]$.

\textbf{3. Exactitude en $H_n(C_*)$ : $\im p_* \subset \ker \partial_*$.}
Pour $[b] \in H_n(B_*)$, on a $\partial_*(p_*[b]) = \partial_*[p(b)]$. On peut relever
$p(b)$ à $b$ même, et $\partial^B(b) = 0$, donc $a = 0$ dans la construction, d'où $\partial_*[p(b)] = 0$.

\textbf{4. Exactitude en $H_n(C_*)$ : $\ker \partial_* \subset \im p_*$.}
Soit $[c] \in H_n(C_*)$ avec $\partial_*[c] = 0$. On relève $c$ à $b \in B_n$, et $\partial^B(b) = i(a)$
avec $[a] = 0$ dans $H_{n-1}(A_*)$. Alors $a = \partial^A(a')$ pour un certain $a' \in A_n$,
d'où $\partial^B(b - i(a')) = i(a) - i(\partial^A(a')) = 0$. Donc $b - i(a')$ est un cycle et
$p(b - i(a')) = p(b) = c$, d'où $[c] = p_*[b - i(a')]$.

\textbf{5. Exactitude en $H_{n-1}(A_*)$ : $\im \partial_* \subset \ker i_*$.}
Si $\partial_*[c] = [a]$, alors $i(a) = \partial^B(b)$ par construction, donc $i_*[a] = [\partial^B(b)] = 0$.

\textbf{6. Exactitude en $H_{n-1}(A_*)$ : $\ker i_* \subset \im \partial_*$.}
Soit $[a] \in H_{n-1}(A_*)$ avec $i_*[a] = 0$. Alors $i(a) = \partial^B(b)$ pour un certain
$b \in B_n$. Posons $c = p(b) \in C_n$. Alors $\partial^C(c) = \partial^C(p(b)) = p(\partial^B(b)) = p(i(a)) = 0$,
donc $c$ est un cycle. Par construction, $\partial_*[c] = [a]$.
\end{proof}

% ===================================================================
\section{Suite exacte longue de la paire}\label{sec:suite-exacte-paire}
% ===================================================================

\begin{theoreme}[Suite exacte longue de la paire]\label{thm:suite-exacte-paire}
\index{suite exacte longue!de la paire}
Soit $A \subset X$ un sous-espace. La suite exacte courte de complexes
$0 \to C_*(A) \xrightarrow{\iota_\sharp} C_*(X) \xrightarrow{\pi} C_*(X,A) \to 0$
induit une suite exacte longue :
\[
  \cdots \xrightarrow{\partial_*} H_n(A) \xrightarrow{\iota_*} H_n(X) \xrightarrow{j_*} H_n(X,A) \xrightarrow{\partial_*} H_{n-1}(A) \xrightarrow{\iota_*} \cdots
\]
\[
  \cdots \xrightarrow{j_*} H_1(X,A) \xrightarrow{\partial_*} H_0(A) \xrightarrow{\iota_*} H_0(X) \xrightarrow{j_*} H_0(X,A) \to 0.
\]
\end{theoreme}

\begin{proof}
C'est une application directe du 
ef{thm:suite-exacte-longue} à la suite exacte courte de
complexes $0 \to C_*(A) \to C_*(X) \to C_*(X,A) \to 0$.
\end{proof}

\begin{remarque}[Naturalité]\label{rem:naturalite-paire}
\index{naturalité}
Si $f \colon (X, A) \to (Y, B)$ est une application de paires (c'est-à-dire $f(A) \subset B$),
alors le diagramme suivant commute :
\[
\begin{tikzcd}
  \cdots \ar[r] & H_n(A) \ar[r, "\iota_*"] \ar[d, "f_*"] & H_n(X) \ar[r, "j_*"] \ar[d, "f_*"] & H_n(X,A) \ar[r, "\partial_*"] \ar[d, "f_*"] & H_{n-1}(A) \ar[r] \ar[d, "f_*"] & \cdots \\
  \cdots \ar[r] & H_n(B) \ar[r, "\iota_*"'] & H_n(Y) \ar[r, "j_*"'] & H_n(Y,B) \ar[r, "\partial_*"'] & H_{n-1}(B) \ar[r] & \cdots
\end{tikzcd}
\]
\end{remarque}

\begin{exemple}[La paire $(D^n, S^{n-1})$]\label{ex:paire-disque-sphere}
Considérons la paire $(D^n, S^{n-1})$ pour $n \geq 1$. Comme $D^n$ est contractile,
$\widetilde{H}_k(D^n) = 0$ pour tout $k$. La suite exacte longue réduite donne :
\[
  \cdots \to \widetilde{H}_k(D^n) \to H_k(D^n, S^{n-1}) \xrightarrow{\partial_*} \widetilde{H}_{k-1}(S^{n-1}) \to \widetilde{H}_{k-1}(D^n) \to \cdots
\]
Puisque les termes en $D^n$ sont nuls, on obtient des isomorphismes
\[
  H_k(D^n, S^{n-1}) \cong \widetilde{H}_{k-1}(S^{n-1})
\]
pour tout $k$. En particulier, $H_k(D^n, S^{n-1}) \cong \Z$ si $k = n$, et $0$ sinon.
\end{exemple}

% ===================================================================
\section{Théorème d'excision}\label{sec:excision}
% ===================================================================

Le théorème d'excision est l'un des résultats techniques les plus importants de l'homologie
singulière. Il affirme que l'on peut «~exciser~» un sous-espace qui est loin du bord sans
changer l'homologie relative.

\begin{theoreme}[Excision]\label{thm:excision}
\index{excision!théorème}\index{théorème d'excision}
Soient $Z \subset A \subset X$ des sous-espaces tels que $\overline{Z} \subset \mathring{A}$
(c'est-à-dire l'adhérence de $Z$ est contenue dans l'intérieur de $A$). Alors l'inclusion
$(X \setminus Z, A \setminus Z) \hookrightarrow (X, A)$ induit un isomorphisme
\[
  H_n(X \setminus Z, A \setminus Z) \xrightarrow{\;\sim\;} H_n(X, A)
\]
pour tout $n \geq 0$.

De manière équivalente, si $X = \mathring{A} \cup \mathring{B}$ (c'est-à-dire $A$ et $B$ sont
des sous-espaces dont les intérieurs recouvrent $X$), alors l'inclusion
$(B, A \cap B) \hookrightarrow (X, A)$ induit un isomorphisme $H_n(B, A \cap B) \cong H_n(X, A)$
pour tout $n$.
\end{theoreme}

La preuve repose sur la technique de subdivision barycentrique, que nous développons ci-dessous.

\begin{remarque}[Formes équivalentes de l'excision]\label{rem:excision-equivalence}
\index{excision!formes équivalentes}
Les deux formulations du 
ef{thm:excision} sont équivalentes. Pour passer de la première à la
seconde, on pose $B = X \setminus Z$, de sorte que $X \setminus Z = B$ et $A \setminus Z = A \cap B$.
La condition $\overline{Z} \subset \mathring{A}$ est équivalente à $X = \mathring{A} \cup \mathring{B}$,
car $X \setminus \mathring{A} \subset X \setminus \overline{Z} = \mathring{B}$.
\end{remarque}

\begin{remarque}[L'excision ne vaut pas pour l'homotopie]\label{rem:excision-pas-homotopie}
Il est important de noter que le théorème d'excision est propre à l'homologie et \emph{ne vaut
pas} pour les groupes d'homotopie. L'absence d'une propriété d'excision pour les groupes
d'homotopie supérieurs est l'une des raisons principales de la difficulté du calcul de $\pi_n(S^m)$
pour $n > m$. Le théorème de Blakers--Massey fournit un substitut partiel (excision homotopique)
sous des hypothèses de connectivité.
\end{remarque}

% ===================================================================
\section{Subdivision barycentrique}\label{sec:subdivision-barycentrique}
% ===================================================================

\begin{definition}[Barycentre]\label{def:barycentre}
\index{barycentre}\index{subdivision barycentrique}
Le \emph{barycentre} du $n$-simplexe standard est le point
\[
  \hat{b}_n = \frac{1}{n+1}(e_0 + e_1 + \cdots + e_n) = \left(\frac{1}{n+1}, \ldots, \frac{1}{n+1}\right) \in \Delta^n.
\]
\end{definition}

\begin{definition}[Opérateur de subdivision barycentrique]\label{def:operateur-subdivision}
\index{subdivision barycentrique!opérateur}
On définit un opérateur $\operatorname{sd} \colon C_n(X) \to C_n(X)$, appelé \emph{subdivision
barycentrique}, de la manière suivante. Pour un $n$-simplexe singulier $\sigma \colon \Delta^n \to X$ :
\begin{itemize}
  \item Si $n = 0$ : $\operatorname{sd}(\sigma) = \sigma$.
  \item Si $n \geq 1$ : on itère la construction. On considère le barycentre $\hat{b}$ de $\Delta^n$
  et on décompose $\Delta^n$ en $(n+1)!$ petits simplexes obtenus en joignant $\hat{b}$ aux
  subdivisions barycentriques des faces.
\end{itemize}
Formellement, on utilise l'\emph{opérateur cône}\index{opérateur cône}. Si $b$ est un point intérieur et
$\alpha = \sum_j n_j \tau_j$ est une $(n-1)$-chaîne dont les simplexes sont à valeurs dans un
convexe contenant $b$, on pose $b \cdot \tau_j$ le simplexe affine qui envoie
$e_0 \mapsto b$ et $e_{i+1} \mapsto \tau_j(e_i)$, puis $b \cdot \alpha = \sum_j n_j (b \cdot \tau_j)$.
La subdivision est alors
\[
  \operatorname{sd}(\sigma) = \sigma(\hat{b}_n) \cdot \operatorname{sd}(\partial \lambda_n),
\]
où $\lambda_n \colon \Delta^n \to \Delta^n$ est l'identité vue comme $n$-simplexe et l'opérateur cône
utilise le barycentre $\sigma(\hat{b}_n)$.
\end{definition}

\begin{lemme}[Propriétés de la subdivision]\label{lem:proprietes-subdivision}
\begin{enumerate}[label=(\roman*)]
  \item $\operatorname{sd}$ est un morphisme de chaînes : $\partial \circ \operatorname{sd} = \operatorname{sd} \circ \partial$.
  \item $\operatorname{sd}$ est homotope à l'identité comme morphisme de chaînes : il existe
  $T_n \colon C_n(X) \to C_{n+1}(X)$ tel que $\partial T + T \partial = \operatorname{sd} - \Id$.
  \item En itérant la subdivision, le diamètre des simplexes tend vers zéro :
  le diamètre de chaque simplexe de $\operatorname{sd}^k(\sigma)$ est au plus
  $\left(\frac{n}{n+1}\right)^k \operatorname{diam}(\sigma(\Delta^n))$.
\end{enumerate}
\end{lemme}

\begin{proof}[Esquisse de preuve]
(i) se vérifie par récurrence, en utilisant la formule $\partial(b \cdot \alpha) = \alpha - b \cdot \partial\alpha$
pour l'opérateur cône.
(ii) L'homotopie $T$ se construit également par récurrence, en utilisant l'opérateur cône.
(iii) C'est un résultat classique de géométrie convexe : la subdivision barycentrique d'un
$n$-simplexe produit des simplexes dont le diamètre est au plus $\frac{n}{n+1}$ fois le diamètre original.
\end{proof}

% ===================================================================
\section{Preuve du théorème d'excision}\label{sec:preuve-excision}
% ===================================================================

\begin{proof}[Preuve du 
ef{thm:excision}]
Soit $\mathcal{U} = \{A, B\}$ avec $B = X \setminus Z$. La condition $\overline{Z} \subset \mathring{A}$
équivaut à $X = \mathring{A} \cup \mathring{B}$. Notons $C_n^\mathcal{U}(X)$ le sous-groupe de
$C_n(X)$ engendré par les simplexes singuliers dont l'image est contenue soit dans $A$, soit dans $B$.

\textbf{Étape 1 : L'inclusion $C_*^\mathcal{U}(X) \hookrightarrow C_*(X)$ est une équivalence d'homotopie de chaînes.}

Par le lemme de Lebesgue\index{lemme de Lebesgue}, tout simplexe singulier $\sigma \colon \Delta^n \to X$
peut être «~raffiné~» par subdivision barycentrique itérée : il existe $k \geq 0$ tel que chaque
simplexe de $\operatorname{sd}^k(\sigma)$ a son image dans $A$ ou dans $B$.
Le 
ef{lem:proprietes-subdivision}(ii) fournit une homotopie de chaînes entre $\operatorname{sd}^k$
et l'identité.

\textbf{Étape 2 : Isomorphisme.}

La suite exacte courte
$0 \to C_*(A) \to C_*^\mathcal{U}(X) \to C_*^\mathcal{U}(X)/C_*(A) \to 0$
et l'identification $C_*^\mathcal{U}(X)/C_*(A) \cong C_*(B)/C_*(A \cap B) = C_*(B, A \cap B)$
montrent, via l'équivalence d'homotopie de l'étape 1, que
l'inclusion $(B, A \cap B) \hookrightarrow (X, A)$ induit un isomorphisme en homologie.

En posant $B = X \setminus Z$, on obtient
$H_n(X \setminus Z, A \setminus Z) \xrightarrow{\sim} H_n(X, A)$.
\end{proof}

\begin{remarque}[Rôle du lemme de Lebesgue]\label{rem:lebesgue-excision}
\index{lemme de Lebesgue}
Le lemme de Lebesgue intervient comme suit. Le recouvrement ouvert
$\{\mathring{A}, \mathring{B}\}$ de $X$ induit, par image réciproque via $\sigma \colon \Delta^n \to X$,
un recouvrement ouvert de $\Delta^n$. Comme $\Delta^n$ est compact métrique, le lemme de Lebesgue
garantit l'existence d'un nombre de Lebesgue $\lambda > 0$ tel que toute partie de $\Delta^n$ de
diamètre $< \lambda$ soit contenue dans l'un des ouverts du recouvrement. La propriété
(iii) du 
ef{lem:proprietes-subdivision} assure qu'après un nombre fini $k$ de subdivisions
barycentriques, chaque simplexe a un diamètre inférieur à $\lambda$, donc son image par $\sigma$
est contenue dans $A$ ou dans $B$.
\end{remarque}

% ===================================================================
\section{Théorème de Mayer--Vietoris (mention)}\label{sec:mayer-vietoris-mention}
% ===================================================================

\begin{theoreme}[Mayer--Vietoris]\label{thm:mayer-vietoris-mention}
\index{Mayer--Vietoris!théorème}
Soient $A, B \subset X$ des sous-espaces tels que $X = \mathring{A} \cup \mathring{B}$.
Il existe une suite exacte longue :
\[
  \cdots \to H_n(A \cap B) \xrightarrow{(\iota_{A*}, \iota_{B*})} H_n(A) \oplus H_n(B) \xrightarrow{j_{A*} - j_{B*}} H_n(X) \xrightarrow{\partial_*} H_{n-1}(A \cap B) \to \cdots
\]
où $\iota_A, \iota_B$ sont les inclusions de $A \cap B$ dans $A$ et $B$, et $j_A, j_B$ les
inclusions de $A$ et $B$ dans $X$.
\end{theoreme}

\begin{remarque}\label{rem:mv-preuve-ch7}
La preuve complète du théorème de Mayer--Vietoris sera donnée au chapitre~7.
Elle repose sur l'excision et la suite exacte longue. Le théorème de
Mayer--Vietoris est l'outil de calcul par excellence en homologie singulière.
\end{remarque}

\begin{remarque}[Mayer--Vietoris en homologie réduite]\label{rem:mv-reduite}
\index{Mayer--Vietoris!homologie réduite}
Il existe une version réduite de la suite de Mayer--Vietoris :
\[
  \cdots \to \widetilde{H}_n(A \cap B) \to \widetilde{H}_n(A) \oplus \widetilde{H}_n(B) \to \widetilde{H}_n(X) \xrightarrow{\partial_*} \widetilde{H}_{n-1}(A \cap B) \to \cdots
\]
valable sous les mêmes hypothèses. Cette version est souvent plus commode car elle
évite les termes $\Z$ parasites en degré $0$.
\end{remarque}

\begin{exemple}[Application : homologie de la suspension]\label{ex:homologie-suspension}
\index{suspension!homologie}
Soit $SX = (X \times [-1,1])/(X \times \{-1\} \cup X \times \{1\})$ la \emph{suspension} de $X$.
On décompose $SX$ en les deux cônes $A = (X \times [-1, 0])/(X \times \{-1\})$ et
$B = (X \times [0, 1])/(X \times \{1\})$, tous deux contractiles. Leur intersection est
$A \cap B \cong X$ (identifié à $X \times \{0\}$). Par Mayer--Vietoris réduit :
\[
  \widetilde{H}_n(SX) \cong \widetilde{H}_{n-1}(X) \quad \text{pour tout } n.
\]
Comme $S^n \cong S(S^{n-1})$ (la sphère $S^n$ est la suspension de $S^{n-1}$), on retrouve
par récurrence $\widetilde{H}_k(S^n) \cong \Z$ si $k = n$ et $0$ sinon.
\end{exemple}

% ===================================================================
\section{Équivalence entre homologie relative et homologie réduite du quotient}\label{sec:equiv-relative-quotient}
% ===================================================================

\begin{definition}[Bonne paire]\label{def:bonne-paire}
\index{bonne paire}
Une paire $(X, A)$ est dite \emph{bonne} (ou \emph{paire CW}, ou paire vérifiant la
propriété d'extension d'homotopie) si $A$ est un sous-espace fermé et possède un
voisinage $V$ dans $X$ tel que $A$ soit un rétracte par déformation de $V$.
\end{definition}

\begin{theoreme}\label{thm:relative-quotient}
\index{homologie relative!et quotient}
Si $(X, A)$ est une bonne paire avec $A \neq \varnothing$, alors la projection
$q \colon (X, A) \to (X/A, A/A)$ induit un isomorphisme
\[
  q_* \colon H_n(X, A) \xrightarrow{\;\sim\;} H_n(X/A, \{*\}) \cong \widetilde{H}_n(X/A)
\]
pour tout $n \geq 0$.
\end{theoreme}

\begin{proof}[Esquisse de preuve]
On considère le voisinage $V$ de $A$ tel que $A \hookrightarrow V$ soit une équivalence d'homotopie.
Par excision, $H_n(X, A) \cong H_n(X, V)$ (en excisant un sous-ensemble approprié).
Le morphisme de projection $q$ induit un homéomorphisme $X \setminus A \to X/A \setminus \{*\}$.
En utilisant l'invariance homotopique et l'excision dans $X/A$, on conclut.
\end{proof}

\begin{exemple}\label{ex:quotient-sn}
La paire $(D^n, S^{n-1})$ est une bonne paire, et $D^n/S^{n-1} \cong S^n$.
Donc $H_k(D^n, S^{n-1}) \cong \widetilde{H}_k(S^n)$, ce qui est cohérent avec
l'
ef{ex:paire-disque-sphere}.
\end{exemple}

% ===================================================================
\section{Applications}\label{sec:applications-excision}
% ===================================================================

\subsection{Calcul de $H_n(S^n)$ par excision}\label{subsec:calcul-Sn-excision}

\begin{proposition}\label{prop:Sn-excision}
Pour tout $n \geq 1$, on a $\widetilde{H}_k(S^n) \cong \widetilde{H}_{k-1}(S^{n-1})$
pour tout $k$.
\end{proposition}

\begin{proof}
Écrivons $S^n = D^n_+ \cup D^n_-$ comme réunion des hémisphères supérieur et inférieur,
avec $D^n_+ \cap D^n_- = S^{n-1}$ (l'équateur). Les deux hémisphères sont contractiles.
Par la suite exacte longue de la paire $(D^n_+, S^{n-1})$ :
\[
  \widetilde{H}_k(D^n_+) \to H_k(D^n_+, S^{n-1}) \xrightarrow{\partial_*} \widetilde{H}_{k-1}(S^{n-1}) \to \widetilde{H}_{k-1}(D^n_+).
\]
Comme $D^n_+$ est contractile, on obtient $H_k(D^n_+, S^{n-1}) \cong \widetilde{H}_{k-1}(S^{n-1})$.
Par le 
ef{thm:relative-quotient}, $H_k(D^n_+, S^{n-1}) \cong \widetilde{H}_k(D^n_+/S^{n-1}) \cong \widetilde{H}_k(S^n)$.
\end{proof}

\subsection{Degré d'une application}\label{subsec:degre}

\begin{definition}[Degré]\label{def:degre}
\index{degré!d'une application}
Soit $f \colon S^n \to S^n$ une application continue, $n \geq 1$. Le morphisme induit
$f_* \colon H_n(S^n) \to H_n(S^n)$ est un endomorphisme de $\Z \cong H_n(S^n)$, donc
de la forme $f_*([\alpha]) = d \cdot [\alpha]$ pour un unique entier $d \in \Z$.
Cet entier est le \emph{degré} de $f$, noté $\deg(f)$.
\end{definition}

\begin{proposition}[Propriétés du degré]\label{prop:proprietes-degre}
\index{degré!propriétés}
\begin{enumerate}[label=(\roman*)]
  \item $\deg(\Id_{S^n}) = 1$.
  \item $\deg(g \circ f) = \deg(g) \cdot \deg(f)$.
  \item Si $f \simeq g$, alors $\deg(f) = \deg(g)$.
  \item Si $f$ n'est pas surjective, alors $\deg(f) = 0$.
  \item La réflexion $r \colon S^n \to S^n$ définie par $r(x_0, x_1, \ldots, x_n) = (-x_0, x_1, \ldots, x_n)$
  a pour degré $-1$.
  \item L'antipodie $a \colon S^n \to S^n$, $a(x) = -x$, a pour degré $(-1)^{n+1}$.
\end{enumerate}
\end{proposition}

\begin{proof}
(i)--(iii) découlent directement de la fonctorialité et de l'invariance homotopique.

(iv) Si $f$ n'est pas surjective, disons $y_0 \notin \im(f)$, alors $f$ se factorise par
$S^n \setminus \{y_0\} \cong \R^n$, qui est contractile. Donc $f_* = 0$.

(v) On décompose $S^n$ en hémisphères séparés par l'hyperplan $x_0 = 0$.
La réflexion $r$ permute les deux hémisphères en inversant l'orientation,
d'où $\deg(r) = -1$ par un calcul direct en homologie relative.

(vi) L'antipodie est la composée de $n+1$ réflexions (une par coordonnée),
donc $\deg(a) = (-1)^{n+1}$.
\end{proof}

\begin{corollaire}\label{cor:antipodie-homotopie}
L'antipodie $a \colon S^n \to S^n$ est homotope à l'identité si et seulement si $n$ est impair.
\end{corollaire}

\begin{proof}
Si $n$ est impair, disons $n = 2m+1$, on peut écrire $S^n \subset \R^{2m+2} = \C^{m+1}$ et
l'homotopie $F(z, t) = e^{i\pi t} z$ relie continûment l'identité ($t=0$) à l'antipodie ($t=1$).
Réciproquement, si $n$ est pair, $\deg(a) = (-1)^{n+1} = -1 \neq 1 = \deg(\Id)$, donc $a$
n'est pas homotope à l'identité par invariance homotopique du degré.
\end{proof}

\begin{remarque}\label{rem:degre-condition-necessaire}
La condition $\deg(f) = \deg(g)$ est nécessaire mais en général pas suffisante pour que
$f \simeq g$. Le théorème de Hopf affirme toutefois que pour les sphères, le degré classifie
complètement les classes d'homotopie : $[S^n, S^n] \cong \Z$ via le degré.
\end{remarque}

\subsection{Théorème du point fixe de Brouwer}\label{subsec:brouwer}

\begin{theoreme}[Point fixe de Brouwer]\label{thm:brouwer-point-fixe}
\index{Brouwer!théorème du point fixe}\index{théorème du point fixe}
Toute application continue $f \colon D^n \to D^n$ du disque fermé dans lui-même possède au
moins un point fixe, pour tout $n \geq 1$.
\end{theoreme}

\begin{proof}
Raisonnons par l'absurde. Supposons que $f(x) \neq x$ pour tout $x \in D^n$.
On définit alors $r \colon D^n \to S^{n-1}$ en envoyant $x$ sur le point d'intersection de
la demi-droite issue de $f(x)$ passant par $x$ avec la sphère $S^{n-1}$.
Cette application $r$ est continue et vérifie $r(x) = x$ pour tout $x \in S^{n-1}$
(car si $x \in S^{n-1}$, la demi-droite de $f(x)$ à $x$ atteint $S^{n-1}$ en $x$ même).
Ainsi $r$ est une rétraction de $D^n$ sur $S^{n-1}$.

Notons $\iota \colon S^{n-1} \hookrightarrow D^n$ l'inclusion. Alors $r \circ \iota = \Id_{S^{n-1}}$,
d'où $r_* \circ \iota_* = \Id \colon H_{n-1}(S^{n-1}) \to H_{n-1}(S^{n-1})$.
En particulier, $r_*$ est surjectif. Or $H_{n-1}(S^{n-1}) \cong \Z$ (pour $n \geq 2$) et
$H_{n-1}(D^n) = 0$ (car $D^n$ est contractile), de sorte que $\iota_*$ se factorise par $0$,
contradiction. Pour $n=1$, l'argument utilise $H_0$.
\end{proof}

\begin{lemme}[Pas de rétraction]\label{lem:pas-retraction}
\index{rétraction!inexistence}
Il n'existe aucune rétraction $r \colon D^n \to S^{n-1}$ pour $n \geq 1$.
\end{lemme}

\begin{proof}
C'est exactement l'argument de la preuve ci-dessus, sans l'hypothèse sur $f$.
La suite exacte longue de la paire $(D^n, S^{n-1})$ montre que
$\iota_* \colon H_{n-1}(S^{n-1}) \to H_{n-1}(D^n)$ est le morphisme nul (car
$H_{n-1}(D^n) = 0$), mais si une rétraction $r$ existait, $r_* \circ \iota_* = \Id$
imposerait $\iota_*$ injectif, contradiction puisque $H_{n-1}(S^{n-1}) \cong \Z \neq 0$.
\end{proof}

\subsection{Invariance du domaine}\label{subsec:invariance-domaine}

\begin{theoreme}[Invariance du domaine de Brouwer]\label{thm:invariance-domaine}
\index{Brouwer!invariance du domaine}\index{invariance du domaine}
Si $U \subset \R^n$ est un ouvert et $f \colon U \to \R^n$ est une injection continue, alors
$f(U)$ est un ouvert de $\R^n$ et $f$ est un homéomorphisme de $U$ sur $f(U)$.
\end{theoreme}

\begin{proof}[Esquisse de preuve]
Il suffit de montrer que $f(U)$ est ouvert. Soit $x \in U$. On choisit une boule fermée
$\overline{B}(x, \varepsilon) \subset U$. Par compacité, $f|_{\overline{B}(x,\varepsilon)}$
est un homéomorphisme sur son image. Il faut montrer que $f(x)$ est intérieur à $f(\overline{B}(x,\varepsilon))$.

On utilise l'homologie. Notons $B = \overline{B}(x,\varepsilon)$ et $S = \partial B \cong S^{n-1}$.
Supposons par l'absurde que $f(x) \notin \mathring{f(B)}$. Alors il existe un rétracte de
$\R^n \setminus \{f(x)\}$ sur $f(S)$. Cela contredit le fait que l'inclusion
$f(S) \hookrightarrow \R^n \setminus \{f(x)\}$ n'est pas nulle en $H_{n-1}$, ce qui se
démontre en utilisant l'excision et le calcul de $H_{n-1}(S^{n-1})$.
\end{proof}

\begin{corollaire}[Invariance de la dimension]\label{cor:invariance-dimension}
\index{invariance de la dimension}
Si $\R^m$ est homéomorphe à $\R^n$, alors $m = n$.
\end{corollaire}

\begin{proof}
Un homéomorphisme $\R^m \cong \R^n$ donne par restriction un homéomorphisme $\R^m \setminus \{0\} \cong \R^n \setminus \{0\}$. Or $\R^k \setminus \{0\}$ a le type d'homotopie de $S^{k-1}$,
donc $S^{m-1} \simeq S^{n-1}$, d'où $H_{m-1}(S^{m-1}) \cong H_{n-1}(S^{n-1})$.
Puisque $\widetilde{H}_k(S^j) \neq 0$ seulement pour $k = j$, on conclut $m - 1 = n - 1$.
\end{proof}

\begin{remarque}[Portée du théorème d'invariance du domaine]\label{rem:portee-invariance}
Le théorème d'invariance du domaine est remarquable en ce qu'il ne suppose \emph{que} la
continuité et l'injectivité de $f$ --- aucune hypothèse de différentiabilité n'est requise.
En dimension $1$, le résultat est élémentaire (toute injection continue d'un intervalle
ouvert dans $\R$ est un homéomorphisme sur son image, par le théorème des valeurs intermédiaires).
En dimension supérieure, des outils d'homologie sont indispensables.
\end{remarque}

\subsection{Théorème de Borsuk--Ulam}\label{subsec:borsuk-ulam}

Le théorème de Borsuk--Ulam est l'une des applications les plus spectaculaires de l'homologie.

\begin{theoreme}[Borsuk--Ulam]\label{thm:borsuk-ulam}
\index{Borsuk--Ulam!théorème}
Pour toute application continue $f \colon S^n \to \R^n$, il existe un point $x \in S^n$
tel que $f(x) = f(-x)$.
\end{theoreme}

\begin{proof}[Esquisse de preuve]
Supposons par l'absurde que $f(x) \neq f(-x)$ pour tout $x$. On définit alors
$g \colon S^n \to S^{n-1}$ par
\[
  g(x) = \frac{f(x) - f(-x)}{\abs{f(x) - f(-x)}}.
\]
Cette application est continue et vérifie $g(-x) = -g(x)$ (elle est \emph{impaire} ou
\emph{antipodal-équivariante}). On peut montrer, par un argument de degré utilisant
l'homologie, qu'une telle application impaire $S^n \to S^{n-1}$ ne peut pas exister,
car elle induirait un morphisme non trivial en homologie qui est incompatible
avec les degrés. La preuve complète utilise le degré modulo $2$ ou la cohomologie
de $\Z/2\Z$ des espaces projectifs.
\end{proof}

\begin{corollaire}[Corollaire géographique]\label{cor:borsuk-ulam-geo}
À tout instant, il existe deux points diamétralement opposés sur la surface de la Terre
ayant la même température et la même pression atmosphérique.
\end{corollaire}

\begin{proof}
La surface de la Terre est (approximativement) $S^2$. L'application
$f \colon S^2 \to \R^2$ définie par $f(x) = (\text{température}(x), \text{pression}(x))$
est continue. Par le 
ef{thm:borsuk-ulam}, il existe $x$ tel que $f(x) = f(-x)$.
\end{proof}

\begin{corollaire}[Recouvrement de $S^n$]\label{cor:lusternik-schnirelmann}
\index{Lusternik--Schnirelmann}
Si $S^n = F_1 \cup F_2 \cup \cdots \cup F_{n+1}$ est un recouvrement par $n+1$ ensembles
fermés, alors l'un des $F_i$ contient une paire de points antipodaux $\{x, -x\}$.
\end{corollaire}

% ===================================================================
\section{Exercices}\label{sec:exo-ch6}
% ===================================================================

\begin{exercice}[$\star\star$]\label{exo:suite-exacte-triade}
Soit $X = A \cup B$ avec $A, B$ des sous-espaces. Écrire la suite exacte longue de la paire
$(X, A)$ et identifier les termes en utilisant l'excision lorsque c'est possible.
\end{exercice}

\begin{exercice}[$\star\star$]\label{exo:degre-calculs}
Calculer le degré des applications suivantes $S^1 \to S^1$ :
\begin{enumerate}[label=(\alph*)]
  \item $z \mapsto z^n$ (où $S^1 \subset \C$),
  \item $z \mapsto \bar{z}$ (conjugaison complexe),
  \item $z \mapsto z^2 \bar{z}$.
\end{enumerate}
\end{exercice}

\begin{exercice}[$\star\star$]\label{exo:brouwer-consequence}
Soit $f \colon S^{2n} \to S^{2n}$ continue. Montrer qu'il existe $x \in S^{2n}$ tel que
$f(x) = x$ ou $f(x) = -x$.
\emph{Indication : utiliser le degré de l'antipodie sur les sphères de dimension paire.}
\end{exercice}

\begin{exercice}[$\star\star\star$]\label{exo:hairy-ball}
\index{théorème de la boule chevelue}
(\emph{Théorème de la boule chevelue.}) Montrer qu'il n'existe pas de champ de vecteurs continu
non nul sur $S^{2n}$. Précisément, montrer qu'il n'existe pas d'application continue
$v \colon S^{2n} \to \R^{2n+1}$ telle que $v(x) \neq 0$ et $\langle v(x), x \rangle = 0$
pour tout $x \in S^{2n}$.
\emph{Indication : si un tel $v$ existe, construire une homotopie entre l'identité et
l'antipodie.}
\end{exercice}

\begin{exercice}[$\star\star$]\label{exo:suite-mv-calcul}
Utiliser la suite de Mayer--Vietoris pour calculer les groupes d'homologie du tore
$T^2 = S^1 \times S^1$, en décomposant $T^2$ en deux cylindres.
\end{exercice}

\begin{exercice}[$\star\star\star$]\label{exo:borsuk-ulam-dim2}
\index{Borsuk--Ulam}
(\emph{Théorème de Borsuk--Ulam en dimension 2.})
Montrer que pour toute application continue $f \colon S^2 \to \R^2$, il existe $x \in S^2$ tel
que $f(x) = f(-x)$.
\emph{Indication : supposer le contraire et en déduire une application $S^2 \to S^1$ de degré
impair, ce qui est impossible.}
\end{exercice}

\begin{exercice}[$\star\star\star$]\label{exo:excision-verification}
Vérifier directement l'énoncé d'excision dans le cas suivant : $X = [0,2]$, $A = [0,1]$,
$Z = [0, 1/2[$. Identifier $H_n(X \setminus Z, A \setminus Z)$ et $H_n(X, A)$, et construire
explicitement l'isomorphisme.
\end{exercice}

\begin{exercice}[$\star\star$]\label{exo:sel-triple}
(\emph{Suite exacte longue d'un triple.})
Soient $B \subset A \subset X$. Montrer qu'il existe une suite exacte longue :
\[
  \cdots \to H_n(A, B) \to H_n(X, B) \to H_n(X, A) \xrightarrow{\partial_*} H_{n-1}(A, B) \to \cdots
\]
\emph{Indication : considérer la suite exacte courte
$0 \to C_*(A)/C_*(B) \to C_*(X)/C_*(B) \to C_*(X)/C_*(A) \to 0$.}
\end{exercice}

\begin{exercice}[$\star\star\star$]\label{exo:homologie-bouquet}
Soit $X = \bigvee_{i=1}^k S^{n_i}$ un bouquet de sphères ($n_i \geq 1$). Montrer, en utilisant
l'excision et le 
ef{thm:relative-quotient}, que
\[
  \widetilde{H}_m(X) \cong \bigoplus_{i=1}^{k} \widetilde{H}_m(S^{n_i})
\]
pour tout $m \geq 0$.
\end{exercice}

\begin{exercice}[$\star\star$]\label{exo:degre-suspension}
\index{degré!et suspension}
Soit $f \colon S^n \to S^n$ de degré $d$. La \emph{suspension} de $f$ est l'application
$Sf \colon S^{n+1} \to S^{n+1}$ définie par $Sf(x, t) = (f(x)\sqrt{1-t^2}, t)$ pour
$(x, t) \in S^n \times [-1, 1] \subset \R^{n+2}$ (en identifiant $S^{n+1}$ à la suspension de $S^n$).
Montrer que $\deg(Sf) = \deg(f) = d$.
\end{exercice}

\begin{exercice}[$\star\star\star$]\label{exo:homologie-projectif}
\index{espace projectif!homologie}
Utiliser la suite exacte longue de la paire $(\R P^n, \R P^{n-1})$ et le fait que
$\R P^n / \R P^{n-1} \cong S^n$ pour calculer par récurrence les groupes $H_k(\R P^2; \Z)$.
On montrera que $H_0(\RP^2) = \Z$, $H_1(\RP^2) = \Z/2\Z$ et $H_k(\RP^2) = 0$ pour $k \geq 2$.
\end{exercice}

% ===================================================================
\section{Remarques complémentaires}\label{sec:remarques-ch6}
% ===================================================================

\begin{remarque}[Théorie des faisceaux et excision]\label{rem:faisceaux-excision}
\index{faisceau}
Le théorème d'excision peut être vu comme une manifestation de la nature «~locale~» de
l'homologie singulière. En langage moderne (cohomologie des faisceaux, développée par
Leray, Cartan et Godement), l'excision est essentiellement équivalente à la propriété que le
préfaisceau des chaînes singulières est un \emph{faisceau mou}. Cette perspective sera
approfondie dans les chapitres sur la cohomologie.
\end{remarque}

\begin{remarque}[Axiomes d'Eilenberg--Steenrod]\label{rem:eilenberg-steenrod}
\index{Eilenberg--Steenrod!axiomes}
Les résultats de ce chapitre et du précédent établissent que l'homologie singulière satisfait
les axiomes d'Eilenberg--Steenrod pour une théorie d'homologie :
\begin{enumerate}[label=(ES\arabic*)]
  \item \textbf{Fonctorialité} : $H_n$ est un foncteur covariant de la catégorie des paires
  d'espaces vers celle des groupes abéliens.
  \item \textbf{Invariance homotopique} : des applications de paires homotopes induisent les
  mêmes morphismes en homologie.
  \item \textbf{Suite exacte longue} : toute paire $(X, A)$ donne lieu à une suite exacte longue
  naturelle reliant $H_*(A)$, $H_*(X)$ et $H_*(X, A)$.
  \item \textbf{Excision} : si $\overline{Z} \subset \mathring{A}$, l'inclusion induit un
  isomorphisme $H_n(X \setminus Z, A \setminus Z) \cong H_n(X, A)$.
  \item \textbf{Axiome de la dimension} : $H_n(\{*\}) = 0$ pour $n \neq 0$.
\end{enumerate}
Le théorème d'Eilenberg--Steenrod affirme que ces cinq axiomes caractérisent l'homologie
singulière de manière unique sur la catégorie des CW-complexes finis (et plus généralement
des paires compactes).
\end{remarque}

% ===================================================================
\section*{Résumé du chapitre}\label{sec:resume-ch6}
\addcontentsline{toc}{section}{Résumé du chapitre}
% ===================================================================

\begin{enumerate}[label=\textbf{\arabic*.}]
  \item Toute suite exacte courte de complexes de chaînes
  $0 \to A_* \to B_* \to C_* \to 0$ induit une suite exacte longue en homologie
  $\cdots \to H_n(A_*) \to H_n(B_*) \to H_n(C_*) \xrightarrow{\partial_*} H_{n-1}(A_*) \to \cdots$
  \item Le morphisme de connexion $\partial_*$ est construit par chasse au diagramme et est naturel.
  \item La suite exacte longue de la paire $(X, A)$ relie $H_*(A)$, $H_*(X)$ et $H_*(X,A)$.
  \item Le théorème d'excision affirme que l'on peut retirer un sous-espace «~intérieur~»
  à $A$ sans changer l'homologie relative $H_*(X, A)$. La preuve utilise la subdivision barycentrique.
  \item Pour une bonne paire, $H_n(X, A) \cong \widetilde{H}_n(X/A)$.
  \item Le degré d'une application $f \colon S^n \to S^n$ est l'entier $d$ tel que $f_* = d \cdot \Id$
  sur $H_n(S^n) \cong \Z$.
  \item Le théorème du point fixe de Brouwer ($D^n \to D^n$ a un point fixe) se déduit de
  l'inexistence d'une rétraction $D^n \to S^{n-1}$, elle-même conséquence de l'homologie.
  \item L'invariance du domaine de Brouwer montre que toute injection continue d'un ouvert
  de $\R^n$ dans $\R^n$ est un homéomorphisme sur son image.
\end{enumerate}

% === FIN DU CHUNK ch5-6 ===% === DEBUT DU CHUNK ch7-8 ===

%%%%%%%%%%%%%%%%%%%%%%%%%%%%%%%%%%%%%%%%%%%%%%%%%%%%%%%%%%%%%%%%%%%%
%  CHAPITRE 7 — Théorème de Mayer-Vietoris
%%%%%%%%%%%%%%%%%%%%%%%%%%%%%%%%%%%%%%%%%%%%%%%%%%%%%%%%%%%%%%%%%%%%
\chapter{Théorème de Mayer-Vietoris}\label{ch:mayer-vietoris}
\index{Mayer-Vietoris, théorème de}

\section*{Introduction}
\addcontentsline{toc}{section}{Introduction}

La suite exacte de Mayer-Vietoris\index{suite exacte!de Mayer-Vietoris} est l'un des
outils de calcul les plus puissants en homologie singulière. Elle permet de déterminer
l'homologie d'un espace topologique~$X$ à partir de celles de deux sous-espaces
dont~$X$ est la réunion, à condition de connaître l'homologie de leur intersection.
Analogue en homologie du théorème de Seifert--van~Kampen\index{Seifert--van Kampen, théorème de}
pour le groupe fondamental, cette suite exacte longue repose de manière essentielle sur
le théorème d'excision.

Dans ce chapitre, nous établissons la suite de Mayer-Vietoris, puis l'appliquons
systématiquement au calcul de l'homologie des sphères, des surfaces compactes, des
espaces projectifs et d'autres espaces classiques. Nous terminons par une version
simplifiée de la formule de Künneth et ses conséquences pour la caractéristique d'Euler.

%% ================================================================
\section{Rappels : excision et suites exactes}\label{sec:mv-rappels}
\index{excision, théorème d'}

Nous rappelons brièvement les résultats fondamentaux qui servent de fondement à la
construction de Mayer-Vietoris.

\begin{theoreme}[Excision]\label{thm:excision-rappel}
\index{excision, théorème d'}
Soit~$X$ un espace topologique et soient~$Z \subset A \subset X$ tels
que~$\overline{Z} \subset \mathring{A}$. Alors l'inclusion~$(X \setminus Z, A \setminus Z)
\hookrightarrow (X, A)$ induit des isomorphismes
\[
  H_n(X \setminus Z,\, A \setminus Z) \xrightarrow{\;\cong\;} H_n(X, A)
  \qquad \text{pour tout } n \geqslant 0.
\]
\end{theoreme}

\begin{theoreme}[Suite exacte longue d'une paire]\label{thm:sel-paire-rappel}
\index{suite exacte!longue d'une paire}
Soit~$(X, A)$ une paire d'espaces topologiques. Il existe une suite exacte longue
\[
  \cdots \longrightarrow H_n(A) \xrightarrow{\;i_*\;} H_n(X)
  \xrightarrow{\;j_*\;} H_n(X, A) \xrightarrow{\;\partial_n\;}
  H_{n-1}(A) \longrightarrow \cdots
\]
où~$i \colon A \hookrightarrow X$ et~$j \colon (X, \varnothing) \hookrightarrow (X, A)$
sont les inclusions, et~$\partial_n$ est l'opérateur bord de connexion.\index{opérateur bord}
\end{theoreme}

\begin{remarque}\label{rem:excision-variante}
De manière équivalente, l'excision peut s'énoncer ainsi : si~$A, B \subset X$ sont
tels que~$X = \mathring{A} \cup \mathring{B}$, alors l'inclusion
$(B, A \cap B) \hookrightarrow (X, A)$ induit un isomorphisme en homologie.
C'est cette formulation qui sera la plus commode pour dériver Mayer-Vietoris.
\end{remarque}

\begin{definition}[Paire excisive]\label{def:paire-excisive}
\index{paire excisive}
On dit que le couple~$(A, B)$ est une \emph{paire excisive} pour~$X = A \cup B$ si
les intérieurs~$\mathring{A}$ et~$\mathring{B}$ recouvrent~$X$ :
\[
  X = \mathring{A} \cup \mathring{B}.
\]
\end{definition}

%% ================================================================
\section{La suite exacte de Mayer-Vietoris}\label{sec:mv-enonce}
\index{Mayer-Vietoris, suite exacte de}

\begin{theoreme}[Mayer-Vietoris]\label{thm:mayer-vietoris}
Soit~$X$ un espace topologique et soient~$A, B \subset X$ tels que~$X = \mathring{A} \cup
\mathring{B}$. On note~$i_A \colon A \cap B \hookrightarrow A$,
$i_B \colon A \cap B \hookrightarrow B$, $j_A \colon A \hookrightarrow X$,
$j_B \colon B \hookrightarrow X$ les inclusions. Alors il existe une suite exacte longue
\[
\cdots \longrightarrow H_n(A \cap B)
\xrightarrow{\;\Phi_n\;} H_n(A) \oplus H_n(B)
\xrightarrow{\;\Psi_n\;} H_n(X)
\xrightarrow{\;\Delta_n\;} H_{n-1}(A \cap B) \longrightarrow \cdots
\]
qui se termine par
\[
\cdots \longrightarrow H_0(A) \oplus H_0(B) \xrightarrow{\;\Psi_0\;} H_0(X)
\longrightarrow 0,
\]
où les morphismes sont définis par
\begin{align}
  \Phi_n(\alpha) &= \bigl((i_A)_*(\alpha),\; -(i_B)_*(\alpha)\bigr),
    \label{eq:mv-phi}\\
  \Psi_n(\alpha, \beta) &= (j_A)_*(\alpha) + (j_B)_*(\beta),
    \label{eq:mv-psi}
\end{align}
et~$\Delta_n \colon H_n(X) \to H_{n-1}(A \cap B)$ est l'opérateur de connexion.\index{opérateur de connexion!de Mayer-Vietoris}
\end{theoreme}

Le diagramme suivant résume la suite de Mayer-Vietoris :
\[
\begin{tikzcd}[column sep=1.6em, row sep=2.5em]
  \cdots \ar[r]
  & H_n(A \cap B) \ar[r, "\Phi_n"]
  & H_n(A) \oplus H_n(B) \ar[r, "\Psi_n"]
  & H_n(X) \ar[dll, "\Delta_n"', out=-10, in=170, overlay]\\
  & H_{n-1}(A \cap B) \ar[r, "\Phi_{n-1}"]
  & H_{n-1}(A) \oplus H_{n-1}(B) \ar[r, "\Psi_{n-1}"]
  & H_{n-1}(X) \ar[r]
  & \cdots
\end{tikzcd}
\]

\begin{remarque}[Description de l'opérateur de connexion]\label{rem:mv-connexion}
\index{opérateur de connexion!description explicite}
L'opérateur de connexion~$\Delta_n \colon H_n(X) \to H_{n-1}(A \cap B)$ admet la
description suivante. Soit~$[\gamma] \in H_n(X)$ une classe d'homologie, représentée
par un cycle~$\gamma \in C_n(X)$. Par le théorème des petites chaînes, on peut supposer
(quitte à subdiviser) que~$\gamma = \alpha + \beta$ avec~$\alpha \in C_n(A)$
et~$\beta \in C_n(B)$. Alors~$\partial \alpha = -\partial \beta$ est un $(n-1)$-cycle
dans~$C_{n-1}(A \cap B)$ (car~$\partial \alpha \in C_{n-1}(A)$
et~$-\partial \beta \in C_{n-1}(B)$, et ces deux chaînes sont égales). On pose
\[
  \Delta_n([\gamma]) = [\partial \alpha] = -[\partial \beta] \in H_{n-1}(A \cap B).
\]
La vérification que cette classe est bien définie (indépendante du choix de~$\gamma$
et de la décomposition~$\gamma = \alpha + \beta$) est un exercice standard
de \emph{diagram chasing}.
\end{remarque}

\begin{remarque}[Naturalité de Mayer-Vietoris]\label{rem:mv-naturalite}
\index{Mayer-Vietoris!naturalité}
La suite de Mayer-Vietoris est \emph{naturelle} au sens suivant. Si~$f \colon X \to X'$
est une application continue telle que~$f(A) \subset A'$ et~$f(B) \subset B'$
(où~$X' = A' \cup B'$ est une autre décomposition excisive), alors il existe un
morphisme de suites exactes longues :
\[
\begin{tikzcd}[column sep=1.4em]
  H_n(A \cap B) \ar[r, "\Phi_n"] \ar[d, "f_*"']
  & H_n(A) \oplus H_n(B) \ar[r, "\Psi_n"] \ar[d, "f_* \oplus f_*"']
  & H_n(X) \ar[r, "\Delta_n"] \ar[d, "f_*"']
  & H_{n-1}(A \cap B) \ar[d, "f_*"']\\
  H_n(A' \cap B') \ar[r, "\Phi_n'"]
  & H_n(A') \oplus H_n(B') \ar[r, "\Psi_n'"]
  & H_n(X') \ar[r, "\Delta_n'"]
  & H_{n-1}(A' \cap B')
\end{tikzcd}
\]
Ce diagramme commute dans chaque carré.
\end{remarque}

%% ================================================================
\section{Preuve du théorème de Mayer-Vietoris}\label{sec:mv-preuve}

\begin{proof}
La démonstration procède en deux étapes.

\medskip
\noindent\textbf{Étape 1 : la suite exacte courte de complexes de chaînes.}
\index{suite exacte courte!de complexes de chaînes}

On considère la suite
\begin{equation}\label{eq:mv-sec}
  0 \longrightarrow C_*(A \cap B)
  \xrightarrow{\;\varphi\;} C_*(A) \oplus C_*(B)
  \xrightarrow{\;\psi\;} C_*(A) + C_*(B) \longrightarrow 0,
\end{equation}
où, pour tout $n$-simplexe singulier~$\sigma$ de~$A \cap B$,
\[
  \varphi(\sigma) = (\sigma, -\sigma),
  \qquad
  \psi(\sigma_A, \sigma_B) = \sigma_A + \sigma_B.
\]
Ici $C_*(A) + C_*(B)$ désigne le sous-complexe de~$C_*(X)$ engendré par les chaînes
singulières à image dans~$A$ ou dans~$B$.

\emph{Injectivité de~$\varphi$.} Si~$\varphi(\sigma) = (0, 0)$, alors~$\sigma = 0$.

\emph{Surjectivité de~$\psi$.} Tout élément de~$C_*(A) + C_*(B)$ s'écrit par
définition comme~$\sigma_A + \sigma_B$ avec~$\sigma_A \in C_*(A)$ et~$\sigma_B \in C_*(B)$.

\emph{Exactitude en~$C_*(A) \oplus C_*(B)$.} On a~$\psi \circ \varphi = 0$ car
$\psi(\sigma, -\sigma) = \sigma - \sigma = 0$. Réciproquement,
si~$\psi(\sigma_A, \sigma_B) = \sigma_A + \sigma_B = 0$, alors~$\sigma_A = -\sigma_B$
est une chaîne dans~$C_*(A) \cap C_*(B) = C_*(A \cap B)$, et
$(\sigma_A, \sigma_B) = \varphi(-\sigma_B)$.

La suite~\eqref{eq:mv-sec} est donc exacte courte.

\medskip
\noindent\textbf{Étape 2 : passage à l'homologie et rôle de l'excision.}

Le théorème du serpent\index{théorème du serpent} (ou théorème de la suite exacte longue
en homologie) appliqué à~\eqref{eq:mv-sec} fournit une suite exacte longue
\[
  \cdots \to H_n(A \cap B) \xrightarrow{\Phi_n} H_n(A) \oplus H_n(B)
  \xrightarrow{\Psi_n} H_n(C_*(A) + C_*(B)) \xrightarrow{\Delta_n}
  H_{n-1}(A \cap B) \to \cdots
\]

Il reste à identifier~$H_n(C_*(A) + C_*(B))$ avec~$H_n(X)$. L'inclusion
$\iota \colon C_*(A) + C_*(B) \hookrightarrow C_*(X)$ induit un morphisme
$\iota_* \colon H_n(C_*(A) + C_*(B)) \to H_n(X)$.

Or la condition~$X = \mathring{A} \cup \mathring{B}$ signifie exactement que le
recouvrement~$\{A, B\}$ est admissible. Le \emph{théorème des petites chaînes}\index{petites chaînes, théorème des}
(conséquence directe de l'excision) affirme que~$\iota_*$ est un isomorphisme.
Plus précisément, toute chaîne singulière de~$X$ est homologue à une chaîne dans
$C_*(A) + C_*(B)$ par subdivision barycentrique\index{subdivision barycentrique}
itérée, car les intérieurs de~$A$ et~$B$ recouvrent~$X$.

En composant avec cet isomorphisme, on obtient la suite exacte longue de
Mayer-Vietoris avec les morphismes~$\Phi_n$, $\Psi_n$ et~$\Delta_n$ tels que
décrits dans l'énoncé du théorème.
\end{proof}

%% ================================================================
\section{Version réduite}\label{sec:mv-reduite}
\index{Mayer-Vietoris!version réduite}
\index{homologie réduite!Mayer-Vietoris}

\begin{theoreme}[Mayer-Vietoris réduite]\label{thm:mv-reduite}
Sous les mêmes hypothèses que le 
ef{thm:mayer-vietoris}, et en supposant de plus
que~$A \cap B \neq \varnothing$, on dispose d'une suite exacte longue en homologie
réduite\index{homologie réduite} :
\[
  \cdots \longrightarrow \widetilde{H}_n(A \cap B)
  \xrightarrow{\;\Phi_n\;}
  \widetilde{H}_n(A) \oplus \widetilde{H}_n(B)
  \xrightarrow{\;\Psi_n\;}
  \widetilde{H}_n(X)
  \xrightarrow{\;\Delta_n\;}
  \widetilde{H}_{n-1}(A \cap B)
  \longrightarrow \cdots
\]
qui se termine par
\[
  \cdots \longrightarrow \widetilde{H}_0(A) \oplus \widetilde{H}_0(B)
  \xrightarrow{\;\Psi_0\;} \widetilde{H}_0(X) \longrightarrow 0.
\]
\end{theoreme}

\begin{proof}
On utilise la chaîne augmentée~$\widetilde{C}_*(-)$ et le fait que l'augmentation
$\varepsilon \colon C_0(-) \to \Z$ est compatible avec les inclusions. La suite
exacte courte augmentée
\[
  0 \to \widetilde{C}_*(A \cap B) \to \widetilde{C}_*(A) \oplus \widetilde{C}_*(B)
  \to \widetilde{C}_*(A) + \widetilde{C}_*(B) \to 0
\]
donne, par le même argument que précédemment, la suite exacte longue annoncée.
\end{proof}

\begin{remarque}\label{rem:mv-reduite-utilite}
La version réduite est souvent plus commode pour les calculs car elle évite les
copies parasites de~$\Z$ en degré~$0$ qui proviennent des composantes connexes.
\end{remarque}

%% ================================================================
\section{Applications}\label{sec:mv-applications}

\subsection{Homologie des sphères}\label{subsec:homologie-spheres}
\index{sphère!homologie de la}

\begin{theoreme}[Homologie de~$S^n$]\label{thm:homologie-spheres}
Pour tout~$n \geqslant 0$, l'homologie réduite de la sphère~$S^n$ est donnée par
\[
  \widetilde{H}_k(S^n) \cong
  \begin{cases}
    \Z & \text{si } k = n,\\
    0  & \text{sinon.}
  \end{cases}
\]
\end{theoreme}

\begin{proof}
On procède par récurrence sur~$n$.

\medskip
\noindent\textit{Cas de base :~$n = 0$.}
La sphère~$S^0 = \{-1, +1\}$ est un espace discret à deux points. On a
$\widetilde{H}_0(S^0) \cong \Z$ et~$\widetilde{H}_k(S^0) = 0$ pour~$k \neq 0$.

\medskip
\noindent\textit{Étape de récurrence.}
Supposons le résultat vrai pour~$S^{n-1}$. On décompose~$S^n$ en
\[
  A = S^n \setminus \{N\}, \qquad B = S^n \setminus \{S\},
\]
où~$N = (0, \ldots, 0, 1)$ et~$S = (0, \ldots, 0, -1)$ sont les pôles nord et sud.
Alors~$A$ et~$B$ sont homéomorphes à~$\R^n$ (par projection stéréographique),
donc contractiles, et
\[
  A \cap B = S^n \setminus \{N, S\} \simeq S^{n-1}
\]
par une rétraction par déformation radiale sur l'équateur.

On a~$\mathring{A} = A$ et~$\mathring{B} = B$ (car~$A$ et~$B$ sont ouverts),
donc~$S^n = \mathring{A} \cup \mathring{B}$ et la suite de Mayer-Vietoris
réduite s'applique :
\[
  \cdots \to \widetilde{H}_k(A) \oplus \widetilde{H}_k(B)
  \to \widetilde{H}_k(S^n)
  \xrightarrow{\Delta_k}
  \widetilde{H}_{k-1}(A \cap B)
  \to \widetilde{H}_{k-1}(A) \oplus \widetilde{H}_{k-1}(B) \to \cdots
\]
Puisque~$A$ et~$B$ sont contractiles,~$\widetilde{H}_k(A) = \widetilde{H}_k(B) = 0$
pour tout~$k$. La suite donne donc des isomorphismes
\[
  \Delta_k \colon \widetilde{H}_k(S^n) \xrightarrow{\;\cong\;} \widetilde{H}_{k-1}(S^{n-1})
  \qquad \text{pour tout } k.
\]
Par hypothèse de récurrence,~$\widetilde{H}_{k-1}(S^{n-1}) \cong \Z$ si~$k-1 = n-1$
(soit~$k = n$) et~$0$ sinon. D'où le résultat.
\end{proof}

\subsection{Homologie du tore}\label{subsec:homologie-tore}
\index{tore!homologie du}

\begin{theoreme}[Homologie de~$T^2$]\label{thm:homologie-tore}
L'homologie du tore~$T^2 = S^1 \times S^1$ est
\[
  H_k(T^2) \cong
  \begin{cases}
    \Z   & \text{si } k = 0 \text{ ou } k = 2,\\
    \Z^2 & \text{si } k = 1,\\
    0    & \text{si } k \geqslant 3.
  \end{cases}
\]
\end{theoreme}

\begin{proof}
On représente le tore comme un quotient du carré~$[0,1]^2$ par les identifications
habituelles. On choisit la décomposition suivante : soit~$p \in T^2$ un point et posons
\[
  A = T^2 \setminus \{p\}, \qquad B = D
\]
où~$D$ est un petit disque ouvert autour de~$p$. Alors~$B \simeq \ast$ est contractile,
et~$A \cap B = D \setminus \{p\} \simeq S^1$.

L'espace~$A = T^2 \setminus \{p\}$ se rétracte par déformation sur un bouquet
$S^1 \vee S^1$\index{bouquet}, donc~$\widetilde{H}_1(A) \cong \Z^2$,
$\widetilde{H}_k(A) = 0$ pour~$k \neq 1$.

La suite de Mayer-Vietoris réduite donne :
\[
\begin{tikzcd}[column sep=1.4em]
  \widetilde{H}_2(A) \oplus \widetilde{H}_2(B) \ar[r]
  & \widetilde{H}_2(T^2) \ar[r, "\Delta_2"]
  & \widetilde{H}_1(A \cap B) \ar[r, "\Phi_1"]
  & \widetilde{H}_1(A) \oplus \widetilde{H}_1(B)
\end{tikzcd}
\]
soit
\[
  0 \longrightarrow \widetilde{H}_2(T^2) \xrightarrow{\Delta_2} \Z
  \xrightarrow{\Phi_1} \Z^2 \oplus 0.
\]
Le morphisme~$\Phi_1$ envoie le générateur de~$\widetilde{H}_1(S^1) \cong \Z$ sur la
classe de l'inclusion du cercle~$S^1 \simeq \partial D \setminus \{p\}$ dans~$A \simeq
S^1 \vee S^1$. Ce cercle est homotope au commutateur~$[a, b] = aba^{-1}b^{-1}$ des
deux générateurs de~$\pi_1(S^1 \vee S^1)$, qui est trivial en~$H_1$ (car~$H_1$ est
l'abélianisé du groupe fondamental). Ainsi~$\Phi_1 = 0$.

Par exactitude,~$\Delta_2$ est un isomorphisme, donc~$\widetilde{H}_2(T^2) \cong \Z$.

Pour~$H_1$, la suite donne
\[
  0 \to \widetilde{H}_1(A) \oplus \widetilde{H}_1(B) \to \widetilde{H}_1(T^2) \to
  \widetilde{H}_0(A \cap B) \to \cdots
\]
soit~$0 \to \Z^2 \to \widetilde{H}_1(T^2) \to 0$,
d'où~$\widetilde{H}_1(T^2) \cong \Z^2$.
\end{proof}

\subsection{Homologie des surfaces orientables}\label{subsec:homologie-surfaces}
\index{surface!orientable!homologie}
\index{surface de genre~$g$}

\begin{theoreme}[Homologie de~$\Sigma_g$]\label{thm:homologie-sigma-g}
Soit~$\Sigma_g$\index{$\Sigma_g$} la surface compacte orientable de genre~$g \geqslant 0$.
Alors
\[
  H_k(\Sigma_g) \cong
  \begin{cases}
    \Z      & \text{si } k = 0 \text{ ou } k = 2,\\
    \Z^{2g} & \text{si } k = 1,\\
    0       & \text{si } k \geqslant 3.
  \end{cases}
\]
\end{theoreme}

\begin{proof}
On procède par récurrence sur~$g$.

\noindent\textit{Cas~$g = 0$:} $\Sigma_0 = S^2$, et le résultat découle du

ef{thm:homologie-spheres}.

\noindent\textit{Cas~$g = 1$:} $\Sigma_1 = T^2$, traité au 
ef{thm:homologie-tore}.

\noindent\textit{Cas~$g \geqslant 2$:} On décompose~$\Sigma_g$ en retirant un
point~$p$ et en prenant un petit disque~$D$ autour de~$p$. L'espace
$A = \Sigma_g \setminus \{p\}$ se rétracte par déformation sur un bouquet de~$2g$
cercles
\[
  A \simeq \bigvee_{i=1}^{2g} S^1,
\]
obtenu à partir du polygone fondamental à~$4g$ côtés\index{polygone fondamental}
en identifiant un point sur le bord. On a~$B = D \simeq \ast$ et~$A \cap B \simeq S^1$.

La suite de Mayer-Vietoris réduite donne, comme dans le cas du tore :
\[
  0 \to \widetilde{H}_2(\Sigma_g) \xrightarrow{\Delta_2} \Z
  \xrightarrow{\Phi_1} \Z^{2g}.
\]
Le morphisme~$\Phi_1$ est nul car le bord du disque~$D$ représente le
mot de relateur~$[a_1, b_1] \cdots [a_g, b_g]$, qui est trivial en homologie.
Donc~$\widetilde{H}_2(\Sigma_g) \cong \Z$.

De même,~$\widetilde{H}_1(\Sigma_g) \cong \Z^{2g}$ par le reste de la suite.
\end{proof}

\subsection{Homologie de l'espace projectif réel}\label{subsec:homologie-rpn}
\index{espace projectif réel!homologie}
\index{$\RP^n$!homologie}

\begin{theoreme}[Homologie de~$\RP^2$]\label{thm:homologie-rp2}
L'homologie de l'espace projectif réel~$\RP^2$ est donnée par
\[
  H_k(\RP^2) \cong
  \begin{cases}
    \Z      & \text{si } k = 0,\\
    \Z/2\Z  & \text{si } k = 1,\\
    0       & \text{si } k \geqslant 2.
  \end{cases}
\]
\end{theoreme}

\begin{proof}
On utilise la structure CW de~$\RP^2$\index{structure CW} obtenue en identifiant les
points antipodaux du bord de~$D^2$. On décompose~$\RP^2 = A \cup B$ avec :
\begin{itemize}
  \item $B$~un voisinage ouvert du point central (homéomorphe à un disque),
    donc~$B \simeq \ast$~;
  \item $A = \RP^2 \setminus \{p\}$, qui se rétracte par déformation sur le
    cercle~$\RP^1 \cong S^1$~;
  \item $A \cap B \simeq S^1$.
\end{itemize}

La suite de Mayer-Vietoris réduite en basse dimension donne :
\[
  0 \to \widetilde{H}_2(\RP^2) \xrightarrow{\Delta_2} \widetilde{H}_1(S^1)
  \xrightarrow{\Phi_1} \widetilde{H}_1(A) \oplus 0
  \xrightarrow{\Psi_1} \widetilde{H}_1(\RP^2) \to 0.
\]
Ici~$\widetilde{H}_1(S^1) \cong \Z$ et~$\widetilde{H}_1(A) \cong \widetilde{H}_1(S^1) \cong \Z$.
Le morphisme~$\Phi_1$ envoie le générateur du cercle~$A \cap B$ vers~$A \simeq S^1$.
Le cercle~$\partial D$ dans~$\RP^2$ parcourt le cercle~$\RP^1$ \emph{deux fois}
(car l'identification antipodale double le tour), donc~$\Phi_1$ est la
multiplication par~$2$ :
\[
  \Phi_1 \colon \Z \xrightarrow{\times 2} \Z.
\]
Ce morphisme est injectif, donc~$\widetilde{H}_2(\RP^2) = \ker(\Phi_1) = 0$,
soit~$H_2(\RP^2) = 0$. De plus,
$\widetilde{H}_1(\RP^2) \cong \coker(\Phi_1) = \Z/2\Z$.
\end{proof}

\begin{proposition}[Homologie de~$\RP^n$]\label{prop:homologie-rpn}
\index{$\RP^n$!homologie}
Pour~$n \geqslant 1$, l'homologie de~$\RP^n$ à coefficients dans~$\Z$ est
\[
  H_k(\RP^n) \cong
  \begin{cases}
    \Z       & \text{si } k = 0,\\
    \Z       & \text{si } k = n \text{ et } n \text{ est impair},\\
    \Z/2\Z   & \text{si } k \text{ est impair, } 0 < k < n,\\
    0        & \text{sinon.}
  \end{cases}
\]
\end{proposition}

\begin{proof}
Ce résultat se démontre par récurrence via l'homologie cellulaire\index{homologie cellulaire}
de la structure CW standard de~$\RP^n$ (une cellule en chaque dimension~$0, 1, \ldots, n$),
en calculant les opérateurs bord cellulaires.
\end{proof}

\subsection{Homologie de la bouteille de Klein}\label{subsec:homologie-klein}
\index{bouteille de Klein!homologie}

\begin{theoreme}[Homologie de la bouteille de Klein]\label{thm:homologie-klein}
Soit~$K$\index{$K$, bouteille de Klein} la bouteille de Klein. Alors
\[
  H_k(K) \cong
  \begin{cases}
    \Z            & \text{si } k = 0,\\
    \Z \oplus \Z/2\Z & \text{si } k = 1,\\
    0             & \text{si } k \geqslant 2.
  \end{cases}
\]
\end{theoreme}

\begin{proof}
La bouteille de Klein admet une structure CW avec une $0$-cellule, deux $1$-cellules~$a$ et~$b$,
et une $2$-cellule attachée le long du mot~$aba^{-1}b$\index{mot d'attachement}.

L'homologie cellulaire donne les complexes de chaînes :
\[
  0 \to \Z \xrightarrow{d_2} \Z^2 \xrightarrow{d_1} \Z \to 0.
\]
L'opérateur~$d_1$ est nul (car l'espace a une seule $0$-cellule). Pour~$d_2$, le mot
d'attachement~$aba^{-1}b$ fournit :
\[
  d_2(1) = (1 - 1, 1 + 1) = (0, 2) \in \Z^2
\]
(la cellule~$a$ apparaît avec exposants~$+1$ et~$-1$, la cellule~$b$ avec~$+1$ et~$+1$).
Ainsi~$\im(d_2) = \{(0, 2m) : m \in \Z\}$ et
\[
  H_2(K) = \ker(d_2) = 0, \quad
  H_1(K) = \Z^2 / \langle (0, 2) \rangle \cong \Z \oplus \Z/2\Z, \quad
  H_0(K) = \Z.
\]
L'absence de classe fondamentale en degré~$2$ reflète le fait que la bouteille de Klein
est non orientable\index{non orientable}.
\end{proof}

%% ================================================================
\section{Formule de Künneth}\label{sec:kunneth}
\index{Künneth, formule de}
\index{produit d'espaces!formule de Künneth}

Pour calculer l'homologie d'un produit d'espaces, on dispose du puissant outil
qu'est la formule de Künneth. Elle exprime~$H_*(X \times Y)$ en termes de
$H_*(X)$ et~$H_*(Y)$ via le produit tensoriel.

\begin{definition}[Produit tensoriel de complexes]\label{def:produit-tensoriel-complexes}
\index{produit tensoriel!de complexes de chaînes}
Si~$(C_*, \partial^C)$ et~$(D_*, \partial^D)$ sont deux complexes de chaînes de groupes
abéliens, leur \emph{produit tensoriel} est le complexe~$(C \otimes D)_*$ défini par
\[
  (C \otimes D)_n = \bigoplus_{p+q=n} C_p \otimes_\Z D_q,
\]
muni de la différentielle~$\partial(c \otimes d) = \partial^C c \otimes d + (-1)^p c \otimes \partial^D d$
pour~$c \in C_p$.
\end{definition}

\begin{theoreme}[Eilenberg--Zilber]\label{thm:eilenberg-zilber}
\index{Eilenberg--Zilber, théorème d'}
Il existe des équivalences d'homotopie de chaînes, naturelles en~$X$ et~$Y$,
\[
  \mathrm{EZ} \colon C_*(X) \otimes C_*(Y) \rightleftarrows C_*(X \times Y) :\! \mathrm{AW},
\]
appelées respectivement \emph{shuffle map} (Eilenberg--Zilber)\index{shuffle map} et
\emph{Alexander--Whitney map}\index{Alexander--Whitney}. En particulier,
$H_*(C_*(X) \otimes C_*(Y)) \cong H_*(X \times Y)$.
\end{theoreme}

\begin{theoreme}[Formule de Künneth pour l'homologie]\label{thm:kunneth-homologie}
Soient~$X$ et~$Y$ des espaces topologiques. Supposons que les groupes
d'homologie~$H_n(Y;\Z)$ sont libres de type fini pour tout~$n$. Alors il existe un
isomorphisme naturel
\[
  H_n(X \times Y;\Z) \cong
  \bigoplus_{p+q=n} H_p(X;\Z) \otimes_\Z H_q(Y;\Z).
\]
\end{theoreme}

\begin{remarque}\label{rem:kunneth-torsion}
\index{terme de torsion de Künneth}
Lorsque l'hypothèse de liberté n'est pas satisfaite, il faut ajouter un terme de
torsion. La formule complète est :
\[
  0 \to \bigoplus_{p+q=n} H_p(X) \otimes H_q(Y)
  \to H_n(X \times Y)
  \to \bigoplus_{p+q=n-1} \operatorname{Tor}_1^\Z(H_p(X), H_q(Y))
  \to 0,
\]
et cette suite exacte courte est scindée (non naturellement).
\end{remarque}

\begin{exemple}\label{ex:kunneth-tore}
\index{tore!via Künneth}
Retrouvons l'homologie de~$T^2 = S^1 \times S^1$ par Künneth. Puisque~$H_0(S^1) \cong \Z$
et~$H_1(S^1) \cong \Z$ (tous deux libres), on peut appliquer le 
ef{thm:kunneth-homologie} :
\begin{align*}
  H_0(T^2) &\cong H_0(S^1) \otimes H_0(S^1) \cong \Z,\\
  H_1(T^2) &\cong \bigl(H_0(S^1) \otimes H_1(S^1)\bigr)
    \oplus \bigl(H_1(S^1) \otimes H_0(S^1)\bigr) \cong \Z^2,\\
  H_2(T^2) &\cong H_1(S^1) \otimes H_1(S^1) \cong \Z.
\end{align*}
Ce résultat coïncide avec le 
ef{thm:homologie-tore}.
\end{exemple}

\begin{exemple}\label{ex:kunneth-tore-n}
\index{tore!$n$-dimensionnel}
Plus généralement, le tore~$T^n = (S^1)^n$ satisfait
\[
  H_k(T^n) \cong \Z^{\binom{n}{k}},
\]
ce qui se montre par récurrence à l'aide de la formule de Künneth.
\end{exemple}

%% ================================================================
\section{Caractéristique d'Euler}\label{sec:euler}
\index{caractéristique d'Euler}

\begin{definition}[Caractéristique d'Euler]\label{def:euler}
Soit~$X$ un espace topologique dont l'homologie~$H_*(X;\Z)$ est de type fini et
nulle à partir d'un certain rang. La \emph{caractéristique d'Euler} de~$X$ est
l'entier
\[
  \chi(X) = \sum_{n=0}^{\infty} (-1)^n \operatorname{rang} H_n(X;\Z).
\]
\end{definition}

\begin{proposition}[Additivité de la caractéristique d'Euler]\label{prop:euler-additivite}
\index{caractéristique d'Euler!additivité}
Si~$X = A \cup B$ avec~$(A, B)$ paire excisive et~$A$, $B$, $A \cap B$ d'homologie
de type fini, alors
\[
  \chi(X) = \chi(A) + \chi(B) - \chi(A \cap B).
\]
\end{proposition}

\begin{proof}
Par la suite exacte de Mayer-Vietoris, la somme alternée des rangs est nulle dans
toute suite exacte longue de groupes abéliens de type fini (propriété d'Euler--Poincaré\index{Euler--Poincaré}).
On obtient :
\[
  \sum_{n}(-1)^n\bigl[\operatorname{rang} H_n(A \cap B) - \operatorname{rang}(H_n(A)
  \oplus H_n(B)) + \operatorname{rang} H_n(X)\bigr] = 0,
\]
ce qui donne~$\chi(X) - \chi(A) - \chi(B) + \chi(A \cap B) = 0$.
\end{proof}

\begin{proposition}[Multiplicativité]\label{prop:euler-multiplicativite}
\index{caractéristique d'Euler!multiplicativité}
Si~$X$ et~$Y$ sont d'homologie de type fini avec~$H_*(Y)$ libre, alors
\[
  \chi(X \times Y) = \chi(X) \cdot \chi(Y).
\]
\end{proposition}

\begin{proof}
Par la formule de Künneth (
ef{thm:kunneth-homologie}),
\begin{align*}
  \chi(X \times Y) &= \sum_n (-1)^n \operatorname{rang} H_n(X \times Y)
  = \sum_n (-1)^n \sum_{p+q=n} \operatorname{rang}(H_p(X) \otimes H_q(Y))\\
  &= \sum_n (-1)^n \sum_{p+q=n} \operatorname{rang} H_p(X) \cdot \operatorname{rang} H_q(Y)\\
  &= \biggl(\sum_p (-1)^p \operatorname{rang} H_p(X)\biggr)
    \biggl(\sum_q (-1)^q \operatorname{rang} H_q(Y)\biggr) = \chi(X)\,\chi(Y).
    \qedhere
\end{align*}
\end{proof}

\begin{exemple}\label{ex:euler-surfaces}
\index{caractéristique d'Euler!des surfaces}
Pour les surfaces compactes :
\begin{itemize}
  \item $\chi(S^2) = 1 - 0 + 1 = 2$.
  \item $\chi(T^2) = 1 - 2 + 1 = 0$, et plus généralement~$\chi(\Sigma_g) = 2 - 2g$.
  \item $\chi(\RP^2) = 1 - 0 + 0 = 1$ (en utilisant les rangs libres uniquement).
  \item $\chi(K) = 1 - 1 + 0 = 0$ (la bouteille de Klein).
\end{itemize}
\end{exemple}

\begin{exemple}\label{ex:euler-verification}
\index{caractéristique d'Euler!vérification}
Vérifions l'additivité pour le tore. On a~$T^2 = A \cup B$ avec~$A = T^2 \setminus \{p\}
\simeq S^1 \vee S^1$ et~$B = D \simeq \ast$, $A \cap B \simeq S^1$. Alors
\begin{align*}
  \chi(A) &= \chi(S^1 \vee S^1) = 1 - 2 = -1,\\
  \chi(B) &= 1,\\
  \chi(A \cap B) &= \chi(S^1) = 0,
\end{align*}
et~$\chi(A) + \chi(B) - \chi(A \cap B) = -1 + 1 - 0 = 0 = \chi(T^2)$.

Pour la multiplicativité, $\chi(S^1 \times S^1) = \chi(S^1) \cdot \chi(S^1) = 0 \cdot 0 = 0$.
Plus intéressant : $\chi(\Sigma_g \times S^1) = (2-2g) \cdot 0 = 0$ pour tout~$g$.
\end{exemple}

\begin{remarque}\label{rem:euler-cellulaire}
\index{caractéristique d'Euler!et décomposition cellulaire}
Si~$X$ admet une structure CW finie avec~$c_k$ cellules de dimension~$k$, alors
\[
  \chi(X) = \sum_{k \geqslant 0} (-1)^k c_k.
\]
Ce résultat, qui se démontre par récurrence sur le nombre de cellules en utilisant
l'additivité, justifie le lien classique entre la caractéristique d'Euler et la formule
combinatoire~$V - E + F$ d'Euler\index{formule de Euler!$V-E+F$} pour les polyèdres
(où~$V$, $E$, $F$ désignent les nombres de sommets, arêtes et faces).
\end{remarque}

%% ================================================================
\section{Exercices}\label{sec:mv-exercices}

\begin{exercice}[$\bigstar\bigstar$]\label{exo:mv-s2}
\index{sphère!$S^2$}
En utilisant la suite de Mayer-Vietoris avec la décomposition~$S^2 = A \cup B$ où~$A$ et~$B$
sont les hémisphères nord et sud (légèrement agrandis), retrouver directement
l'homologie de~$S^2$.
\end{exercice}

\begin{exercice}[$\bigstar\bigstar$]\label{exo:mv-wedge}
\index{bouquet}
Soient~$X$ et~$Y$ des espaces topologiques bien pointés\index{bien pointé} (c'est-à-dire
que le point de base admet un voisinage contractile). Montrer que
\[
  \widetilde{H}_n(X \vee Y) \cong \widetilde{H}_n(X) \oplus \widetilde{H}_n(Y)
  \qquad \text{pour tout } n.
\]
\emph{Indication :} utiliser Mayer-Vietoris avec des voisinages appropriés.
\end{exercice}

\begin{exercice}[$\bigstar\bigstar$]\label{exo:mv-suspension}
\index{suspension}
Soit~$\Sigma X$ la suspension\index{suspension!homologie} de~$X$. En utilisant
Mayer-Vietoris, montrer que
\[
  \widetilde{H}_n(\Sigma X) \cong \widetilde{H}_{n-1}(X)
  \qquad \text{pour tout } n.
\]
En déduire une autre preuve de l'homologie de~$S^n$ (puisque~$S^n \cong \Sigma S^{n-1}$).
\end{exercice}

\begin{exercice}[$\bigstar\bigstar$]\label{exo:mv-genus2}
\index{surface!de genre~2}
Calculer l'homologie de la surface~$\Sigma_2$ de genre~$2$ directement par Mayer-Vietoris,
en décomposant~$\Sigma_2$ comme union de deux tores épointés recollés le long d'un cercle.
\end{exercice}

\begin{exercice}[$\bigstar\bigstar$]\label{exo:euler-produit}
Vérifier la formule~$\chi(X \times Y) = \chi(X)\chi(Y)$ pour~$X = S^1$, $Y = \Sigma_g$.
\end{exercice}

\begin{exercice}[$\bigstar\bigstar\bigstar$]\label{exo:mv-tore-plein}
\index{tore plein}
Soit~$T = S^1 \times D^2$ le tore plein. Calculer l'homologie de~$T$ et celle de son bord
$\partial T = S^1 \times S^1 = T^2$, puis vérifier la compatibilité avec la suite exacte
de la paire~$(T, \partial T)$.
\end{exercice}

\begin{exercice}[$\bigstar\bigstar\bigstar$]\label{exo:mv-revetement}
\index{revêtement double}
Soit~$p \colon S^n \to \RP^n$ le revêtement double. En utilisant la suite exacte longue
associée à cette application de revêtement et le théorème de transfert, retrouver
l'homologie de~$\RP^n$ pour~$n = 2, 3$.
\end{exercice}

\begin{exercice}[$\bigstar\bigstar\bigstar$]\label{exo:mv-mapping-torus}
\index{tore d'application}
Soit~$f \colon X \to X$ un homéomorphisme et soit~$T_f$ le tore d'application\index{tore d'application}
(mapping torus) :
\[
  T_f = \frac{X \times [0,1]}{(x, 0) \sim (f(x), 1)}.
\]
En décomposant~$T_f$ en deux exemplaires de~$X \times (0,1)$ qui se recouvrent, établir
une suite exacte de Mayer-Vietoris reliant~$H_*(T_f)$ à~$H_*(X)$ et au morphisme
$\Id_* - f_* \colon H_*(X) \to H_*(X)$.
\end{exercice}

\begin{exercice}[$\bigstar\bigstar\bigstar$]\label{exo:mv-klein-mv}
En utilisant une décomposition de la bouteille de Klein~$K$ en bande de Möbius et
cylindre, retrouver~$H_*(K)$ par la suite de Mayer-Vietoris.
\end{exercice}

%% ================================================================
\section*{Résumé du chapitre}
\addcontentsline{toc}{section}{Résumé du chapitre}

\begin{itemize}
  \item La \textbf{suite de Mayer-Vietoris} est une suite exacte longue associée à une
    décomposition~$X = A \cup B$ avec~$X = \mathring{A} \cup \mathring{B}$ :
    \[
      \cdots \to H_n(A \cap B) \xrightarrow{\Phi} H_n(A) \oplus H_n(B) \xrightarrow{\Psi}
      H_n(X) \xrightarrow{\Delta} H_{n-1}(A \cap B) \to \cdots
    \]
  \item Elle se déduit de l'excision via le théorème des petites chaînes.
  \item \textbf{Applications principales} : homologie des sphères ($\widetilde{H}_n(S^n) \cong \Z$),
    du tore ($H_1(T^2) \cong \Z^2$), des surfaces~$\Sigma_g$, des espaces projectifs, de la
    bouteille de Klein.
  \item La \textbf{formule de Künneth} : $H_n(X \times Y) \cong \bigoplus_{p+q=n} H_p(X) \otimes H_q(Y)$
    lorsque~$H_*(Y)$ est libre.
  \item La \textbf{caractéristique d'Euler} $\chi(X) = \sum (-1)^n \operatorname{rang} H_n(X)$
    est additive (via Mayer-Vietoris) et multiplicative (via Künneth).
\end{itemize}

\newpage

%%%%%%%%%%%%%%%%%%%%%%%%%%%%%%%%%%%%%%%%%%%%%%%%%%%%%%%%%%%%%%%%%%%%
%  CHAPITRE 8 — Cohomologie Singulière et Cup-Produit
%%%%%%%%%%%%%%%%%%%%%%%%%%%%%%%%%%%%%%%%%%%%%%%%%%%%%%%%%%%%%%%%%%%%
\chapter{Cohomologie Singulière et Cup-Produit}\label{ch:cohomologie}
\index{cohomologie singulière}
\index{cup-produit}

\section*{Introduction}
\addcontentsline{toc}{section}{Introduction}

Si l'homologie singulière capture des informations fondamentales sur la topologie d'un
espace, la \emph{cohomologie} la enrichit d'une structure algébrique supplémentaire :
le \emph{cup-produit}\index{cup-produit}. Grâce à cette opération, les groupes de
cohomologie~$H^n(X; R)$ s'assemblent en un \emph{anneau gradué}, l'anneau de
cohomologie~$H^*(X; R)$, qui est un invariant strictement plus fin que les groupes
d'homologie seuls.

Ce chapitre développe la théorie de la cohomologie singulière à partir des cochaînes,
établit le théorème des coefficients universels, puis construit le cup-produit et ses
propriétés. Nous calculons l'anneau de cohomologie de plusieurs espaces fondamentaux et
montrons comment cette structure distingue des espaces que l'homologie seule ne peut
séparer.

%% ================================================================
\section{Cochaînes et opérateur cobord}\label{sec:cochaines}
\index{cochaîne singulière}
\index{cobord, opérateur}

\begin{definition}[Groupe de cochaînes singulières]\label{def:cochaines}
Soit~$X$ un espace topologique et~$G$ un groupe abélien. Le \emph{groupe des $n$-cochaînes
singulières} de~$X$ à coefficients dans~$G$ est le groupe
\[
  C^n(X; G) = \Hom(C_n(X), G),
\]
où~$C_n(X)$ est le groupe des $n$-chaînes singulières de~$X$.\index{chaîne singulière}
Un élément~$\varphi \in C^n(X; G)$ est donc un homomorphisme de groupes abéliens
$\varphi \colon C_n(X) \to G$, autrement dit une fonction qui attribue à chaque
$n$-simplexe singulier~$\sigma \colon \Delta^n \to X$ un élément~$\varphi(\sigma) \in G$.
\end{definition}

\begin{definition}[Opérateur cobord]\label{def:cobord}
\index{cobord, opérateur!définition}
L'\emph{opérateur cobord} $\delta^n \colon C^n(X; G) \to C^{n+1}(X; G)$ est le dual
de l'opérateur bord~$\partial_{n+1} \colon C_{n+1}(X) \to C_n(X)$, défini par
\[
  (\delta^n \varphi)(\sigma) = \varphi(\partial_{n+1} \sigma)
\]
pour tout $(n+1)$-simplexe singulier~$\sigma$ et toute $n$-cochaîne~$\varphi$.

De manière explicite, si~$\sigma \colon \Delta^{n+1} \to X$, alors
\[
  (\delta^n \varphi)(\sigma) = \sum_{i=0}^{n+1} (-1)^i \varphi(\sigma \circ d_i),
\]
où~$d_i \colon \Delta^n \to \Delta^{n+1}$ est la $i$-ème coface\index{coface}
(face du simplexe standard).
\end{definition}

\begin{lemme}\label{lem:cobord-carre}
On a~$\delta^{n+1} \circ \delta^n = 0$ pour tout~$n$.
\end{lemme}

\begin{proof}
Pour tout~$\varphi \in C^n(X; G)$ et tout $(n+2)$-simplexe~$\sigma$,
\[
  (\delta^{n+1} \delta^n \varphi)(\sigma)
  = (\delta^n \varphi)(\partial_{n+2} \sigma)
  = \varphi(\partial_{n+1} \partial_{n+2} \sigma) = \varphi(0) = 0,
\]
car~$\partial_{n+1} \circ \partial_{n+2} = 0$.
\end{proof}

Ainsi~$(C^*(X; G), \delta)$ est un complexe de cochaînes\index{complexe de cochaînes} :
\[
  0 \to C^0(X; G) \xrightarrow{\delta^0} C^1(X; G) \xrightarrow{\delta^1}
  C^2(X; G) \xrightarrow{\delta^2} \cdots
\]

%% ================================================================
\section{Groupes de cohomologie}\label{sec:groupes-cohomologie}
\index{cohomologie singulière!groupes de}

\begin{definition}[Groupes de cohomologie singulière]\label{def:cohomologie}
Le $n$-ième \emph{groupe de cohomologie singulière} de~$X$ à coefficients dans~$G$
est le quotient
\[
  H^n(X; G) = \frac{\ker(\delta^n \colon C^n(X;G) \to C^{n+1}(X;G))}
    {\im(\delta^{n-1} \colon C^{n-1}(X;G) \to C^n(X;G))}
  = \frac{Z^n(X; G)}{B^n(X; G)},
\]
où~$Z^n(X; G) = \ker \delta^n$ est le groupe des \emph{$n$-cocycles}\index{cocycle}
et~$B^n(X; G) = \im \delta^{n-1}$ est le groupe des \emph{$n$-cobords}\index{cobord}.
\end{definition}

\begin{remarque}\label{rem:cohomologie-fonctorialite}
\index{fonctorialité!de la cohomologie}
La cohomologie est un foncteur \emph{contravariant} : si~$f \colon X \to Y$ est continue,
on obtient un morphisme induit~$f^* \colon H^n(Y; G) \to H^n(X; G)$ défini par
$f^*(\varphi) = \varphi \circ f_\#$, où~$f_\# \colon C_n(X) \to C_n(Y)$ est le
morphisme de chaînes induit par~$f$.
\end{remarque}

\begin{proposition}\label{prop:cohomologie-homotopie}
\index{invariance par homotopie!cohomologie}
Si~$f, g \colon X \to Y$ sont homotopes, alors~$f^* = g^* \colon H^n(Y; G) \to H^n(X; G)$
pour tout~$n$. En particulier, les groupes de cohomologie sont des invariants d'homotopie.
\end{proposition}

\begin{proof}
C'est le dual de l'invariance par homotopie de l'homologie. L'homotopie de chaînes
$P \colon C_n(X) \to C_{n+1}(Y)$ satisfaisant~$\partial P + P \partial = g_\# - f_\#$
induit, par dualisation,
$P^* \colon C^{n+1}(Y; G) \to C^n(X; G)$ vérifiant~$\delta P^* + P^* \delta = g^* - f^*$
au niveau des cochaînes, ce qui implique~$f^* = g^*$ en cohomologie.
\end{proof}

%% ================================================================
\section{Théorème des coefficients universels}\label{sec:coefficients-universels}
\index{coefficients universels, théorème des!pour la cohomologie}

\begin{theoreme}[Coefficients universels pour la cohomologie]\label{thm:coeff-universels}
Soit~$X$ un espace topologique et~$G$ un groupe abélien. Pour tout~$n \geqslant 0$,
il existe une suite exacte courte naturelle
\[
  0 \longrightarrow \operatorname{Ext}^1_\Z(H_{n-1}(X), G)
  \longrightarrow H^n(X; G)
  \xrightarrow{\;h\;}
  \Hom(H_n(X), G) \longrightarrow 0.
\]
De plus, cette suite est scindée (non naturellement).
\end{theoreme}

\index{Ext@$\operatorname{Ext}$}

\begin{remarque}\label{rem:ext-calcul}
Rappelons les valeurs utiles de~$\operatorname{Ext}^1_\Z$ :
\begin{itemize}
  \item $\operatorname{Ext}^1_\Z(\Z, G) = 0$ pour tout~$G$~;
  \item $\operatorname{Ext}^1_\Z(\Z/m\Z, G) \cong G/mG$ pour tout~$G$~;
  \item $\operatorname{Ext}^1_\Z$ est additif en le premier argument.
\end{itemize}
En particulier, si~$H_*(X)$ est sans torsion, le terme $\operatorname{Ext}$ s'annule
et~$H^n(X; G) \cong \Hom(H_n(X), G)$.
\end{remarque}

\begin{proof}[Esquisse de preuve]
On considère la suite exacte courte de complexes de cochaînes
\[
  0 \to Z_n \hookrightarrow C_n \xrightarrow{\partial} B_{n-1} \to 0,
\]
où~$Z_n = \ker \partial_n$ est le groupe des cycles et~$B_{n-1} = \im \partial_n$
est le groupe des bords. En appliquant le foncteur~$\Hom(-, G)$, on obtient une
suite exacte longue en cohomologie faisant intervenir les foncteurs dérivés
$\operatorname{Ext}^k$. Puisque~$B_{n-1}$ est un sous-groupe de~$C_{n-1}$, qui est
libre, $B_{n-1}$ est libre, et donc~$\operatorname{Ext}^k(B_{n-1}, G) = 0$
pour~$k \geqslant 1$. L'analyse de la suite exacte longue donne alors la suite
exacte courte annoncée. Le scindage provient du fait que~$\Hom(H_n(X), G)$
est un quotient du groupe~$\Hom(Z_n, G)$.
\end{proof}

\begin{corollaire}\label{cor:coeff-univ-Z}
\index{cohomologie!à coefficients entiers}
Si~$H_n(X)$ est libre de type fini pour tout~$n$, alors
\[
  H^n(X; \Z) \cong \Hom(H_n(X), \Z) \cong H_n(X).
\]
\end{corollaire}

\begin{proof}
Le terme $\operatorname{Ext}^1_\Z(H_{n-1}(X), \Z) = 0$ car~$H_{n-1}(X)$ est libre.
De plus, $\Hom(\Z^r, \Z) \cong \Z^r$ pour tout~$r$.
\end{proof}

\begin{corollaire}\label{cor:coeff-univ-corps}
Si~$k$ est un corps, alors
\[
  H^n(X; k) \cong \Hom_k(H_n(X; k), k),
\]
c'est-à-dire que la cohomologie à coefficients dans un corps est le dual de l'homologie.
\end{corollaire}

\begin{exemple}\label{ex:cohomologie-rp2}
\index{$\RP^2$!cohomologie}
Pour~$\RP^2$, on a~$H_0 = \Z$, $H_1 = \Z/2\Z$, $H_k = 0$ pour~$k \geqslant 2$.
Le théorème des coefficients universels donne :
\begin{align*}
  H^0(\RP^2; \Z) &\cong \Hom(\Z, \Z) = \Z,\\
  H^1(\RP^2; \Z) &\cong \Hom(\Z/2\Z, \Z) \oplus \operatorname{Ext}^1(\Z, \Z)
    = 0 \oplus 0 = 0,\\
  H^2(\RP^2; \Z) &\cong \Hom(0, \Z) \oplus \operatorname{Ext}^1(\Z/2\Z, \Z)
    = 0 \oplus \Z/2\Z = \Z/2\Z.
\end{align*}
On observe que les groupes de cohomologie ne coïncident pas avec ceux d'homologie
en présence de torsion.
\end{exemple}

%% ================================================================
\section{Cup-produit}\label{sec:cup-produit}
\index{cup-produit!définition}

Nous arrivons à la structure algébrique centrale de ce chapitre. Pour simplifier la
présentation, nous travaillons avec un anneau commutatif~$R$ comme coefficients.

\begin{definition}[Cup-produit sur les cochaînes]\label{def:cup-cochaines}
Soient~$\varphi \in C^p(X; R)$ et~$\psi \in C^q(X; R)$. Le \emph{cup-produit}
$\varphi \cup \psi \in C^{p+q}(X; R)$ est défini, pour tout $(p+q)$-simplexe
singulier~$\sigma \colon \Delta^{p+q} \to X$, par
\[
  (\varphi \cup \psi)(\sigma)
  = \varphi(\sigma|_{[v_0, \ldots, v_p]}) \cdot \psi(\sigma|_{[v_p, \ldots, v_{p+q}]}),
\]
où~$\sigma|_{[v_{i_0}, \ldots, v_{i_k}]}$ désigne la restriction de~$\sigma$ à la face
engendrée par les sommets~$v_{i_0}, \ldots, v_{i_k}$, et le produit est celui de~$R$.
\end{definition}

\begin{lemme}[Formule de Leibniz]\label{lem:leibniz-cup}
\index{cup-produit!formule de Leibniz}
Pour toutes cochaînes~$\varphi \in C^p(X; R)$ et~$\psi \in C^q(X; R)$, on a
\[
  \delta(\varphi \cup \psi) = (\delta\varphi) \cup \psi
  + (-1)^p \varphi \cup (\delta\psi).
\]
\end{lemme}

\begin{proof}
Soit~$\sigma \colon \Delta^{p+q+1} \to X$ un $(p+q+1)$-simplexe singulier. On a
\begin{align*}
  \delta(\varphi \cup \psi)(\sigma)
  &= (\varphi \cup \psi)(\partial \sigma)
  = \sum_{i=0}^{p+q+1} (-1)^i (\varphi \cup \psi)(\sigma \circ d_i).
\end{align*}
En séparant les termes selon que~$i \leqslant p$ ou~$i > p$ et en réarrangeant, on
retrouve exactement
\[
  (\delta\varphi)(\sigma|_{[v_0,\ldots,v_{p+1}]}) \cdot \psi(\sigma|_{[v_{p+1},\ldots,v_{p+q+1}]})
  + (-1)^p \varphi(\sigma|_{[v_0,\ldots,v_p]}) \cdot (\delta\psi)(\sigma|_{[v_p,\ldots,v_{p+q+1}]}),
\]
ce qui donne la formule annoncée.
\end{proof}

\begin{corollaire}\label{cor:cup-bien-defini}
Le cup-produit passe au quotient et définit une opération bilinéaire
\[
  \cup \colon H^p(X; R) \times H^q(X; R) \longrightarrow H^{p+q}(X; R).
\]
\end{corollaire}

\begin{proof}
Soient~$[\varphi] \in H^p(X; R)$ et~$[\psi] \in H^q(X; R)$ des classes de cocycles,
donc~$\delta\varphi = 0$ et~$\delta\psi = 0$. Par le 
ef{lem:leibniz-cup},
$\delta(\varphi \cup \psi) = 0$, donc~$\varphi \cup \psi$ est un cocycle.

De plus, si~$\varphi = \delta\alpha$, alors
\[
  \varphi \cup \psi = (\delta\alpha) \cup \psi = \delta(\alpha \cup \psi)
  - (-1)^{p-1} \alpha \cup (\delta\psi) = \delta(\alpha \cup \psi),
\]
donc~$[\varphi \cup \psi] = 0$. Le même argument vaut si~$\psi$ est un cobord.
\end{proof}

%% ================================================================
\section{Propriétés du cup-produit}\label{sec:cup-proprietes}

\begin{theoreme}[Propriétés du cup-produit]\label{thm:cup-proprietes}
\index{cup-produit!propriétés}
Le cup-produit possède les propriétés suivantes :
\begin{enumerate}[label=(\roman*)]
  \item \textbf{Associativité.}\index{associativité du cup-produit}
    Pour~$\alpha \in H^p$, $\beta \in H^q$, $\gamma \in H^r$ :
    \[
      (\alpha \cup \beta) \cup \gamma = \alpha \cup (\beta \cup \gamma).
    \]
  \item \textbf{Élément unité.}\index{cup-produit!élément unité}
    La classe~$1 \in H^0(X; R) \cong R$ (envoyant chaque $0$-simplexe sur~$1_R$)
    satisfait
    \[
      1 \cup \alpha = \alpha \cup 1 = \alpha
      \qquad \text{pour tout } \alpha \in H^*(X; R).
    \]
  \item \textbf{Anti-commutativité graduée.}\index{anti-commutativité graduée}
    Pour~$\alpha \in H^p(X; R)$ et~$\beta \in H^q(X; R)$ :
    \[
      \alpha \cup \beta = (-1)^{pq}\, \beta \cup \alpha.
    \]
  \item \textbf{Fonctorialité.}\index{cup-produit!fonctorialité}
    Si~$f \colon X \to Y$ est continue, alors
    \[
      f^*(\alpha \cup \beta) = f^*(\alpha) \cup f^*(\beta)
      \qquad \text{pour tous } \alpha, \beta \in H^*(Y; R).
    \]
\end{enumerate}
\end{theoreme}

\begin{proof}[Esquisse de preuve]
L'associativité~(i) et l'élément unité~(ii) se vérifient directement au niveau des
cochaînes.

Pour l'anti-commutativité graduée~(iii), la preuve au niveau des cochaînes ne donne
\emph{pas} directement le signe~$(-1)^{pq}$, car la définition du cup-produit n'est
pas symétrique (on utilise la face avant et la face arrière du simplexe). On construit
une homotopie de chaînes explicite entre les deux ordres, ce qui requiert l'approximation
d'Eilenberg--Zilber\index{Eilenberg--Zilber, approximation d'} ou la méthode des modèles
acycliques\index{modèles acycliques}. Cette homotopie introduit le signe~$(-1)^{pq}$
en cohomologie.

La fonctorialité~(iv) découle immédiatement de la définition.
\end{proof}

%% ================================================================
\section{L'anneau de cohomologie}\label{sec:anneau-cohomologie}
\index{anneau de cohomologie}

\begin{definition}[Anneau de cohomologie]\label{def:anneau-cohomologie}
Soit~$X$ un espace topologique et~$R$ un anneau commutatif. L'\emph{anneau de cohomologie}
de~$X$ à coefficients dans~$R$ est la $R$-algèbre graduée
\[
  H^*(X; R) = \bigoplus_{n \geqslant 0} H^n(X; R),
\]
munie du cup-produit comme multiplication. C'est une $R$-algèbre graduée
\emph{anti-commutative} (au sens gradué)\index{algèbre graduée anti-commutative} :
\[
  \alpha \cup \beta = (-1)^{|\alpha|\,|\beta|}\, \beta \cup \alpha,
\]
où~$|\alpha|$ et~$|\beta|$ désignent les degrés de~$\alpha$ et~$\beta$.
\end{definition}

\begin{remarque}\label{rem:anneau-vs-groupe}
L'anneau de cohomologie est un invariant \emph{strictement plus fin} que la collection
des groupes de cohomologie. Deux espaces peuvent avoir les mêmes groupes~$H^n$ pour
tout~$n$ mais des structures d'anneau différentes, comme nous le verrons dans les calculs
qui suivent.
\end{remarque}

%% ================================================================
\section{Calculs de l'anneau de cohomologie}\label{sec:calculs-cohomologie}

\subsection{Cohomologie de la sphère~$S^n$}\label{subsec:cohomologie-sphere}
\index{sphère!cohomologie de la}

\begin{theoreme}\label{thm:cohomologie-sphere}
L'anneau de cohomologie de~$S^n$ ($n \geqslant 1$) à coefficients dans~$\Z$ est
\[
  H^*(S^n; \Z) \cong \Z[\alpha]/(\alpha^2),
\]
où~$\alpha \in H^n(S^n; \Z)$ est un générateur, de sorte que~$\alpha^2 = 0$ pour des
raisons de degré (puisque~$H^{2n}(S^n) = 0$ pour~$n \geqslant 1$).

En termes de groupes :~$H^0(S^n) \cong \Z$, $H^n(S^n) \cong \Z$, et le cup-produit
est trivial (c'est-à-dire que tout produit de classes de degré strictement positif est nul).
\end{theoreme}

\begin{proof}
Les groupes de cohomologie se déduisent de l'homologie de~$S^n$ (
ef{thm:homologie-spheres})
et du théorème des coefficients universels. L'homologie étant libre, on a
$H^k(S^n;\Z) \cong \Hom(H_k(S^n), \Z)$. Pour le cup-produit, si~$n \geqslant 1$, le seul
produit non trivial possible serait~$\alpha \cup \alpha \in H^{2n}(S^n) = 0$.
\end{proof}

\subsection{Cohomologie du tore}\label{subsec:cohomologie-tore}
\index{tore!cohomologie du}

\begin{theoreme}[Anneau de cohomologie de~$T^2$]\label{thm:cohomologie-tore}
L'anneau de cohomologie du tore~$T^2$ à coefficients dans~$\Z$ est
\[
  H^*(T^2; \Z) \cong \Lambda_\Z[\alpha, \beta],
\]
l'algèbre extérieure\index{algèbre extérieure} sur deux générateurs~$\alpha, \beta \in H^1(T^2;\Z)$.
Explicitement :
\begin{itemize}
  \item $H^0(T^2;\Z) \cong \Z$, engendré par~$1$~;
  \item $H^1(T^2;\Z) \cong \Z^2$, engendré par~$\alpha$ et~$\beta$~;
  \item $H^2(T^2;\Z) \cong \Z$, engendré par~$\alpha \cup \beta$~;
\end{itemize}
avec les relations~$\alpha^2 = \beta^2 = 0$ et~$\alpha \cup \beta = -\beta \cup \alpha$.
\end{theoreme}

\begin{proof}
Puisque~$T^2 = S^1 \times S^1$ et~$H^*(S^1;\Z) \cong \Z[\gamma]/(\gamma^2)$ avec
$\gamma \in H^1$, la formule de Künneth en cohomologie (
ef{thm:kunneth-cohomologie}
ci-dessous) donne
\[
  H^*(T^2;\Z) \cong H^*(S^1;\Z) \otimes_\Z H^*(S^1;\Z)
\]
en tant qu'algèbre graduée. Posons~$\alpha = \gamma \otimes 1$ et~$\beta = 1 \otimes \gamma$.
Alors~$\alpha, \beta \in H^1$ et
\[
  \alpha \cup \beta = (\gamma \otimes 1) \cup (1 \otimes \gamma) = \gamma \otimes \gamma,
\]
qui est un générateur de~$H^2(T^2;\Z) \cong \Z$. De plus,
$\alpha^2 = \gamma^2 \otimes 1 = 0$ et~$\beta^2 = 1 \otimes \gamma^2 = 0$, et
l'anti-commutativité graduée donne~$\beta \cup \alpha = (-1)^{1 \cdot 1} \alpha \cup \beta
= -\alpha \cup \beta$. C'est bien l'algèbre extérieure sur~$\{\alpha, \beta\}$.
\end{proof}

\subsection{Cohomologie de l'espace projectif réel}\label{subsec:cohomologie-rpn}
\index{$\RP^n$!cohomologie}
\index{anneau de polynômes tronqué}

\begin{theoreme}[Anneau de cohomologie de~$\RP^n$]\label{thm:cohomologie-rpn}
L'anneau de cohomologie de~$\RP^n$ à coefficients dans~$\Z/2\Z$ est
\[
  H^*(\RP^n; \Z/2\Z) \cong (\Z/2\Z)[\alpha]/(\alpha^{n+1}),
\]
où~$\alpha \in H^1(\RP^n; \Z/2\Z)$ est un générateur. Autrement dit, les groupes de
cohomologie sont
\[
  H^k(\RP^n; \Z/2\Z) \cong
  \begin{cases}
    \Z/2\Z & \text{si } 0 \leqslant k \leqslant n,\\
    0      & \text{sinon,}
  \end{cases}
\]
et le cup-produit est donné par~$\alpha^k \cup \alpha^\ell = \alpha^{k+\ell}$ si
$k + \ell \leqslant n$, et~$0$ sinon.
\end{theoreme}

\begin{proof}[Esquisse de preuve]
On procède par récurrence sur~$n$ en utilisant la suite de Gysin\index{Gysin, suite de}
associée au fibré en sphères~$S^0 \to S^n \to \RP^n$, ou bien par la méthode de
l'homologie cellulaire sur la structure CW de~$\RP^n$\index{$\RP^n$!structure CW}
(une cellule en chaque dimension~$0, 1, \ldots, n$). L'identification des
opérateurs cobords cellulaires à coefficients dans~$\Z/2\Z$ montre que tous sont
nuls, d'où~$H^k \cong \Z/2\Z$ pour~$0 \leqslant k \leqslant n$.

La structure multiplicative~$\alpha^k \neq 0$ se démontre par récurrence sur~$k$
en utilisant la naturalité du cup-produit par rapport aux inclusions
$\RP^{n-1} \hookrightarrow \RP^n$.
\end{proof}

\begin{remarque}\label{rem:rpn-Z2-essentiel}
Le choix des coefficients~$\Z/2\Z$ est crucial ici. À coefficients entiers, la
cohomologie de~$\RP^n$ est plus subtile (avec de la $2$-torsion) et la structure
d'anneau est plus complexe.
\end{remarque}

\subsection{Cohomologie de l'espace projectif complexe}\label{subsec:cohomologie-cpn}
\index{$\CP^n$!cohomologie}
\index{espace projectif complexe!cohomologie}

\begin{theoreme}[Anneau de cohomologie de~$\CP^n$]\label{thm:cohomologie-cpn}
L'anneau de cohomologie de~$\CP^n$ à coefficients dans~$\Z$ est
\[
  H^*(\CP^n; \Z) \cong \Z[\beta]/(\beta^{n+1}),
\]
où~$\beta \in H^2(\CP^n; \Z)$ est un générateur. Les groupes de cohomologie sont
\[
  H^k(\CP^n; \Z) \cong
  \begin{cases}
    \Z & \text{si } k \in \{0, 2, 4, \ldots, 2n\},\\
    0  & \text{sinon.}
  \end{cases}
\]
\end{theoreme}

\begin{proof}[Esquisse de preuve]
L'espace~$\CP^n$ possède une structure CW avec une cellule en chaque dimension
paire~$0, 2, 4, \ldots, 2n$. Puisqu'il n'y a pas de cellules en dimensions impaires
adjacentes, tous les opérateurs bords sont nuls et l'homologie cellulaire donne les
groupes annoncés.

Pour la structure multiplicative, on utilise la suite de Gysin\index{Gysin, suite de}
du fibré~$S^1 \to S^{2n+1} \to \CP^n$ et la classe d'Euler\index{classe d'Euler}
$e \in H^2(\CP^n;\Z)$. Le cup-produit avec la classe d'Euler réalise l'isomorphisme
de Gysin~$H^k(\CP^n) \xrightarrow{\cong} H^{k+2}(\CP^n)$ pour~$0 \leqslant k \leqslant 2n-2$,
ce qui force~$\beta^k \neq 0$ pour~$k \leqslant n$ et~$\beta^{n+1} = 0$.
\end{proof}

Le diagramme suivant illustre les anneaux de cohomologie calculés :
\[
\begin{tikzcd}[row sep=1.2em, column sep=0.8em]
  & H^0 & H^1 & H^2 & H^3 & H^4 & \text{Structure} \\
  S^n \ar[dash]{r} & \Z & 0 & \cdots & \Z_{\text{(deg }n\text{)}} & \cdots
    & \Z[\alpha]/(\alpha^2) \\
  T^2 \ar[dash]{r} & \Z & \Z^2 & \Z & 0 &
    & \Lambda_\Z[\alpha,\beta] \\
  \RP^n\; (\Z/2) \ar[dash]{r} & \Z/2 & \Z/2 & \cdots & \Z/2 & \cdots
    & (\Z/2)[\alpha]/(\alpha^{n+1}) \\
  \CP^n \ar[dash]{r} & \Z & 0 & \Z & 0 & \Z
    & \Z[\beta]/(\beta^{n+1})
\end{tikzcd}
\]

%% ================================================================
\section{Cross-product et formule de Künneth en cohomologie}\label{sec:kunneth-cohomologie}
\index{cross-product}
\index{Künneth, formule de!en cohomologie}

\begin{definition}[Cross-product]\label{def:cross-product}
\index{cross-product!définition}
Soient~$X$ et~$Y$ des espaces topologiques. Le \emph{cross-product} (ou produit
extérieur) est le morphisme
\[
  \times \colon H^p(X; R) \otimes_R H^q(Y; R) \longrightarrow H^{p+q}(X \times Y; R)
\]
défini par
\[
  \alpha \times \beta = \pi_X^*(\alpha) \cup \pi_Y^*(\beta),
\]
où~$\pi_X \colon X \times Y \to X$ et~$\pi_Y \colon X \times Y \to Y$ sont les
projections canoniques.
\end{definition}

\begin{theoreme}[Künneth en cohomologie]\label{thm:kunneth-cohomologie}
\index{Künneth, formule de!en cohomologie}
Supposons que~$R$ est un corps, ou que~$R = \Z$ et~$H^*(Y;\Z)$ est sans torsion et
de type fini. Alors le cross-product induit un isomorphisme d'algèbres graduées
\[
  H^*(X; R) \otimes_R H^*(Y; R) \xrightarrow{\;\cong\;} H^*(X \times Y; R),
\]
où le membre de gauche est muni du produit tensoriel gradué\index{produit tensoriel gradué} :
\[
  (\alpha \otimes \beta) \cdot (\alpha' \otimes \beta')
  = (-1)^{|\beta|\,|\alpha'|} (\alpha \cup \alpha') \otimes (\beta \cup \beta').
\]
\end{theoreme}

\begin{remarque}\label{rem:kunneth-corps}
Lorsque~$R = k$ est un corps, le théorème de Künneth en cohomologie est toujours
valable sans condition sur l'homologie de~$Y$.
\end{remarque}

\begin{corollaire}\label{cor:cohomologie-tore-n}
\index{tore!$n$-dimensionnel!cohomologie}
L'anneau de cohomologie du tore~$T^n$ est l'algèbre extérieure
\[
  H^*(T^n; \Z) \cong \Lambda_\Z[\alpha_1, \ldots, \alpha_n],
  \qquad \alpha_i \in H^1, \quad \dim_\Z H^k(T^n;\Z) = \binom{n}{k}.
\]
\end{corollaire}

%% ================================================================
\section{Applications : distinction par le cup-produit}\label{sec:distinction-cup}
\index{cup-produit!applications}

\begin{theoreme}\label{thm:tore-vs-bouquet}
Le tore~$T^2$ et le bouquet~$S^1 \vee S^1 \vee S^2$ ne sont pas homotopiquement
équivalents, bien qu'ils aient les mêmes groupes d'homologie :
\[
  H_0 \cong \Z, \qquad H_1 \cong \Z^2, \qquad H_2 \cong \Z, \qquad H_k = 0 \; (k \geqslant 3).
\]
\end{theoreme}

\begin{proof}
Calculons les anneaux de cohomologie et comparons-les.

\noindent\emph{Pour~$T^2$ :} on a~$H^*(T^2;\Z) \cong \Lambda_\Z[\alpha,\beta]$
(
ef{thm:cohomologie-tore}), et en particulier~$\alpha \cup \beta \neq 0$
dans~$H^2$.

\noindent\emph{Pour~$W = S^1 \vee S^1 \vee S^2$ :} les générateurs de~$H^1(W;\Z) \cong \Z^2$
proviennent des deux facteurs~$S^1$, et le générateur de~$H^2(W;\Z) \cong \Z$
provient du facteur~$S^2$.

Or, si~$\alpha', \beta' \in H^1(W;\Z)$ sont les générateurs provenant des deux cercles,
alors~$\alpha' \cup \beta' = 0 \in H^2(W;\Z)$. En effet, par fonctorialité du
cup-produit, l'inclusion~$S^1 \vee S^1 \hookrightarrow W$ est un facteur par lequel
$\alpha'$ et~$\beta'$ se factorisent, et dans~$H^2(S^1 \vee S^1) = 0$, tout
cup-produit est nul.

Puisque tout cup-produit de classes de degré~$1$ est nul dans~$H^*(W)$ mais pas
dans~$H^*(T^2)$, les deux anneaux de cohomologie ne sont pas isomorphes, et
$T^2 \not\simeq W$.
\end{proof}

Le diagramme suivant illustre la distinction :
\[
\begin{tikzcd}[row sep=2em]
  H^1(T^2) \otimes H^1(T^2) \ar[r, "\cup"] \ar[d, "\cong"']
  & H^2(T^2) \ar[d, "\cong"]\\
  \Z \otimes \Z \ar[r, "\neq 0"]
  & \Z
\end{tikzcd}
\qquad\qquad
\begin{tikzcd}[row sep=2em]
  H^1(W) \otimes H^1(W) \ar[r, "\cup"] \ar[d, "\cong"']
  & H^2(W) \ar[d, "\cong"]\\
  \Z \otimes \Z \ar[r, "= 0"]
  & \Z
\end{tikzcd}
\]

\begin{corollaire}\label{cor:cp2-vs-bouquet}
\index{$\CP^2$!vs bouquet}
De même,~$\CP^2$ n'est pas homotopiquement équivalent à~$S^2 \vee S^4$, bien que
leurs groupes d'homologie coïncident ($H_0 = H_2 = H_4 = \Z$, le reste nul).
En effet, dans~$H^*(\CP^2;\Z) \cong \Z[\beta]/(\beta^3)$, on a~$\beta^2 \neq 0$,
tandis que dans~$H^*(S^2 \vee S^4;\Z)$, le cup-produit de deux classes de degré~$2$
est nul.
\end{corollaire}

\begin{exemple}\label{ex:sigma-g-cohomologie}
\index{surface!orientable!cohomologie}
L'anneau de cohomologie de la surface~$\Sigma_g$ à coefficients dans~$\Z$ est
\[
  H^*(\Sigma_g; \Z) \cong \Lambda_\Z[\alpha_1, \beta_1, \ldots, \alpha_g, \beta_g]
  \Big/ \mathcal{I},
\]
où~$\alpha_i, \beta_i \in H^1$ et~$\mathcal{I}$ est l'idéal engendré par les
relations~$\alpha_i \cup \alpha_j = \beta_i \cup \beta_j = 0$ pour tous~$i, j$
et~$\alpha_i \cup \beta_j = \delta_{ij} \mu$ où~$\mu \in H^2(\Sigma_g;\Z) \cong \Z$
est la classe fondamentale. La forme d'intersection\index{forme d'intersection}
\[
  Q \colon H^1(\Sigma_g;\Z) \times H^1(\Sigma_g;\Z) \to H^2(\Sigma_g;\Z) \cong \Z
\]
est la forme symplectique standard de rang~$2g$.
\end{exemple}

%% ================================================================
\section{Connexion avec la cohomologie de de~Rham}\label{sec:de-rham}
\index{de Rham, cohomologie de}
\index{théorème de de Rham}

\begin{remarque}[Théorème de de~Rham]\label{rem:de-rham}
Pour une variété différentiable~$M$, la cohomologie de de~Rham~$H^*_{\mathrm{dR}}(M)$,
définie comme le quotient des formes différentielles fermées par les formes exactes,
est reliée à la cohomologie singulière par un isomorphisme naturel d'algèbres graduées :
\[
  H^*_{\mathrm{dR}}(M) \cong H^*(M; \R).
\]
Cet isomorphisme est réalisé par l'intégration\index{intégration des formes}
\[
  I \colon \Omega^p(M) \longrightarrow C^p(M; \R), \qquad
  I(\omega)(\sigma) = \int_\sigma \omega = \int_{\Delta^p} \sigma^* \omega,
\]
qui envoie une $p$-forme fermée~$\omega$ sur la $p$-cochaîne singulière qui à un
$p$-simplexe singulier lisse~$\sigma \colon \Delta^p \to M$ associe~$\int_\sigma \omega$.
Le théorème de Stokes\index{Stokes, théorème de} assure que~$I$ est un morphisme de
complexes :
\[
  I(d\omega)(\sigma) = \int_\sigma d\omega = \int_{\partial \sigma} \omega
  = (I(\omega))(\partial \sigma) = (\delta I(\omega))(\sigma),
\]
de sorte que~$I \circ d = \delta \circ I$.
\end{remarque}

\begin{remarque}[Compatibilité des produits]\label{rem:de-rham-produit}
\index{de Rham!compatibilité des produits}
Le produit extérieur de formes~$\omega \wedge \eta$\index{produit extérieur de formes}
correspond au cup-produit~$[\omega] \cup [\eta]$ via l'isomorphisme de de~Rham. Plus
précisément, l'application~$I$ n'est pas strictement multiplicative au niveau des
cochaînes, mais elle l'est à homotopie de chaînes près, ce qui suffit pour obtenir
un isomorphisme d'algèbres graduées en cohomologie :
\[
  I \colon (H^*_{\mathrm{dR}}(M), \wedge)
  \xrightarrow{\;\cong\;} (H^*(M; \R), \cup).
\]
Ce pont entre analyse (formes différentielles) et topologie algébrique (cochaînes
singulières) est l'un des résultats les plus profonds de la géométrie différentielle.
Il fournit notamment des outils de calcul puissants : on peut déterminer la structure
multiplicative de~$H^*(M;\R)$ en calculant des intégrales de formes différentielles.
\end{remarque}

\begin{exemple}\label{ex:de-rham-tore}
\index{de Rham!cohomologie du tore}
Sur le tore~$T^2 = \R^2/\Z^2$, les $1$-formes~$d\theta_1$ et~$d\theta_2$ (projections
des formes coordonnées de~$\R^2$) sont fermées et non exactes. Elles engendrent
$H^1_{\mathrm{dR}}(T^2) \cong \R^2$, et leur produit extérieur
$d\theta_1 \wedge d\theta_2$ engendre~$H^2_{\mathrm{dR}}(T^2) \cong \R$.
Via l'isomorphisme de de~Rham, cela correspond au fait que le cup-produit des
classes~$\alpha$ et~$\beta$ dans~$H^1(T^2;\R)$ est non nul.
\end{exemple}

%% ================================================================
\section{Exercices}\label{sec:cohomologie-exercices}

\begin{exercice}[$\bigstar\bigstar$]\label{exo:cohomologie-s1}
Calculer directement~$H^*(S^1; \Z)$ à partir de la définition des cochaînes et
du cobord, en utilisant la structure simpliciale minimale de~$S^1$ (un sommet, une arête).
\end{exercice}

\begin{exercice}[$\bigstar\bigstar$]\label{exo:coeff-univ-klein}
\index{bouteille de Klein!cohomologie}
En utilisant le théorème des coefficients universels et l'homologie de la bouteille
de Klein, calculer~$H^*(K; \Z)$ et~$H^*(K; \Z/2\Z)$.
\end{exercice}

\begin{exercice}[$\bigstar\bigstar$]\label{exo:cup-tore}
Soit~$T^2$ le tore, représenté comme quotient du carré~$[0,1]^2$. Construire
explicitement des $1$-cocycles simpliciales~$\alpha$ et~$\beta$ (sur une triangulation
du tore) dont le cup-produit~$\alpha \cup \beta$ est un $2$-cocycle non nul.
\end{exercice}

\begin{exercice}[$\bigstar\bigstar$]\label{exo:cup-s1-wedge-s1}
Montrer que tout cup-produit~$H^1(X;\Z) \otimes H^1(X;\Z) \to H^2(X;\Z)$ est nul
lorsque~$X = S^1 \vee S^1$. En déduire que le cup-produit est nul dans~$H^*(X;\Z)$
pour tout bouquet de cercles.
\end{exercice}

\begin{exercice}[$\bigstar\bigstar$]\label{exo:cohomologie-kunneth}
En utilisant la formule de Künneth en cohomologie, calculer l'anneau de cohomologie
de~$S^2 \times S^3$ à coefficients dans~$\Z$ et montrer que le cup-produit est
entièrement trivial en degrés positifs.
\end{exercice}

\begin{exercice}[$\bigstar\bigstar\bigstar$]\label{exo:rp2-coeff-Z}
Montrer que la cohomologie de~$\RP^2$ à coefficients dans~$\Z$ est :
\[
  H^0(\RP^2;\Z) \cong \Z, \quad H^1(\RP^2;\Z) = 0, \quad
  H^2(\RP^2;\Z) \cong \Z/2\Z, \quad H^k(\RP^2;\Z) = 0 \; (k \geqslant 3).
\]
Expliquer pourquoi~$H^1 = 0$ alors que~$H_1 = \Z/2\Z$.
\end{exercice}

\begin{exercice}[$\bigstar\bigstar\bigstar$]\label{exo:rp-infini}
\index{$\RP^\infty$!cohomologie}
Montrer que l'espace projectif réel infini~$\RP^\infty = \varinjlim \RP^n$ satisfait
\[
  H^*(\RP^\infty; \Z/2\Z) \cong (\Z/2\Z)[\alpha],
\]
l'anneau de polynômes en une variable~$\alpha$ de degré~$1$ (sans troncation).
\end{exercice}

\begin{exercice}[$\bigstar\bigstar\bigstar$]\label{exo:cp-infini}
\index{$\CP^\infty$!cohomologie}
Montrer que~$H^*(\CP^\infty; \Z) \cong \Z[\beta]$, l'anneau de polynômes en un
générateur~$\beta$ de degré~$2$.
\emph{Indication :} utiliser~$\CP^\infty = \varinjlim \CP^n$ et la compatibilité
des cup-produits avec les inclusions.
\end{exercice}

\begin{exercice}[$\bigstar\bigstar\bigstar$]\label{exo:distinction-lens}
\index{espace lenticulaire}
Les espaces lenticulaires~$L(5,1)$ et~$L(5,2)$ ont les mêmes groupes d'homologie.
Montrer que leurs anneaux de cohomologie à coefficients dans~$\Z/5\Z$ permettent
de les distinguer.
\emph{Indication :} étudier le cup-produit~$H^1 \otimes H^1 \to H^2$ et identifier
le carré du générateur.
\end{exercice}

\begin{exercice}[$\bigstar\bigstar\bigstar$]\label{exo:hopf}
\index{fibration de Hopf}
Soit~$\eta \colon S^3 \to S^2$ la fibration de Hopf\index{Hopf, fibration de}. Le cône
d'application (mapping cone)~$C_\eta$ a les mêmes groupes d'homologie que~$\CP^2$.
Montrer que l'anneau de cohomologie de~$C_\eta$ est isomorphe à celui de~$\CP^2$.
Comparer avec le cône d'application de l'application constante~$S^3 \to S^2$.
\end{exercice}

%% ================================================================
\section*{Résumé du chapitre}
\addcontentsline{toc}{section}{Résumé du chapitre}

\begin{itemize}
  \item La \textbf{cohomologie singulière} est définie par dualisation :
    $C^n(X;G) = \Hom(C_n(X), G)$, le cobord~$\delta$ est le dual de~$\partial$,
    et~$H^n(X;G) = \ker\delta^n / \im\delta^{n-1}$.
  \item Le \textbf{théorème des coefficients universels} :
    $0 \to \operatorname{Ext}^1(H_{n-1}(X), G) \to H^n(X;G) \to \Hom(H_n(X), G) \to 0$.
  \item Le \textbf{cup-produit} $\cup \colon H^p(X;R) \times H^q(X;R) \to H^{p+q}(X;R)$
    munit~$H^*(X;R)$ d'une structure de $R$-algèbre graduée anti-commutative.
  \item \textbf{Calculs fondamentaux} :
    \begin{itemize}
      \item $H^*(S^n;\Z) \cong \Z[\alpha]/(\alpha^2)$, cup-produit trivial.
      \item $H^*(T^2;\Z) \cong \Lambda_\Z[\alpha,\beta]$, cup-produit non trivial.
      \item $H^*(\RP^n;\Z/2\Z) \cong (\Z/2\Z)[\alpha]/(\alpha^{n+1})$.
      \item $H^*(\CP^n;\Z) \cong \Z[\beta]/(\beta^{n+1})$.
    \end{itemize}
  \item Le cup-produit permet de \textbf{distinguer des espaces} ayant les mêmes groupes
    d'homologie, par exemple~$T^2 \neq S^1 \vee S^1 \vee S^2$ et~$\CP^2 \neq S^2 \vee S^4$.
  \item La \textbf{formule de Künneth en cohomologie} :
    $H^*(X \times Y; R) \cong H^*(X;R) \otimes_R H^*(Y;R)$ comme algèbre graduée
    (sous hypothèses de liberté ou lorsque~$R$ est un corps).
  \item La cohomologie singulière à coefficients réels s'identifie à la \textbf{cohomologie
    de de~Rham} pour les variétés différentiables, le cup-produit correspondant au produit
    extérieur de formes.
\end{itemize}

% === FIN DU CHUNK ch7-8 ===
% === DEBUT DU CHUNK ch9-11+end ===
% ============================================================
%  Topologie Algébrique — Chapitres 9–11 + Matière finale
%  Master / Doctorat, Mathématiques
% ============================================================

% ============================================================
%  CHAPITRE 9 : Dualité de Poincaré
% ============================================================
\chapter{Dualité de Poincaré}\label{chap:poincare}
\index{dualité de Poincaré|(}

La dualité de Poincaré constitue l'un des résultats les plus profonds et les plus
élégants de la topologie algébrique. Elle établit un isomorphisme remarquable entre
la cohomologie et l'homologie d'une variété orientable compacte, mettant en lumière
une symétrie cachée dans la structure algébrique de ces espaces. Ce chapitre développe
la théorie complète de cette dualité, depuis la notion d'orientabilité jusqu'aux
applications en topologie des variétés de dimension~$4$.

% ============================================================
\section{Variétés topologiques}\label{sec:varietes-topo}
\index{variété topologique}

\begin{definition}[Variété topologique]\label{def:variete-topo}
  Une \emph{variété topologique de dimension~$n$}\index{variété topologique!de dimension $n$}
  est un espace topologique~$M$ séparé, à base dénombrable, tel que tout point
  $x \in M$ possède un voisinage ouvert homéomorphe à~$\R^n$. Un tel homéomorphisme
  $\varphi \colon U \to \R^n$ est appelé une \emph{carte}\index{carte} au
  voisinage de~$x$.
\end{definition}

\begin{remarque}\label{rem:variete-localement-compacte}
  Toute variété topologique de dimension~$n$ est localement compacte, localement
  connexe par arcs, et métrisable (par le théorème de Smirnov, puisqu'elle est
  paracompacte et localement métrisable). De plus, elle est localement contractile,
  ce qui garantit que tout revêtement admet un relèvement local.
\end{remarque}

\begin{exemple}\label{ex:varietes-exemples}
  \begin{enumerate}[label=(\roman*)]
    \item La sphère $S^n$\index{sphère} est une variété compacte de dimension~$n$.
    \item Le tore $T^n = (S^1)^n$\index{tore} est une variété compacte de dimension~$n$.
    \item L'espace projectif réel $\RP^n$\index{espace projectif réel} est une variété
          compacte de dimension~$n$.
    \item L'espace projectif complexe $\CP^n$\index{espace projectif complexe} est une
          variété compacte de dimension~$2n$.
    \item Toute surface compacte connexe est une variété de dimension~$2$.
  \end{enumerate}
\end{exemple}

% ============================================================
\section{Orientabilité et classes fondamentales}\label{sec:orientabilite}
\index{orientabilité|(}

Pour définir l'orientabilité, nous devons comprendre la structure de l'homologie
locale d'une variété.

\begin{proposition}[Homologie locale]\label{prop:homologie-locale}
  Soit $M$ une variété topologique de dimension~$n$ et $x \in M$. Alors :
  \[
    H_k(M, M \setminus \set{x}; \Z) \cong
    \begin{cases}
      \Z & \text{si } k = n, \\
      0  & \text{sinon.}
    \end{cases}
  \]
  Un générateur $\mu_x \in H_n(M, M \setminus \set{x}; \Z)$ est appelé une
  \emph{orientation locale}\index{orientation locale} en~$x$.
\end{proposition}

\begin{proof}
  Par excision, si $U$ est un voisinage ouvert de~$x$ homéomorphe à~$\R^n$, on a
  \[
    H_k(M, M \setminus \set{x}) \cong H_k(U, U \setminus \set{x})
    \cong H_k(\R^n, \R^n \setminus \set{0}).
  \]
  Par la suite exacte longue de la paire $(\R^n, \R^n \setminus \set{0})$ et le fait
  que $\R^n$ est contractile, on obtient
  $H_k(\R^n, \R^n \setminus \set{0}) \cong \widetilde{H}_{k-1}(S^{n-1})$,
  d'où le résultat.
\end{proof}

\begin{definition}[Orientation]\label{def:orientation}
  \index{orientation}
  Soit $M$ une variété topologique de dimension~$n$. Une \emph{orientation} de~$M$
  est le choix, pour tout $x \in M$, d'un générateur
  $\mu_x \in H_n(M, M \setminus \set{x}; \Z)$,
  vérifiant la condition de \emph{cohérence locale} suivante : pour tout $x \in M$,
  il existe un voisinage ouvert $U \ni x$ homéomorphe à~$\R^n$ et une classe
  $\mu_U \in H_n(M, M \setminus U; \Z)$
  telle que, pour tout $y \in U$, l'image de~$\mu_U$ par le morphisme naturel
  $H_n(M, M \setminus U) \to H_n(M, M \setminus \set{y})$
  soit~$\mu_y$.
\end{definition}

\begin{definition}[Variété orientable]\label{def:variete-orientable}
  \index{variété!orientable}
  Une variété topologique~$M$ est dite \emph{orientable} si elle admet une orientation.
  Une variété munie d'une orientation est dite \emph{orientée}.
\end{definition}

\begin{exemple}\label{ex:orientabilite}
  \begin{enumerate}[label=(\roman*)]
    \item Les sphères $S^n$ sont orientables pour tout $n \geq 0$.
    \item Le tore $T^n$ est orientable.
    \item L'espace projectif réel $\RP^n$ est orientable si et seulement si $n$ est
          impair.
    \item La bouteille de Klein\index{bouteille de Klein} n'est pas orientable.
    \item Toute variété simplement connexe est orientable.
  \end{enumerate}
\end{exemple}

\begin{theoreme}[Classe fondamentale]\label{thm:classe-fondamentale}
  \index{classe fondamentale}
  Soit $M$ une variété topologique compacte, connexe et orientée de dimension~$n$.
  Alors :
  \begin{enumerate}[label=(\roman*)]
    \item $H_n(M; \Z) \cong \Z$.
    \item Il existe une unique classe $[M] \in H_n(M; \Z)$, appelée
          \emph{classe fondamentale}\index{classe fondamentale} de~$M$, telle que
          pour tout $x \in M$, l'image de~$[M]$ par le morphisme
          $H_n(M; \Z) \to H_n(M, M \setminus \set{x}; \Z)$
          soit l'orientation locale~$\mu_x$.
    \item Si $M$ n'est pas orientable, alors $H_n(M; \Z) = 0$, mais
          $H_n(M; \Z/2\Z) \cong \Z/2\Z$ et il existe une classe fondamentale à
          coefficients dans~$\Z/2\Z$.
  \end{enumerate}
\end{theoreme}

\begin{proof}
  La preuve repose sur une analyse détaillée du faisceau d'orientations
  $\mathcal{O}_M$, dont la fibre en~$x$ est $H_n(M, M \setminus \set{x}; \Z) \cong \Z$.
  Ce faisceau est localement constant de fibre~$\Z$, et $M$ est orientable si et
  seulement si $\mathcal{O}_M$ est trivial. On montre alors que
  $H_n(M; \Z) \cong \Gamma_c(M; \mathcal{O}_M)$
  (sections à support compact du faisceau d'orientations). Lorsque $M$ est orientable
  et compacte, on obtient $H_n(M; \Z) \cong \Z$, et la classe fondamentale est la
  section globale déterminée par l'orientation. Lorsque $M$ n'est pas orientable,
  $\Gamma(M; \mathcal{O}_M) = 0$, d'où $H_n(M; \Z) = 0$. Pour les
  coefficients~$\Z/2\Z$, le faisceau d'orientations est toujours trivial, ce qui
  donne l'assertion~(iii).
\end{proof}

\begin{proposition}[Critère d'orientabilité]\label{prop:critere-orientabilite}
  \index{orientabilité!critère}
  Soit $M$ une variété topologique connexe de dimension~$n$. Les conditions suivantes
  sont équivalentes :
  \begin{enumerate}[label=(\roman*)]
    \item $M$ est orientable.
    \item Le faisceau d'orientations $\mathcal{O}_M$ est trivial.
    \item Le revêtement d'orientations $\widetilde{M} \to M$ (revêtement double
          associé à l'homomorphisme $\pi_1(M) \to \Aut(\Z) = \set{\pm 1}$
          déterminé par l'action de la monodromie sur les orientations locales)
          est trivial.
    \item La première classe de Stiefel--Whitney $w_1(M) \in H^1(M; \Z/2\Z)$
          est nulle.
  \end{enumerate}
\end{proposition}

\begin{proof}
  L'équivalence (i)$\Leftrightarrow$(ii) résulte de la définition : une orientation
  est une section globale non nulle d'un faisceau localement constant de fibre~$\Z$,
  et un tel faisceau admet une section non nulle si et seulement s'il est trivial
  (puisque $\Z$ n'a pas d'automorphisme d'ordre fini préservant un générateur, si ce
  n'est l'identité). L'équivalence (ii)$\Leftrightarrow$(iii) est la correspondance
  classique entre faisceaux localement constants et revêtements. L'équivalence
  (iii)$\Leftrightarrow$(iv) traduit le fait que $w_1(M)$ est précisément la classe
  caractéristique du revêtement double d'orientations, vue dans
  $H^1(M; \Z/2\Z) \cong \Hom(\pi_1(M), \Z/2\Z)$.
\end{proof}

\begin{lemme}[Orientabilité des revêtements]\label{lem:orientabilite-revetement}
  Si $M$ est une variété et $p \colon \widetilde{M} \to M$ est un revêtement fini
  de degré impair, alors $M$ est orientable si et seulement si $\widetilde{M}$ est
  orientable.
\end{lemme}

\index{orientabilité|)}

% ============================================================
\section{Cup-produit et cap-produit}\label{sec:cap-produit}
\index{cup-produit}

Avant de définir le cap-produit, rappelons brièvement le cup-produit qui joue un
rôle central en cohomologie.

\begin{definition}[Cup-produit]\label{def:cup-produit}
  \index{cup-produit}
  Le \emph{cup-produit} est le morphisme bilinéaire
  \[
    \smile \colon H^p(X; R) \times H^q(X; R) \longrightarrow H^{p+q}(X; R)
  \]
  défini au niveau des cochaînes par
  $(\varphi \smile \psi)(\sigma) = \varphi(\sigma \circ [e_0, \ldots, e_p])
  \cdot \psi(\sigma \circ [e_p, \ldots, e_{p+q}])$
  pour $\sigma \colon \Delta^{p+q} \to X$.
\end{definition}

\begin{proposition}[Anneau de cohomologie]\label{prop:anneau-cohomologie}
  \index{anneau de cohomologie}
  Le cup-produit munit $H^*(X; R) = \bigoplus_{n \geq 0} H^n(X; R)$ d'une
  structure d'anneau gradué, associatif et unitaire. De plus, il est
  \emph{gradué-commutatif} :
  \[
    \alpha \smile \beta = (-1)^{pq} \, \beta \smile \alpha
  \]
  pour $\alpha \in H^p(X; R)$ et $\beta \in H^q(X; R)$.
\end{proposition}

\index{cap-produit|(}

Le cap-produit est l'opération bilinéaire fondamentale qui relie la cohomologie
et l'homologie, et qui permet de formuler la dualité de Poincaré.
\index{cap-produit|(}

Le cap-produit est l'opération bilinéaire fondamentale qui relie la cohomologie
et l'homologie, et qui permet de formuler la dualité de Poincaré.

\begin{definition}[Cap-produit]\label{def:cap-produit}
  Soit $R$ un anneau commutatif. Le \emph{cap-produit} est le morphisme bilinéaire
  \[
    \frown \colon H^p(X; R) \times H_n(X; R) \longrightarrow H_{n-p}(X; R)
  \]
  défini au niveau des chaînes et cochaînes comme suit. Soit
  $\varphi \in C^p(X; R)$ et $\sigma \colon \Delta^n \to X$ un $n$-simplexe
  singulier. On pose
  \[
    \varphi \frown \sigma = \varphi(\sigma \circ [e_0, \ldots, e_p]) \cdot
    (\sigma \circ [e_p, \ldots, e_n]),
  \]
  où $[e_0, \ldots, e_p]$ désigne la face avant et $[e_p, \ldots, e_n]$ la face
  arrière. Cette opération passe aux classes de (co)homologie.
\end{definition}

\begin{proposition}[Propriétés du cap-produit]\label{prop:cap-proprietes}
  Le cap-produit satisfait les propriétés suivantes :
  \begin{enumerate}[label=(\roman*)]
    \item \textbf{Associativité :}\index{cap-produit!associativité}
          $(\alpha \smile \beta) \frown \sigma = \alpha \frown (\beta \frown \sigma)$
          pour $\alpha \in H^p$, $\beta \in H^q$, $\sigma \in H_n$ ($p + q \leq n$).
    \item \textbf{Unitarité :} $1 \frown \sigma = \sigma$ où $1 \in H^0(X; R)$ est
          l'élément augmenté.
    \item \textbf{Naturalité :} si $f \colon X \to Y$ est continue, alors
          $f_*(f^*(\alpha) \frown \sigma) = \alpha \frown f_*(\sigma)$.
    \item \textbf{Formule du bord :}
          $\partial(\varphi \frown \sigma) = (-1)^p(\delta\varphi \frown \sigma
          - \varphi \frown \partial\sigma)$ au niveau des chaînes.
  \end{enumerate}
\end{proposition}

\begin{proof}
  Les vérifications sont directes au niveau des chaînes et cochaînes. Pour~(i),
  on utilise la formule d'Alexander--Whitney pour le cup-produit et le cap-produit.
  Pour~(iii), on observe que
  $f^*(\alpha)(\sigma \circ [e_0, \ldots, e_p]) = \alpha(f \circ \sigma \circ [e_0, \ldots, e_p])$.
  Les détails sont laissés en exercice.
\end{proof}

\index{cap-produit|)}

% ============================================================
\section{Théorème de dualité de Poincaré}\label{sec:dualite-poincare}
\index{dualité de Poincaré!théorème}

\begin{theoreme}[Dualité de Poincaré]\label{thm:dualite-poincare}
  Soit $M$ une variété topologique compacte, connexe et orientée de dimension~$n$,
  et soit $[M] \in H_n(M; R)$ l'image de la classe fondamentale. Alors, pour tout
  $0 \leq p \leq n$, le morphisme
  \[
    D_M \colon H^p(M; R) \longrightarrow H_{n-p}(M; R), \qquad
    \alpha \longmapsto \alpha \frown [M]
  \]
  est un isomorphisme.
\end{theoreme}

\begin{remarque}\label{rem:dualite-non-orientable}
  Si $M$ n'est pas orientable, la dualité de Poincaré subsiste à coefficients
  dans~$\Z/2\Z$, en utilisant la classe fondamentale $[M]_2 \in H_n(M; \Z/2\Z)$.
\end{remarque}

\subsection{Esquisse de la preuve}\label{subsec:preuve-poincare}

La démonstration procède par induction sur le nombre de cartes d'un recouvrement
ouvert fini de~$M$, exploitant la suite de Mayer--Vietoris.

\begin{proof}[Preuve (esquisse rigoureuse)]
  \textbf{Étape 1 : Cas $M = \R^n$.}
  Pour une variété ouverte homéomorphe à~$\R^n$, on utilise la cohomologie à support
  compact. L'isomorphisme $H^p_c(\R^n; R) \cong H_{n-p}(\R^n; R)$ se vérifie
  directement : les deux côtés ne sont non nuls que pour $p = n$ (à gauche) et
  $n-p = 0$ (à droite), et l'application cap-produit donne l'isomorphisme
  $R \cong R$.

  \textbf{Étape 2 : Mayer--Vietoris.}
  Supposons que $M = U \cup V$ avec $U, V$ ouverts, et que le théorème vaut pour
  $U$, $V$ et $U \cap V$. Le diagramme suivant, où les lignes horizontales sont les
  suites de Mayer--Vietoris, est commutatif :
  \[
    \begin{tikzcd}[column sep=small, font=\small]
      \cdots \ar[r] & H^p(M) \ar[r] \ar[d, "D_M"] &
      H^p(U) \oplus H^p(V) \ar[r] \ar[d, "D_U \oplus D_V"] &
      H^p(U \cap V) \ar[r] \ar[d, "D_{U\cap V}"] &
      H^{p+1}(M) \ar[r] \ar[d, "D_M"] & \cdots \\
      \cdots \ar[r] & H_{n-p}(M) \ar[r] &
      H_{n-p}(U) \oplus H_{n-p}(V) \ar[r] &
      H_{n-p}(U \cap V) \ar[r] &
      H_{n-p-1}(M) \ar[r] & \cdots
    \end{tikzcd}
  \]
  Par le lemme des cinq, $D_M$ est un isomorphisme.

  \textbf{Étape 3 : Induction.}
  Soit $\set{U_1, \ldots, U_k}$ un recouvrement ouvert fini de~$M$ par des ouverts
  homéomorphes à~$\R^n$ (ou plus généralement tels que toutes les intersections
  finies vérifient la dualité). On pose $A = U_1 \cup \cdots \cup U_{k-1}$ et
  $B = U_k$. Par hypothèse d'induction, la dualité vaut pour~$A$, pour~$B$
  (étape~1) et pour $A \cap B = (U_1 \cap U_k) \cup \cdots \cup (U_{k-1} \cap U_k)$
  (induction avec $k-1$ ouverts). L'étape~2 conclut.

  \textbf{Étape 4 : Compacité.}
  La compacité de~$M$ garantit l'existence d'un recouvrement fini, ce qui rend
  l'induction possible. La preuve complète nécessite en outre de travailler avec la
  cohomologie à support compact pour les ouverts non compacts et de vérifier la
  compatibilité des classes fondamentales dans la suite de Mayer--Vietoris.
\end{proof}

% ============================================================
\section{Applications}\label{sec:applications-poincare}

\subsection{Dualité pour les variétés classiques}\label{subsec:dualite-classiques}

\begin{exemple}[Sphères]\label{ex:dualite-spheres}
  \index{dualité de Poincaré!pour les sphères}
  Pour $S^n$ (orientable, compacte, de dimension~$n$), la dualité donne
  $H^p(S^n; \Z) \cong H_{n-p}(S^n; \Z)$. On retrouve le fait que
  $H^0(S^n) \cong H^n(S^n) \cong \Z$ et $H^p(S^n) = 0$ pour $0 < p < n$.
\end{exemple}

\begin{exemple}[Tores]\label{ex:dualite-tores}
  \index{dualité de Poincaré!pour les tores}
  Pour $T^n = (S^1)^n$ (orientable, compacte, de dimension~$n$), on a
  $H^p(T^n; \Z) \cong \Z^{\binom{n}{p}}$, et la dualité
  $H^p(T^n) \cong H_{n-p}(T^n)$ se traduit par la symétrie
  $\binom{n}{p} = \binom{n}{n-p}$ des coefficients binomiaux.
\end{exemple}

\begin{exemple}[Espaces projectifs réels]\label{ex:dualite-rpn}
  \index{dualité de Poincaré!pour $\RP^n$}
  L'espace $\RP^n$ est orientable si et seulement si $n$ est impair. Pour $n = 2k+1$,
  la dualité de Poincaré à coefficients dans~$\Z/2\Z$ donne
  $H^p(\RP^n; \Z/2\Z) \cong H_{n-p}(\RP^n; \Z/2\Z) \cong \Z/2\Z$
  pour tout $0 \leq p \leq n$.
\end{exemple}

\subsection{Caractéristique d'Euler des variétés de dimension impaire}
\label{subsec:euler-impair}
\index{caractéristique d'Euler!dimension impaire}

\begin{corollaire}\label{cor:euler-dim-impaire}
  Soit $M$ une variété orientable compacte de dimension impaire~$n = 2k+1$.
  Alors la caractéristique d'Euler de~$M$ est nulle : $\chi(M) = 0$.
\end{corollaire}

\begin{proof}
  Par la dualité de Poincaré (à coefficients dans un corps~$\mathbb{F}$),
  $\dim H^p(M; \mathbb{F}) = \dim H^{n-p}(M; \mathbb{F})$, c'est-à-dire
  que les nombres de Betti vérifient $b_p = b_{n-p}$. Donc
  \[
    \chi(M) = \sum_{p=0}^{n} (-1)^p b_p
    = \sum_{p=0}^{n} (-1)^p b_{n-p}
    = (-1)^n \sum_{p=0}^{n} (-1)^{n-p} b_{n-p}
    = (-1)^n \chi(M)
    = -\chi(M),
  \]
  puisque $n$ est impair. Donc $2\chi(M) = 0$, d'où $\chi(M) = 0$ (car
  $\chi(M) \in \Z$).
\end{proof}

\subsection{Forme d'intersection en dimension~$4$}\label{subsec:forme-intersection}
\index{forme d'intersection|(}

\begin{definition}[Forme d'intersection]\label{def:forme-intersection}
  Soit $M$ une variété orientable compacte de dimension~$4$. La \emph{forme
  d'intersection} est la forme bilinéaire symétrique
  \[
    Q_M \colon H^2(M; \Z) \times H^2(M; \Z) \longrightarrow \Z
  \]
  définie par $Q_M(\alpha, \beta) = \langle \alpha \smile \beta, [M] \rangle$,
  où $\langle \cdot, \cdot \rangle$ désigne l'évaluation (kronecker)
  $H^4(M; \Z) \times H_4(M; \Z) \to \Z$.
\end{definition}

\begin{theoreme}\label{thm:forme-intersection-unimod}
  \index{forme d'intersection!unimodularité}
  Si $M$ est une $4$-variété orientable compacte connexe, la forme d'intersection
  $Q_M$ est unimodulaire, c'est-à-dire que l'application
  $H^2(M; \Z)/\mathrm{torsion} \to \Hom(H^2(M; \Z)/\mathrm{torsion}, \Z)$
  induite par $Q_M$ est un isomorphisme.
\end{theoreme}

\begin{proof}
  C'est une conséquence directe de la dualité de Poincaré. Le morphisme cap-produit
  $D_M \colon H^2(M; \Z) \to H_2(M; \Z)$ est un isomorphisme. Modulo la torsion,
  l'identification de $\Hom(H_2(M; \Z)/\mathrm{torsion}, \Z)$ avec
  $H^2(M; \Z)/\mathrm{torsion}$ via le théorème des coefficients universels montre
  que $Q_M$ définit un isomorphisme entre $H^2(M; \Z)/\mathrm{torsion}$ et son dual,
  c'est-à-dire une forme unimodulaire.
\end{proof}

\begin{exemple}\label{ex:formes-intersection}
  \begin{enumerate}[label=(\roman*)]
    \item Pour $M = S^4$ : $H^2(S^4; \Z) = 0$, donc $Q_{S^4} = 0$.
    \item Pour $M = \CP^2$ : $H^2(\CP^2; \Z) \cong \Z$, engendré par la classe
          $\alpha$ d'un hyperplan. On a $Q_{\CP^2} = (1)$ (la forme de rang~$1$
          et de signature~$+1$).
    \item Pour $M = S^2 \times S^2$ : $H^2(M; \Z) \cong \Z^2$, et la forme
          d'intersection est $Q_M = \begin{pmatrix} 0 & 1 \\ 1 & 0 \end{pmatrix}$,
          la forme hyperbolique.
    \item Pour $M = \CP^2 \mathbin{\#} \overline{\CP^2}$ :
          $Q_M = \begin{pmatrix} 1 & 0 \\ 0 & -1 \end{pmatrix}$.
  \end{enumerate}
\end{exemple}

\begin{remarque}\label{rem:donaldson-freedman}
  \index{Donaldson, théorème de}
  \index{Freedman, théorème de}
  La forme d'intersection joue un rôle central dans la classification des
  $4$-variétés. Le théorème de Freedman (1982) classifie les $4$-variétés topologiques
  simplement connexes par leur forme d'intersection. Le théorème de Donaldson (1983)
  impose des contraintes supplémentaires pour les variétés différentielles : si
  $Q_M$ est définie positive, alors elle est diagonalisable sur~$\Z$.
\end{remarque}

\index{forme d'intersection|)}

% ============================================================
\section{Dualité de Poincaré--Lefschetz}\label{sec:poincare-lefschetz}
\index{dualité de Poincaré--Lefschetz}

\begin{theoreme}[Dualité de Poincaré--Lefschetz]\label{thm:poincare-lefschetz}
  Soit $M$ une variété orientable compacte de dimension~$n$, à bord~$\partial M$.
  Alors il existe un isomorphisme
  \[
    D \colon H^p(M, \partial M; R) \xrightarrow{\;\sim\;} H_{n-p}(M; R)
  \]
  donné par le cap-produit avec la classe fondamentale relative
  $[M, \partial M] \in H_n(M, \partial M; R)$. De même, on a un isomorphisme
  \[
    D' \colon H^p(M; R) \xrightarrow{\;\sim\;} H_{n-p}(M, \partial M; R).
  \]
\end{theoreme}

Ces isomorphismes se combinent avec les suites exactes longues de la paire
$(M, \partial M)$ pour produire le diagramme commutatif suivant :
\[
  \begin{tikzcd}[column sep=small, font=\small]
    \cdots \ar[r] & H^p(\partial M) \ar[r] \ar[d, "\cong"', "D_{\partial M}"]
    & H^{p+1}(M, \partial M) \ar[r] \ar[d, "\cong"', "D"]
    & H^{p+1}(M) \ar[r] \ar[d, "\cong"', "D'"]
    & H^{p+1}(\partial M) \ar[r] \ar[d, "\cong"', "D_{\partial M}"]
    & \cdots \\
    \cdots \ar[r] & H_{n-1-p}(\partial M) \ar[r]
    & H_{n-p-1}(M) \ar[r]
    & H_{n-p-1}(M, \partial M) \ar[r]
    & H_{n-2-p}(\partial M) \ar[r]
    & \cdots
  \end{tikzcd}
\]

\begin{exemple}[Dualité pour le disque]\label{ex:dualite-disque}
  Soit $D^2$ le disque fermé, avec $\partial D^2 = S^1$. On a
  $H^0(D^2, S^1; \Z) \cong \Z$ et $H^p(D^2, S^1; \Z) = 0$ pour $p \neq 0, 2$,
  $H^2(D^2, S^1; \Z) \cong \Z$. La dualité donne :
  \begin{align*}
    H^0(D^2, S^1) &\cong H_2(D^2) = 0, \\
    H^1(D^2, S^1) &\cong H_1(D^2) = 0, \\
    H^2(D^2, S^1) &\cong H_0(D^2) \cong \Z.
  \end{align*}
  On vérifie bien que les deux lignes à gauche et à droite sont cohérentes.
\end{exemple}

\begin{proposition}[Dualité et suite exacte longue]\label{prop:dualite-sel}
  Pour une variété à bord $(M, \partial M)$, le diagramme de dualité de
  Poincaré--Lefschetz est compatible avec les suites exactes longues de cohomologie
  et d'homologie de la paire, ce qui fournit un puissant outil de calcul. En
  combinant ce diagramme avec le théorème des coefficients universels, on peut
  souvent déterminer complètement les groupes d'homologie et de cohomologie de
  variétés à bord.
\end{proposition}

\begin{remarque}[Connexion avec la dualité de Hodge]\label{rem:hodge}
  \index{dualité de Hodge}
  Lorsque $M$ est une variété riemannienne compacte orientée de dimension~$n$,
  l'opérateur étoile de Hodge $\star \colon \Omega^p(M) \to \Omega^{n-p}(M)$
  induit un isomorphisme entre les espaces de formes harmoniques
  $\mathcal{H}^p(M) \cong \mathcal{H}^{n-p}(M)$. Par le théorème de Hodge,
  $\mathcal{H}^p(M) \cong H^p_{\mathrm{dR}}(M; \R)$, de sorte que la dualité de
  Hodge fournit une réalisation analytique de la dualité de Poincaré à coefficients
  réels.
\end{remarque}

% ============================================================
\section{Exercices}\label{sec:exo-poincare}

\begin{exercice}\label{exo:poincare-1}
  Soit $M$ une variété compacte orientée de dimension~$n$. Montrer que les nombres
  de Betti satisfont $b_p(M) = b_{n-p}(M)$ pour tout $0 \leq p \leq n$.
\end{exercice}

\begin{exercice}\label{exo:poincare-2}
  Calculer la forme d'intersection $Q_M$ pour $M = \CP^2 \mathbin{\#} \CP^2$
  (somme connexe de deux copies de~$\CP^2$). Quelle est sa signature ?
\end{exercice}

\begin{exercice}\label{exo:poincare-3}
  Montrer que toute variété compacte orientée de dimension~$3$ a
  $b_0 = b_3 = 1$ et $b_1 = b_2$. En déduire que $\chi(M) = 0$.
\end{exercice}

\begin{exercice}\label{exo:poincare-4}
  Soit $\Sigma_g$ la surface orientable compacte de genre~$g$. Vérifier directement
  la dualité de Poincaré : $H^0(\Sigma_g) \cong H_2(\Sigma_g) \cong \Z$,
  $H^1(\Sigma_g) \cong H_1(\Sigma_g) \cong \Z^{2g}$,
  $H^2(\Sigma_g) \cong H_0(\Sigma_g) \cong \Z$.
\end{exercice}

\begin{exercice}\label{exo:poincare-5}
  \index{bouteille de Klein!dualité}
  Soit $K$ la bouteille de Klein. Montrer que la dualité de Poincaré échoue à
  coefficients entiers, mais vaut à coefficients dans~$\Z/2\Z$.
\end{exercice}

\begin{exercice}\label{exo:poincare-6}
  Soit $M^4$ une $4$-variété compacte orientée simplement connexe avec
  $Q_M \cong E_8$ (la forme d'intersection exceptionnelle de rang~$8$).
  Calculer tous les nombres de Betti et la caractéristique d'Euler de~$M$.
\end{exercice}

\begin{exercice}\label{exo:poincare-7}
  Montrer que pour une variété compacte orientée de dimension paire~$n = 2k$, la
  forme bilinéaire $Q \colon H^k(M; \R) \times H^k(M; \R) \to \R$ définie par
  $Q(\alpha, \beta) = \langle \alpha \smile \beta, [M] \rangle$ est symétrique
  si $k$ est pair et antisymétrique si $k$ est impair.
\end{exercice}

\begin{exercice}\label{exo:poincare-8}
  \textbf{(Dualité de Poincaré--Lefschetz.)}
  Soit $D^2$ le disque fermé, vu comme variété à bord $\partial D^2 = S^1$.
  Vérifier la dualité $H^p(D^2, S^1) \cong H_{2-p}(D^2)$ pour tout~$p$.
\end{exercice}

\begin{exercice}\label{exo:poincare-9}
  Soit $M$ une $3$-variété compacte orientée à bord~$\Sigma_g$ (surface de genre~$g$).
  En utilisant la dualité de Poincaré--Lefschetz et la suite exacte longue de la
  paire, montrer que $\mathrm{rang}\, H_1(M; \Z) \geq g$.
\end{exercice}

% ============================================================
\section*{Résumé du chapitre}
\addcontentsline{toc}{section}{Résumé du chapitre}

\begin{itemize}
  \item Une variété topologique de dimension~$n$ est orientable lorsqu'on peut choisir
        de manière cohérente des orientations locales $\mu_x \in H_n(M, M \setminus \set{x})$.
  \item Une variété orientable compacte connexe de dimension~$n$ possède une unique
        classe fondamentale $[M] \in H_n(M; \Z)$.
  \item Le cap-produit $\frown \colon H^p(M) \times H_n(M) \to H_{n-p}(M)$ est
        l'opération clé reliant cohomologie et homologie.
  \item \textbf{Dualité de Poincaré :} $\alpha \mapsto \alpha \frown [M]$ donne
        $H^p(M) \cong H_{n-p}(M)$ pour $M$ variété orientable compacte de dimension~$n$.
  \item La forme d'intersection $Q_M$ sur $H^2(M^4)$ est unimodulaire et joue un
        rôle central dans la classification des $4$-variétés.
  \item La dualité de Poincaré--Lefschetz étend le résultat aux variétés à bord.
\end{itemize}

\index{dualité de Poincaré|)}


% ============================================================
%  CHAPITRE 10 : Groupes d'Homotopie Supérieurs — Introduction
% ============================================================
\chapter{Groupes d'homotopie supérieurs --- Introduction}\label{chap:homotopie-sup}
\index{groupes d'homotopie supérieurs|(}

Le groupe fondamental $\pi_1(X, x_0)$, étudié dans les chapitres précédents,
capture les propriétés de lacets dans un espace. La généralisation naturelle consiste
à remplacer les lacets (images de~$S^1$) par des sphères de dimension
supérieure. Les \emph{groupes d'homotopie supérieurs} $\pi_n(X, x_0)$, introduits
par \v{C}ech en 1932 et développés par Hurewicz, forment une famille d'invariants
algébriques d'une richesse considérable, dont le calcul explicite demeure l'un des
problèmes centraux de la topologie algébrique.

% ============================================================
\section{Définitions et premières propriétés}\label{sec:def-homotopie-sup}

\begin{definition}[Groupes d'homotopie]\label{def:pi-n}
  \index{groupes d'homotopie!définition}
  Soit $(X, x_0)$ un espace topologique pointé et $n \geq 0$ un entier. Le
  \emph{$n$-ième groupe d'homotopie} de~$(X, x_0)$ est l'ensemble
  \[
    \pi_n(X, x_0) = [S^n, X]_0 = \set{f \colon (S^n, s_0) \to (X, x_0)} / {\simeq}
  \]
  des classes d'homotopie d'applications pointées de $(S^n, s_0)$ dans $(X, x_0)$,
  où $s_0 \in S^n$ est un point base fixé.
\end{definition}

\begin{remarque}\label{rem:pi-n-equivalences}
  De manière équivalente, on peut définir $\pi_n(X, x_0)$ comme l'ensemble des
  classes d'homotopie d'applications $(I^n, \partial I^n) \to (X, x_0)$, ou encore
  comme l'ensemble des classes d'homotopie d'applications
  $(D^n, S^{n-1}) \to (X, x_0)$.
\end{remarque}

\begin{definition}[Structure de groupe]\label{def:pi-n-groupe}
  \index{groupes d'homotopie!structure de groupe}
  Pour $n \geq 1$, on munit $\pi_n(X, x_0)$ d'une loi de groupe comme suit. Étant
  données $f, g \colon (I^n, \partial I^n) \to (X, x_0)$, on définit
  $f \cdot g \colon (I^n, \partial I^n) \to (X, x_0)$ par
  \[
    (f \cdot g)(t_1, t_2, \ldots, t_n) =
    \begin{cases}
      f(2t_1, t_2, \ldots, t_n) & \text{si } 0 \leq t_1 \leq \tfrac{1}{2}, \\
      g(2t_1 - 1, t_2, \ldots, t_n) & \text{si } \tfrac{1}{2} \leq t_1 \leq 1.
    \end{cases}
  \]
  L'élément neutre est la classe de l'application constante, et l'inverse de $[f]$ est
  la classe de $\bar{f}(t_1, \ldots, t_n) = f(1-t_1, t_2, \ldots, t_n)$.
\end{definition}

\begin{theoreme}[Commutativité]\label{thm:pi-n-abelien}
  \index{groupes d'homotopie!commutativité}
  \index{argument d'Eckmann--Hilton}
  Pour $n \geq 2$, le groupe $\pi_n(X, x_0)$ est abélien.
\end{theoreme}

\begin{proof}
  La preuve utilise l'\emph{argument d'Eckmann--Hilton}\index{Eckmann--Hilton, argument}.
  Pour $n \geq 2$, on peut définir deux lois de composition sur $\pi_n(X, x_0)$ :
  la concaténation selon la première coordonnée (notée~$\cdot_1$) et la concaténation
  selon la deuxième coordonnée (notée~$\cdot_2$). Ces deux opérations satisfont la
  relation d'échange
  \[
    (f \cdot_1 g) \cdot_2 (h \cdot_1 k) = (f \cdot_2 h) \cdot_1 (g \cdot_2 k)
  \]
  (car la composition dans les deux directions est indépendante). En prenant
  $f = \alpha$, $g = e$, $h = e$, $k = \beta$ (où $e$ est l'élément neutre), on
  obtient
  \[
    \alpha \cdot_2 \beta = (\alpha \cdot_1 e) \cdot_2 (e \cdot_1 \beta)
    = (\alpha \cdot_2 e) \cdot_1 (e \cdot_2 \beta)
    = \alpha \cdot_1 \beta.
  \]
  Ainsi $\cdot_1 = \cdot_2$. En prenant ensuite $f = e$, $g = \beta$,
  $h = \alpha$, $k = e$ :
  \[
    \beta \cdot_1 \alpha = (e \cdot_1 \beta) \cdot_2 (\alpha \cdot_1 e)
    = (e \cdot_2 \alpha) \cdot_1 (\beta \cdot_2 e)
    = \alpha \cdot_1 \beta.
  \]
  Donc la loi de groupe est commutative.
\end{proof}

% ============================================================
\section{Fonctorialité et invariance homotopique}\label{sec:fonctorialite-pi}
\index{groupes d'homotopie!fonctorialité}

\begin{proposition}\label{prop:pi-n-fonctoriel}
  Pour tout $n \geq 0$, l'assignation $(X, x_0) \mapsto \pi_n(X, x_0)$ est un
  foncteur de la catégorie des espaces pointés vers la catégorie des groupes (abéliens
  pour $n \geq 2$, ensembles pointés pour $n = 0$). Plus précisément, toute
  application continue pointée $f \colon (X, x_0) \to (Y, y_0)$ induit un
  homomorphisme $f_* \colon \pi_n(X, x_0) \to \pi_n(Y, y_0)$ par
  $f_*([\alpha]) = [f \circ \alpha]$.
\end{proposition}

\begin{theoreme}[Invariance homotopique]\label{thm:pi-n-inv-homo}
  \index{groupes d'homotopie!invariance homotopique}
  Si $f, g \colon (X, x_0) \to (Y, y_0)$ sont homotopes relativement à~$x_0$,
  alors $f_* = g_* \colon \pi_n(X, x_0) \to \pi_n(Y, y_0)$. En particulier, une
  équivalence d'homotopie pointée induit un isomorphisme sur tous les groupes
  d'homotopie.
\end{theoreme}

\begin{remarque}\label{rem:pi-n-changement-base}
  Pour un espace connexe par arcs, le choix du point base est inoffensif : un chemin
  de $x_0$ à~$x_1$ induit un isomorphisme
  $\pi_n(X, x_0) \cong \pi_n(X, x_1)$. Pour $n \geq 2$, cet isomorphisme est
  canonique (ne dépend que de la classe d'homotopie du chemin), tandis que pour
  $n = 1$, il dépend du chemin à conjugaison près.
\end{remarque}

\begin{proposition}[Groupes d'homotopie et produits]\label{prop:pi-produit}
  Pour tout $n \geq 0$ et tous espaces pointés $(X, x_0)$, $(Y, y_0)$, on a
  un isomorphisme naturel
  \[
    \pi_n(X \times Y, (x_0, y_0)) \cong \pi_n(X, x_0) \times \pi_n(Y, y_0),
  \]
  induit par les projections $\mathrm{pr}_X$ et $\mathrm{pr}_Y$.
\end{proposition}

\begin{proof}
  Une application $f \colon S^n \to X \times Y$ est déterminée par ses composantes
  $\mathrm{pr}_X \circ f$ et $\mathrm{pr}_Y \circ f$, et une homotopie de~$f$
  est déterminée par les homotopies de ses composantes. Donc l'application
  $(\mathrm{pr}_{X*}, \mathrm{pr}_{Y*}) \colon \pi_n(X \times Y) \to
  \pi_n(X) \times \pi_n(Y)$ est bijective. C'est un homomorphisme de groupes
  pour $n \geq 1$ car les projections sont des morphismes de groupes.
\end{proof}

\begin{corollaire}\label{cor:pi-tore}
  Pour le $k$-tore $T^k = (S^1)^k$, on a
  $\pi_1(T^k) \cong \Z^k$ et $\pi_n(T^k) = 0$ pour tout $n \geq 2$.
\end{corollaire}

% ============================================================
\section{Suite exacte longue de la paire}\label{sec:suite-exacte-paire-homotopie}
\index{suite exacte longue!d'une paire (homotopie)}

\begin{definition}[Groupes d'homotopie relatifs]\label{def:pi-n-relatif}
  \index{groupes d'homotopie!relatifs}
  Soit $(X, A, x_0)$ un triplet topologique pointé avec $x_0 \in A \subset X$.
  Pour $n \geq 1$, le \emph{groupe d'homotopie relatif} est
  \[
    \pi_n(X, A, x_0) = \set{f \colon (D^n, S^{n-1}, s_0) \to (X, A, x_0)} / {\simeq}.
  \]
  Pour $n \geq 2$, c'est un groupe ; pour $n \geq 3$, il est abélien.
\end{definition}

\begin{theoreme}[Suite exacte longue]\label{thm:suite-exacte-homotopie}
  \index{suite exacte longue!homotopie}
  Soit $x_0 \in A \subset X$. Il existe une suite exacte longue
  \[
    \cdots \to \pi_n(A, x_0) \xrightarrow{i_*} \pi_n(X, x_0)
    \xrightarrow{j_*} \pi_n(X, A, x_0)
    \xrightarrow{\partial} \pi_{n-1}(A, x_0) \to \cdots
    \to \pi_0(A, x_0) \xrightarrow{i_*} \pi_0(X, x_0)
  \]
  où $i \colon A \hookrightarrow X$ est l'inclusion, $j$ est induit par
  l'inclusion des paires, et $\partial$ est l'opérateur bord (restriction au bord
  de la boule).
\end{theoreme}

\begin{proof}
  L'exactitude en $\pi_n(X)$ signifie que $\ker j_* = \im i_*$, c'est-à-dire que
  $[\alpha] \in \pi_n(X)$ provient de~$A$ si et seulement si $\alpha$ est homotope,
  relativement au bord, à une application à valeurs dans~$A$. L'exactitude en
  $\pi_n(X, A)$ signifie que $\ker \partial = \im j_*$ : la restriction de~$f$
  à~$S^{n-1}$ est homotopiquement triviale dans~$A$ si et seulement si $f$ s'étend
  (à homotopie près) depuis~$X$ tout entier. Les vérifications sont analogues au
  cas $n = 1$, en utilisant les homotopies dans~$I^n$ plutôt que dans~$I$.
\end{proof}

% ============================================================
\section{Suite exacte de la fibration}\label{sec:suite-fibration}
\index{suite exacte longue!d'une fibration}
\index{fibration|(}

\begin{definition}[Fibration de Serre]\label{def:fibration-serre}
  \index{fibration!de Serre}
  Une application continue $p \colon E \to B$ est une \emph{fibration de Serre} si
  elle possède la propriété de relèvement des homotopies pour tous les cubes~$I^n$
  ($n \geq 0$). La fibre au-dessus de $b_0 \in B$ est $F = p^{-1}(b_0)$.
\end{definition}

\begin{theoreme}[Suite exacte d'une fibration]\label{thm:suite-exacte-fibration}
  Soit $F \hookrightarrow E \xrightarrow{p} B$ une fibration de Serre avec $B$
  connexe par arcs. Alors il existe une suite exacte longue
  \[
    \cdots \to \pi_n(F) \xrightarrow{i_*} \pi_n(E) \xrightarrow{p_*} \pi_n(B)
    \xrightarrow{\partial} \pi_{n-1}(F) \to \cdots
    \to \pi_1(B) \xrightarrow{\partial} \pi_0(F) \to \pi_0(E) \to 0.
  \]
\end{theoreme}

\begin{proof}
  On identifie $\pi_n(B, b_0)$ avec $\pi_n(E, F, e_0)$ via la propriété de
  relèvement : toute application $(D^n, S^{n-1}) \to (B, b_0)$ se relève en une
  application $(D^n, S^{n-1}) \to (E, F)$, et cette correspondance passe aux classes
  d'homotopie. Le théorème découle alors de la suite exacte de la paire $(E, F)$.
\end{proof}

\begin{exemple}[Fibration de Hopf]\label{ex:hopf}
  \index{fibration de Hopf}
  La \emph{fibration de Hopf} est la fibration $S^1 \hookrightarrow S^3 \xrightarrow{p} S^2$
  définie par $p(z_0, z_1) = z_0/z_1 \in \CP^1 \cong S^2$ (en coordonnées
  homogènes). La suite exacte longue associée donne
  \[
    \cdots \to \pi_n(S^1) \to \pi_n(S^3) \to \pi_n(S^2) \xrightarrow{\partial}
    \pi_{n-1}(S^1) \to \cdots
  \]
  Puisque $\pi_n(S^1) = 0$ pour $n \geq 2$ (car $\R$ est le revêtement universel
  de~$S^1$), on obtient $\pi_n(S^3) \cong \pi_n(S^2)$ pour $n \geq 3$.
  En particulier, $\pi_3(S^2) \cong \pi_3(S^3) \cong \Z$.
\end{exemple}

\begin{exemple}\label{ex:fibrations-classiques}
  D'autres fibrations classiques incluent :
  \begin{enumerate}[label=(\roman*)]
    \item Le revêtement $\Z \hookrightarrow \R \to S^1$ (fibration triviale à fibre
          discrète).
    \item La fibration $S^0 \hookrightarrow S^n \to \RP^n$ (revêtement double).
    \item La fibration $S^3 \hookrightarrow S^7 \to S^4$ (fibration de Hopf
          quaternionique)\index{fibration de Hopf!quaternionique}.
    \item La fibration des espaces de lacets :
          $\Omega B \hookrightarrow PB \to B$ où $PB$ est l'espace des chemins.
  \end{enumerate}
\end{exemple}

\index{fibration|)}

% ============================================================
\section{Calculs fondamentaux}\label{sec:calculs-pi-n}

\subsection{Groupes d'homotopie des sphères}\label{subsec:homotopie-spheres}
\index{groupes d'homotopie!des sphères}

\begin{theoreme}[Approximation cellulaire]\label{thm:approx-cellulaire}
  \index{approximation cellulaire}
  Toute application continue entre CW-complexes est homotope à une application
  cellulaire (envoyant le $n$-squelette dans le $n$-squelette).
\end{theoreme}

\begin{corollaire}\label{cor:pi-n-spheres-nul}
  Pour $n < m$, on a $\pi_n(S^m) = 0$.
\end{corollaire}

\begin{proof}
  Par le théorème d'approximation cellulaire, toute application $f \colon S^n \to S^m$
  est homotope à une application cellulaire $g \colon S^n \to S^m$. Or, la structure
  CW standard de~$S^m$ consiste en un $0$-squelette réduit à un point et une seule
  $m$-cellule. Puisque $n < m$, l'image de~$g$ est contenue dans le
  $(m-1)$-squelette, qui est réduit à un point. Donc $g$, et donc $f$, est
  homotope à l'application constante.
\end{proof}

\begin{theoreme}\label{thm:pi-n-sn}
  \index{Hurewicz, théorème de!pour les sphères}
  Pour tout $n \geq 1$, on a $\pi_n(S^n) \cong \Z$, engendré par la classe de
  l'identité $[\Id_{S^n}]$.
\end{theoreme}

L'isomorphisme est donné par le \emph{degré}\index{degré d'une application} :
à toute application $f \colon S^n \to S^n$ on associe l'entier $\deg(f) \in \Z$
tel que $f_* \colon H_n(S^n) \to H_n(S^n)$ soit la multiplication par $\deg(f)$.

\begin{remarque}[Degré et homotopie]\label{rem:degre-homotopie}
  \index{degré d'une application!et homotopie}
  Le degré fournit un isomorphisme $\deg \colon \pi_n(S^n) \to \Z$.
  Deux applications continues $f, g \colon S^n \to S^n$ sont homotopes si et
  seulement si $\deg(f) = \deg(g)$. Pour $n = 1$, ceci se démontre par relèvement
  au revêtement universel $\R \to S^1$. Pour $n \geq 2$, on utilise le théorème
  de Hurewicz : l'isomorphisme $\pi_n(S^n) \to H_n(S^n; \Z) \cong \Z$ est
  précisément le degré.

  Rappelons quelques propriétés du degré :
  \begin{enumerate}[label=(\roman*)]
    \item $\deg(\Id_{S^n}) = 1$.
    \item $\deg(f \circ g) = \deg(f) \cdot \deg(g)$.
    \item L'antipode $a \colon S^n \to S^n$, $a(x) = -x$, a pour degré $(-1)^{n+1}$.
    \item Si $f$ n'est pas surjective, $\deg(f) = 0$.
  \end{enumerate}
\end{remarque}

\begin{theoreme}[$\pi_3(S^2) \cong \Z$]\label{thm:pi3-s2}
  Le groupe $\pi_3(S^2)$ est isomorphe à~$\Z$, engendré par la classe de la
  fibration de Hopf $\eta \colon S^3 \to S^2$.
\end{theoreme}

\begin{proof}
  C'est une conséquence de la suite exacte de la fibration de Hopf
  (Exemple~\ref{ex:hopf}) :
  $\pi_3(S^1) \to \pi_3(S^3) \to \pi_3(S^2) \to \pi_2(S^1)$.
  Puisque $\pi_3(S^1) = \pi_2(S^1) = 0$, on obtient
  $\pi_3(S^2) \cong \pi_3(S^3) \cong \Z$.
\end{proof}

\begin{remarque}\label{rem:homotopie-spheres-difficile}
  Le calcul des groupes d'homotopie des sphères $\pi_n(S^m)$ pour $n > m$ est un
  problème extrêmement difficile, loin d'être résolu en toute généralité. Quelques
  résultats remarquables :
  \begin{center}
    \renewcommand{\arraystretch}{1.3}
    \begin{tabular}{c|cccccc}
      $\pi_n(S^m)$ & $m=1$ & $m=2$ & $m=3$ & $m=4$ & $m=5$ & $m=6$ \\ \hline
      $n = 1$ & $\Z$ & $0$ & $0$ & $0$ & $0$ & $0$ \\
      $n = 2$ & $0$ & $\Z$ & $0$ & $0$ & $0$ & $0$ \\
      $n = 3$ & $0$ & $\Z$ & $\Z$ & $0$ & $0$ & $0$ \\
      $n = 4$ & $0$ & $\Z/2$ & $\Z/2$ & $\Z$ & $0$ & $0$ \\
      $n = 5$ & $0$ & $\Z/2$ & $\Z/2$ & $\Z/2$ & $\Z$ & $0$ \\
      $n = 6$ & $0$ & $\Z/{12}$ & $\Z/{12}$ & $\Z/2$ & $\Z/2$ & $\Z$
    \end{tabular}
  \end{center}
\end{remarque}

% ============================================================
\section{Théorème de Hurewicz}\label{sec:hurewicz}
\index{Hurewicz, théorème de|(}

Le théorème de Hurewicz établit un lien fondamental entre les groupes d'homotopie et
les groupes d'homologie.

\begin{definition}[Morphisme de Hurewicz]\label{def:morphisme-hurewicz}
  \index{morphisme de Hurewicz}
  Le \emph{morphisme de Hurewicz} est l'homomorphisme
  \[
    h \colon \pi_n(X, x_0) \longrightarrow H_n(X; \Z)
  \]
  défini par $h([f]) = f_*(\iota_n)$, où $\iota_n \in H_n(S^n; \Z) \cong \Z$ est
  le générateur fondamental.
\end{definition}

\begin{theoreme}[Hurewicz, cas $n = 1$]\label{thm:hurewicz-n1}
  Si $X$ est connexe par arcs, le morphisme de Hurewicz
  $h \colon \pi_1(X, x_0) \to H_1(X; \Z)$
  est surjectif, et son noyau est le sous-groupe des commutateurs
  $[\pi_1, \pi_1]$. Autrement dit,
  \[
    H_1(X; \Z) \cong \pi_1(X, x_0)^{\mathrm{ab}}.
  \]
\end{theoreme}

\begin{proof}
  Pour un CW-complexe, $H_1(X; \Z)$ est engendré par les $1$-cellules modulo les
  relations données par les $2$-cellules. Le morphisme de Hurewicz envoie un lacet
  $\gamma$ sur le cycle singulier correspondant. La surjectivité résulte du fait que
  tout $1$-cycle est une combinaison de lacets. Pour le noyau, si $h([\gamma]) = 0$
  dans $H_1$, le cycle $\gamma$ est un bord, ce qui signifie que $[\gamma]$ est
  un produit de commutateurs dans $\pi_1$. La preuve détaillée utilise la
  présentation de $\pi_1$ par générateurs et relations via le théorème de
  Van Kampen et la comparaison avec la présentation de~$H_1$.
\end{proof}

\begin{theoreme}[Hurewicz, cas général]\label{thm:hurewicz-general}
  \index{Hurewicz, théorème de!cas général}
  Soit $X$ un espace connexe par arcs tel que $\pi_i(X) = 0$ pour tout $1 \leq i < n$
  (on dit que $X$ est \emph{$(n{-}1)$-connexe})\index{espace $(n-1)$-connexe}. Alors
  $H_i(X; \Z) = 0$ pour $1 \leq i < n$, et le morphisme de Hurewicz
  \[
    h \colon \pi_n(X) \xrightarrow{\;\sim\;} H_n(X; \Z)
  \]
  est un isomorphisme.
\end{theoreme}

\begin{proof}[Esquisse de preuve]
  On procède par induction sur~$n$. Le cas $n = 1$ est traité ci-dessus (dans ce
  cas, $0$-connexe signifie connexe par arcs, et $h$ est l'abélianisation, qui est
  un isomorphisme si $\pi_1$ est déjà abélien --- mais ici la $0$-connexité ne
  suffit pas pour un isomorphisme ; on a besoin de $1$-connexité, qui est la
  simple connexité). Pour $n \geq 2$, on utilise la fibration des chemins
  $\Omega X \hookrightarrow PX \to X$, où $PX$ est contractile. La suite exacte
  longue donne $\pi_i(\Omega X) \cong \pi_{i+1}(X)$, de sorte que $\Omega X$ est
  $(n{-}2)$-connexe. Par hypothèse d'induction,
  $\pi_{n-1}(\Omega X) \cong H_{n-1}(\Omega X)$. On conclut par la
  \emph{suspension de l'homologie} (théorème de suspension de Freudenthal) et des
  arguments de théorie spectrale.
\end{proof}

\index{Hurewicz, théorème de|)}

% ============================================================
\section{Théorème de Whitehead}\label{sec:whitehead}
\index{Whitehead, théorème de}

\begin{theoreme}[Whitehead]\label{thm:whitehead}
  Soit $f \colon X \to Y$ une application continue entre CW-complexes connexes.
  Si $f$ induit un isomorphisme $f_* \colon \pi_n(X) \xrightarrow{\sim} \pi_n(Y)$
  pour tout $n \geq 0$, alors $f$ est une équivalence d'homotopie.
\end{theoreme}

\begin{proof}[Esquisse de preuve]
  On construit une inverse homotopique $g \colon Y \to X$ par induction sur les
  cellules de~$Y$. L'hypothèse que $f$ est un isomorphisme sur tous les $\pi_n$
  garantit que chaque étape de la construction est possible : pour relever une
  $n$-cellule de~$Y$, on a besoin de la surjectivité de
  $f_* \colon \pi_{n-1}(X) \to \pi_{n-1}(Y)$ (pour l'existence) et de
  l'injectivité (pour l'unicité à homotopie près). Les détails techniques
  utilisent le cylindre d'application (\emph{mapping cylinder}) et la théorie
  de l'obstruction.
\end{proof}

\begin{remarque}\label{rem:whitehead-cw}
  L'hypothèse que $X$ et $Y$ sont des CW-complexes est essentielle. Il existe des
  espaces topologiques qui ont les mêmes groupes d'homotopie sans être homotopiquement
  équivalents (par exemple, le <<~peigne de Varsovie~>>\index{peigne de Varsovie}
  et le segment).
\end{remarque}

\begin{corollaire}\label{cor:whitehead-contractile}
  Un CW-complexe connexe dont tous les groupes d'homotopie sont triviaux est
  contractile.
\end{corollaire}

\begin{proof}
  L'application constante $X \to \set{*}$ induit un isomorphisme sur tous les
  $\pi_n$ (les deux côtés étant triviaux). Par le théorème de Whitehead, c'est
  une équivalence d'homotopie, donc $X$ est contractile.
\end{proof}

\begin{corollaire}[Équivalences faibles entre CW-complexes]\label{cor:equi-faible}
  Soit $f \colon X \to Y$ une application entre CW-complexes connexes. Les
  conditions suivantes sont équivalentes :
  \begin{enumerate}[label=(\roman*)]
    \item $f$ est une équivalence d'homotopie.
    \item $f_* \colon \pi_n(X) \to \pi_n(Y)$ est un isomorphisme pour tout $n$.
    \item $f_* \colon H_n(X; \Z) \to H_n(Y; \Z)$ est un isomorphisme pour tout $n$,
          et $f$ induit un isomorphisme sur~$\pi_1$.
  \end{enumerate}
\end{corollaire}

\begin{proof}
  (i)$\Rightarrow$(ii) est clair. (ii)$\Rightarrow$(i) est le théorème de Whitehead.
  (ii)$\Rightarrow$(iii) résulte du théorème de Hurewicz appliqué au cylindre
  d'application. (iii)$\Rightarrow$(ii) se démontre par une version relative du
  théorème de Hurewicz : en passant aux revêtements universels et en utilisant
  l'hypothèse sur~$\pi_1$, l'isomorphisme en homologie entraîne l'isomorphisme
  sur les groupes d'homotopie supérieurs.
\end{proof}

\begin{exemple}[Application du théorème de Whitehead]\label{ex:whitehead-application}
  Soit $f \colon S^1 \vee S^3 \to S^1 \times S^3$ l'inclusion naturelle.
  Cette application induit un isomorphisme sur $\pi_1$ (les deux sont isomorphes
  à~$\Z$), mais \emph{pas} sur~$\pi_3$ : en effet, $\pi_3(S^1 \vee S^3)$ contient
  des éléments provenant de l'attachement de Whitehead, tandis que
  $\pi_3(S^1 \times S^3) \cong \pi_3(S^3) \cong \Z$. Cet exemple illustre le fait
  que le théorème de Whitehead requiert un isomorphisme sur \emph{tous} les~$\pi_n$.
\end{exemple}

% ============================================================
\section{Exercices}\label{sec:exo-homotopie-sup}

\begin{exercice}\label{exo:homotopie-1}
  Montrer que $\pi_n(X \times Y) \cong \pi_n(X) \times \pi_n(Y)$ pour tout $n \geq 0$.
\end{exercice}

\begin{exercice}\label{exo:homotopie-2}
  Calculer $\pi_n(T^k)$ pour tout $n \geq 1$ et $k \geq 1$, où $T^k = (S^1)^k$
  est le $k$-tore.
\end{exercice}

\begin{exercice}\label{exo:homotopie-3}
  \index{espace de lacets}
  Montrer que $\pi_n(X, x_0) \cong \pi_{n-1}(\Omega X, c_{x_0})$ pour tout $n \geq 1$,
  où $\Omega X = \set{\gamma \colon [0,1] \to X \mid \gamma(0) = \gamma(1) = x_0}$
  est l'espace de lacets\index{espace de lacets} et $c_{x_0}$ le lacet constant.
\end{exercice}

\begin{exercice}\label{exo:homotopie-4}
  Utiliser la fibration de Hopf quaternionique
  $S^3 \hookrightarrow S^7 \to S^4$
  pour calculer $\pi_n(S^4)$ pour $n \leq 6$.
\end{exercice}

\begin{exercice}\label{exo:homotopie-5}
  Soit $X = S^1 \vee S^2$ le bouquet d'un cercle et d'une $2$-sphère.
  \begin{enumerate}[label=(\alph*)]
    \item Calculer $\pi_1(X)$.
    \item Montrer que le revêtement universel $\widetilde{X}$ est homotopiquement
          équivalent à $\bigvee_{k \in \Z} S^2$ (un bouquet dénombrable de $2$-sphères).
    \item En déduire $\pi_2(X) \cong \Z[\Z] = \bigoplus_{k \in \Z} \Z$ comme
          $\Z[\pi_1(X)]$-module.
  \end{enumerate}
\end{exercice}

\begin{exercice}\label{exo:homotopie-6}
  Donner une preuve directe (sans utiliser la suite exacte de la fibration) que
  $\pi_2(S^2) \cong \Z$ en utilisant le théorème de Hurewicz.
\end{exercice}

\begin{exercice}\label{exo:homotopie-7}
  Soit $f \colon S^3 \to S^2$ la fibration de Hopf. Montrer que $f$ n'est pas
  homotope à l'application constante, c'est-à-dire que $[f] \neq 0$ dans
  $\pi_3(S^2)$.
  \emph{Indication :} considérer l'espace obtenu en attachant une $4$-cellule
  à~$S^2$ via~$f$.
\end{exercice}

\begin{exercice}\label{exo:homotopie-8}
  \index{suspension de Freudenthal}
  \textbf{(Théorème de suspension de Freudenthal.)}
  Le morphisme de suspension $\Sigma \colon \pi_n(S^k) \to \pi_{n+1}(S^{k+1})$
  est un isomorphisme pour $n < 2k - 1$ et un épimorphisme pour $n = 2k - 1$.
  Vérifier ce résultat pour les premiers cas $k = 1, 2, 3$ à l'aide de la table de
  la Remarque~\ref{rem:homotopie-spheres-difficile}.
\end{exercice}

\begin{exercice}\label{exo:homotopie-9}
  Soit $X$ un CW-complexe simplement connexe. Montrer que si $H_n(X; \Z) = 0$ pour
  tout $n \geq 1$, alors $X$ est contractile.
  \emph{Indication :} utiliser le théorème de Hurewicz et le théorème de Whitehead.
\end{exercice}

% ============================================================
\section*{Résumé du chapitre}
\addcontentsline{toc}{section}{Résumé du chapitre}

\begin{itemize}
  \item Le $n$-ième groupe d'homotopie $\pi_n(X, x_0) = [S^n, X]_0$ est un groupe
        pour $n \geq 1$, abélien pour $n \geq 2$ (argument d'Eckmann--Hilton).
  \item La suite exacte longue de la paire
        $\cdots \to \pi_n(A) \to \pi_n(X) \to \pi_n(X, A) \to \pi_{n-1}(A) \to \cdots$
        est l'outil central de calcul.
  \item La suite exacte d'une fibration $F \to E \to B$ relie les groupes d'homotopie
        de la fibre, de l'espace total et de la base.
  \item Résultats clés : $\pi_n(S^m) = 0$ pour $n < m$ ; $\pi_n(S^n) \cong \Z$ ;
        $\pi_3(S^2) \cong \Z$ (fibration de Hopf).
  \item \textbf{Hurewicz :} si $\pi_i(X) = 0$ pour $i < n$, alors
        $\pi_n(X) \cong H_n(X; \Z)$.
  \item \textbf{Whitehead :} une application entre CW-complexes induisant des
        isomorphismes sur tous les $\pi_n$ est une équivalence d'homotopie.
\end{itemize}

\index{groupes d'homotopie supérieurs|)}


% ============================================================
%  CHAPITRE 11 : Théorie de l'Obstruction et Applications
% ============================================================
\chapter{Théorie de l'obstruction et applications}\label{chap:obstruction}
\index{théorie de l'obstruction|(}

La théorie de l'obstruction fournit un cadre systématique pour répondre aux
problèmes d'extension et de relèvement d'applications continues. L'idée
fondamentale est que les obstacles à l'extension d'une application sont
détectés par des classes de cohomologie à valeurs dans les groupes d'homotopie
de l'espace cible. Ce chapitre présente les résultats principaux de cette
théorie, puis en développe les applications, des classes caractéristiques aux
espaces d'Eilenberg--MacLane.

% ============================================================
\section{Le problème d'extension}\label{sec:probleme-extension}
\index{problème d'extension}

\begin{definition}[Problème d'extension]\label{def:probleme-extension}
  Soit $A \subset X$ un sous-espace d'un CW-complexe et $Y$ un espace topologique.
  Étant donnée une application continue $f \colon A \to Y$, le \emph{problème
  d'extension} consiste à déterminer si $f$ admet une extension continue
  $\tilde{f} \colon X \to Y$ telle que $\tilde{f}|_A = f$.
  \[
    \begin{tikzcd}
      A \ar[r, "f"] \ar[d, hookrightarrow] & Y \\
      X \ar[ur, dashed, "\tilde{f}"']
    \end{tikzcd}
  \]
\end{definition}

La stratégie de la théorie de l'obstruction est d'étendre $f$ cellule par cellule.
Supposons que $f$ a été étendue au $n$-squelette $X^{(n)}$. Pour chaque
$(n+1)$-cellule $e_\alpha^{n+1}$ de $X \setminus A$, l'application d'attachement
$\varphi_\alpha \colon S^n \to X^{(n)}$ composée avec l'extension partielle donne
un élément
\[
  [f \circ \varphi_\alpha] \in \pi_n(Y).
\]
L'extension à $e_\alpha^{n+1}$ est possible si et seulement si cet élément est nul.

% ============================================================
\section{Classes d'obstruction}\label{sec:classes-obstruction}
\index{classes d'obstruction}

\begin{definition}[Cochaîne d'obstruction]\label{def:cochaine-obstruction}
  Soit $f^{(n)} \colon X^{(n)} \cup A \to Y$ une extension partielle de~$f$
  au $n$-squelette. La \emph{cochaîne d'obstruction} est l'élément
  \[
    \mathfrak{o}^{n+1}(f^{(n)}) \in C^{n+1}(X, A; \pi_n(Y))
  \]
  défini sur chaque $(n+1)$-cellule orientée $e_\alpha^{n+1}$ par
  $\mathfrak{o}^{n+1}(f^{(n)})(e_\alpha^{n+1}) = [f^{(n)} \circ \varphi_\alpha]
  \in \pi_n(Y)$.
\end{definition}

\begin{theoreme}\label{thm:obstruction-cocycle}
  La cochaîne d'obstruction $\mathfrak{o}^{n+1}(f^{(n)})$ est un cocycle :
  $\delta \mathfrak{o}^{n+1}(f^{(n)}) = 0$ dans
  $C^{n+2}(X, A; \pi_n(Y))$.
\end{theoreme}

\begin{proof}
  Pour toute $(n+2)$-cellule $e^{n+2}$, le bord $\partial e^{n+2}$ est une
  $(n+1)$-sphère, et l'image de $\mathfrak{o}^{n+1}$ sur ce bord est la classe
  d'homotopie de la restriction de $f^{(n)}$ à la sphère $\partial e^{n+2}$. Or
  cette restriction s'étend au disque $D^{n+2}$ (qui est contractile), donc sa
  classe dans $\pi_{n+1}$ est triviale. Le calcul de $\delta$ reproduit exactement
  cette situation.
\end{proof}

\begin{definition}[Classe d'obstruction]\label{def:classe-obstruction}
  La \emph{classe d'obstruction} est la classe de cohomologie
  \[
    [\mathfrak{o}^{n+1}(f^{(n)})] \in H^{n+1}(X, A; \pi_n(Y)).
  \]
  Elle ne dépend pas du choix de l'extension $f^{(n)}$ au $n$-squelette (parmi
  toutes celles qui coïncident avec une extension fixée au $(n{-}1)$-squelette).
\end{definition}

\begin{theoreme}[Théorème d'obstruction principal]\label{thm:obstruction-principal}
  \index{théorème d'obstruction principal}
  Soit $A \subset X$ une paire CW et $Y$ un espace connexe par arcs. Soit
  $f \colon A \to Y$ une application continue, et supposons que $f$ a été étendue
  en $f^{(n)} \colon X^{(n)} \cup A \to Y$.
  \begin{enumerate}[label=(\roman*)]
    \item L'extension $f^{(n)}$ se prolonge au $(n{+}1)$-squelette si et seulement si
          $\mathfrak{o}^{n+1}(f^{(n)}) = 0$ dans $C^{n+1}(X, A; \pi_n(Y))$.
    \item Il existe une extension $g^{(n)}$ de $f$ au $n$-squelette, coïncidant avec
          $f^{(n)}$ sur le $(n{-}1)$-squelette, et se prolongeant au
          $(n{+}1)$-squelette, si et seulement si la classe de cohomologie
          $[\mathfrak{o}^{n+1}(f^{(n)})] = 0$ dans $H^{n+1}(X, A; \pi_n(Y))$.
    \item En particulier, si $H^{n+1}(X, A; \pi_n(Y)) = 0$ pour tout $n$, toute
          application $f \colon A \to Y$ s'étend à~$X$.
  \end{enumerate}
\end{theoreme}

\begin{remarque}\label{rem:obstruction-itere}
  Dans la pratique, on rencontre des \emph{obstructions successives} :
  même si la première obstruction s'annule, il peut exister des obstructions
  dans les degrés de cohomologie supérieurs. L'analyse complète nécessite
  parfois l'utilisation de suites spectrales ou de la théorie de Postnikov.
\end{remarque}

\begin{corollaire}\label{cor:extension-n-connexe}
  Soit $X$ un CW-complexe de dimension~$\leq n$ et $Y$ un espace $(n{-}1)$-connexe
  (c'est-à-dire $\pi_k(Y) = 0$ pour $k < n$). Alors toute application continue
  $f \colon A \to Y$ définie sur un sous-complexe $A \subset X$ admet une extension
  $\tilde{f} \colon X \to Y$.
\end{corollaire}

\begin{proof}
  L'extension de $f$ au $(k{+}1)$-squelette rencontre une obstruction dans
  $H^{k+1}(X, A; \pi_k(Y))$. Pour $k < n$, $\pi_k(Y) = 0$, donc ces obstructions
  s'annulent trivialement. Pour $k \geq n$, le groupe $H^{k+1}(X, A; \pi_k(Y)) = 0$
  puisque $\dim X \leq n \leq k$. Donc toutes les obstructions sont nulles et
  l'extension existe.
\end{proof}

\begin{proposition}[Unicité de l'extension]\label{prop:unicite-extension}
  \index{problème d'extension!unicité}
  Dans le cadre du théorème d'obstruction, l'\emph{obstruction à l'unicité} de
  l'extension au $n$-squelette (c'est-à-dire l'obstruction à ce que deux extensions
  soient homotopes) vit dans $H^n(X, A; \pi_n(Y))$. Ainsi, si
  $H^n(X, A; \pi_n(Y)) = 0$ pour tout~$n$, l'extension est unique à homotopie
  près.
\end{proposition}

\begin{exemple}[Extension des applications vers les sphères]\label{ex:extension-spheres}
  Considérons le problème d'étendre une application $f \colon S^n \to S^n$ à la boule
  $D^{n+1}$. L'unique obstruction vit dans
  $H^{n+1}(D^{n+1}, S^n; \pi_n(S^n)) \cong \pi_n(S^n) \cong \Z$,
  et elle est précisément le degré de~$f$. Ainsi, $f$ s'étend à $D^{n+1}$ si et
  seulement si $\deg(f) = 0$.
\end{exemple}

% ============================================================
\section{Applications classiques}\label{sec:applications-obstruction}

\subsection{Théorème de Borsuk--Ulam}\label{subsec:borsuk-ulam}
\index{Borsuk--Ulam, théorème de}

\begin{theoreme}[Borsuk--Ulam]\label{thm:borsuk-ulam}
  Pour toute application continue $f \colon S^n \to \R^n$, il existe un point
  $x \in S^n$ tel que $f(x) = f(-x)$.
\end{theoreme}

\begin{proof}
  Supposons par l'absurde que $f(x) \neq f(-x)$ pour tout~$x$. Alors l'application
  $g \colon S^n \to S^{n-1}$ définie par
  \[
    g(x) = \frac{f(x) - f(-x)}{\abs{f(x) - f(-x)}}
  \]
  est bien définie, continue, et vérifie $g(-x) = -g(x)$ (elle est impaire).
  Une application impaire $S^n \to S^{n-1}$ induit un morphisme non trivial en
  cohomologie modulo~$2$ : la classe $g^*(u) = v$ où $u$ engendre
  $H^1(\RP^{n-1}; \Z/2\Z)$ et $v$ engendre $H^1(\RP^n; \Z/2\Z)$. Mais alors
  $v^n = g^*(u^n) \neq 0$, ce qui est impossible car $u^n = 0$ dans
  $H^n(\RP^{n-1}; \Z/2\Z)$ (puisque $\RP^{n-1}$ est de dimension~$n-1$).
  Contradiction.
\end{proof}

\subsection{Théorème de Lusternik--Schnirelman}\label{subsec:lusternik-schnirelman}
\index{Lusternik--Schnirelman, théorème de}

\begin{theoreme}[Lusternik--Schnirelman]\label{thm:lusternik-schnirelman}
  Si $S^n = F_1 \cup F_2 \cup \cdots \cup F_{n+1}$, où les $F_i$ sont des fermés,
  alors au moins un $F_i$ contient une paire de points antipodaux $(x, -x)$.
\end{theoreme}

\begin{proof}
  Supposons que chaque $F_i$ ne contient aucune paire antipodale. Posons
  $U_i = S^n \setminus F_i$. Alors $\set{U_1, \ldots, U_{n+1}}$ est un recouvrement
  ouvert de~$S^n$ (car $F_1 \cup \cdots \cup F_{n+1} = S^n$ implique que les $U_i$
  ne recouvrent pas, contradiction). On rectifie : on définit
  $g_i(x) = d(x, F_i)$ et $g(x) = (g_1(x), \ldots, g_n(x)) \in \R^n$, où l'on
  utilise seulement les $n$ premiers $F_i$. Par le théorème de Borsuk--Ulam, il
  existe $x$ tel que $g(x) = g(-x)$, c'est-à-dire $d(x, F_i) = d(-x, F_i)$ pour
  $i = 1, \ldots, n$. Puisque ni $x$ ni $-x$ n'est dans le même $F_i$ et que
  $x \in S^n = \bigcup F_i$, on a $x \in F_{n+1}$, et de même $-x \in F_{n+1}$,
  ce qui contredit l'hypothèse sur $F_{n+1}$.
\end{proof}

\subsection{Classes caractéristiques}\label{subsec:classes-caract}
\index{classes caractéristiques|(}

\begin{remarque}\label{rem:classes-caract}
  La théorie de l'obstruction fournit un cadre naturel pour la classification des
  fibrés vectoriels. Étant donné un fibré vectoriel réel $\xi$ de rang~$k$ sur un
  CW-complexe~$X$, les obstructions à la trivialisation de~$\xi$ vivent dans les
  groupes $H^n(X; \pi_{n-1}(\mathrm{GL}_k(\R)))$, ou de manière équivalente,
  dans $H^n(X; \pi_{n-1}(O(k)))$ (en choisissant une métrique riemannienne).
\end{remarque}

Les principales classes caractéristiques sont :

\begin{enumerate}[label=(\roman*)]
  \item \textbf{Classes de Stiefel--Whitney}\index{classes de Stiefel--Whitney}
        $w_i(\xi) \in H^i(X; \Z/2\Z)$ pour les fibrés vectoriels réels. La classe
        $w_1$ mesure l'orientabilité ($w_1 = 0$ si et seulement si $\xi$ est
        orientable), et $w_2$ est l'obstruction à l'existence d'une structure spin.

  \item \textbf{Classes de Chern}\index{classes de Chern}
        $c_i(\xi) \in H^{2i}(X; \Z)$ pour les fibrés vectoriels complexes.
        La première classe de Chern $c_1$ classifie les fibrés en droites complexes.

  \item \textbf{Classes de Pontriaguine}\index{classes de Pontriaguine}
        $p_i(\xi) \in H^{4i}(X; \Z)$ pour les fibrés vectoriels réels, définies par
        $p_i(\xi) = (-1)^i c_{2i}(\xi \otimes_{\R} \C)$.

  \item \textbf{Classe d'Euler}\index{classe d'Euler}
        $e(\xi) \in H^k(X; \Z)$ pour un fibré orienté de rang~$k$, reliée à la
        caractéristique d'Euler par $\langle e(TM), [M] \rangle = \chi(M)$
        (théorème de Gauss--Bonnet--Chern).
\end{enumerate}

\begin{proposition}[Axiomes des classes de Chern]\label{prop:axiomes-chern}
  \index{classes de Chern!axiomes}
  Les classes de Chern sont caractérisées de manière unique par les axiomes suivants :
  \begin{enumerate}[label=(\roman*)]
    \item \textbf{Naturalité :} $c_i(f^*\xi) = f^*(c_i(\xi))$ pour toute application
          continue $f$.
    \item \textbf{Formule de Whitney :}\index{formule de Whitney}
          $c(\xi \oplus \eta) = c(\xi) \smile c(\eta)$, où
          $c(\xi) = 1 + c_1(\xi) + c_2(\xi) + \cdots$ est la \emph{classe de Chern
          totale}.
    \item \textbf{Normalisation :} pour le fibré tautologique $\gamma^1$ sur
          $\CP^1 \cong S^2$, on a $\langle c_1(\gamma^1), [\CP^1] \rangle = -1$
          (ou $+1$, selon la convention).
    \item \textbf{Dimension :} $c_i(\xi) = 0$ pour $i > \mathrm{rang}(\xi)$.
  \end{enumerate}
\end{proposition}

\begin{exemple}\label{ex:classes-chern-calculs}
  \begin{enumerate}[label=(\roman*)]
    \item Le fibré tangent $T\CP^n$ satisfait $c(T\CP^n) = (1 + h)^{n+1}$
          dans $H^*(\CP^n; \Z) \cong \Z[h]/(h^{n+1})$, où $h$ est le générateur
          de $H^2(\CP^n; \Z)$.
    \item Pour le fibré de Hopf $\eta$ sur $S^2 \cong \CP^1$ :
          $c_1(\eta) = h \in H^2(S^2; \Z) \cong \Z$.
    \item La classe de Chern totale d'un fibré trivial est $c(\varepsilon^k) = 1$.
  \end{enumerate}
\end{exemple}

\begin{theoreme}[Classification des fibrés]\label{thm:classification-fibres}
  \index{fibré vectoriel!classification}
  Les fibrés vectoriels réels de rang~$k$ sur un espace paracompact~$X$ sont
  classifiés par les classes d'homotopie d'applications $X \to BO(k)$, où $BO(k)$
  est l'espace classifiant\index{espace classifiant} du groupe orthogonal $O(k)$.
  Les classes de Stiefel--Whitney sont les images réciproques des classes universelles
  dans $H^*(BO(k); \Z/2\Z)$.
\end{theoreme}

\index{classes caractéristiques|)}

% ============================================================
\section{Espaces d'Eilenberg--MacLane}\label{sec:eilenberg-maclane}
\index{Eilenberg--MacLane, espaces d'|(}

\begin{definition}[Espace d'Eilenberg--MacLane]\label{def:eilenberg-maclane}
  Soit $G$ un groupe (abélien si $n \geq 2$) et $n \geq 1$. Un espace
  topologique~$X$ est un \emph{espace d'Eilenberg--MacLane de type $K(G, n)$} si
  \[
    \pi_k(X) \cong
    \begin{cases}
      G & \text{si } k = n, \\
      0 & \text{si } k \neq n.
    \end{cases}
  \]
\end{definition}

\begin{theoreme}[Existence et unicité]\label{thm:em-existence}
  \index{Eilenberg--MacLane, espaces d'!existence et unicité}
  Pour tout groupe~$G$ (abélien si $n \geq 2$) et tout $n \geq 1$, il existe un
  CW-complexe de type $K(G, n)$, unique à équivalence d'homotopie près.
\end{theoreme}

\begin{proof}[Esquisse de preuve]
  \emph{Existence :} on construit $K(G, n)$ par étapes.
  \begin{enumerate}[label=\arabic*.]
    \item On part d'un bouquet $\bigvee_\alpha S^n$ indexé par un système de
          générateurs de~$G$.
    \item On attache des $(n{+}1)$-cellules pour tuer les relations dans~$G$,
          obtenant un espace~$X$ avec $\pi_n(X) \cong G$ et $\pi_k(X) = 0$
          pour $k < n$.
    \item On attache des cellules de dimension $\geq n+2$ pour tuer les groupes
          d'homotopie $\pi_k(X)$ pour $k > n$, ce qui est possible par construction.
  \end{enumerate}
  \emph{Unicité :} si $X$ et $Y$ sont deux $K(G, n)$, on construit une application
  $f \colon X \to Y$ induisant un isomorphisme sur~$\pi_n$, et donc (par Hurewicz
  et Whitehead) une équivalence d'homotopie.
\end{proof}

\begin{exemple}\label{ex:em-espaces}
  \begin{enumerate}[label=(\roman*)]
    \item $K(\Z, 1) \simeq S^1$\index{$K(\Z,1)$} (car $\pi_1(S^1) = \Z$ et
          $\pi_n(S^1) = 0$ pour $n \geq 2$).
    \item $K(\Z, 2) \simeq \CP^\infty$\index{$K(\Z,2)$}, l'espace projectif complexe
          de dimension infinie.
    \item $K(\Z/2\Z, 1) \simeq \RP^\infty$\index{$K(\Z/2\Z,1)$}.
    \item $K(\Z/n\Z, 1)$ est un espace de lentilles\index{espace de lentilles} de
          dimension infinie.
  \end{enumerate}
\end{exemple}

\begin{remarque}[Relation avec les espaces de lacets]\label{rem:em-lacets}
  \index{espace de lacets!et $K(G,n)$}
  Pour tout groupe abélien~$G$ et tout $n \geq 1$, on a une équivalence d'homotopie
  $\Omega K(G, n+1) \simeq K(G, n)$. En effet, la suite exacte d'homotopie de la
  fibration des chemins $\Omega K(G, n+1) \hookrightarrow PK(G, n+1) \to K(G, n+1)$,
  combinée avec la contractibilité de $PK(G, n+1)$, donne
  $\pi_k(\Omega K(G, n+1)) \cong \pi_{k+1}(K(G, n+1))$. Or
  $\pi_{k+1}(K(G, n+1)) = G$ si $k+1 = n+1$ (c'est-à-dire $k = n$) et $0$ sinon.
  Donc $\Omega K(G, n+1)$ est un $K(G, n)$.
\end{remarque}

\begin{proposition}[Propriétés de la cohomologie des espaces $K(G,n)$]
  \label{prop:cohomologie-KGn}
  \begin{enumerate}[label=(\roman*)]
    \item $H^n(K(G, n); G) \cong \Hom(G, G)$ contient un élément distingué~$\iota_n$
          (correspondant à l'identité de~$G$), appelé \emph{classe fondamentale}.
    \item L'anneau $H^*(K(\Z, 2); \Z) \cong \Z[c]$ est un anneau de polynômes
          en un générateur $c$ de degré~$2$ (cohomologie de $\CP^\infty$).
    \item L'algèbre $H^*(K(\Z/2\Z, 1); \Z/2\Z) \cong (\Z/2\Z)[w]$ est un anneau
          de polynômes en un générateur $w$ de degré~$1$ (cohomologie de $\RP^\infty$).
  \end{enumerate}
\end{proposition}

% ============================================================
\section{Cohomologie et classes d'homotopie}\label{sec:representabilite}
\index{représentabilité de la cohomologie}

\begin{theoreme}[Représentabilité de la cohomologie]\label{thm:representabilite-brown}
  \index{Brown, théorème de représentabilité de}
  Soit $G$ un groupe abélien et $n \geq 1$. Pour tout CW-complexe~$X$, il existe
  une bijection naturelle
  \[
    [X, K(G, n)] \xrightarrow{\;\sim\;} H^n(X; G)
  \]
  qui, à la classe d'homotopie d'une application $f \colon X \to K(G, n)$,
  associe $f^*(\iota_n)$, où $\iota_n \in H^n(K(G,n); G)$ est la \emph{classe
  fondamentale}\index{classe fondamentale!d'un espace $K(G,n)$}.
\end{theoreme}

\begin{proof}[Esquisse de preuve]
  La surjectivité résulte de la construction d'une application $f \colon X \to K(G,n)$
  cellule par cellule, à partir d'un cocycle représentant une classe de cohomologie
  donnée. L'injectivité se démontre en montrant que deux applications induisant la
  même classe en cohomologie sont homotopes, par un argument d'obstruction
  (les obstructions vivent dans des groupes de cohomologie à coefficients dans
  $\pi_k(K(G,n)) = 0$ pour $k \neq n$). Ce théorème est un cas particulier du
  \emph{théorème de représentabilité de Brown}\index{Brown, théorème de
  représentabilité de}, qui s'applique à tout foncteur cohomologique
  vérifiant les axiomes de Mayer--Vietoris et de wedge.
\end{proof}

\begin{remarque}\label{rem:operations-cohomologiques}
  \index{opérations cohomologiques}
  Le théorème de représentabilité implique que les \emph{opérations cohomologiques}
  naturelles $H^n(-; G) \to H^m(-; G')$ correspondent bijectivement aux éléments de
  $H^m(K(G,n); G') = [K(G,n), K(G',m)]$. Par exemple, les \emph{carrés de
  Steenrod}\index{carrés de Steenrod} $\mathrm{Sq}^i \colon H^n(-; \Z/2\Z) \to
  H^{n+i}(-; \Z/2\Z)$ correspondent à des classes dans
  $H^{n+i}(K(\Z/2\Z, n); \Z/2\Z)$.
\end{remarque}

% ============================================================
\section{Ouvertures}\label{sec:ouvertures}

\subsection{Homotopie stable et spectres}\label{subsec:homotopie-stable}
\index{homotopie stable}
\index{spectre (topologie)}

Le théorème de suspension de Freudenthal montre que les morphismes de suspension
$\pi_n(S^k) \to \pi_{n+1}(S^{k+1})$ se stabilisent pour $n < 2k - 1$.
Les \emph{groupes d'homotopie stables} $\pi_n^s = \varinjlim_k \pi_{n+k}(S^k)$
forment un anneau gradué (via le produit de composition), dont le calcul est l'un
des problèmes centraux de la topologie algébrique moderne.

\begin{definition}[Spectre]\label{def:spectre}
  Un \emph{spectre} est une suite d'espaces pointés $(E_n)_{n \geq 0}$ munis
  d'applications $\sigma_n \colon \Sigma E_n \to E_{n+1}$ (ou de manière
  équivalente, d'applications $\tilde{\sigma}_n \colon E_n \to \Omega E_{n+1}$).
  Un spectre définit une théorie de cohomologie généralisée par
  $E^n(X) = \varinjlim_k [\Sigma^k X, E_{n+k}]$.
\end{definition}

\begin{exemple}\label{ex:spectres}
  \begin{enumerate}[label=(\roman*)]
    \item Le \emph{spectre de sphères}\index{spectre de sphères}
          $\mathbb{S} = (S^0, S^1, S^2, \ldots)$ représente l'homotopie stable.
    \item Le \emph{spectre d'Eilenberg--MacLane}\index{spectre d'Eilenberg--MacLane}
          $HG = (K(G,0), K(G,1), K(G,2), \ldots)$ représente la cohomologie singulière.
    \item Le \emph{spectre de Thom}\index{spectre de Thom}
          $MO = (MO(0), MO(1), MO(2), \ldots)$ représente le cobordisme non orienté.
  \end{enumerate}
\end{exemple}

\subsection{Connexions avec la physique mathématique}\label{subsec:physique}
\index{physique mathématique}

\begin{remarque}\label{rem:physique}
  Les concepts de ce chapitre trouvent des applications profondes en physique
  théorique, notamment en \emph{théorie de jauge}\index{théorie de jauge} :
  \begin{enumerate}[label=(\roman*)]
    \item Un champ de jauge (connexion) sur un espace-temps~$M$ est une connexion
          sur un fibré principal $P \to M$ de groupe de structure~$G$
          (par exemple $G = SU(2)$ ou $G = SU(3)$).
    \item Les classes caractéristiques (classes de Chern) du fibré mesurent les
          \emph{charges topologiques} ; par exemple, la deuxième classe de Chern
          $c_2(P)$ est le \emph{nombre d'instantons}\index{instanton}.
    \item Les \emph{anomalies}\index{anomalie (physique)} en théorie quantique des
          champs correspondent à des obstructions topologiques à la trivialisation
          de certains fibrés déterminants.
    \item La classification des phases topologiques de la matière utilise la
          $K$-théorie\index{K-théorie} et le cobordisme.
  \end{enumerate}
\end{remarque}

% ============================================================
\section{Exercices}\label{sec:exo-obstruction}

\begin{exercice}\label{exo:obstruction-1}
  Soit $X$ un CW-complexe de dimension $\leq n$ et $Y$ un espace $(n{-}1)$-connexe.
  Montrer que toute application $f \colon A \to Y$ (où $A$ est un sous-complexe
  de~$X$) s'étend à~$X$.
\end{exercice}

\begin{exercice}\label{exo:obstruction-2}
  Montrer que tout fibré vectoriel de rang~$k$ sur une sphère~$S^n$ est trivial si
  $n < k$. \emph{Indication :} utiliser le fait que la fibre du fibré en repères
  est $(k{-}2)$-connexe.
\end{exercice}

\begin{exercice}\label{exo:obstruction-3}
  Soit $\xi$ un fibré vectoriel réel de rang~$2$ sur un CW-complexe de dimension~$2$.
  Montrer que $\xi$ est trivial si et seulement si $w_1(\xi) = 0$ et $e(\xi) = 0$
  (où $e$ est la classe d'Euler).
\end{exercice}

\begin{exercice}\label{exo:obstruction-4}
  En utilisant le théorème de représentabilité de Brown, montrer que pour tout
  groupe abélien~$G$, la première classe de cohomologie $H^1(X; G)$ est en bijection
  avec $\Hom(\pi_1(X), G)$ lorsque $X$ est connexe par arcs.
\end{exercice}

\begin{exercice}\label{exo:obstruction-5}
  Montrer que $K(\Z, n)$ a le type d'homotopie de $\Omega K(\Z, n+1)$.
  En déduire que la cohomologie singulière vérifie l'axiome de la suite de
  Mayer--Vietoris.
\end{exercice}

\begin{exercice}\label{exo:obstruction-6}
  Calculer les premiers groupes de cohomologie de $K(\Z, 2) \simeq \CP^\infty$ et
  retrouver l'anneau de cohomologie $H^*(\CP^\infty; \Z) \cong \Z[c]$, $\deg c = 2$.
\end{exercice}

\begin{exercice}\label{exo:obstruction-7}
  \textbf{(Application du théorème de Borsuk--Ulam.)}
  Montrer qu'il n'existe pas de plongement topologique de $\RP^2$ dans~$\R^3$.
  \emph{Indication :} supposer qu'un tel plongement existe, considérer le complémentaire
  et utiliser la dualité d'Alexander.
\end{exercice}

\begin{exercice}\label{exo:obstruction-8}
  \index{catégorie de Lusternik--Schnirelman}
  La \emph{catégorie de Lusternik--Schnirelman} $\mathrm{cat}(X)$ est le plus petit
  entier~$k$ tel que $X$ puisse être recouvert par $k+1$ ouverts contractiles
  dans~$X$. Montrer que $\mathrm{cat}(S^n) = 1$ et que
  $\mathrm{cat}(T^n) = n$.
\end{exercice}

\begin{exercice}\label{exo:obstruction-9}
  Montrer que pour tout CW-complexe fini~$X$ et tout groupe abélien~$G$, l'ensemble
  $[X, K(G, n)]$ est un groupe abélien, et que ce groupe est isomorphe à~$H^n(X; G)$.
  Expliciter cet isomorphisme pour $X = S^n$ et $G = \Z$.
\end{exercice}

\begin{exercice}\label{exo:obstruction-10}
  \textbf{(Tour de Postnikov.)}\index{tour de Postnikov}
  Soit $X$ un CW-complexe connexe. Montrer qu'il existe des espaces $X_n$ et
  des fibrations $X_n \to X_{n-1}$ de fibre $K(\pi_n(X), n)$, tels que
  $\pi_k(X_n) \cong \pi_k(X)$ pour $k \leq n$ et $\pi_k(X_n) = 0$ pour $k > n$.
  Les obstructions à relever des applications vers $X_{n-1}$ en applications vers
  $X_n$ vivent dans $H^{n+1}(-; \pi_n(X))$.
\end{exercice}

% ============================================================
\section*{Résumé du chapitre}
\addcontentsline{toc}{section}{Résumé du chapitre}

\begin{itemize}
  \item La théorie de l'obstruction analyse l'extension d'applications cellule par
        cellule. L'obstruction à l'extension au $(n{+}1)$-squelette est un cocycle
        à valeurs dans~$\pi_n(Y)$.
  \item La classe de cohomologie $[\mathfrak{o}^{n+1}] \in H^{n+1}(X, A; \pi_n(Y))$
        est l'obstruction fondamentale.
  \item Le théorème de Borsuk--Ulam et le théorème de Lusternik--Schnirelman sont
        des applications classiques de ces méthodes.
  \item Les classes caractéristiques (Stiefel--Whitney, Chern, Pontriaguine, Euler)
        mesurent les obstructions à la trivialisation des fibrés vectoriels.
  \item Les espaces $K(G,n)$ représentent la cohomologie :
        $[X, K(G,n)] \cong H^n(X; G)$.
  \item La théorie des spectres fournit le cadre moderne pour les théories
        cohomologiques généralisées.
\end{itemize}

\index{théorie de l'obstruction|)}


% ============================================================
%  APPENDICE : Rappels d'algèbre homologique
% ============================================================
\appendix
\chapter{Rappels d'algèbre homologique}\label{app:algebre-homologique}
\index{algèbre homologique|(}

Cet appendice rassemble les notions fondamentales d'algèbre homologique utilisées
tout au long du cours. Les preuves sont omises ou seulement esquissées ; le lecteur
pourra se reporter aux ouvrages de Weibel~\cite{weibel} ou de
Cartan--Eilenberg~\cite{cartan-eilenberg} pour les détails.

\section{Complexes de chaînes}\label{sec:complexes-chaines}
\index{complexe de chaînes}

\begin{definition}[Complexe de chaînes]\label{def:complexe-chaines}
  Un \emph{complexe de chaînes} $(C_\bullet, \partial)$ sur un anneau~$R$ est une
  suite de $R$-modules et de morphismes
  \[
    \cdots \xrightarrow{\partial_{n+1}} C_n \xrightarrow{\partial_n}
    C_{n-1} \xrightarrow{\partial_{n-1}} \cdots
    \xrightarrow{\partial_1} C_0 \xrightarrow{\partial_0} 0
  \]
  telle que $\partial_n \circ \partial_{n+1} = 0$ pour tout~$n$. Dualement, un
  \emph{complexe de cochaînes}\index{complexe de cochaînes} est une suite
  $(C^\bullet, \delta)$ avec $\delta^{n+1} \circ \delta^n = 0$.
\end{definition}

\begin{definition}[Homologie]\label{def:homologie-complexe}
  Le $n$-ième \emph{groupe d'homologie}\index{homologie!d'un complexe} du complexe
  $(C_\bullet, \partial)$ est
  \[
    H_n(C_\bullet) = \ker \partial_n / \im \partial_{n+1}
    = Z_n(C_\bullet) / B_n(C_\bullet),
  \]
  où $Z_n = \ker \partial_n$ est le module des \emph{cycles}\index{cycle} et
  $B_n = \im \partial_{n+1}$ celui des \emph{bords}\index{bord}.
\end{definition}

\begin{definition}[Morphisme de complexes]\label{def:morphisme-complexes}
  Un \emph{morphisme de complexes} $f \colon C_\bullet \to D_\bullet$ est une
  famille de morphismes $f_n \colon C_n \to D_n$ commutant avec les différentielles :
  $\partial^D_n \circ f_n = f_{n-1} \circ \partial^C_n$. Un tel morphisme induit
  des morphismes en homologie $f_* \colon H_n(C_\bullet) \to H_n(D_\bullet)$.
\end{definition}

\begin{definition}[Homotopie de chaînes]\label{def:homotopie-chaines}
  \index{homotopie de chaînes}
  Deux morphismes $f, g \colon C_\bullet \to D_\bullet$ sont \emph{homotopes}
  (noté $f \simeq g$) s'il existe des morphismes $s_n \colon C_n \to D_{n+1}$ tels
  que $f_n - g_n = \partial^D_{n+1} \circ s_n + s_{n-1} \circ \partial^C_n$.
  Dans ce cas, $f_* = g_*$ en homologie.
\end{definition}

\section{Suites exactes}\label{sec:suites-exactes}
\index{suite exacte}

\begin{definition}[Suite exacte]\label{def:suite-exacte}
  Une suite de morphismes de $R$-modules
  $A \xrightarrow{f} B \xrightarrow{g} C$
  est dite \emph{exacte en~$B$} si $\ker g = \im f$. Une \emph{suite exacte courte}
  \index{suite exacte!courte} est une suite exacte
  $0 \to A \xrightarrow{f} B \xrightarrow{g} C \to 0$.
\end{definition}

\begin{theoreme}[Suite exacte longue en homologie]\label{thm:sel-homologie}
  \index{suite exacte longue!en homologie}
  Toute suite exacte courte de complexes de chaînes
  $0 \to A_\bullet \to B_\bullet \to C_\bullet \to 0$
  induit une suite exacte longue en homologie :
  \[
    \cdots \to H_n(A) \to H_n(B) \to H_n(C) \xrightarrow{\partial_*}
    H_{n-1}(A) \to H_{n-1}(B) \to \cdots
  \]
  Le morphisme $\partial_* \colon H_n(C) \to H_{n-1}(A)$ est appelé
  \emph{morphisme connectant}\index{morphisme connectant}. Sa construction repose
  sur le \emph{lemme du serpent}\index{lemme du serpent}.
\end{theoreme}

\begin{theoreme}[Lemme des cinq]\label{thm:lemme-cinq}
  \index{lemme des cinq}
  Soit le diagramme commutatif à lignes exactes :
  \[
    \begin{tikzcd}
      A_1 \ar[r] \ar[d, "f_1"] & A_2 \ar[r] \ar[d, "f_2"] & A_3 \ar[r] \ar[d, "f_3"]
      & A_4 \ar[r] \ar[d, "f_4"] & A_5 \ar[d, "f_5"] \\
      B_1 \ar[r] & B_2 \ar[r] & B_3 \ar[r] & B_4 \ar[r] & B_5
    \end{tikzcd}
  \]
  Si $f_1, f_2, f_4, f_5$ sont des isomorphismes, alors $f_3$ est un isomorphisme.
\end{theoreme}

\section{Formule de Künneth}\label{sec:kunneth-app}
\index{formule de Künneth}

\begin{theoreme}[Künneth]\label{thm:kunneth-app}
  Soient $C_\bullet$ et $D_\bullet$ des complexes de chaînes de $R$-modules libres,
  où $R$ est un anneau principal. Alors il existe une suite exacte courte naturelle
  \[
    0 \to \bigoplus_{p+q=n} H_p(C) \otimes_R H_q(D) \to H_n(C \otimes_R D)
    \to \bigoplus_{p+q=n-1} \mathrm{Tor}_1^R(H_p(C), H_q(D)) \to 0,
  \]
  scindée (non naturellement). En particulier, si $R$ est un corps, on a
  $H_n(C \otimes_R D) \cong \bigoplus_{p+q=n} H_p(C) \otimes_R H_q(D)$.
\end{theoreme}

\begin{corollaire}[Künneth topologique]\label{cor:kunneth-topo}
  Pour des espaces topologiques~$X$ et~$Y$ et un corps~$\mathbb{F}$ :
  \[
    H_n(X \times Y; \mathbb{F}) \cong
    \bigoplus_{p+q=n} H_p(X; \mathbb{F}) \otimes_\mathbb{F} H_q(Y; \mathbb{F}).
  \]
\end{corollaire}

\section{Foncteurs dérivés}\label{sec:foncteurs-derives}
\index{foncteurs dérivés|(}

\begin{definition}[Module projectif, module injectif]\label{def:projectif-injectif}
  \index{module projectif}\index{module injectif}
  Un $R$-module~$P$ est \emph{projectif} si pour tout épimorphisme $B \twoheadrightarrow C$
  et tout morphisme $P \to C$, il existe un relèvement $P \to B$. Un
  $R$-module~$I$ est \emph{injectif} si pour tout monomorphisme $A \hookrightarrow B$
  et tout morphisme $A \to I$, il existe une extension $B \to I$.
\end{definition}

\begin{definition}[Résolution projective]\label{def:resolution-projective}
  \index{résolution projective}
  Une \emph{résolution projective} d'un $R$-module~$M$ est une suite exacte
  $\cdots \to P_1 \to P_0 \to M \to 0$
  où chaque $P_n$ est projectif.
\end{definition}

\begin{definition}[Foncteurs Ext et Tor]\label{def:ext-tor}
  \index{Ext}\index{Tor}
  Soit $M$ un $R$-module et $P_\bullet \to M \to 0$ une résolution projective.
  \begin{enumerate}[label=(\roman*)]
    \item $\mathrm{Tor}_n^R(M, N) = H_n(P_\bullet \otimes_R N)$.
    \item $\mathrm{Ext}_R^n(M, N) = H^n(\Hom_R(P_\bullet, N))$.
  \end{enumerate}
  Ces foncteurs sont indépendants du choix de la résolution (à isomorphisme
  canonique près).
\end{definition}

\begin{theoreme}[Théorème des coefficients universels]\label{thm:coeff-universels-app}
  \index{théorème des coefficients universels}
  Pour un espace topologique~$X$ et un groupe abélien~$G$ :
  \[
    0 \to \mathrm{Ext}^1_\Z(H_{n-1}(X; \Z), G) \to H^n(X; G)
    \to \Hom(H_n(X; \Z), G) \to 0.
  \]
  Cette suite exacte est scindée (non naturellement).
\end{theoreme}

\begin{exemple}\label{ex:ext-tor-calculs}
  Quelques calculs fondamentaux :
  \begin{enumerate}[label=(\roman*)]
    \item $\mathrm{Tor}_1^\Z(\Z/m\Z, \Z/n\Z) \cong \Z/\gcd(m,n)\Z$.
    \item $\mathrm{Ext}^1_\Z(\Z/m\Z, \Z/n\Z) \cong \Z/\gcd(m,n)\Z$.
    \item $\mathrm{Tor}_1^\Z(\Z/m\Z, G) \cong \set{g \in G \mid mg = 0}$
          pour tout groupe abélien~$G$.
    \item $\mathrm{Ext}^1_\Z(\Z, G) = 0$ pour tout groupe abélien~$G$
          (car $\Z$ est libre, donc projectif).
    \item Si $P$ est un module projectif, $\mathrm{Ext}^n_R(P, M) = 0$ pour tout
          $n \geq 1$ et tout $R$-module~$M$.
  \end{enumerate}
\end{exemple}

\begin{remarque}[Suites spectrales]\label{rem:suites-spectrales}
  \index{suites spectrales}
  Pour les calculs avancés en topologie algébrique, les \emph{suites spectrales}
  constituent un outil indispensable. Une suite spectrale est une suite de pages
  $(E_r^{p,q}, d_r)_{r \geq 0}$ de bimodules différentiels, avec
  $E_{r+1} = H(E_r, d_r)$, qui <<~converge~>> vers un objet gradué associé à une
  filtration de l'objet final. Les plus utilisées sont :
  \begin{enumerate}[label=(\roman*)]
    \item La suite spectrale de \emph{Serre}\index{suite spectrale!de Serre}
          pour une fibration $F \to E \to B$ :
          $E_2^{p,q} = H^p(B; H^q(F)) \Rightarrow H^{p+q}(E)$.
    \item La suite spectrale de \emph{Mayer--Vietoris}\index{suite spectrale!de Mayer--Vietoris}
          pour un recouvrement ouvert.
    \item La suite spectrale des \emph{coefficients universels} pour le calcul
          de la cohomologie à coefficients.
  \end{enumerate}
  Nous renvoyons à McCleary~\cite{weibel} pour une introduction détaillée.
\end{remarque}

\index{foncteurs dérivés|)}
\index{algèbre homologique|)}


% ============================================================
%  BIBLIOGRAPHIE
% ============================================================
\begin{thebibliography}{99}

\bibitem{hatcher}
  A.~Hatcher,
  \emph{Algebraic Topology},
  Cambridge University Press, 2002.
  \index{Hatcher}

\bibitem{spanier}
  E.~H.~Spanier,
  \emph{Algebraic Topology},
  Springer-Verlag, 1966 ; réédition corrigée, 1981.
  \index{Spanier}

\bibitem{bredon}
  G.~E.~Bredon,
  \emph{Topology and Geometry},
  Graduate Texts in Mathematics~139, Springer, 1993.
  \index{Bredon}

\bibitem{may}
  J.~P.~May,
  \emph{A Concise Course in Algebraic Topology},
  Chicago Lectures in Mathematics, University of Chicago Press, 1999.
  \index{May}

\bibitem{tomdieck}
  T.~tom~Dieck,
  \emph{Algebraic Topology},
  EMS Textbooks in Mathematics, European Mathematical Society, 2008.
  \index{tom Dieck}

\bibitem{eilenberg-steenrod}
  S.~Eilenberg et N.~E.~Steenrod,
  \emph{Foundations of Algebraic Topology},
  Princeton University Press, 1952.
  \index{Eilenberg}\index{Steenrod}

\bibitem{godement}
  R.~Godement,
  \emph{Topologie algébrique et théorie des faisceaux},
  Hermann, Paris, 1958.
  \index{Godement}

\bibitem{serre}
  J.-P.~Serre,
  <<\,Homologie singulière des espaces fibrés. Applications\,>>,
  \emph{Annals of Mathematics}, 54 (1951), 425--505.
  \index{Serre}

\bibitem{bott-tu}
  R.~Bott et L.~W.~Tu,
  \emph{Differential Forms in Algebraic Topology},
  Graduate Texts in Mathematics~82, Springer, 1982.
  \index{Bott}\index{Tu}

\bibitem{milnor-stasheff}
  J.~W.~Milnor et J.~D.~Stasheff,
  \emph{Characteristic Classes},
  Annals of Mathematics Studies~76, Princeton University Press, 1974.
  \index{Milnor}\index{Stasheff}

\bibitem{weibel}
  C.~A.~Weibel,
  \emph{An Introduction to Homological Algebra},
  Cambridge Studies in Advanced Mathematics~38, Cambridge University Press, 1994.
  \index{Weibel}

\bibitem{cartan-eilenberg}
  H.~Cartan et S.~Eilenberg,
  \emph{Homological Algebra},
  Princeton University Press, 1956.
  \index{Cartan}

\bibitem{whitehead}
  G.~W.~Whitehead,
  \emph{Elements of Homotopy Theory},
  Graduate Texts in Mathematics~61, Springer, 1978.
  \index{Whitehead, G.\,W.}

\bibitem{switzer}
  R.~M.~Switzer,
  \emph{Algebraic Topology: Homotopy and Homology},
  Classics in Mathematics, Springer, 2002.
  \index{Switzer}

\bibitem{dold}
  A.~Dold,
  \emph{Lectures on Algebraic Topology},
  Classics in Mathematics, Springer, 1995.
  \index{Dold}

\end{thebibliography}


% ============================================================
%  INDEX
% ============================================================
\printindex

\end{document}
% === FIN DU CHUNK ch9-11+end ===
